\documentclass[11pt,a4paper]{article}
\ifdefined\pdfmapfile
\pdfmapfile{+cm.map}
\pdfmapfile{+cmextra.map}
\pdfmapfile{+symbols.map}
\pdfmapfile{+euler.map}
\pdfmapfile{+lm.map}
\pdfmapfile{+rsfs.map}
\fi
\usepackage{lmodern}
\usepackage[margin=25mm]{geometry}
\usepackage{amsmath,amssymb,amsthm,mathtools,mathrsfs}
\usepackage{booktabs,tabularx,longtable,array}
\usepackage[hyphens]{url}
\usepackage[colorlinks=true,linkcolor=blue!45!black,citecolor=blue!45!black,urlcolor=blue!45!black]{hyperref}
\usepackage{xcolor}
\usepackage{enumitem}
\setlist{itemsep=3pt,topsep=5pt}
\setlength{\emergencystretch}{3em}
\setlength{\parskip}{4pt}
\setlength{\parindent}{0pt}
\renewcommand{\arraystretch}{1.17}
\newtheorem{proposition}{Proposition}[section]
\newtheorem{corollary}[proposition]{Corollary}
\newtheorem{theorem}[proposition]{Theorem}
\newtheorem{lemma}[proposition]{Lemma}
\hypersetup{pdftitle={Purity, positive Hilbert series, and the GCT flip},pdfauthor={Research analysis},pdfsubject={Checked deductions, exact source scopes, and open gates}}
\title{Purity, positive Hilbert series, and the GCT flip\\[5pt]\large Checked deductions and remaining gates}
\author{}
\date{21 November 2026}
\begin{document}
\maketitle
\begin{abstract}
We examine what the supplied RH/Dirichlet-GRH and subexponential counting
theorems, and their proposed positivity methods, yield for geometric
complexity theory. For ordinary cubic plethysm, a complete classification
of extremal two-power stabilizer supports makes the chain curve the unique
nonzero fixed component of the invariant quotient. Its Hilbert series
has actual simple poles at roots of unity of order $2^{n-1}$,
although the conjectured size
parameter is $O(n^6)$. Thus both the numerical positive-form clause and
the stronger canonical clause of GCT VI PH3 are false. An explicit
polynomial bound for the coefficient
positivity index of a positive rational Hilbert series proves the revised
GCT implication PH3 $\Rightarrow$ PH2. Fixed-rank torus embeddings refute
the printed rank-only bounds in the broader subgroup-restriction problem.
Fixed inner dimension nevertheless gives a finite uniform bound; its
optimal canonical envelope is therefore finite but not polynomially
bounded. The pole theorem does not settle PH2 or saturation; an
exterior-power reduction gives a separate square-root lower bound for
their indices on an infinite cubic family. Positive Hilbert data yield shifted polytope models; the
standard plethysm counterexample has an optimal two-dimensional model
and an actual normal multiplicity algebra with a positive multiplicative
basis and pure even Tate intersection cohomology. These results separate qualitative purity
from the additional grading, uniformity and quotient-rank information
needed for an obstruction. A named regularized theta construction
identifies the degree-$80$ cubic discriminant in the actual five-variable
plethysm algebra. Its exact input has a specified definite-unitary theta
factor and Hilbert-scheme cohomology models for all its coefficients.
For the nonstandard cubic example, an explicit positive central word
realizes the quadratic operator inside the actual algebra; an alternative
normalization proves that its positive action and Lefschetz trace and
norm do not characterize the source scale. An unconditional cubic average-polynomial estimate
and a signed global cone decomposition further delimit the remaining
index problem. Explicit Hessian contractions give polynomial Fermat
nonvanishing and conductor bounds, sharpened by the supplied Dirichlet
theorem. These are actual-support statements, separate from the strict
constituent indices. We give checked deductions for all seven
requested Grothendieck bridges, distinguish Dirichlet from number-field
GRH in complexity applications, and track the precise consequences and
limitations of the supplied ETH theorem. No resolution of $\mathsf P$
versus $\mathsf{NP}$ or the usual $\mathsf{VP}$ versus $\mathsf{VNP}$
problem is claimed; unresolved targets are listed separately.
\end{abstract}
{\small\setlength{\parskip}{0pt}\tableofcontents}
\clearpage


\section*{Scope and research status}
\addcontentsline{toc}{section}{Scope and research status}
This paper treats the supplied Alp\"oge and OpenAI theorems as cited
results, as requested. It distinguishes those premises from deductions
proved in this run. ``Finding'' means such a deduction, including a
counterexample or an elementary sharpened form of a known result; it does
not assert historical priority. Exact source locators accompany the
statements used. The accompanying check register describes what was
actually checked and does not promote a text scan or a numerical test to
a formal proof certificate.

The main negative conclusion is P-F31: both clauses of the ordinary
plethysm PH3 in GCT VI, Hypothesis~1.6.7, are false. The family has
the single inner shape $(3)$. Its invariant Hilbert series has an
actual simple pole at every primitive $2^{n-1}$st root, although the
source's parameter is $O(n^6)$. Every rational expression with
polynomial numerator and factors $1-t^a$ in its denominator must
therefore use some $a\ge2^{n-1}$. P-F30 supplies the complete
support classification that excludes every competing nonzero fixed
component; P-F23, P-F24 and P-F28 then prove the pole and its leading
coefficient. P-F19 had already refuted the canonical parameter clause.
These results do not refute the narrower PH2 or saturation bounds.

A central positive result obtained here is P-F11: a global
quasi-polynomial with nonnegative rational Hilbert numerator and
$r$ denominator factors of degrees at most $A$ has positivity index
less than $rA+128r^3A^3$. P-F12 applies this to show that the revised
PH3 of GCT VI implies its revised PH2. This is an implication from
bounded positive Hilbert data, not a proof that the requested
denominator bounds exist.
P-F14--P-F15 separately refute the rank-only PH3, PH2 and saturation
bounds in the broader subgroup-restriction formulation: faithful
polynomial torus embeddings have fixed group ranks and unbounded
indices as their character weights vary. The narrower plethysm
formulation is not refuted by those examples.
P-F16--P-F18 prove finite uniform bounds for a fixed semisimple
embedding and, for ordinary $GL_n$ plethysm, for fixed inner dimension.
P-F19 then gives an exponential lower bound for that optimal canonical
modular-index envelope, using the single inner shape $(3)$.
It refutes the stronger canonical clause of ordinary plethysm PH3.
P-F31 additionally refutes the assertion about some positive rational
Hilbert expression. The narrower PH2 and saturation bounds remain
separate questions.
P-F20 distinguishes a large universal denominator from an ordinary
Hilbert-series pole. P-F21--P-F22 refute general transfers from large
parameter degrees even for rational Gorenstein domains with one fixed
Hilbert series. P-F23 supplies a precise noncancellation criterion;
P-F24 proves that the chain gives a smooth one-dimensional fixed
component in the actual cubic quotient. P-F26 confines every competing component
to the discriminant boundary. P-F27 proves that any such component
requires three-distinct-variable monomials in every finite-lift
eigenbasis; even a connected graph part forces the smooth chain.
P-F29 excludes the complementary cube-and-squarefree support class,
so any remaining supports would have to mix the two types with
disconnected graph part. P-F30 closes that final support case,
including rank-deficient supports and arbitrary two-power lift order.
P-F28 gives the exact local chain coefficient and its cyclotomic
valuation; P-F31 identifies it with the actual leading coefficient.
P-F32 supplies a separate exterior-duality reduction and an exact
first nonzero stretch on an infinite cubic family, giving a
$\Theta(\sqrt n)$ lower bound compatible with polynomial PH2 and SH.
P-F33 prevents an inference from large period and Gorenstein
reciprocity alone to large indices. P-F34 proves a large Fourier
amplitude for the actual cubic family and reduces one route to
an exponential index lower bound to a specified, still unproved
bound on the residue-average polynomial. P-F35 computes both
indices exactly as $23$ in the four-variable case.
Positive Hilbert forms also yield
shifted polytope models (P-F2), and the main counterexample ray admits
an optimal two-dimensional model (P-F1).

The converse direction from qualitative purity to the required counting
and complexity properties fails in concrete examples. PLE-F7 supplies a
normal surface with pure even Tate ordinary and intersection cohomology
whose Hilbert function violates homogeneous PH1. PLE-F8 identifies it
with the actual canonical plethysm multiplicity algebra on that ray,
including multiplication, and supplies a basis with nonnegative
multiplication constants. P-F10 keeps the smooth
underlying variety fixed while its grading indices grow arbitrarily.
P-F7 shows that finite smooth proper pure Tate schemes can encode all
\#3SAT counts when specified by equations. These examples identify
missing hypotheses; they do not prove a logical independence theorem
between RH and any GCT assertion.

No separation or equality of $\mathsf P$ and $\mathsf{NP}$, or of the
usual complex algebraic classes $\mathsf{VP}$ and $\mathsf{VNP}$, is
proved. The remaining research targets are displayed separately in
Part~II and in \texttt{QUESTIONS.tsv}. A statement is not marked proved
merely because its missing input has been named.

References to numbered GCT VI hypotheses concern the supplied
arXiv:0704.0229v4 text. The broad subgroup-restriction formulation
of Problem 1.1.3 and Sections 3.3--3.4 is distinguished throughout
from irreducible-inner Problem 1.1.2 and its ordinary $GL_n$
plethysm subfamily. Introductory notational inconsistencies about
the canonical invariant algebra are resolved using the explicit
construction and modular-index definition in Sections 4.1.1--4.1.2.

\subsection*{The supplied analytic statements}
The five theorem declarations in \texttt{ALPOGE\_Main.lean} are the
precise analytic premises. Their types are imported from the supplied
\texttt{ALPOGE\_Statements.lean}, except for Mathlib's
\texttt{RiemannHypothesis} itself.
\begin{longtable}{@{}p{.48\linewidth}p{.46\linewidth}@{}}
\toprule
Theorem name (statement identifier) & Scope used here\\
\midrule
\endhead
\path{riemannHypothesis} (A-RH) & Mathlib's named Riemann-hypothesis proposition.\\[3pt]
\path{riemannZeta_eq_zero_re} (A-ZETA) & A zero of $\zeta(s)$, other than a negative even integer and with $s\ne1$, has $\Re s=1/2$.\\[3pt]
\path{LFunction_ne_zero_of_half_lt_re} (A-DIR-NV) & For every positive modulus and Dirichlet character, $L(\chi,s)\ne0$ if $\Re s>1/2$ and $(\chi\ne1\text{ or }s\ne1)$.\\[3pt]
\path{LFunction_eq_zero_re} (A-DIR-Z) & All Dirichlet $L$-zeros in $0<\Re s<1$ have real part $1/2$.\\[3pt]
\path{LFunction_eq_zero_re_of_isPrimitive} (A-DIR-PRIM) & The corresponding statement restricted to primitive Dirichlet characters.\\
\bottomrule
\end{longtable}
These declarations quantify over Dirichlet characters, not arbitrary
number fields or their Hecke characters. C-GRH-01 proves the exact
cyclotomic-field extension needed below. The present working set contains
the theorem wrappers, their statement file and a supplied scan report;
it does not contain all imported proof modules. No Lean kernel rebuild
was performed. This records the scope of this run's verification, without
reopening the requested treatment of the theorems as cited results.

\subsection*{The positivity method in the supplied note}
The determinant sign $W(1_S)\ge0$ is in \emph{Remark}, p.~3, numbered
item~4. The sign-only discussion of finite-field purity is in footnote~3
on that page; numbered item~5 concerns switching to the von Mangoldt
sum. GGP periods and interlacing appear in item~7, iterated theta lifts in
item~8, and the Bessel-period seesaw in item~12 on p.~4. We use these
locators rather than conflating the footnote with numbered item~5.

The transfer question is mathematical: identify the representation
multiplicity or graded coordinate-ring quotient with the object to which
a sign or purity theorem applies, then supply quantitative control.
P-F3 constructs a positive determinant for each coefficient
unconditionally, exposing the dimension cost. P-F4 separates pointwise
sign from character positivity. P-F5 distinguishes standard interlacing
from restriction along a Schur-functor homomorphism. P-F6 states the
precise categorical data sufficient for positivity of a based product.
Neither the four-page note nor the five theorem types supplies those GCT
identifications automatically.
P-F36 supplies a more specific requirement in the actual nonstandard
cubic example: a quadratic coefficient extension survives every
finite power regrading. Rational projectors cannot isolate its
conjugate blocks, and finite graded Laurent characters cannot
produce the displayed raw coefficient. A proposed transfer must
account for this descent datum; the relaxed conjectures themselves
are not refuted.

\subsection*{Reading the package}
Part~I contains the checked findings and the exact cited background on
which they rely. Part~II contains unresolved targets and unsuccessful
transfer attempts; it asserts no unchecked theorem. The source archive
is identified by a SHA-256 manifest. The audit register records model,
effort disclosure, agent/run, date, method and outcome for each check.
The finite computations are sanity checks, separately labelled from
whole-proof readings and independent re-proofs.


\clearpage\part{Checked findings and cited background}

\section{The main question: which positivity can actually pass?}
\label{sec:positivity}

\subsection{Six distinct assertions}
Throughout, representations are finite dimensional over $\mathbb C$ unless
specified otherwise. Write
\[
 g(\alpha,\beta,\gamma)
 =\dim([\alpha]\otimes[\beta]\otimes[\gamma])^{\mathfrak S_d},
 \qquad
 p_{\lambda,\mu}^{\nu}
 =\dim\operatorname{Hom}_{GL(V)}
     (S_\nu V,S_\lambda(S_\mu V)).
\]
Partitions in the first formula have size $d$; in the second,
$|\nu|=|\lambda||\mu|$ and $V$ is large enough. Neither nonnegativity is
conjectural: these numbers are dimensions.

\begin{center}\small
\begin{tabularx}{\linewidth}{@{}p{.18\linewidth}X@{}}
\toprule
Assertion & Exact additional content\\
\midrule
Numerical positivity & A specified real number, determinant or norm is nonnegative.\\
Schur positivity & Expansion in the irreducible-character basis has
nonnegative integer coefficients.\\
Positive counting rule & The coefficient counts explicit combinatorial
objects; a complexity-theoretic $\#P$ rule also bounds witness length and
verification time in the stated input encoding.\\
Homogeneous PH1 & A stretching function is exactly
$\#(sP\cap\mathbb Z^D)$ for a rational polytope $P$, with the size and
separation-oracle bounds in GCT VI, Hypothesis 1.6.4.\\
Shifted counting model & The count is $\#\{z\in\mathbb Z^D:Az\le sb+c\}$.
The fixed shift $c$ changes the mathematical assertion.\\
Saturation/PH2 & GCT VI, Definitions 1.2.2 and 1.2.4 and
Hypotheses 1.6.5--1.6.6 concern constituent quasi-polynomials and bounds
on their saturation or coefficient-positivity indices. They are stronger
than values being nonnegative.\\
\bottomrule
\end{tabularx}
\end{center}

The exact stretching convention in GCT VI, Section 1.6, is
$p_{s\lambda,\mu}^{s\nu}$: the inner weight $\mu$ is fixed.
In its revised version the small-index conjectures replace a disproved
zero-index conjecture. A counterexample to naive dilation saturation must
not be presented as a counterexample to every revised hypothesis.

\subsection{Known refutations that a new purity proof cannot reverse}
\paragraph{P-K1: the literal homogeneous PH1 is false.}
Kahle--Micha\l ek, \emph{Obstructions to combinatorial formulas for
plethysm}, Theorem 5 and Remark 6, give
\[
 \phi(s)=p_{(4s),(3)}^{(7s,5s)}
 =\frac{s+r(s)}3,\qquad
 (r(0),\ldots,r(5))=(3,-1,1,0,2,-2),
 \quad r(s+6)=r(s).
 \tag{P-K1}\label{eq:km}
\]
Their Question 3 explicitly identifies this as the stretching convention
of GCT VI, Hypothesis 1.6.4. The function is not an Ehrhart function:
$\phi(1)=0$ but $-\phi(-1)=1$, contrary to Ehrhart reciprocity.
The same paper, Theorem 5, excludes one-dimensional shifted models too.
These are cited results, not new results of this investigation.

\paragraph{P-K2: the strong, zero-index Kronecker conjectures are false.}
Briand--Orellana--Rosas, \emph{Reduced Kronecker coefficients and
counter-examples to Mulmuley's strong saturation conjecture SH},
Theorem 2.4, give, in particular,
\[
 g\bigl((6s,6s),(7s,5s),(6s,4s,2s)\bigr)
 =
 \begin{cases}s/2+1&s\text{ even},\\(s-1)/2&s\text{ odd}.\end{cases}
 \tag{P-K2}
\]
The odd constituent is nonzero but vanishes at $s=1$, and has a negative
constant coefficient. This refutes the old strong SH and PH2, not the
revised polynomial index bounds. The distinction is explicit in that
paper's Section 2.5 and Mulmuley's appendix.

\paragraph{P-K3: unrestricted exact positivity is computationally hard.}
Ikenmeyer--Mulmuley--Walter, \emph{On vanishing of Kronecker coefficients},
Theorem 1.3, prove NP-hardness of Kronecker positivity, including unary
partition encoding. Fischer--Ikenmeyer, \emph{The computational complexity
of plethysm coefficients}, Theorem 3.1, prove NP-hardness for plethysm
positivity, already for fixed inner symmetric degree $3$.
Consequently, a proposed polynomial-time algorithm for all of either
family would imply $\mathsf P=\mathsf{NP}$. This is an implication;
it does not assume $\mathsf P\ne\mathsf{NP}$ and is compatible with the
supplied subexponential SAT theorem.

\subsection{A sharp repair on the counterexample ray}

\begin{proposition}[P-F1: optimal shifted model on the PH1 counterexample ray]
\label{prop:shiftedray}
For every integer $s\ge0$, the function in \eqref{eq:km} counts the integer
points of the single rational polytope
\[
 P_s=\{(t,z)\in\mathbb R^2:
 0\le z\le1,\quad 3t+4z\ge s,\quad 2t+4z\le s\}.
 \tag{P-F1}
\]
It has dimension two for $s>4$, even though $\phi$ has degree one.
Two is the smallest possible eventual dimension of an inhomogeneous
rational-polytope model for this function.
\end{proposition}
\begin{proof}
Inspection of the six residue classes gives
$\phi(s)=\lfloor s/3\rfloor+\mathbf1_{2\mid s}$, including $s=0$.
As formal power series,
\[
 \sum_{s\ge0}\phi(s)T^s
 =\frac{T^3}{(1-T)(1-T^3)}+\frac1{1-T^2}
 =\frac{1+T^4}{(1-T^2)(1-T^3)}.
\]
Expanding the last expression shows that $\phi(s)$ is the number of
integer triples
\[
 a,b\ge0,\quad z\in\{0,1\},\quad 2a+3b+4z=s.
\]
The integer coordinate change
$(a,b,z)\mapsto(k,t,z)=(2a+3b+4z,a+b,z)$ has determinant $-1$.
On $k=s$ its inverse is
$a=3t+4z-s$, $b=s-2t-4z$.
It is therefore an exact bijection of the relevant integer lattices,
and nonnegativity becomes the two displayed inequalities for $P_s$.
For $s>4$, at each $0<z<1$ the interval for $t$ has length
$(s-4z)/6>0$. The area is $(s-2)/6$, so the dimension is two.
A zero-dimensional convex polytope has at most one integer point,
whereas $\phi$ is unbounded. For completeness, the exclusion of
one-dimensional models also holds when an affine line is embedded in
a larger ambient lattice, as follows. Suppose
$Q_s=\{x\in\mathbb R^N:Ax\le sb+c\}$ is a bounded, eventually
one-dimensional model. Candidate vertices are the finitely many affine
functions of $s$ obtained by solving independent sets of $N$ active
constraints. Their feasibility and the equality of each constraint at
each vertex stabilize for sufficiently large $s$, since all these tests
are signs of affine functions. The constraints tight at every vertex
therefore give a fixed affine-hull system $Mx=sb'+c'$, whose kernel is
a rational line.

The count $\phi(s)$ is positive for every integer $s\ge2$. Choose
integer support points $x_{s_0},x_{s_0+1}$ at two consecutive sufficiently
large parameters, put $w=x_{s_0+1}-x_{s_0}$ and $v_0=x_{s_0}-s_0w$.
Then $v_0+sw$ lies in the affine support for every integer $s$.
If $v$ is a primitive integer generator of $\ker M$, all integer
points in that support are uniquely $v_0+sw+tv$, $t\in\mathbb Z$.
Substitution turns the inequalities into a bounded interval for $t$.
On a tail its endpoints are affine, say $as+\alpha$ and $bs+\beta$,
because maxima and minima of finitely many affine functions stabilize.
Its integer count differs from its length by a bounded quantity,
so $b-a=1/3$. Translate $t$ by an integer multiple of $s$ to arrange
$0\le a<1$.
If $b\le1$, at most one integer can enter at the right endpoint in a
unit step, and the nondecreasing left endpoint cannot add any.
But $\phi(s+1)-\phi(s)=2$ at arbitrarily large $s\equiv5\pmod6$.
Hence $b>1$. Then at least one integer enters at the right in every
step and at most one leaves at the left, so the count is eventually
nondecreasing. This contradicts its drop by one at every large
$s\equiv0\pmod6$. Thus no eventual dimension-one model exists.
This extends the one-variable exclusion in the cited theorem to an
arbitrary fixed ambient lattice and completes the minimality proof.
\end{proof}

This repairs this ray, not homogeneous PH1. It also gives a positive
formula for the equal rectangular Kronecker ray
$g((4s,4s,4s),(4s,4s,4s),(7s,5s))$; the equality of the two coefficient
families is recorded by Kahle--Micha\l ek on page 6.
No assertion of historical priority is made for this elementary
reconstruction. Positive formulas in the two-row range have independent
literature, including Pak--Panova--Swanson (2026 version).

\begin{proposition}[P-F2: from an explicit positive Hilbert series to one shifted polytope]
\label{prop:selector}
Let
\[
 F(T)=\frac{\sum_{i\in I}h_iT^i}{\prod_{j=1}^{r}(1-T^{a_j})},
 \qquad I\subset\mathbb Z_{\ge0}\text{ finite},\quad
 h_i,a_j\in\mathbb Z_{>0}.
\]
Every coefficient $[T^s]F(T)$, for $s\ge0$, is the integer-point count
of one bounded rational polytope whose inequalities have right sides
affine in $s$. Its description length is polynomial in the explicit
sparse numerator, the expanded list of denominator factors, and the
binary encoding of $s$.
\end{proposition}
\begin{proof}
Use real variables $x_1,\ldots,x_r,u$ and $e_i$ for $i\in I$ and impose
\[
 x_j,u,e_i\ge0,\qquad
 \sum_{i\in I}e_i=1,\qquad
 \sum_{j=1}^{r}a_jx_j+\sum_{i\in I}i e_i=s,\qquad
 u+1\le\sum_{i\in I}h_i e_i.
 \tag{P-F2}
\]
These are linear equalities and inequalities, hence define a convex
rational polytope (equalities can be replaced by two inequalities).
At an integer point, the nonnegative $e_i$ summing to one select exactly
one $i$. For that choice there are exactly $h_i$ possibilities
$u=0,\ldots,h_i-1$. The remaining equation counts nonnegative integer
solutions of $\sum_j a_jx_j=s-i$, which is the coefficient of
$T^{s-i}$ in the denominator's reciprocal. Summing the disjoint integer
selector slices gives exactly $[T^s]F$.
Boundedness holds even for real points: $0\le e_i\le1$,
$0\le u\le\max_i h_i-1$, and $0\le x_j\le s/a_j$ whenever feasible.
There are $r+|I|+1$ variables and linearly many constraints; their
coefficients are the explicitly supplied integers $a_j,i,h_i$, together
with $0,1,-1$. This proves the description bound. If the numerator is
zero, use the fixed inconsistent inequality $0\le-1$.
\end{proof}

\paragraph{What P-F2 settles.}
GCT VI, Theorem 1.6.1(d), supplies positive rational Hilbert forms for
its canonical models. P-F2 supplies an ordinary, possibly
large-dimensional shifted counting model from any such form.
The gate for an \emph{efficient uniform} model is now concrete:
construct sufficiently short positive numerator and denominator data
uniformly, or replace their explicit enumeration by an effective
description with the required separation bounds. A mere existence
theorem for an h.s.o.p.\ or a nonnegative Hilbert numerator does not give
those bounds. Neither P-F2 nor a polynomial description proves
saturation: integer feasibility in a rational polytope remains a
separate issue. This is a sufficient route to shifted counting, not a
claim that every possible route must use Hilbert series.

\subsection{Positive determinants already exist for the coefficients}

\begin{proposition}[P-F3: an unconditional determinant certificate, with its cost exposed]
\label{prop:projector}
Let $K$ be a compact group and $W$ a finite-dimensional continuous
unitary representation. With normalized Haar measure put
$P=\int_K\rho(k)\,dk$ and $m=\dim W^K$. Then
\[
 P=P^*=P^2,\qquad
 \det(I+tP)=(1+t)^m
 =\sum_{j=0}^{m}\binom mj t^j.
\]
For a finite group, the integral is the finite average.
Consequently every Kronecker coefficient and every plethysm coefficient
has a positive-semidefinite determinant representation of this form.
\end{proposition}
\begin{proof}
For every $h\in K$, invariance of Haar measure gives $\rho(h)P=P$,
so the image of $P$ lies in $W^K$. If $w\in W^K$, then $Pw=w$.
Thus $P$ is a projection onto $W^K$.
Unitarity and invariance under inversion give
$P^*=\int_K\rho(k^{-1})\,dk=P$. It has $m$ eigenvalues one
and all remaining eigenvalues zero, which proves the determinant formula.
For Kronecker take $K=\mathfrak S_d$ and
$W=[\alpha]\otimes[\beta]\otimes[\gamma]$.
For plethysm take $K=U(V)$ and
$W=(S_\nu V)^*\otimes S_\lambda(S_\mu V)$.
An algebraic map commuting with $U(V)$ commutes with $GL(V)$, because
$U(V)$ is Zariski dense in $GL(V)$, or equivalently because its complexified
Lie algebra is $\mathfrak{gl}(V)$ and $GL(V)$ is connected.
Hence $W^K$ is precisely the required Hom space.
\end{proof}

The matrix in P-F3 acts on the \emph{whole} tensor/Hom space. Its
dimension and the direct averaging cost are not polynomially bounded
in a compressed coefficient input. The determinant formula contains the
unknown rank; it is not a noncircular efficient algorithm for that rank.
It shows exactly why an additional positive determinant, without a
size bound and a coefficient-identification theorem, does not supply
PH1 or a multiplicity comparison.

\begin{proposition}[P-F4: scalar positivity does not imply character positivity]
\label{prop:scalar}
A real nonnegative class function, even an integer virtual character,
need not have nonnegative irreducible multiplicities.
\end{proposition}
\begin{proof}
On $C_2=\{1,\tau\}$ let $\mathbf1$ and $\varepsilon$ be the trivial
and sign characters. The virtual character
$f=\mathbf1-\varepsilon$ has values $f(1)=0$, $f(\tau)=2$,
both nonnegative, but its sign-character coefficient is $-1$.
Moreover the kernel $(f(x^{-1}y))_{x,y\in C_2}$ is
$\left(\begin{smallmatrix}0&2\\2&0\end{smallmatrix}\right)$,
with eigenvalues $2,-2$. Pointwise positivity is therefore different
even from positive type. The scalar inequality $W(1_S)\ge0$ cannot
be substituted for positivity of all representation projectors without
an additional identity.
\end{proof}

\begin{proposition}[P-F5: ordinary interlacing is not plethysm restriction]
\label{prop:interlacing}
The homomorphism $GL_2\to GL_3$ defined by $\operatorname{Sym}^2$
is not conjugate to the standard block embedding
$g\mapsto\operatorname{diag}(g,1)$. Their branching multiplicities
already disagree on the defining $GL_3$ representation.
\end{proposition}
\begin{proof}
The pullback of $\mathbb C^3$ by the first map is
$\operatorname{Sym}^2\mathbb C^2$, an irreducible three-dimensional
module. The pullback by the second is $\mathbb C^2\oplus\mathbb C$,
which is reducible. More explicitly, $zI_2$ acts in the first as $z^2I_3$,
and in the second with eigenvalues $z,z,1$; for example $z=2$ gives
different spectra. Conjugate homomorphisms would give isomorphic pullbacks,
so they cannot be conjugate.
\end{proof}

Thus GGP interlacing and a chain of standard unitary restrictions do not
by themselves compute restriction along the representation homomorphism
$GL(V)\to GL(S_\mu V)$. A theta-lift route would have to identify these
specific multiplicity spaces, prove compatibility of its compositions
and specializations, and control sizes. An order of vanishing gained
by an analytic lift is not, by definition, the dimension of one of
these Hom spaces. No such identifying theorem is stated in the supplied
four-page note or in the five supplied Lean declarations.

\subsection{Purity: the exact categorical gate and a complexity barrier}

\paragraph{P-K4: the established geometric template.}
Deligne's \emph{Weil II}, Theorem 3.3.1 and Corollaries 3.3.4--3.3.6,
control weights under direct image; Corollary 3.3.9 gives the expected
pure weights for smooth proper cohomology. BBD, \emph{Faisceaux pervers}, Theorems 5.3.8 and 5.4.5,
give geometric semisimplicity and decomposition for pure complexes,
and Theorem 6.2.5 gives the characteristic-zero decomposition theorem.
For Schubert varieties of a split connected reductive group over
$\mathbb F_q$, Haines, \emph{A proof of the Kazhdan--Lusztig purity
theorem via the decomposition theorem of BBD}, Theorem 2.0.1,
gives odd-stalk vanishing and the eigenvalues $q^{i/2}$ in even degree
$i$ for the unshifted intersection complex. His Theorem 5.0.4 gives
$P_{y,w}(q)=\sum_i\dim H^{2i}_{yB}(IC(X_w))q^i$, identifying the
Kazhdan--Lusztig polynomials with local intersection-cohomology
dimensions. These statements have hypotheses and identifications;
purity alone is not that identification.

\begin{proposition}[P-F6: a sufficient categorical transfer criterion]
\label{prop:categorical}
Let $\mathcal C$ be an additive semisimple graded category whose
indecomposables are $T^jL_\lambda$, with unique finite decompositions.
Suppose a bilinear product $\star$ preserves $\mathcal C$ and satisfies
$(T^iA)\star(T^jB)\simeq T^{i+j}(A\star B)$.
Suppose its split Grothendieck ring is identified with a based
$\mathbb Z[v,v^{-1}]$-algebra $A$, with $[L_\lambda]\mapsto b_\lambda$
and $[TM]\mapsto v[M]$. Then the structure constants of the
$b_\lambda$ lie in $\mathbb Z_{\ge0}[v,v^{-1}]$.
If specialization at $v=1$ identifies these constants with specified
representation multiplicities, those specializations are nonnegative
sums of actual multiplicities. This does not assert an efficient
combinatorial rule.
\end{proposition}
\begin{proof}
For each pair $\alpha,\beta$, the assumptions give an actual
decomposition
$L_\alpha\star L_\beta\simeq
\bigoplus_{\lambda,j}(T^jL_\lambda)^{\oplus m_{\alpha\beta}^{\lambda,j}}$
with $m_{\alpha\beta}^{\lambda,j}\in\mathbb Z_{\ge0}$.
Pass to the split Grothendieck group and apply the specified ring
identification. The coefficient of $b_\lambda$ is
$\sum_jm_{\alpha\beta}^{\lambda,j}v^j$.
Specialization at $v=1$ sums the nonnegative multiplicities, once the
asserted comparison with representations has independently been proved.
\end{proof}

The use of the \emph{split} Grothendieck group matters. In a triangulated
Grothendieck group a cohomological shift introduces a minus sign.
The rank-one constant Frobenius module $\mathbb Q_\ell$, with Frobenius
one, placed in cohomological degree one has alternating trace $-1$,
despite having a pure eigenvalue. Parity, basis identification, product
compatibility and geometric semisimplicity must all be supplied in the
intended realization. P-F6 does not produce an efficient description of
its multiplicity spaces.

\begin{proposition}[P-F7: even pure Tate geometry can encode all \#3SAT counts]
\label{prop:tatesat}
Fix a finite field $\mathbb F_q$ and a prime $\ell\ne\operatorname{char}\mathbb F_q$.
From any 3-CNF formula $F$ with $n$
variables and $m$ clauses one can construct, in time polynomial in
$n+m$, equations for a finite smooth proper $\mathbb F_q$-scheme $X_F$
such that
\[
 \#X_F(\mathbb F_{q^r})=\#\operatorname{SAT}(F)\quad(r\ge1),\qquad
 H^0_{\mathrm{\acute et}}(X_{F,\overline{\mathbb F}_q},\mathbb Q_\ell)
 \simeq \mathbb Q_\ell^{\,\#\operatorname{SAT}(F)}
\]
with Frobenius the identity, and all other cohomology zero.
On this explicitly presented family, nonemptiness is NP-complete and
the count is \#P-complete even under this promised purity, parity and
smoothness.
\end{proposition}
\begin{proof}
Impose $x_i^2-x_i=0$ for $1\le i\le n$.
For a positive literal $x_i$ let its false-value factor be $1-x_i$;
for a negative literal $\neg x_i$ let it be $x_i$.
For each clause impose the product of its at most three false-value
factors equal to zero.
The algebra
$\mathbb F_q[x_1,\ldots,x_n]/(x_i^2-x_i:1\le i\le n)$
is, by repeated Chinese remainder, the product of one copy of
$\mathbb F_q$ for every Boolean assignment. This remains true in
characteristic two, since $0$ and $1$ are distinct and
$x_i(x_i-1)$ has distinct roots.
In a factor corresponding to an assignment, each clause product is
zero exactly when the clause is satisfied, and otherwise equals one.
Quotienting by all clause products therefore deletes precisely the
unsatisfying factors. The resulting coordinate algebra is
$\prod_{\sigma\models F}\mathbb F_q$ (the zero algebra for no assignments).
Consequently $X_F$ is a disjoint union of rational points, is finite
\'etale, and hence smooth and proper. The displayed point counts and
cohomology follow point by point.
There are $n+m$ equations of degree at most three with coefficients of
bounded size. The construction is parsimonious, so the
NP-completeness and \#P-completeness of 3SAT transfer (a constructive
verifier proof is given in C-COUNT-03 below). Membership follows
by guessing and checking an assignment.
\end{proof}

This is a reduction, not a proof that $\mathsf P\ne\mathsf{NP}$.
It establishes an NP-hardness boundary for any proposed recipe using
a polynomial list of equations and pure even Tate cohomology:
a polynomial algorithm from those data would solve 3SAT.
The supplied ETH theorem improves the subexponential computation of
these examples, without changing the exact polynomial-time implication.

\subsection{Nonstandard quantum groups: what remains to be constructed}
GCT VII, Theorem 1.2, constructs the nonstandard quantization.
Its Conjectures 1.3 and 1.4 are duality and reciprocity assertions.
GCT VIII's Conjectures 2.6--2.9 and 3.4--3.7 impose cell, positivity,
saturation and compatibility properties. The general setting even
permits signed factors and a relaxed positivity condition; it is not
the assertion that every raw coefficient is nonnegative.
The intended geometric realization and specialization remain additional
input. P-F6 supplies a sufficient transfer criterion once those data
exist; P-F3--P-F7 identify why determinant signs, interlacing or purity
alone do not supply them. No theorem of nonexistence of such a future
realization is claimed.

\subsection{Occurrence, multiplicity and actual equations}

\paragraph{P-K5: the no-occurrence theorem has a precise range.}
B\"urgisser--Ikenmeyer--Panova, \emph{No occurrence obstructions in
geometric complexity theory}, Theorem 1.4, prove that for
$n\ge m^{25}$ every irreducible occurring in the coordinate ring of
the padded permanent orbit closure also occurs in the determinant
orbit closure. This is a theorem about these closures, not merely an
upper bound using rectangular Kronecker coefficients. It does not
exclude multiplicity obstructions or all other uses of representation
theory in complexity.

\begin{proposition}[P-F8: the surviving multiplicity test and its logical limit]
\label{prop:multiplicity}
Let $X,Y$ be $G$-stable affine cones in one representation of a complex
reductive group, and let $m_X(d,\lambda)$, $m_Y(d,\lambda)$ be their
graded coordinate-ring multiplicities. If
$m_Y(d,\lambda)>m_X(d,\lambda)$ for some $(d,\lambda)$ then
$Y\not\subseteq X$. The reverse inference from all weak inequalities
to containment is false, even when all the multiplicities agree.
In the no-occurrence range above, any multiplicity obstruction for the
specified determinant and permanent closures has both multiplicities
strictly positive.
\end{proposition}
\begin{proof}
If $Y\subseteq X$, restriction is a surjective graded $G$-map
$\mathbb C[X]\twoheadrightarrow\mathbb C[Y]$.
Finite-dimensional rational $G$-modules in characteristic zero are
semisimple, so taking the $\lambda$-multiplicity in degree $d$ gives
$m_Y(d,\lambda)\le m_X(d,\lambda)$. This proves the first assertion.
For the failure of the converse, let $G=\mathbb G_m$ act on
$\mathbb A^2$ by scalar multiplication, let $X$ be the first coordinate
axis and $Y$ the second. Both are stable irreducible cones and both
coordinate rings have exactly one copy of the character of weight
$-d$ in each degree $d$. Nevertheless neither axis contains the other.
Finally a strict inequality implies $m_Y>0$; P-K5 then gives $m_X>0$.
\end{proof}

The theorem sought for determinant versus permanent is consequently
a strict comparison of the \emph{actual quotient} multiplicities, or
an explicit equation separating the closures. Ambient plethysm
multiplicities and determinant-orbit Kronecker bounds do not determine
those quotient multiplicities. The Chow/Foulkes--Howe map in the
plethysm section makes the missing rank information explicit.



\subsection{Quantitative bounds from positive Hilbert data}

\begin{proposition}[P-F9: a common-period bound for the positivity index]
\label{prop:periodbound}
Let $f(s)$, $s\ge0$, be a quasi-polynomial whose generating function has
an explicit positive form
\[
 \sum_{s\ge0}f(s)T^s=\frac{H(T)}{\prod_{j=1}^r(1-T^{a_j})},
 \quad H(T)\in\mathbb Z_{\ge0}[T]\setminus\{0\},\quad a_j>0.
\]
Put $h=\deg H$, let $L$ be any positive common multiple of the $a_j$, and suppose
$r\ge2$. In the definitions of GCT VI, Definitions 1.2.2 and 1.2.4,
\[
 s(f)\le p(f)\le
 \max\{0,h+(r-1)L-\textstyle\sum_j a_j\}
 \le (r-1)L-1.
 \tag{P-F9}
\]
For $r=1$, $f$ is periodic with nonnegative constant constituents and
$p(f)=s(f)=0$. Thus a positive form with polynomially bounded $r$ and
common multiple $L$ supplies polynomial positivity and saturation
indices. A bound on the largest $a_j$ alone does not make this particular
estimate polynomial, because their least common multiple can be large.
\end{proposition}
\begin{proof}
First $h<\sum_j a_j$. To see this without an assumption on cancellations,
write the quasi-polynomial on each residue class modulo a period as a
polynomial in its nonnegative quotient variable. Express that polynomial
in the binomial basis $\binom{k}{j}$. Its generating function is a linear
combination of $T^{c+jD}/(1-T^D)^{j+1}$ with $0\le c<D$, all proper
rational functions. Their sum tends to zero at infinity. The displayed
positive rational expression for the same function therefore also has
numerator degree strictly below its denominator degree.

Multiply numerator and denominator by geometric sums to obtain
\[
 \frac{B(T)}{(1-T^L)^r},\qquad
 B(T)=H(T)\prod_{j=1}^r
       (1+T^{a_j}+\cdots+T^{L-a_j}).
\]
All coefficients $b_k$ of $B$ are nonnegative and
$D_B:=\deg B=h+rL-\sum_j a_j\le rL-1$.
For a fixed residue $c$ modulo $L$, the eventual constituent of $f$ is
\[
 F_c(s)=\sum_{k\equiv c\ (L)} b_k
 \binom{(s-k)/L+r-1}{r-1}
 =\frac1{L^{r-1}(r-1)!}
 \sum_{k\equiv c\ (L)} b_k
 \prod_{j=1}^{r-1}(s-k+jL).
\]
Indeed this follows by expanding $(1-T^L)^{-r}$ for all sufficiently
large $s$. The constituents coincide with those of the given
quasi-polynomial everywhere: two polynomials agreeing on an infinite
arithmetic progression agree identically (use a common period if
necessary). If $p\ge\max\{0,D_B-L\}$, each factor of $F_c(s+p)$ is
$s+(p-k+jL)$ with nonnegative constant term. Every constituent thus has
nonnegative coefficients. Translation merely permutes the residue
labels, so this proves the bound on $p(f)$. Any nonzero polynomial with
nonnegative coefficients is positive at every real $s\ge1$, which gives
$s(f)\le p(f)$. For $r=1$ the constituents are nonnegative constants.
Finally $D_B\le rL-1$ gives the last bound when $r\ge2$.
\end{proof}

This is an explicit sufficient condition for the revised PH2, not a proof
of that hypothesis for the GCT families. Taking $L$ to be the least common
multiple always gives a bound whose bit length is polynomial in the
expanded positive-form data. GCT VI, Lemma 3.1.6, supplies a different
polynomial modular-index bound for saturation in its stated setting;
we do not replace that known result by a new claim. The remaining
quantitative issue is to control the appropriate parameters for the
specific multiplicity rings, uniformly in the GCT input.

\begin{proposition}[P-F10: purity and normality do not control a grading's indices]
\label{prop:weightedplane}
Let $a,b>1$ be relatively prime. Give the polynomial ring $k[x,y]$ the
positive grading $\deg x=a$, $\deg y=b$, and put
$f(n)=\dim_k k[x,y]_n$. Then $f$ is a quasi-polynomial and its GCT
positivity and saturation indices are exactly
\[
 p(f)=s(f)=ab-a-b.
 \tag{P-F10}
\]
The underlying variety is always the same smooth normal affine plane.
Over a finite field its only compactly supported cohomology is
$H_c^4=\mathbb Q_\ell(-2)$; its intersection complex is the shifted
constant sheaf. Thus even the same pure geometry can support arbitrarily
large indices when its grading varies.
\end{proposition}
\begin{proof}
The monomial basis gives
$f(n)=\#\{(u,v)\in\mathbb Z_{\ge0}^2:au+bv=n\}$ and generating function
$1/((1-T^a)(1-T^b))$. Write
\[
 \frac1{(1-T^a)(1-T^b)}
 =\frac{\displaystyle\sum_{i=0}^{b-1}\sum_{j=0}^{a-1}T^{ai+bj}}
 {(1-T^{ab})^2}.
\]
The $ab$ exponents $r=ai+bj$ represent distinct residues modulo $ab$:
reduction modulo $b$ fixes $i$, and modulo $a$ fixes $j$.
Their minimum is zero and their maximum is
$r_{\max}=a(b-1)+b(a-1)=2ab-a-b<2ab$.
For the residue represented by $r$, its constituent is therefore
\[
 F_r(n)=\frac{n-r}{ab}+1.
\]
For $n\ge r$ this is the geometric-series coefficient.
If $0\le n<r$ in the same residue, necessarily $n=r-ab$;
the displayed formula is zero, as the coefficient is. Hence the
constituent formula is valid for every nonnegative integer $n$.

All constituents have positive slope $1/(ab)$. After a shift by $p$,
the smallest constant coefficient is
$1+(p-r_{\max})/(ab)=(p-(ab-a-b))/(ab)$.
It is nonnegative precisely when $p\ge ab-a-b$, proving the positivity
index. Every constituent is nonzero. Strict saturation after a shift
$s$ requires its value at $n=1$ to be strictly positive for the worst
constituent; the condition is $1+s>ab-a-b$, or, for integer $s$,
$s\ge ab-a-b$. This proves the second index.

The ring is a polynomial ring independently of its grading, hence its
spectrum is smooth and normal. The cohomology statement is the standard
affine-space calculation and duality (Milne, \emph{Lectures on \'Etale
Cohomology}, Example 16.3 and Theorem 24.1). Smoothness identifies its
perverse intersection complex with $\mathbb Q_\ell[2]$.
Taking $a=2$ and odd $b$ arbitrarily large makes the indices $b-2$
unbounded while this underlying geometry stays fixed.
\end{proof}

P-F10 by itself makes no representation-theoretic identification.
P-F15 below supplies that identification for the broad subgroup
restriction problem, where it refutes the printed rank-only index bounds.
It does not refute the narrower irreducible-inner plethysm bounds.
Independently, it refutes a proposed derivation of quantitative bounds
from purity, smoothness or normality alone. The action and its
grading, and a bound relating their weights to the input, are essential
additional data. Normality of the original ring also must not be
confused with saturation of the image of its exponent semigroup under
the degree map.


\subsection{A polynomial modular-index bound for coefficient positivity}

\begin{proposition}[P-F11: positive rational form implies a polynomial positivity-index bound]
\label{prop:modularbound}
Let $f(n)$, $n\ge0$, be a quasi-polynomial with nonzero generating function
\[
 F(T)=\frac{H(T)}{\prod_{i=1}^r(1-T^{a_i})},\qquad
 H\in\mathbb Z_{\ge0}[T]\setminus\{0\},\quad a_i>0.
\]
Here $r\ge1$. Set $A=\max_i a_i$ and $h=\deg H$. Then
\[
 p(f)\le h+128r^3A^3 < rA+128r^3A^3.
 \tag{P-F11}
\]
The bound concerns the positivity index of every constituent polynomial,
not just eventual positivity of the values in its own congruence class.
In particular, polynomial bounds for $r$ and the modular index $A$
suffice for a polynomial bound for the positivity index; a polynomial
bound on the common period is unnecessary.
\end{proposition}
\begin{proof}
The inequality $h<\sum_i a_i\le rA$ was proved in P-F9 from global
quasi-polynomiality. It remains to prove the first bound. Begin with
\[
 G(T)=\prod_{i=1}^r(1-T^{a_i})^{-1},\qquad g(n)=[T^n]G(T),
\]
assuming $\gcd(a_1,\ldots,a_r)=1$ and $r\ge2$. Put
$D=16rA^2$ and $\rho=1/D$. Every pole $\zeta$ of $G$ is a root of unity
of order at most $A$. Its order as a pole is
$m=m(\zeta)=\#\{i:\zeta^{a_i}=1\}$.
There are at most $N:=\sum_i a_i\le rA$ distinct poles. The only pole
of order $r$ is $\zeta=1$, by the greatest-common-divisor assumption.
Write the partial fractions as
\[
 G(T)=\sum_{\zeta}\sum_{j=1}^{m(\zeta)}
 C_{\zeta,j}(1-T/\zeta)^{-j}.
\]
There is no polynomial part because $G$ is proper, and
$C_{1,r}=c:=1/\prod_i a_i>0$.

We first establish, for every displayed term,
\[
 |C_{\zeta,j}|/c\le 2D^{r-j}.                 \tag{*}
\]
Use the local variable $z=1-T/\zeta$ and put
$u_a(z)=(1-(1-z)^a)/(az)=a^{-1}\sum_{k=0}^{a-1}(1-z)^k$.
For $|z|\le\rho$,
\[
 |u_a(z)-1|\le\tfrac12(a-1)\rho e^{a\rho}\le a\rho,
 \qquad
 |(1-z)^a-1|\le a\rho e^{a\rho}\le2a\rho.
\]
These inequalities use $a\rho\le1/16$ and the finite geometric sum;
the last constants are deliberately loose. If $\zeta$ has order $d$
and $\zeta^{a_i}\ne1$, then
\[
 |1-\zeta^{a_i}|\ge2\sin(\pi/d)\ge4/d\ge4/A.
\]
Consequently
$|1-\zeta^{a_i}(1-z)^{a_i}|\ge4/A-2A\rho>3/A$.
For the $m$ factors vanishing at $z=0$, the product of the inverses of
$u_{a_i}$ is at most
\[
 \prod_{\zeta^{a_i}=1}(1-a_i\rho)^{-1}
 \le\exp(2\rho\textstyle\sum_i a_i)
 \le\exp(1/(8A))<2.
\]
Thus the analytic function $z^mG(\zeta(1-z))/c$ has absolute value at
most $2A^{2(r-m)}$ on $|z|\le\rho$: each of the other factors contributes
at most $a_i/(3/A)\le A^2$.
Cauchy's coefficient bound gives
\[
 |C_{\zeta,j}|/c
 \le2A^{2(r-m)}\rho^{-(m-j)}
 \le2D^{r-j},
\]
which is (*). The denominator estimates also show explicitly that this
analytic function has no further pole in that disk.

Choose a period divisible by the orders of all poles. On a fixed residue
class $b$, partial fractions give the constituent polynomial
\[
 g_b(n)=c\binom{n+r-1}{r-1}
 +\sum_{(\zeta,j)\ne(1,r)}
     C_{\zeta,j}\zeta^{-b}\binom{n+j-1}{j-1}.
\]
This identity holds for all nonnegative integers on the class, hence as
a polynomial identity for the constituent. Let $p>0$ be real and let
$[n^k]$ denote a polynomial coefficient. All coefficients of
$B_j(n):=\binom{n+p+j-1}{j-1}$ are nonnegative. For $j<r$ and
$0\le k\le j-1$,
\[
 \frac{[n^k]B_j(n)}{[n^k]B_r(n)}
 \le \frac{(r-1)!}{(j-1)!p^{r-j}}
 \le (r/p)^{r-j}.                              \tag{**}
\]
Indeed, multiply the $j-1$ factors in $B_j$ by the additional factors
$(n+p+j),\ldots,(n+p+r-1)$, whose constant product is at least
$p^{r-j}$, and compare coefficients. For $k>j-1$ the numerator is zero.

Using (*) and (**), the absolute value of the total error coefficient,
divided by the corresponding positive coefficient of $cB_r$, is at most
\[
 2N\sum_{\ell=1}^{r-1}(Dr/p)^\ell.
\]
At $p_0=128r^3A^3$ one has $Dr/p_0=1/(8rA)\le1/(8N)$.
The last display is bounded by
$2N(1/(8N))/(1-1/(8N))\le2/7<1$.
Each constituent has real coefficients, since it agrees with integer
values at infinitely many real integers. The strict bound on the
absolute value of its complex error therefore shows that every
coefficient of every $g_b(n+p_0)$ is strictly positive.
The same holds for every real $p\ge p_0$, either by the same estimates
or by a further nonnegative translation.

Now let $d=\gcd(a_1,\ldots,a_r)$ be arbitrary and set $a_i'=a_i/d$.
The coefficients of $G$ vanish outside multiples of $d$, and on those
multiples equal $g'(n/d)$ for the relatively prime weights $a_i'$.
A nonzero constituent is correspondingly a constituent of $g'$ evaluated
at $n/d$. It has nonnegative coefficients after every shift
$p\ge d\cdot128r^3(A/d)^3\le128r^3A^3$; zero constituents remain zero.
If $r=1$, the coefficient sequence is simply the indicator of divisibility
by $a_1$, and already has nonnegative constant constituents.

Finally, $F(T)=\sum_i H_iT^iG(T)$ with $H_i\ge0$.
On sufficiently large indices its constituent polynomials are the sums
of $H_i$ times the appropriately relabeled constituents of $g(n-i)$.
The same polynomial identities hold everywhere because $f$ is globally
a quasi-polynomial. After shifting $n$ by $h+p_0$, each summand has
shift $h+p_0-i\ge p_0$, so every coefficient is nonnegative.
This proves the asserted positivity-index bound and completes the proof.
\end{proof}

P-F11 is an elementary analytic bound on rational generating functions;
its complex analysis is Cauchy's coefficient estimate on a fixed small
disk, and it uses neither RH nor a purity theorem. It strengthens the
common-period sufficient condition P-F9. A historical novelty assertion
would require a separate literature review; the statement is included
here as a fully proved deduction, not as a claim of priority.


\subsection{PH3 implies PH2 for the representation-theoretic stretching families}

\begin{proposition}[P-F12: the denominator length is controlled by Weyl growth]
\label{prop:ph3ph2}
Let $G,H$ be connected reductive complex algebraic groups, let
$\rho:H\to G$ be an algebraic homomorphism, and let $\lambda,\pi$ be
dominant integral weights. Put
\[
 f(n)=\dim\operatorname{Hom}_H
       (V_{n\pi}(H),V_{n\lambda}(G)|_H).
\]
Suppose this is a nonzero global quasi-polynomial on $n\ge0$ and has a
positive rational form
\[
 \sum_{n\ge0}f(n)T^n
 =\frac{N(T)}{\prod_{j=1}^{r}(1-T^{a_j})},
 \quad N\in\mathbb Z_{\ge0}[T]\setminus\{0\},\quad r\ge1,
 \quad a_j\ge1.
\]
Let $A=\max_j a_j$, and let
\[
 d_\lambda
 =\#\{\alpha\in\Phi_G^+:\langle\lambda,\alpha^\vee\rangle>0\}.
\]
Then
\[
 r=\deg f+1\le d_\lambda+1\le|\Phi_G^+|+1\le\dim G+1,
\]
and P-F11 gives
\[
 p(f)<(d_\lambda+1)A+128(d_\lambda+1)^3A^3.
 \tag{P-F12a}
\]
In particular, for the plethysm and generalized plethysm problem of
GCT VI, Problem 1.1.2, Hypothesis 1.6.7 (PH3) implies Hypothesis 1.6.6
(PH2), with their stated parameter
$h=\dim G+\operatorname{ht}(\lambda)+\operatorname{ht}(\pi)$.
The same implication holds for connected subgroup restriction with
classical ambient $G$, using $\operatorname{rank}(x)$ as in GCT VI,
Section 3.4, and Hypotheses 3.3.3--3.3.4. This includes the usual
Kronecker restriction realization of Problem 1.1.1.

For $G=GL_m$, $m\ge1$, and a partition $\lambda$, the same positive form has a
binary description of length polynomial in
$m,A,\log(\lambda_1+1)$, with $\lambda_1=0$ for the empty partition.
Consequently PH3 also implies the existence of a polynomial-size shifted
polytope description through P-F2 in this expanded group-size parameter.
This last assertion is an existence and encoding bound, not a uniform
algorithm for constructing the positive form.
\end{proposition}

\begin{proof}
Complete reducibility gives
\[
 0\le f(n)\le\dim V_{n\lambda}(G).
\]
The Weyl dimension formula is
\[
 \dim V_{n\lambda}(G)
 =\prod_{\alpha\in\Phi_G^+}
   \frac{n\langle\lambda,\alpha^\vee\rangle
             +\langle\varrho_G,\alpha^\vee\rangle}
        {\langle\varrho_G,\alpha^\vee\rangle},
 \tag{P-F12b}
\]
where $\varrho_G$ is half the sum of positive roots. See
Borcherds--Haiman--Reshetikhin--Serganova--Johnson-Freyd,
\emph{Lectures in Lie Groups and Quantum Groups}, Part I (2016),
Theorem 6.1.1.30, printed pp.~102--103.
For reductive $G$, apply that formula to its semisimple derived Lie
algebra; the central character does not change the dimension.
Every nonconstant factor in (P-F12b) has positive slope, and their
number is $d_\lambda$. Hence $f(n)=O((n+1)^{d_\lambda})$.
The implied constant may depend on $\lambda$; its size is irrelevant to
this degree comparison.

We next bound the number of denominator factors without assuming that
the positive form came from a particular h.s.o.p. Choose $i_0$ with
$N_{i_0}>0$ and put $L=\operatorname{lcm}(a_1,\ldots,a_r)$.
For every nonnegative tuple $(y_1,\ldots,y_r)$ with $\sum_jy_j=k$,
the tuple $x_j=(L/a_j)y_j$ is integral and satisfies
$\sum_ja_jx_j=Lk$. Expanding the positive rational form therefore gives
\[
 f(Lk+i_0)\ge N_{i_0}\binom{k+r-1}{r-1}.
\]
Comparison with the Weyl bound forces $r-1\le d_\lambda$.
The possibly large integer $L$ is used only to prove this inequality;
it is not part of a required computational construction.

For completeness, expanding the same form shows
$f(n)=O((n+1)^{r-1})$: after fixing the numerator term and the first
$r-1$ nonnegative variables, the last variable is determined, and each
of the first $r-1$ variables is at most $n$ whenever the term contributes.
Together with the preceding lower bound on an infinite progression,
this proves $\deg f=r-1$. Finally
$|\Phi_G^+|=(\dim G-\operatorname{rank}G)/2\le\dim G$.
Global quasi-polynomiality gives
$\deg N<\sum_ja_j\le rA$ as in P-F9. Apply P-F11 and substitute
$r\le d_\lambda+1$ to obtain (P-F12a).

In Problem 1.1.2, $G=GL(V_\mu(H))$ and the inner weight $\mu$ is held
fixed while $\lambda$ and $\pi$ are stretched. Thus the preceding
multiplicity and the parameter $h$ are exactly the ones in
Hypotheses 1.6.6--1.6.7. If $A$ is bounded by a fixed polynomial in $h$,
then (P-F12a) bounds $p(f)$ by a fixed polynomial in $h$.
For a product of classical ambient groups, the numbers of positive
roots in types $A_\ell,B_\ell,C_\ell,D_\ell$ are respectively
$\ell(\ell+1)/2,\ell^2,\ell^2,\ell(\ell-1)$, and the sum for a product
is at most $(\operatorname{rank}G)^2$. The general linear central torus
does not add roots. Since the definition of $\operatorname{rank}(x)$
includes $\operatorname{rank}G$, one has
$r\le(\operatorname{rank}(x))^2+1$ in these cases.
The same argument proves the implication of Hypothesis 3.3.4 to 3.3.3.
For Kronecker coefficients use the restriction
$GL_a\times GL_b\to GL(\mathbb C^a\otimes\mathbb C^b)$, with the
third partition as the highest weight of the ambient general linear
group; this is the realization in Problem 1.1.1.
If all positive-degree multiplicities vanish, their constituent
polynomials are zero and their positivity index is zero. Such a ray
is handled separately: the isolated value $f(0)=1$ of the literal
branching formula is not part of a global quasi-polynomial and is
outside the hypothesis above.

It remains to prove the encoding assertion. Let $B=\deg N$ and write
$N(T)=\sum_{i=0}^B N_iT^i$. The reciprocal denominator has nonnegative
coefficients and constant coefficient one, so
\[
 N_i\le f(i)\le\dim V_{i\lambda}(GL_m)
 \le (i\lambda_1+m)^{m(m-1)/2}.
\]
The last inequality follows term by term from the $GL_m$ specialization
of (P-F12b): each factor is
$(i(\lambda_j-\lambda_k)+k-j)/(k-j)$ for $j<k$, and its denominator
is at least one. Since $B<rA$ and $r\le m(m-1)/2+1$, there are fewer
than $rA+1$ numerator entries, each requiring at most
\[
 O\!\left(1+m^2\log\bigl(2+rA\lambda_1+m\bigr)\right)
\]
bits. The denominator list requires $O(r\log(A+1))$ bits.
This is polynomial in $m,A,\log(\lambda_1+1)$.
Applying P-F2 to this explicit data gives the asserted polynomial-size
shifted description whenever such a form is supplied; existentially,
PH3 supplies a form with these bounds. The proof does not find that
form or certify its equality to the multiplicity series in polynomial
time.
\end{proof}

\paragraph{Scope and exact remaining gate.}
GCT VI itself remarks in Section 4.2.2 (printed p.~73) that PH2 may
follow from PH3. P-F11--P-F12 establish this implication for the stated
representation-theoretic families with an explicit bound. They do not
prove PH3. They also do not replace the geometric coefficient and basis
identifications needed to infer PH3 from a proposed purity construction.

The group-size parameter must be preserved. In generalized plethysm,
$\dim G=(\dim V_\mu(H))^2$ can be very large compared with the binary
encoding of $\mu$. GCT VI, Section 3.4, explicitly permits this compact
input; Hypothesis 1.6.4 requires its ambient polytope dimension to be
polynomial in the bitlength of $\mu$ and the heights of $\lambda,\pi$.
Thus the polynomial bounds just proved in $h$ do not automatically give
a uniform shifted model with that more compressed complexity bound.

No estimate $\deg N\le\deg f$ has been used. The numerator index in
the displays of GCT VI, Theorem 1.6.1(d) and equation (4.4), cannot
literally impose such a bound. Indeed the KM example has degree one,
so that bound would require a positive numerator $1+cT$ and two
denominator factors. Its value $\phi(1)=0$ forces $c=0$ and both
denominator weights to exceed one. Then $\phi(2)=\phi(3)=1$ forces
the weights to be two and three, but the resulting coefficient of
$T^4$ is one rather than $\phi(4)=2$, a contradiction. This identifies
an index clash in the printed display; the correct positive-form
consequence of Cohen--Macaulayness has an unrestricted numerator degree.
The required bound here is the proved properness inequality
$\deg N<\sum_j a_j$.


\subsection{A polynomial finite-prefix certificate under a denominator bound}

\begin{proposition}[P-F13: universal denominator and finite reconstruction]
\label{prop:finiteprefix}
Let $R,A\ge1$ be given integers. Suppose a nonzero global
quasi-polynomial $f(n)$, $n\ge0$, has a positive rational form
\[
 F(T)=\sum_{n\ge0}f(n)T^n
 =\frac{N(T)}{Q(T)},\qquad
 Q(T)=\prod_{j=1}^{r}(1-T^{a_j}).
\]
\[
 1\le r\le R,\quad 1\le a_j\le A,
 \quad N\in\mathbb Z_{\ge0}[T]\setminus\{0\}.
\]
Define
\[
 D(T)=\prod_{a=1}^{A}(1-T^a)^R,
 \qquad M=\deg D=\frac{RA(A+1)}2.
\]
Then $F=P/D$ for a unique polynomial $P\in\mathbb Z[T]$ with
$\deg P<M$, and
\[
 P(T)=D(T)\left(\sum_{n=0}^{M-1}f(n)T^n\right)\bmod T^M.
 \tag{P-F13a}
\]
Consequently those first $M$ coefficients determine the entire
generating function. Under the stated denominator promise, a proposed
positive form with at most $R$ denominator factors of weights at most
$A$ and with numerator degree smaller than its denominator degree can
be certified by these $M$ coefficients.

Let $b$ bound the binary lengths of the supplied integers
$f(0),\ldots,f(M-1)$. Reconstruction and verification of any proposed
denominator take bit complexity polynomial in the numeric parameters
$R,A$ and $b$. A positive form can be found by a finite search through
at most $\binom{A+R}{R}-1$ denominators, with polynomial work per
denominator, once those coefficient values are supplied.
Equivalently, positive-form selection for the reconstructed rational
function has certificates verifiable in polynomial time in this data.
No polynomial-time algorithm for evaluating the original multiplicities
or selecting a successful denominator is assumed or proved.
\end{proposition}

\begin{proof}
Let $c_a$ be the number of occurrences of $a$ among the $a_j$.
Then $0\le c_a\le R$, and
\[
 D/Q=\prod_{a=1}^{A}(1-T^a)^{R-c_a}\in\mathbb Z[T].
\]
By the global quasi-polynomial properness argument in P-F9,
$\deg N<\deg Q$. Hence
\[
 P=N(D/Q)\in\mathbb Z[T],\qquad
 \deg P=\deg N+M-\deg Q<M.
\]
Since $D(0)=1$, multiplying the first $M$ coefficients of $F$ by $D$
and discarding terms of degree at least $M$ recovers every coefficient
of $P$, proving (P-F13a). Uniqueness follows immediately. More
explicitly, two strictly proper rational functions whose denominators
divide $D$ and whose first $M$ coefficients agree have a difference
$(P_1-P_2)/D$ with $\deg(P_1-P_2)<M$ and zero of order at least $M$
at zero. Multiplication by $D(0)\ne0$ forces $P_1-P_2=0$.

For an effective verifier, fix a candidate list
$1\le a_1\le\cdots\le a_r\le A$, $1\le r\le R$, and put
$Q'=\prod_j(1-T^{a_j})$ and $d'=\deg Q'\le RA\le M$.
Its only possible strictly proper numerator is
\[
 N'(T)=Q'(T)\left(\sum_{n=0}^{M-1}f(n)T^n\right)\bmod T^{d'}.
 \tag{P-F13b}
\]
Accept this list exactly if all coefficients of $N'$ are nonnegative,
$N'\ne0$, and the polynomial identity
\[
 D(T)N'(T)=Q'(T)P(T)
 \tag{P-F13c}
\]
holds. If it is accepted, (P-F13c) proves $N'/Q'=P/D=F$.
Conversely any positive form with the stated bounds has a strictly
proper numerator, is recovered by (P-F13b), and passes (P-F13c).
Thus the test is both necessary and sufficient under the promise.
The same argument proves the first-$M$-coefficients version of the
certificate.

Here is an explicit bit bound. Write $\|B\|_1$ for the sum of the
absolute values of the coefficients of a polynomial $B$. One has
$\|D\|_1\le2^{AR}$ and $\|Q'\|_1\le2^R$. If
$|f(n)|<2^b$ for $n<M$, formula (P-F13a) bounds the absolute value
of every coefficient of $P$ by $2^{b+AR}$, and (P-F13b) bounds that
of every coefficient of $N'$ by $2^{b+R}$. All degrees in the
construction and in (P-F13c) are at most $M+RA$. Ordinary integer
polynomial multiplication, truncation and equality tests therefore
have bit complexity polynomial in $R,A,b$. This counts $R,A$ by
their numeric values; it does not claim polynomial complexity in
$\log R,\log A$.

There are $\binom{A+r-1}{r}$ weakly increasing lists of length $r$,
so the sum for $1\le r\le R$ is $\binom{A+R}{R}-1$.
Enumerating these and using the verifier must find an accepted list
under the assumed existence of a positive form. The list, together
with the deterministically recovered numerator if desired, is a
polynomial-size witness for successful selection from the supplied
finite data. The number of lists is not asserted to be polynomial.
\end{proof}

\paragraph{Application and scope.}
Under a specified bound $A\le A_0$ in PH3, P-F12 supplies a
representation-theoretic bound $r\le R$. P-F13 reduces construction of
a positive form to the first $RA_0(A_0+1)/2$ exact stretching
multiplicities and finite denominator selection. For $G=GL_m$, the
Weyl estimate in P-F12 bounds the binary lengths of all these values
by a polynomial in $m,R,A_0,\log(\lambda_1+1)$. The number of
required coefficients and the work of verifying a candidate are
therefore polynomial in these expanded parameters. This avoids
verification across a possibly exponential least-common-multiple
period, and it does not provide an efficient coefficient oracle.

The denominator promise is essential. A finite prefix, by itself,
does not prove PH3 for an otherwise unbounded original series.
For example, the two nonnegative global quasi-polynomials
$1$ and $1+\mathbf1_{\{n\equiv M\pmod{M+1}\}}$ agree for
$n=0,\ldots,M-1$ and differ at $n=M$. Thus (P-F13c) certifies the
original multiplicity series only when its membership in the stated
denominator class is known, as it is conditionally under PH3.
Historical priority for this elementary rational reconstruction is
not asserted.


\subsection{Fixed-rank torus embeddings refute the unrestricted index hypotheses}

The next two findings concern the full subgroup restriction problem,
GCT VI, Problem 1.1.3, with the specifications of Section 3.4.
They do not concern its irreducible-inner plethysm subproblem,
Problem 1.1.2. The distinction is material: the broad problem permits
direct sums of polynomial torus characters.

\begin{proposition}[P-F14: fixed-rank modular indices are unbounded]
\label{prop:torusmodular}
For each integer $N\ge2$, let $H=\mathbb G_m$, $G=GL_2$ and
\[
 \rho_N(t)=\operatorname{diag}(t^{N+1},t),\qquad
 \lambda=(1),\qquad \pi=2.
\]
This is a faithful explicit polynomial embedding allowed by
GCT VI, Section 3.4.1. Its stretching multiplicity is
\[
 f_N(n)=\dim\operatorname{Hom}_{\mathbb G_m}
       (\mathbb C_{2n},\operatorname{Sym}^n(\mathbb C^2)|_{\rho_N})
       =\mathbf1_{\{N\mid n\}},\qquad n\ge0,
\]
and its generating function has minimum possible modular index $N$:
\[
 F_N(T)=\frac1{1-T^N}.
 \tag{P-F14}
\]
The input has $\operatorname{rank}(x)=5$ in the exact definition of
GCT VI, Section 3.4. Consequently Hypothesis 3.3.4 (PH3) is false for
the full class of explicit polynomial subgroup restrictions as stated.
Here $p(f_N)=s(f_N)=0$, so this example does not refute PH2 or SH(a).
\end{proposition}

\begin{proof}
The representation defining $\rho_N$ is the direct sum of the
polynomial characters of weights $N+1$ and $1$. Section 3.4.1,
``Basic explicit homomorphisms'', explicitly permits constructions
from irreducible polynomial representations by $\oplus$ and $\otimes$;
it specifies that polynomial representations of a torus are sums
of polynomial one-dimensional characters. Both exponents here are
positive. The map is faithful because its second diagonal entry is
$t$, and the inverse on its image is that entry. Its image is closed:
inside the diagonal torus it is given by $z_1=z_2^{N+1}$.

A monomial $x^k y^{n-k}$, $0\le k\le n$, in the symmetric power
has $H$-weight $(N+1)k+(n-k)=n+Nk$. It has weight $2n$ if and only
if $Nk=n$. There is exactly one such monomial when $N\mid n$ and
none otherwise; when it exists, $k=n/N$ belongs to $[0,n]$.
This proves the multiplicity and generating-function formulas.
Its natural graded multiplicity algebra is
$\mathbb C[xy^{N-1}]$, a polynomial ring with generator of degree $N$:
the weight equation forces every admissible monomial to be a unique
power of $xy^{N-1}$.

Let $\zeta$ be a primitive $N$th root of unity. The rational function
$F_N$ has a pole at $\zeta$. In any expression
\[
 F_N(T)=\frac{B(T)}{\prod_j(1-T^{a_j})},
 \qquad B\in\mathbb C[T],\quad a_j>0,
\]
the denominator must vanish at $\zeta$; otherwise the right side is
regular there. Thus $\zeta^{a_j}=1$ for some $j$, so $N\mid a_j$
and $\max_j a_j\ge N$. This argument does not need positivity of
$B$ or a bound on the number of denominator factors. The displayed
form (P-F14) attains the lower bound, proving minimality. Its
constituent polynomials are the constants zero and one, so both
indices are zero by GCT VI, Definitions 1.2.2 and 1.2.4.

Section 3.4 defines
$\operatorname{rank}(x)=\operatorname{rank}H+\operatorname{rank}G+
\|\lambda\|+\|\pi\|$.
Section 3.4.2 specifies $\|\lambda\|$ as the height for a
$GL_m$ partition, and specifies the combinatorial length of a torus
character as $\dim H$. Therefore the four terms here are
$1,2,1,1$, giving five independently of $N$. A polynomial in this
fixed parameter cannot bound the unbounded modular indices $N$.
This directly contradicts Hypothesis 3.3.4 in the stated broad scope.
\end{proof}

\begin{proposition}[P-F15: fixed-rank positivity and saturation indices are unbounded]
\label{prop:torusindices}
For odd $N\ge3$, let $H=\mathbb G_m$, $G=GL_3$ and
\[
 \sigma_N(t)=\operatorname{diag}(t,t^3,t^{N+1}),
 \qquad \lambda=(1),\qquad\pi=2.
\]
This is again a faithful explicit polynomial embedding. Its stretching
multiplicity $g_N$ satisfies
\[
 g_N(n)=\#\{(v,w)\in\mathbb Z_{\ge0}^2:2v+Nw=n\},
 \qquad
 \sum_{n\ge0}g_N(n)T^n=\frac1{(1-T^2)(1-T^N)}.
\]
In the definitions of GCT VI,
\[
 p(g_N)=s(g_N)=N-2,
 \qquad \operatorname{rank}(x)=6.
 \tag{P-F15}
\]
Its natural graded multiplicity algebra is a polynomial ring on two
generators of degrees $2,N$. Thus the ungraded affine scheme is the
same smooth normal affine plane for every $N$; over finite fields it
has pure Tate ordinary and intersection cohomology, as in P-F10.
Thus Hypothesis 3.3.2(a) (the uniform polynomial saturation-index
bound) and Hypothesis 3.3.3 (PH2) are false for the full subgroup
restriction class with the stated parameter. This finding does not
assert a refutation of the density assertion in Hypothesis 3.3.2(b).
\end{proposition}

\begin{proof}
The three characters have positive weights, so the same exact source
definition as in P-F14 applies. The first diagonal entry is $t$,
giving faithfulness and the inverse on the image. Inside the diagonal
torus, that image is closed by the equations
$z_2=z_1^3$ and $z_3=z_1^{N+1}$.
A monomial $x^u y^v z^w$ of degree $n$ has $u+v+w=n$ and weight
\[
 u+3v+(N+1)w=n+2v+Nw.
\]
It has the required weight $2n$ exactly when $2v+Nw=n$.
Conversely every nonnegative solution of that equation satisfies
$v+w\le n$, because both coefficients are at least two.
It therefore gives the unique admissible value $u=n-v-w$.
This is a bijection with the monomial weight basis and proves the
counting formula and its generating series.
In fact, $0\le w\le\lfloor n/N\rfloor$ and
$w\equiv n\pmod2$, because $N$ is odd. Counting those values gives
the closed formula
\[
 g_N(n)=\left\lfloor
 \frac{\lfloor n/N\rfloor+2-(n\bmod2)}2
 \right\rfloor.
\]
It holds also at $n=0,1$. Integer division and parity computation
evaluate it in time polynomial in the binary lengths of $n,N$.

More precisely, let $S=\mathbb C[x,y,z]$, with ordinary polynomial
degree and torus weights $1,3,N+1$. Its diagonal multiplicity algebra
\[
 T_N=\bigoplus_{n\ge0}(S_n)_{2n}
\]
is closed under multiplication, since degrees and weights both add.
The monomial condition is $u=v+(N-1)w$, so its monomial basis is
\[
 x^{v+(N-1)w}y^v z^w=(xy)^v(x^{N-1}z)^w,
 \qquad v,w\ge0.
\]
The exponent triples for distinct pairs $(v,w)$ are distinct.
Hence these two generators are algebraically independent and generate
all of $T_N$, proving
$T_N=\mathbb C[xy,x^{N-1}z]\simeq\mathbb C[U,V]$ with
$\deg U=2$, $\deg V=N$. The same monomial proof works over every
field. For $k=\mathbb F_q$ and any prime $\ell\nmid q$, P-F10's
affine-space computation therefore gives
$H_c^4(\operatorname{Spec}(T_N)_{\overline k},\mathbb Q_\ell)
=\mathbb Q_\ell(-2)$ and zero in all other degrees, with intersection
complex $\mathbb Q_\ell[2]$. This is an actual graded multiplicity
algebra, rather than only a ring with an equal Hilbert function.

P-F10, already proved in full, applies with $a=2$, $b=N$ since
$\gcd(2,N)=1$, and gives
$p(g_N)=s(g_N)=2N-2-N=N-2$. To make the exact index test explicit,
the constituents can also be read from
\[
 \frac1{(1-T^2)(1-T^N)}
 =\frac{\displaystyle\sum_{i=0}^{N-1}\sum_{j=0}^{1}T^{2i+Nj}}
 {(1-T^{2N})^2}.
\]
The exponents $r=2i+Nj$ represent all residues modulo $2N$ exactly
once: reduction modulo two determines $j$, and reduction modulo $N$
then determines $i$. The constituent belonging to $r$ is
$n\mapsto (n-r)/(2N)+1$. This is valid even when $0\le n<r$,
because the only possible such index in the residue is $r-2N$,
where the formula is zero. The greatest representative is $3N-2$.
After a nonnegative integer shift $c$, the smallest constant term
is $(c-(N-2))/(2N)$. All slopes are positive, so coefficientwise
positivity is equivalent to $c\ge N-2$. For strict saturation
each nonzero constituent must be positive at every integer $n\ge1$;
the smallest value at $n=1$ is $(1+c-(N-2))/(2N)$, which is
strictly positive precisely when $c\ge N-2$. This rechecks both
indices directly with the source conventions.

The four terms in the source's rank parameter are now $1,3,1,1$,
giving six. Since $N-2$ is unbounded, no polynomial in this fixed
parameter bounds either index. These are precisely the universal
bounds asserted in Hypotheses 3.3.2(a) and 3.3.3. A sequence of
counterexamples to a universal bound does not by itself determine
the density appearing in Hypothesis 3.3.2(b).
\end{proof}

\paragraph{Exact source scope and consequences.}
The primary locators are GCT VI, Problem 1.1.3; Hypotheses 3.3.1--3.3.4,
printed pp.~50--51; the definition of $\operatorname{rank}(x)$ in
Section 3.4, printed pp.~52--53; Section 3.4.1, printed pp.~53--54;
and Section 3.4.2, the specifications of $[\rho]$, items (4)--(5),
and of $[\lambda],[\pi]$, printed pp.~55--56.
The source expressly applies Hypotheses 3.3.1--3.3.4 to this
subgroup-restriction input. The maps above use its permitted
direct sums, rather than an irreducible representation of $H$.
Every irreducible representation of the torus is one-dimensional,
so these two- and three-dimensional representations are not inputs
to the irreducible-inner plethysm Problem 1.1.2. No refutation of
Hypotheses 1.6.6--1.6.7 for that narrower family follows.
The conditional implication PH3$\Rightarrow$PH2 in P-F12 is also
unaffected; it does not assert PH3 for these broader inputs.

Both examples even have homogeneous rational-polytope models of
constant ambient dimension. For P-F14 use the singleton
$\{1/N\}\subset\mathbb R$, whose $n$th dilation contains one
integer exactly when $N\mid n$. For P-F15 use
\[
 \{(v,w)\in\mathbb R_{\ge0}^2:2v+Nw=1\}.
\]
Their descriptions have binary length $O(\log N)$, as do the
torus-character inputs. Thus PH1-style efficient descriptions are
compatible with these failures of the index bounds. The indices
in question grow exponentially in this binary length; adding only
the bitlength of $\rho$ to the rank parameter would not make a
polynomial index bound true. A valid replacement must restrict
the family or retain suitable magnitude information about the
embedding weights. These explicit counterexamples do not identify
what occurs for the narrower determinant/permanent obstruction
families, and no historical priority claim is made.
Both examples have easy exact counting formulas---divisibility in
P-F14 and the displayed floor formula in P-F15---so their large
indices do not themselves imply a computational lower bound.

In particular, the broad rank-only bounds cannot follow solely from
purity, smoothness, normality, a positive monomial basis, or an
efficient PH1 polytope: the examples above have all these properties.
Quantitative control of the grading or embedding weights is a
separate requirement even when the graded algebra already realizes
the multiplicities exactly.


\subsection{An embedding-dependent uniform bound for ray diagonals}

The preceding counterexamples vary the embedding while retaining its rank.
The following abstract sufficient theorem records what fixed multigraded
algebra data do control. It makes no claim that an arbitrary proposed
branching algebra already satisfies all the stated hypotheses.

\begin{proposition}[P-F16: fixed positive multigrading gives uniform ray bounds]
\label{prop:fixedmultigrading}
Let
\[
 R=\bigoplus_{\eta\in\mathbb Z^s}R_\eta
\]
be a finitely generated complex domain, generated as a complex algebra by
nonzero homogeneous elements $x_1,\ldots,x_m$, with $m\ge1$ and
$\deg(x_i)=\beta_i\in\mathbb Z^s$.
Assume that a real linear functional $\ell$ satisfies
$\ell(\beta_i)>0$ for every $i$. Write
\[
 M=(\beta_1\ \cdots\ \beta_m),\qquad
 \Delta=\max\bigl(\{1\}\cup
 \{|\det M_{I,J}|:\ |I|=|J|\ge1,\ \det M_{I,J}\ne0\}\bigr),
\]
and set
\[
 D=m\Delta,\qquad L=\operatorname{lcm}(1,2,\ldots,D),
 \qquad d=\dim R.
 \tag{P-F16a}
\]
For every $\gamma\in\mathbb Z^s$ with $\ell(\gamma)>0$, form the ray
diagonal, graded by $n$,
\[
 A_\gamma=\bigoplus_{n\ge0} R_{n\gamma}.
\]
Then $A_\gamma$ is generated by homogeneous elements of ray degree at
most $D$, uniformly in $\gamma$.

Suppose further, for the rays under consideration, that $A_\gamma\ne\mathbb C$,
that $A_\gamma$ is Cohen--Macaulay, and that
$f_\gamma(n)=\dim_\mathbb C R_{n\gamma}$ is a global quasi-polynomial
on all $n\ge0$. If $r_\gamma=\dim A_\gamma$, then $A_\gamma$ has an
h.s.o.p. consisting of $r_\gamma$ elements all of ray degree $L$.
Its Hilbert series therefore has a positive form
\[
 \sum_{n\ge0}f_\gamma(n)T^n
 =\frac{H_\gamma(T)}{(1-T^L)^{r_\gamma}},
 \quad H_\gamma\in\mathbb Z_{\ge0}[T]\setminus\{0\},
 \quad \deg H_\gamma<r_\gamma L.
 \tag{P-F16b}
\]
Consequently its positivity and saturation indices satisfy
\[
 s(f_\gamma)\le p(f_\gamma)
 \le\max\{0,\deg H_\gamma-L\}
 \le(r_\gamma-1)L\le(d-1)L.
 \tag{P-F16c}
\]
The last bound depends on the fixed presentation of $R$ and its weight
matrix, and is independent of $\gamma$. It is an existence bound;
it is not asserted to be polynomial in the binary encoding of $M$.
\end{proposition}

\begin{proof}
The positivity of $\ell$ first gives $R_0=\mathbb C$ and ensures that
every multihomogeneous component of the polynomial algebra
$P=\mathbb C[X_1,\ldots,X_m]$, with $\deg X_i=\beta_i$, is finite
dimensional. Indeed, if $M u=\eta$ with $u_i\ge0$, then
\[
 \sum_i u_i\ell(\beta_i)=\ell(\eta),
\]
which bounds every $u_i$ when such a solution exists. The homogeneous
surjection $P\longrightarrow R$, $X_i\mapsto x_i$, restricts to a
surjection of ray diagonals: a preimage of a multihomogeneous element
can be replaced by its component of the same multidegree. Thus it
suffices first to bound the generating degrees of the semigroup
\[
 S_\gamma=\{(u,n)\in\mathbb Z_{\ge0}^{m+1}:Mu=n\gamma\}.
 \tag{P-F16d}
\]
If this semigroup has no element with $n>0$, its algebra and $A_\gamma$
are just $\mathbb C$, and the first assertion is immediate. Otherwise
the polytope
\[
 Q_\gamma=\{u\in\mathbb R_{\ge0}^m:Mu=\gamma\}
\]
is nonempty. It is bounded, since
$u_i\le\ell(\gamma)/\ell(\beta_i)$.

We give the required rational-polyhedral argument explicitly.
At any vertex $v$ of $Q_\gamma$, the columns of $M$ indexed by
$J=\{i:v_i>0\}$ are linearly independent. Otherwise a nonzero
vector supported on $J$ in the kernel of $M$ would permit sufficiently
small perturbations of $v$ in both directions, contradicting extremality.
If $q=|J|$, select $q$ rows $I$ for which $M_{I,J}$ is nonsingular.
Cramer's rule shows that $|\det M_{I,J}|v$ is integral. Hence the
primitive integral generator of the ray through $(v,1)$ has last
coordinate at most $\Delta$.

For clarity, every point of $Q_\gamma$ is a convex combination of
these vertices without an additional integrality assumption. If its
positive-support columns are dependent, choose a nonzero supported
kernel vector. Such a vector has both positive and negative entries,
since applying $\ell$ to $M u=0$ forbids a nonzero nonnegative or
nonpositive kernel vector. Move in its two directions until a positive
coordinate first reaches zero. The original point is a convex
combination of these two endpoints, each with smaller support.
Induction on support size terminates at points with independent
support columns. Such a point is a vertex: in any convex decomposition,
its zero coordinates remain zero, and independence makes both endpoints
equal to it. There are only finitely many possible independent supports,
so there are finitely many such vertices.

Consider now the real cone
\[
 C_\gamma=\{(u,n)\in\mathbb R_{\ge0}^{m+1}:Mu=n\gamma\}.
\]
Its only point with $n=0$ is zero, by positivity of $\ell$.
Dividing a nonzero point by $n$ and using the preceding vertex
decomposition shows that $C_\gamma$ is generated by the finitely many
primitive integral rays just found. Its dimension $e$ is at most $m$:
the last coordinate is determined by
$n=\ell(Mu)/\ell(\gamma)$, so projection to $u$ is injective.
In particular the positivity assumptions exclude the rank-zero matrix
case.

Every nonzero $w\in S_\gamma$ can be expressed as a nonnegative real
combination of at most $e$ linearly independent primitive ray vectors,
say
\[
 w=\sum_{j=1}^k a_jv_j,\qquad a_j\ge0,\qquad k\le e\le m.
\]
To justify independence, begin with any finite ray expression.
If the active rays are dependent, choose a dependence, which has both
signs because their last coordinates are positive. Subtract a suitable
multiple of it from the coefficients until one active coefficient
becomes zero, retaining nonnegativity. Repeating yields independent
active rays.
Subtract the integer multiples $\lfloor a_j\rfloor v_j$ and obtain
the integral residual
\[
 w_0=\sum_{j=1}^k\{a_j\}v_j\in S_\gamma,
 \qquad n(w_0)<k\Delta\le m\Delta=D.
 \tag{P-F16e}
\]
This is the fundamental-parallelepiped argument, with its degree
estimate made explicit. The primitive vectors themselves have ray
degree at most $\Delta\le D$. Hence $S_\gamma$ is generated by its
elements of degrees at most $D$. There are only finitely many such
elements, by the same positive-functional bounds on their coordinates.
Its semigroup algebra is the polynomial ray diagonal; passing to its
quotient $A_\gamma$ proves the uniform generating-degree assertion.

Fix now a ray satisfying the additional assumptions, and abbreviate
$A=A_\gamma$ and $r=\dim A$. Choose nonzero homogeneous algebra
generators $y_1,\ldots,y_t$ of degrees $1\le b_j\le D$.
Let $B=\mathbb C[A_L]$, the subalgebra generated by its degree-$L$
component. This component is finite dimensional. Since $L$ is divisible
by every $b_j$, the equation
\[
 Y^{L/b_j}-y_j^{L/b_j}=0
\]
is a monic equation for $y_j$ with coefficients in $B$.
Thus $A$ is integral and finite over $B$: adjoining finitely many
integral elements gives a finite module, as is seen by reducing powers
using their monic equations. The domain $B$ has dimension $r$ and is
standard graded when its degree $L$ is renamed degree one.

Here is a direct homogeneous Noether-normalization argument giving
parameters of that same degree. If a standard graded domain $B$ is
generated by $z_1,\ldots,z_q$ of degree one and $q>\dim B$, choose a
nonzero homogeneous polynomial relation $F(z_1,\ldots,z_q)=0$.
Over the infinite field $\mathbb C$, choose $c_1,\ldots,c_{q-1}$ with
$F(c_1,\ldots,c_{q-1},1)\ne0$. Such a choice exists because
dehomogenizing a nonzero homogeneous polynomial at its last variable
gives a nonzero polynomial. Put $z_i'=z_i-c_i z_q$ for $i<q$.
After substituting $z_i'+c_i z_q$ in $F$, its leading coefficient
as a polynomial in $z_q$ is the displayed nonzero scalar.
Rescaling makes this a monic relation for $z_q$ over
$B'=\mathbb C[z_1',\ldots,z_{q-1}']$. Therefore $B$ is finite over
$B'$, while $\dim B'=\dim B$. Induction on the number of generators
terminates with $r$ algebraically independent linear forms
$\theta_1,\ldots,\theta_r$ such that $B$, and hence $A$, is finite
over $\mathbb C[\theta_1,\ldots,\theta_r]$. These forms all have
original ray degree $L$. They are an h.s.o.p. for $A$.

We spell out the Cohen--Macaulay step to keep its hypothesis visible.
Put $\mathfrak m=A_{>0}$. The quotient by all the $\theta_i$ is
finite dimensional, because $A$ is finite over their polynomial
algebra. It is connected and nonnegatively graded; every element of
positive degree is nilpotent. Thus $\mathfrak m$ is the unique maximal
ideal above $(\theta_1,\ldots,\theta_r)$, and the local ring
$A_{\mathfrak m}$ has dimension $r$ (equivalently, lift a maximal
prime chain from the parameter polynomial algebra by going up).
The Cohen--Macaulay regular-parameter criterion therefore says that
$\theta_1,\ldots,\theta_r$ is a regular sequence in $A_{\mathfrak m}$;
an exact primary reference is the Stacks Project,
Lemma 10.104.2, Tag \texttt{02JN},
\url{https://stacks.math.columbia.edu/tag/02JN}.
It is also a regular sequence in $A$. Indeed the kernels of the
successive homogeneous multiplications are graded. A nonzero
homogeneous element of such a kernel cannot disappear on localization
at $\mathfrak m$: multiplying it by an element outside $\mathfrak m$
has a nonzero lowest-degree term, coming from the nonzero scalar
constant coefficient of that element. Local injectivity therefore
implies graded global injectivity.

The exact sequences for multiplication by these regular homogeneous
parameters give
\[
 \operatorname{Hilb}_{A/(\theta_1,\ldots,\theta_r)}(T)
 =(1-T^L)^r\operatorname{Hilb}_{A}(T).
\]
The quotient is finite dimensional and has nonnegative integer graded
dimensions and constant coefficient one. Its Hilbert series, which
is a polynomial, is the required $H_\gamma(T)$. The global quasi-polynomial assumption
and the properness argument of P-F9 give
$\deg H_\gamma<rL$.
In this equal-denominator-degree situation a direct argument improves
the general bound of P-F11. Write $h=\deg H_\gamma$ and
$H_\gamma(T)=\sum_{i=0}^h H_iT^i$.
For each residue $b$ modulo $L$, the corresponding eventual constituent
of $f_\gamma$ is
\[
 q_b(X)=\sum_{\substack{0\le i\le h\\i\equiv b\pmod L}}
 H_i\binom{(X-i)/L+r-1}{r-1}.
\]
This follows by expanding $(1-T^L)^{-r}$ at sufficiently large
indices in that residue. Since $f_\gamma$ is globally
quasi-polynomial, these are its constituents everywhere: any
constituent for a common refinement of its period with $L$ agrees
with the displayed polynomial at infinitely many integers and hence
is identical to it. Put $c=\max\{0,h-L\}$. After the integer shift
$c$, each binomial polynomial appearing in $q_b(X+c)$ is
\[
 \frac{1}{(r-1)!}\prod_{j=1}^{r-1}
 \left(\frac{X}{L}+\frac{c-i}{L}+j\right).
\]
All its coefficients are nonnegative, because
$c-i+L\ge c-h+L\ge0$ and $j\ge1$.
For $r=1$ the empty product is one. Nonnegative linear combinations
retain that property, and shifting merely permutes the residue
labels. Therefore $p(f_\gamma)\le\max\{0,h-L\}\le(r-1)L$,
using $h<rL$; when $r=1$ this bound is zero.
Since $A$ is a finitely generated subdomain of $R$,
$r=\operatorname{trdeg}_\mathbb C\operatorname{Frac}(A)
\le\operatorname{trdeg}_\mathbb C\operatorname{Frac}(R)=d$.
Finally, a nonzero polynomial with nonnegative coefficients is
strictly positive at every positive integer, so coefficient positivity
after a shift implies the source's strict saturation after the same
shift. This proves all assertions.
\end{proof}

\paragraph{Precise use and remaining gate.}
For a fixed embedding, an application requires a single finitely
generated, positively multigraded domain whose $(\lambda,\pi)$
components are the actual branching multiplicity spaces, together with
Cohen--Macaulayness and global quasi-polynomiality for the ray diagonals.
A polynomial branching algebra is a possible input to verify.
The full algebra $\mathbb C[GL_m/U]$ has determinant units and cannot
be inserted into the positivity hypothesis on $\ell$ without addressing
its grading. No such verification is supplied by the abstract theorem.

The quantitative data required here are the generator number, their
multidegrees, the resulting minors, and the parameter degree $L$.
They can change when the embedding changes, as in P-F14--P-F15.
The formula $L=\operatorname{lcm}(1,\ldots,m\Delta)$ is a finite
explicit bound; even a polynomial numerical bound on $m\Delta$
does not by itself turn this particular formula into a polynomial
bound for $L$. A uniform PH3 or PH2 conclusion needs stronger control
of parameter degrees, or another bound avoiding this least-common-multiple
growth. The theorem supplies a corrected sufficient gate and a fixed-data
uniformity statement, not the missing uniform algorithm or a claim of
historical priority.


\subsection{Uniform finite bounds for a fixed semisimple embedding}

\begin{proposition}[P-F17: the fixed-embedding gate closes for semisimple branching]
\label{prop:fixedplethysm}
Fix connected semisimple complex algebraic groups $H,G$ and an algebraic
homomorphism $\rho:H\to G$. There are finite constants $L_\rho\ge1$
and $C_\rho\ge0$, depending on this fixed homomorphism, such that the
following holds for every pair of dominant integral highest weights
$\lambda,\pi$ of $G,H$. Put
\[
 f_{\lambda,\pi}(n)
 =\dim\operatorname{Hom}_H
       (V_{n\pi}(H),V_{n\lambda}(G)|_\rho).
\]
If this ray has a nonzero positive-degree value, its canonical
highest-weight multiplicity algebra has an h.s.o.p. all of whose
degrees are $L_\rho$, and its Hilbert series has a positive form with
all denominator weights $L_\rho$. Moreover
\[
 s(f_{\lambda,\pi})\le p(f_{\lambda,\pi})\le C_\rho.
 \tag{P-F17a}
\]
The largest degree in a lexicographically minimal h.s.o.p. of that
canonical algebra, the modular index of GCT VI, Section 4.1.1,
is also at most $L_\rho$.
For a ray vanishing in all positive degrees, use the zero-constituent
convention and index zero; its isolated value at $n=0$ is not called
a global quasi-polynomial. The trivial pair $(\lambda,\pi)=(0,0)$
has $f(n)=1$ and both indices zero.

Consequently, for fixed $n$ and a fixed partition $\mu$ of length at
most $n$, there is a finite constant $C_{n,\mu}$ bounding the positivity
and saturation indices of all rays
\[
 k\longmapsto
 \dim\operatorname{Hom}_{GL_n}
 \bigl(S^{k\pi}\mathbb C^n,
       S^{k\lambda}(S^\mu\mathbb C^n)\bigr),
 \tag{P-F17b}
\]
where $\lambda,\pi$ range over their admissible partition lengths.
There is likewise one finite bound on the denominator weights in
some positive rational forms, and on the canonical modular indices,
for all nonzero such rays. No polynomial
bound on these constants in $n$, $\mu$, or the group dimension, and
no uniform algorithm producing those forms, is asserted.
\end{proposition}

\begin{proof}
We verify the algebraic inputs to P-F16, including the multiplication
on its diagonals. Choose Borel subgroups of $G,H$, with unipotent
radicals $U_G,U_H$ and maximal tori $T_G,T_H$.
For a dominant weight $\alpha$, write
$\alpha^*=-w_0\alpha$ for the highest weight of the dual representation.
Let
\[
 A_G=\mathbb C[G]^{U_G},\quad
 A_H=\mathbb C[H]^{U_H},\quad
 R_\rho=(A_G\otimes A_H)^H,
 \tag{P-F17c}
\]
where $H$ acts from the left through $\rho$ on the first factor and
by its left regular action on the second. The unipotent invariants
are for the right actions, so the left and right actions commute.

The exact primary construction is GCT VI, proof of Theorem 4.4.1,
printed pp.~75--76, equations (4.7)--(4.8): $A_G,A_H$ are finitely
generated, and their left-group/right-torus decomposition is
\[
 A_G=\bigoplus_{\lambda}V_\lambda(G)
       \quad\text{with right weight }\lambda^*,\qquad
 A_H=\bigoplus_{\nu}V_\nu(H)
       \quad\text{with right weight }\nu^*.
\]
That proof applies Hilbert finite generation for reductive invariants
to $R_\rho$. The same construction works for a homomorphism with
kernel: only the algebraic $H$-action and reductivity are used.
Furthermore $A_G\otimes A_H$ embeds in $\mathbb C[G\times H]$,
which is a domain because the connected complex groups are
irreducible. Thus $R_\rho$ is a domain.

Grade $R_\rho$ by the right $T_G\times T_H$ characters. Its component
of degree $(\lambda^*,\pi)$ is exactly
\[
 (R_\rho)_{(\lambda^*,\pi)}
   =(V_\lambda(G)\otimes V_{\pi^*}(H))^H
   \simeq\operatorname{Hom}_H(V_\pi(H),V_\lambda(G)|_\rho).
 \tag{P-F17d}
\]
All degrees lie in the product of the two dominant-weight cones.
For semisimple groups this cone is pointed: the sum of the pairings
with all simple coroots is nonnegative on it and strictly positive
on each nonzero dominant weight. This supplies the functional $\ell$
in P-F16. The zero-degree component is $\mathbb C$. A finite generating
set can be replaced by its finitely many homogeneous components;
omit scalar components. Choosing lattice bases for the character
lattices gives a fixed integral generator matrix $M_\rho$.

It remains to check the Cohen--Macaulay and global quasi-polynomial
conditions for the actual ray subalgebras, rather than merely for
some rings with equal Hilbert functions. For a dominant weight
$\lambda$ of $G$, let
\[
 C_G(\lambda)=\bigoplus_{k\ge0}(A_G)_{k\lambda^*},
\]
with its external grading by $k$; for $\lambda=0$ this means
$\mathbb C[t]$ with a trivial $G$-action. When $\lambda\ne0$,
this is the Cartan section ring of the highest-weight projective
orbit in $\mathbb P(V_{\lambda^*})$. Here is a direct verification
of its multiplication. Fix a highest vector $v\in V_{\lambda^*}$.
The functions $g\mapsto\varphi(gv)$, for
$\varphi\in V_{\lambda^*}^*=V_\lambda$, form precisely
$(A_G)_{\lambda^*}$. Their products of degree $k$ lie in
$(A_G)_{k\lambda^*}=V_{k\lambda}$. The product map is $G$-equivariant
and nonzero: the $k$th power of any nonzero such function is nonzero
in the domain $\mathbb C[G]$. Its image is therefore all the
irreducible module $V_{k\lambda}$. Thus the actual multiplication
is the Cartan multiplication, and this ray ring is the homogeneous
coordinate ring of that affine highest-weight cone. Equivalently,
the pullback along $g\mapsto gv$ realizes its degree-one coordinate
functions and all their products. A nonzero highest-weight character
takes all values in $\mathbb C^*$, so the orbit closure is the
affine cone over the projective orbit.

Define $C_H(\pi^*)$ in the same way; its degree-$k$ component is
$V_{k\pi}^*$. Give the two factors opposite degree actions of
$\mathbb G_m$. For $\gamma=(\lambda^*,\pi)\ne0$, the componentwise
matrix-coefficient identifications give an isomorphism of graded
algebras
\[
 \bigoplus_{k\ge0}(R_\rho)_{k\gamma}
 \simeq
 \left((C_G(\lambda)\otimes C_H(\pi^*))^{\mathbb G_m}\right)^H
 =\bigoplus_{k\ge0}
      (V_{k\lambda}(G)\otimes V_{k\pi}(H)^*)^H.
 \tag{P-F17e}
\]
In the rightmost expression the multiplication is the tensor product
of the Cartan multiplications just checked. This proves compatibility
with multiplication, not just equality of the displayed vector spaces.
The formula remains valid when one weight is zero: its external
section ring is $\mathbb C[t]$, and the nonzero joint weight
distinguishes different diagonal degrees.

This is exactly the ring displayed explicitly in GCT VI,
Section 4.1.2, equation (4.5), printed p.~69. Its multiplication is
the one constructed in the proof of Theorem 4.1.1,
equations (4.2)--(4.3), printed p.~67, applied to the highest-weight
cone as in Section 3.5.1, printed p.~59, with stretch degree one.
The section rings of the highest-weight cones are normal and have
rational singularities. The proof then uses the opposite-degree
$\mathbb G_m$ quotient and the reductive $H$ quotient to retain
rational singularities. Consequently the algebra in (P-F17e) is
normal and Cohen--Macaulay. These are the exact algebraic assertions
of Theorem 3.4.1(c) and the construction in Theorem 4.1.1(c), not an
inference from its numerical Hilbert series. For a ray whose algebra
is different from $\mathbb C$,
Theorem 3.4.1(b), proved there using Lemma 4.1.3, supplies global
quasi-polynomiality. The constant ring arising when all positive
degrees vanish is excluded as specified in the statement. No bound
on the numerator degree from the erroneous index in display (4.4)
is used.

All hypotheses of P-F16 now hold simultaneously for the
positive-dimensional ray algebras of the one fixed $R_\rho$.
If $R_\rho\ne\mathbb C$,
choose its homogeneous generators as above, put
\[
 L_\rho=\operatorname{lcm}(1,\ldots,m_\rho\Delta_\rho),\qquad
 C_\rho=(\dim R_\rho-1)L_\rho,
\]
and apply P-F16. These choices are independent of $\lambda,\pi$.

We also justify the claim about the lexicographically minimal h.s.o.p.;
existence of one all-$L_\rho$ h.s.o.p. alone is not the comparison
needed for that claim. Fix a ray algebra $A$ of dimension $r$ and
write $L=L_\rho$. The proof of P-F16 shows that every homogeneous
generator has a power in $A_L$. Hence the ideal generated by $A_L$
has radical $A_{>0}$. After quotienting by any partial homogeneous
system of parameters, its image still has radical the positive-degree
maximal ideal of the quotient. If that quotient has positive
dimension, no positive-dimensional minimal prime can contain the
whole image of $A_L$. There are finitely many minimal primes, and
over $\mathbb C$ a finite union of proper linear subspaces does not
cover the finite-dimensional vector space $A_L$. Choose an element
avoiding the minimal primes of maximal positive dimension. Cutting
by it lowers the dimension by one: on every component of maximal
dimension it is nonzero, and the principal ideal dimension theorem
gives the drop by exactly one. Repeat until dimension zero. This
extends any partial parameter sequence by parameters all of degree
$L$. Arrange parameter degrees nondecreasingly; permuting an h.s.o.p.
does not change its parameter property. Now fix a lexicographically
minimal such degree sequence. Its first
entry is at most $L$, by existence of the all-$L$ sequence. Given
any initial part of this minimal sequence, extend the corresponding
parameters by degree-$L$ elements as above. Lexicographic minimality
forces its next entry to be at most $L$. Induction bounds every
entry by $L$, proving the claimed modular-index bound.

The zero joint weight corresponds only to $(0,0)$; its multiplicity
function is one. Its section algebra is $\mathbb C[t]$ and the
parameter $t^{L_\rho}$ has degree $L_\rho$, with positive numerator
$1+T+\cdots+T^{L_\rho-1}$. Thus it obeys the same denominator
bound and index conclusion. If $R_\rho=\mathbb C$, every nonzero
joint-weight ray has zero positive-degree multiplicities; taking
$L_\rho=1$ and $C_\rho=0$ handles that case too. This proves
the semisimple statement.

Finally fix $n,\mu$ as in (P-F17b), put
$W=S^\mu\mathbb C^n$, and write $m=\dim W$.
The representation map $\rho_\mu:GL_n\to GL(W)$ restricts to
\[
 \overline\rho_\mu:SL_n\longrightarrow SL_m.
\]
Indeed, the determinant of its restriction is an algebraic character
of the connected semisimple group $SL_n$ and is therefore trivial.
For clarity, this last fact follows also by differentiating: the
derivative of a character vanishes on commutators, while the
semisimple Lie algebra equals its commutator algebra; a character
with zero derivative on a connected group has trivial image.

If $\alpha$ is a $GL_j$ highest weight, write $\bar\alpha$ for its
$SL_j$ highest weight, obtained by forgetting a determinant twist;
restriction of $V_\alpha(GL_j)$ is the irreducible
$V_{\bar\alpha}(SL_j)$, and $\overline{k\alpha}=k\bar\alpha$.
The source of the Hom space in (P-F17b) has central character
$z^{k|\pi|}$, while its target has central character
$z^{k|\mu||\lambda|}$. If
\[
 |\pi|\ne|\mu||\lambda|,
 \tag{P-F17f}
\]
all positive-degree multiplicities vanish. If the two integers agree,
every $SL_n$-equivariant map between these representations is already
$GL_n$-equivariant, since the central scalar subgroup acts by the
same character on both sides and $GL_n$ is generated by that subgroup
and $SL_n$. Explicitly any $g\in GL_n(\mathbb C)$ can be written
as $zh$ with $h\in SL_n$, after choosing an $n$th root of $\det g$.
Hence, on the balanced rays, the multiplicities in (P-F17b) equal
\[
 \dim\operatorname{Hom}_{SL_n}
     (V_{k\bar\pi},V_{k\bar\lambda}|_{\overline\rho_\mu}).
\]
The fixed semisimple result applies and gives $C_{n,\mu}$ and a
uniform positive-form denominator bound. The cases $n=1$ or $m=1$
reduce to trivial semisimple groups and yield the already handled
constant or zero multiplicity functions.

For the canonical modular index we must also identify multiplication,
rather than infer the claim from this equality of coefficients.
Restriction of matrix-coefficient functions from $GL_j$ to $SL_j$
gives, degree by degree, an isomorphism of the Cartan section ring
$\bigoplus_{k\ge0}V_{k\alpha}(GL_j)$ with
$\bigoplus_{k\ge0}V_{k\bar\alpha}(SL_j)$ as a graded
$SL_j$-algebra. Each degree map is nonzero (choose a matrix coefficient
nonzero at the identity) and maps between the isomorphic irreducible
$SL_j$ modules; it is therefore an isomorphism. Restriction of
functions respects multiplication, so the degree maps are compatible
with all Cartan products. If $\bar\alpha=0$, the external grading
is retained, giving $\mathbb C[t]$ on the $SL_j$ side instead of
collapsing different degrees into constants.
Apply these maps to both factors of the diagonal section algebra
$\bigoplus_k V_{k\lambda}(GL_m)\otimes V_{k\pi}(GL_n)^*$.
On a balanced ray its central $GL_n$ action is trivial in every
degree. Its $GL_n$ and $SL_n$ invariant subalgebras consequently
coincide, and the compatible restriction maps identify this canonical
graded invariant algebra with its semisimple counterpart. H.s.o.p.
degrees and their lexicographic minimum are preserved by a graded
algebra isomorphism. The canonical modular-index bound therefore
passes to (P-F17b) as well. This proves the corollary.
\end{proof}

\paragraph{What has been gained.}
For fixed semisimple branching, the unspecified realization gate in
P-F16 is now verified: the same multigraded algebra contains all the
actual canonical ray algebras, with the required multiplication and
singularity properties. For plethysm both $n$ and $\mu$ are fixed;
fixing only the shape $\mu$ while allowing $n$ to grow is not a
fixed embedding and is not covered by this conclusion.
The remaining quantitative question is how these generator and
parameter bounds grow as the embedding varies. Neither the existence
argument nor its least-common-multiple bound proves the uniform
polynomial growth required in the narrow PH2/PH3 hypotheses. The
torus counterexamples to the unrestricted rank-only formulation
remain consistent with this fixed-embedding result. No historical
priority is asserted.


\subsection{A qualitative bound depending only on the inner dimension}

\begin{proposition}[P-F18: finite dimension envelope for ordinary plethysm]
\label{prop:innerdimensionbound}
For every integer $m\ge1$ there is a finite integer $C(m)\ge1$ with
the following property. Let $n\ge1$, let $\mu$ be a partition of
length at most $n$, and suppose
\[
 W=S^\mu\mathbb C^n,\qquad \dim W=m.
\]
For partitions $\lambda,\pi$ of lengths at most $m,n$, respectively,
put
\[
 f_{n,\mu,\lambda,\pi}(s)
 =\dim\operatorname{Hom}_{GL_n}
   \bigl(S^{s\pi}\mathbb C^n,S^{s\lambda}(W)\bigr).
 \tag{P-F18a}
\]
If this function has a nonzero positive-degree value, the canonical
graded multiplicity algebra has a positive Hilbert form whose
denominator weights are at most $C(m)$; its canonical modular index
$\chi$, defined by a lexicographically minimal h.s.o.p., also satisfies
\[
 \chi\le C(m),\qquad
 s(f_{n,\mu,\lambda,\pi})\le
 p(f_{n,\mu,\lambda,\pi})\le C(m).
 \tag{P-F18b}
\]
For a ray zero in every positive degree, assign indices and $\chi$
equal to zero. Its literal value at degree zero is one and is treated
separately from its zero quasi-polynomial constituents.

No polynomial growth, effective size estimate, or uniform algorithm
for $C(m)$ is asserted. More precisely, define the optimal envelopes
over the ordinary plethysm family (P-F18a) of inner dimension $m$ by
\[
 \Xi(m)=\max\bigl(\{1\}\cup\{\chi\}\bigr),\qquad
 P(m)=\max\bigl(\{0\}\cup\{p(f_{n,\mu,\lambda,\pi})\}\bigr).
 \tag{P-F18c}
\]
These maxima are finite. The restriction of GCT VI,
Hypothesis 1.6.7 (including its canonical modular-index clause), to
ordinary $GL_n$ plethysm is equivalent to polynomial growth of
$\Xi(m)$. The corresponding restriction of Hypothesis 1.6.6 is
equivalent to polynomial growth of $P(m)$.
Here the canonical modular index uses the parameter ring $T$
in the detailed definition of GCT VI, Section 4.1.1, printed p.~68,
and the multiplicity-ring formula (4.5), printed p.~69.
The introductory paragraph before Hypothesis 1.6.7 instead writes
$S$; we follow the detailed $T=S^H$ definition, whose Hilbert
function is the multiplicity, consistently with Theorem 1.6.1(c).
\end{proposition}

\begin{proof}
Assume first $m>1$, pad $\mu$ with zeros to length $n$, and write
\[
 k=\mu_n,\qquad \nu=\mu-k(1^n),\qquad
 \delta_i=\mu_i-\mu_{i+1}\quad(1\le i<n).
\]
Then $\nu_n=0$, the integers $\delta_i$ are nonnegative, and
\[
 S^\mu\mathbb C^n\simeq
 (\det)^k\otimes S^\nu\mathbb C^n.
 \tag{P-F18d}
\]
In particular, the dimension and the restriction to $SL_n$ are
unchanged by replacing $\mu$ with $\nu$.

We give explicit bounds on the possible normalized highest weights.
The type-$A$ Weyl dimension formula is
\[
 m=\prod_{1\le i<j\le n}
       \frac{\nu_i-\nu_j+j-i}{j-i}
   =\prod_{1\le i<j\le n}
       \left(1+\frac{\delta_i+\cdots+\delta_{j-1}}{j-i}\right).
 \tag{P-F18e}
\]
The primary locator is Borcherds--Haiman--Reshetikhin--Serganova--
Johnson-Freyd, \emph{Lectures in Lie Groups and Quantum Groups},
Part I (2016), Theorem 6.1.1.30, printed pp.~102--103;
see also its type-$A$ Example 6.1.1.31, printed p.~103.
(P-F18e) is its $GL_n$ specialization, whose central character
does not affect dimension. Every factor in (P-F18e) is at least
one, and its factor for $(i,j)=(i,i+1)$ is $\delta_i+1$.
Consequently
\[
 0\le\delta_i\le m-1\qquad(1\le i<n).
 \tag{P-F18f}
\]
Moreover $n\ge2$ and at least one $\delta_i$ is nonzero: if all
were zero, (P-F18e) would give $m=1$.
Choose an index $i$ with $\delta_i\ge1$. Every factor in (P-F18e)
is monotone in every $\delta_j$, so comparison with the weight
having only $\delta_i=1$ gives
\[
 m\ge\prod_{a=1}^{i}\prod_{b=i+1}^{n}
                  \frac{b-a+1}{b-a}
   =\prod_{a=1}^{i}\frac{n-a+1}{i-a+1}
   =\binom ni\ge n.
 \tag{P-F18g}
\]
For the last inequality, the binomial coefficients increase from
$\binom n1=n$ up to the middle because
$\binom n{j+1}/\binom nj=(n-j)/(j+1)\ge1$ there, and symmetry
handles the remaining indices $1\le i\le n-1$.

It follows that, for a fixed $m>1$, the pairs $(n,\nu)$ that can
occur belong to the finite set specified by
\[
 2\le n\le m,\qquad \nu_n=0,\qquad
 \nu_i-\nu_{i+1}\in\{0,1,\ldots,m-1\},\qquad
 \dim S^\nu\mathbb C^n=m.
 \tag{P-F18h}
\]
Indeed $\nu_i=\sum_{j=i}^{n-1}\delta_j$ is uniquely determined
by the bounded tuple of differences. There are at most
$\sum_{n=2}^{m}m^{n-1}$ tuples before imposing the dimension
condition. This is a finiteness observation, not a polynomial bound
on the representation-theoretic constants below.

Each pair in (P-F18h) defines one algebraic homomorphism
\[
 \overline\rho_{n,\nu}:SL_n\longrightarrow SL_m,
 \tag{P-F18i}
\]
after a choice of basis. The image lies in $SL_m$ because the
determinant of a representation of a connected semisimple group is
a trivial character. The map is not required to be injective;
it can have a central kernel. Changing the basis conjugates the
homomorphism and gives isomorphic canonical multiplicity algebras.
By (P-F18d), every $\mu$ of inner dimension $m$ restricts to one
of these finitely many $SL_n$ representations.

We spell out why determinant twists and the central characters do
not reintroduce dependence on their unbounded sizes. A scalar
$zI_n$ acts on the two representations in (P-F18a) by the
characters $z^{s|\pi|}$ and $z^{s|\mu||\lambda|}$.
If
\[
 |\pi|\ne |\mu||\lambda|,
 \tag{P-F18j}
\]
then all positive-degree multiplicities are zero. If equality
holds, the central actions agree, and $GL_n$-equivariance is
equivalent to $SL_n$-equivariance. Indeed every complex invertible
matrix is a product of a scalar matrix and a matrix of determinant
one. Write $\bar\lambda,\bar\pi$ for the corresponding highest
weights of $SL_m,SL_n$. On a balanced ray (P-F18a) is therefore
\[
 \dim\operatorname{Hom}_{SL_n}
  \left(V_{s\bar\pi},
       V_{s\bar\lambda}|_{\overline\rho_{n,\nu}}\right).
 \tag{P-F18k}
\]
P-F17 proves this identification as an isomorphism of the actual
graded Cartan multiplicity algebras: matrix-coefficient restriction
respects products, and on the balanced diagonal the scalar action
is trivial, so the $GL_n$ and $SL_n$ invariant subalgebras coincide.
The external grading is retained even when a highest weight is a
determinant character. Thus the conclusion applies to the canonical
modular index as well as to the numerical Hilbert function.

For each of the finitely many homomorphisms (P-F18i), choose the
finite bounds $L_{n,\nu}$ and $C_{n,\nu}$ provided by P-F17.
Take
\[
 C(m)=\max\bigl(\{1\}\cup
       \{L_{n,\nu},C_{n,\nu}:(n,\nu)\text{ satisfies (P-F18h)}\}\bigr).
 \tag{P-F18l}
\]
P-F17 now gives (P-F18b) for every balanced ray with a nonzero
positive-degree value, uniformly over $\lambda,\pi$ and every
determinant twist $k$. It also handles the trivial joint $SL$
weight, whose canonical section algebra is $\mathbb C[t]$.
The remaining rays have the stated zero convention.

If $m=1$, (P-F18e), or the single-coordinate case $n=1$, shows
that $\mu=(k^n)$ and $W=(\det)^k$. Since $\lambda$ has at most
one part,
\[
 S^{s\lambda}(W)=(\det)^{sk|\lambda|}.
\]
Thus every positive-degree multiplicity is one if
$\pi=(k|\lambda|)^n$, and zero otherwise. In the first case the
canonical algebra is $\mathbb C[t]$, its modular index is one,
and its positivity and saturation indices are zero. The second
case has the stated zero convention. Taking $C(1)=1$ completes
the finiteness proof.

The maxima in (P-F18c) exist because their entries are nonnegative
integers bounded by $C(m)$; the sets are nonempty by the added
values one and zero. These envelopes depend on the optimal
canonical modular and positivity indices, not on an arbitrary
choice of the coarse bounds in (P-F18l).
For $m>1$, (P-F18g) and the admissible partition lengths give
\[
 m^2\le
 h=\dim GL(W)+\operatorname{ht}(\lambda)+\operatorname{ht}(\pi)
 \le m^2+m+n\le m^2+2m.
 \tag{P-F18m}
\]
These are exactly the parameters in GCT VI, Hypotheses 1.6.6--1.6.7.
If the canonical modular index is bounded by a fixed polynomial in
$h$, then (P-F18m) bounds $\Xi(m)$ by a fixed polynomial in $m$.
Conversely, a polynomial bound on $\Xi(m)$ bounds every canonical
modular index polynomially in $h$, since $m^2\le h$.
The $m=1$ case satisfies both bounds directly. This proves the
claimed equivalence for the full canonical clause of ordinary
plethysm PH3; the associated positive form is supplied by its
h.s.o.p. The identical argument with $p$ proves the PH2 equivalence
for $P(m)$.
\end{proof}

\paragraph{Quantitative scope.}
This closes a qualitative uniformity question for ordinary
$GL_n$ plethysm: even allowing $n$, the determinant twist of $\mu$,
and the outer and target partitions to vary, fixed inner dimension
permits finite bounds. The canonical clause of PH3 would require
polynomial growth of the optimal envelope $\Xi(m)$. The next
result, P-F19, shows that this requirement fails on an explicit
subsequence, while the finite bound proved here remains valid.
The weaker existence of some bounded positive rational form and the
corresponding PH2 question for $P(m)$ are separate assertions.
P-F31 later refutes the former by an actual ordinary Hilbert-series
pole; the PH2 question remains open here. P-F12 supplies the valid
conditional implication between those two bounds.
The proof here neither estimates the finite bound nor constructs
the required parameters uniformly. Polynomial growth of one particular coarse choice
$C(m)$ would suffice, but is not asserted to be equivalent to PH3.
The unrestricted reductive-group formulation of Problem 1.1.2 is
not reduced to this ordinary $GL_n$ case here. For $m>1$, the parameter
$h\asymp m^2$ is also not the binary encoding length of the inner
partition, so no compressed-input polytope construction follows.
No historical priority is claimed.


\subsection{An exponential canonical modular index inside ordinary plethysm}

\begin{proposition}[P-F19: chain cubics obstruct the canonical clause of PH3]
\label{prop:cubiccanonicalcounterexample}
Let $n\ge2$, put
\[
 g=\gcd(n,3),\qquad d_0=\frac ng,\qquad
 W=\operatorname{Sym}^3\mathbb C^n,\qquad
 m=\dim W=\binom{n+2}{3},
\]
and consider the ordinary plethysm ray with
\[
 \mu=(3),\qquad \lambda=(d_0),\qquad
 \pi=((3/g)^n).
 \tag{P-F19a}
\]
Its canonical multiplicity algebra is
\[
 T^{(n)}=\bigoplus_{s\ge0}
  \operatorname{Hom}_{GL_n}
  \left((\det)^{3s/g},\operatorname{Sym}^{d_0s}W\right),
 \tag{P-F19b}
\]
with the Cartan multiplication. Every homogeneous system of parameters
of $T^{(n)}$ contains a parameter of stretch degree at least
$2^{n-1}$. In particular its canonical modular index satisfies
\[
 \chi(T^{(n)})\ge2^{n-1}.
 \tag{P-F19c}
\]
The parameter of GCT VI, Hypothesis 1.6.7, is
\[
 h=m^2+\operatorname{ht}(\lambda)+\operatorname{ht}(\pi)
   =\binom{n+2}{3}^{\!2}+n+1=O(n^6).
 \tag{P-F19d}
\]
Consequently the canonical modular-index clause of that hypothesis
is false already for ordinary plethysm with the fixed inner shape
$\mu=(3)$. Equivalently, the optimal canonical envelope in P-F18
has the subsequence
\[
 \Xi\!\left(\binom{n+2}{3}\right)\ge2^{n-1}.
 \tag{P-F19e}
\]
This statement concerns the h.s.o.p. of the actual canonical algebra.
The parameter argument below alone does not assert that every positive
rational expression for its Hilbert series has a large denominator
weight. P-F31 later proves that stronger conclusion using the full
support theorem and an actual Hilbert-series pole. Neither argument
gives a superpolynomial positivity or saturation index.
\end{proposition}

\begin{proof}
Write $V=\mathbb C^n$ and identify
\[
 R=\mathbb C[W^*]^{SL_n}=(\operatorname{Sym}W)^{SL_n}.
 \tag{P-F19f}
\]
Thus points of $W^*=\operatorname{Sym}^3V^*$ are cubic forms in
the linear coordinate functions $x_1,\ldots,x_n$ on $V$.
This choice avoids identifying a representation with its dual.
The usual action on such a form is $(a\cdot F)(x)=F(a^{-1}x)$.
The ring $R$ has its ordinary coefficient grading, denoted $R_D$.

First we check the canonical grading and multiplication in (P-F19b).
If $R_D\ne0$, a scalar $\zeta I_n\in SL_n$ with $\zeta^n=1$
acts on $\operatorname{Sym}^D W$ by $\zeta^{3D}$, whence
$n\mid3D$. This is equivalent to $D=d_0s$ for some integer
$s\ge0$. On $R_{d_0s}$ the entire $GL_n$ acts by the character
$(\det)^{3s/g}$. Indeed, write $a=zh$ with $h\in SL_n$ and
$z^n=\det a$. Since $h$ fixes this subspace and $zI_n$ acts by
$z^{3d_0s}$, its action is $(\det a)^{3s/g}$, independently of
the choice of the $n$th root $z$.

Choose a nonzero volume vector $\varepsilon\in\bigwedge^nV$.
Evaluation of an intertwiner at $\varepsilon^{\otimes3s/g}$
identifies the degree-$s$ space in (P-F19b) with $R_{d_0s}$.
The Cartan product of these one-dimensional determinant modules
sends the compatible volume vectors to their tensor product.
Evaluation consequently respects multiplication in
$\operatorname{Sym}W$. We obtain an isomorphism of graded
algebras
\[
 T^{(n)}\simeq\bigoplus_{s\ge0}R_{d_0s},
 \qquad \deg_{T^{(n)}} R_{d_0s}=s.
 \tag{P-F19g}
\]
All nonzero ordinary degrees of $R$ occur in this display, so this
is the whole invariant ring with its grading divided by $d_0$,
not a proper Veronese subring. The canonical ring and its
multiplication are those in GCT VI, Section 4.1.2, equation (4.5),
printed p.~69, and equations (4.2)--(4.3), printed p.~67.
For the canonical modular index we use the detailed $T$-ring
definition of Section 4.1.1, printed p.~68, as clarified in P-F18.

Consider the explicit point
\[
 F_n=x_1^2x_2+x_2^2x_3+\cdots+x_{n-1}^2x_n+x_n^3
       \in W^*.
 \tag{P-F19h}
\]
Its projective hypersurface is smooth. To see this directly, its
partial derivatives are
\[
 \partial_{x_1}F_n=2x_1x_2,\qquad
 \partial_{x_i}F_n=x_{i-1}^2+2x_ix_{i+1}\ (1<i<n),\qquad
 \partial_{x_n}F_n=x_{n-1}^2+3x_n^2.
\]
Suppose they all vanish at a nonzero vector, and let $j$ be the
smallest index with $x_j\ne0$. If $j<n$, the $j$th derivative
gives $x_{j+1}=0$, and the $(j+1)$st derivative is then $x_j^2$,
which is nonzero, a contradiction. If $j=n$, the last derivative
is $3x_n^2\ne0$, also a contradiction. Thus the partials have
no common nonzero zero, proving smoothness over $\mathbb C$.

Let $\Delta$ be the discriminant of cubic forms. It is an
$SL_n$-invariant homogeneous polynomial of coefficient degree
$n2^{n-1}$, and it is nonzero precisely on smooth projective
hypersurfaces. An exact primary locator for these properties is
McCrory--Shifrin--Varley, \emph{Siegel modular forms generated by
invariants of cubic hypersurfaces}, arXiv:alg-geom/9403005v3,
Section 3.1.2(i), printed p.~15: the discriminant is the resultant
of the partial derivatives, with degree $n(k-1)^{n-1}$ for forms
of degree $k$.
For a direct resultant treatment see also Bus\'e--Jouanolou,
\emph{A computational approach to the discriminant of homogeneous
polynomials}, arXiv:1210.4697v1, Definition 4.6 and
Proposition 4.7(i), printed p.~40, and Proposition 4.13,
printed pp.~42--43, for the coordinate-change character.
Hence
\[
 \Delta(F_n)\ne0,\qquad
 \Delta\in R_{n2^{n-1}}
      =T^{(n)}_{g2^{n-1}}.
 \tag{P-F19i}
\]
In particular $F_n$ is outside the null cone. No claim about
stability beyond this nonvanishing is needed.

We next prove the decisive divisibility statement:
\[
 P\in T^{(n)}_s,\quad P(F_n)\ne0
       \quad\Longrightarrow\quad 2^{n-1}\mid s.
 \tag{P-F19j}
\]
We give two independent calculations.

\emph{Coefficient-weight calculation.}
Write $D=d_0s$ for the ordinary coefficient degree. A polynomial
invariant is a sum of coefficient monomials of torus weight zero:
the diagonal torus acts diagonally on the monomial basis, so a
nonzero weight cannot occur in an invariant polynomial.
If $P(F_n)\ne0$, at least one monomial occurring in $P$ uses
only the $n$ coefficients displayed in (P-F19h), because every
other coefficient is zero at $F_n$. Let its nonnegative exponent
tuple be $(a_1,\ldots,a_n)$; then $\sum_i a_i=D$.
Triviality of its character on the determinant-one diagonal torus
means its exponent vector in $x_1,\ldots,x_n$ is constant.
Thus, for one integer $c$,
\[
 2a_1=c,\qquad
 a_{i-1}+2a_i=c\quad(2\le i<n),\qquad
 a_{n-1}+3a_n=c.
 \tag{P-F19k}
\]
Adding these equations gives $nc=3D$, so $c=3s/g$.
The first $n-1$ equations solve to
\[
 a_i=\frac{c\bigl(2^i-(-1)^i\bigr)}{3\,2^i}
       \qquad(1\le i<n),
 \tag{P-F19l}
\]
as is checked by substitution and the first equation.
The numerator factor $2^i-(-1)^i$ is odd. Since $a_{n-1}$
is an integer, (P-F19l) forces $2^{n-1}\mid c$.
But $c=(3/g)s$ and $3/g$ is odd, so $2^{n-1}\mid s$.
This proves (P-F19j). Possible cancellation among the surviving
monomials causes no problem: nonzero evaluation guarantees that
at least one such monomial occurs, and every one obeys the same
divisibility condition.

\emph{Diagonal-stabilizer calculation.}
Let $\zeta$ be a primitive root of unity of order $N=3\,2^{n-1}$
and define
\[
 A=\operatorname{diag}
  (\zeta,\zeta^{-2},\zeta^{(-2)^2},\ldots,
                         \zeta^{(-2)^{n-1}})\in GL_n.
\]
Each chain monomial in (P-F19h) is fixed by $A$, since
$2(-2)^{i-1}+(-2)^i=0$; the last monomial is fixed since
$3(-2)^{n-1}$ is a multiple of $N$. The same is true with
$A^{-1}$, so $A\cdot F_n=F_n$ for the specified action on forms.
Set
\[
 S_n=1-2+(-2)^2+\cdots+(-2)^{n-1}
     =\frac{1-(-2)^n}{3}.
\]
The integer $S_n$ is odd and satisfies $S_n\equiv n\pmod3$.
Consequently
\[
 \operatorname{ord}(\det A)
 =\frac{N}{\gcd(N,S_n)}
 =\frac{3\,2^{n-1}}{g}.
 \tag{P-F19m}
\]
An $SL_n$-invariant polynomial $P$ of coefficient degree
$D=d_0s$ satisfies, as a function on $W^*$,
\[
 P(a\cdot F)=(\det a)^{-3s/g}P(F).
 \tag{P-F19n}
\]
This follows by writing $a=zh$ as above: $h$ does not change
$P$, while a scalar sends $F$ to $z^{-3}F$.
If $P(F_n)\ne0$, evaluation at its stabilizer $A$ therefore
requires $(\det A)^{3s/g}=1$. By (P-F19m),
$(3\,2^{n-1}/g)\mid(3s/g)$, which is again equivalent to
$2^{n-1}\mid s$. This verifies (P-F19j) independently of
the coefficient recurrence.

Finally, let $q_1,\ldots,q_r$ be any homogeneous system of
parameters of $T^{(n)}$. The algebra is finite as a module over
$\mathbb C[q_1,\ldots,q_r]$, so
$T^{(n)}/(q_1,\ldots,q_r)$ is finite-dimensional and graded.
Every positive-degree homogeneous element in this quotient is
nilpotent: its powers eventually have degree beyond the largest
nonzero graded piece. In particular
$\Delta^M\in(q_1,\ldots,q_r)$ for some integer $M\ge1$.
If every $q_i(F_n)$ were zero, evaluation of this ideal membership
would give $\Delta(F_n)^M=0$, contradicting (P-F19i).
Thus some $q_i$ is nonzero at $F_n$. Its positive stretch degree
is divisible by $2^{n-1}$ by (P-F19j), and is at least that
large. This applies to every h.s.o.p., including one with a
lexicographically minimal degree sequence, proving (P-F19c).

The inner module is the single irreducible $S^{(3)}V$, and
the partitions in (P-F19a) have heights one and $n$.
Thus (P-F19d) follows directly from the source's stated
parameter. Exponential growth of $2^{n-1}$ exceeds every
polynomial in $h=O(n^6)$, proving the asserted failure of the
canonical modular-index clause. The same family belongs to the
maximum defining $\Xi(m)$, yielding (P-F19e).
\end{proof}

\paragraph{Exact scope of the refutation.}
GCT VI, Hypothesis 1.6.7 first asks for some positive rational
form of polynomial modular index, and then explicitly asks for
this bound on the minimal form arising from a lexicographically
minimal h.s.o.p. of its canonical ring. P-F19 refutes the latter,
stronger clause. Its proof concerns values of actual invariants
at a particular semistable point. A positive rational expression
for the Hilbert series need not arise from parameters of that
ring, so the first existence clause is not refuted by this
argument alone. P-F30--P-F31 subsequently close that gate by proving
an actual primitive-$2^{n-1}$ pole of the ordinary Hilbert series,
and thereby refute the first clause as well. Neither result shows
that the positivity or saturation index of this ray is superpolynomial. P-F12's implication from
bounded positive-form denominator weights to PH2 remains valid.

This is an irreducible-inner ordinary plethysm family, unlike
the reducible torus embeddings in P-F14--P-F15. Determinant
twisting is absent: the inner shape is always $(3)$, and only
the number of variables grows. The qualitative finite bound
of P-F18 remains true, while its optimal canonical envelope
has the exponential subsequence (P-F19e). No historical
priority is asserted for the chain forms, their stabilizers,
or the resulting invariant-degree obstruction. The same chain
forms, after reversing the variables, already appear in
Bus\'e--Jouanolou, Example 4.12, printed p.~42, with attribution
to Demazure; the statement here gives their explicit consequence
for the specified canonical plethysm parameter bound.


\subsection{What the chain point forces in universal denominators}

\begin{proposition}[P-F20: an exponential universal factor and an actual covariant pole]
\label{prop:chainuniversaldenominator}
Use the notation of P-F19 and put $M=2^{n-1}$.
Let $T=T^{(n)}$ have the stretch grading, and let
$R=(\operatorname{Sym}W)^{SL_n}$ have its ordinary coefficient grading.
Write $\Phi_a(t)$ for the cyclotomic polynomial of order $a$.
Then
\[
 \Phi_M(t)\mid\operatorname{udenom}(T,t),\qquad
 \Phi_{d_0M}(t)\mid\operatorname{udenom}(R,t).
 \tag{P-F20a}
\]
Here the universal denominator is the least common multiple of the
reduced Hilbert denominators of graded ideals, equivalently the
universal denominator in Derksen's Definition 1.5.

There is an irreducible $SL_n$-module $U$ such that the actual Hilbert
series of the covariant module
\[
 \mathcal M_U=(\operatorname{Sym}W\otimes U)^{SL_n}
 \tag{P-F20b}
\]
has a pole at a primitive $(d_0M)$th root of unity.
Consequently every eventual quasi-polynomial period of its coefficient
sequence is divisible by $d_0M$.
This assertion concerns a fixed-target covariant sequence. This
argument alone does not prove such a pole for $\operatorname{Hilb}_T$,
nor for a simultaneously stretched plethysm ray with a nontrivial
target weight. The ordinary pole for this particular cubic invariant
ring is proved separately in P-F31 using additional geometry.
\end{proposition}

\begin{proof}
Let $I^{[M]}\subset T$ be the homogeneous ideal generated by all
$T_s$ with $M\nmid s$. P-F19 proved that every such element
vanishes at the chain point $F_n$. Hence every element of
$I^{[M]}$ vanishes there. The discriminant $\Delta$ satisfies
$\Delta(F_n)\ne0$ and has stretch degree $gM$. Its image in
\[
 Q=T/I^{[M]}
\]
therefore has no zero positive power: otherwise some $\Delta^j$
would lie in $I^{[M]}$ and evaluation would contradict
$\Delta(F_n)^j\ne0$. In particular $Q$ is infinite-dimensional.
It is finitely generated and positively graded, so $\dim Q\ge1$.
Indeed a zero-dimensional finitely generated algebra over
$\mathbb C$ is Artinian and finite-dimensional.

By construction $Q_s=0$ unless $M\mid s$. Rename its grading by
putting $\widetilde Q_j=Q_{Mj}$, so that
\[
 \operatorname{Hilb}_Q(t)
   =\operatorname{Hilb}_{\widetilde Q}(t^M).
 \tag{P-F20c}
\]
The surviving powers of the discriminant occupy distinct degrees
$gj$ of $\widetilde Q$. Thus, for real $0<t<1$,
\[
 \operatorname{Hilb}_{\widetilde Q}(t)
    \ge\sum_{j\ge0}t^{gj}=\frac1{1-t^g}.
\]
The Hilbert series is rational because the algebra is finitely
generated. The displayed divergence as $t$ tends to one forces
a pole there. For any primitive $M$th root $\zeta$, the map
$t\mapsto t^M$ has nonzero derivative $M\zeta^{M-1}$ at
$\zeta$. Equation (P-F20c) consequently has a pole at every
such $\zeta$.

The exact sequence
\[
 0\longrightarrow I^{[M]}\longrightarrow T
     \longrightarrow Q\longrightarrow0
\]
gives
$\operatorname{Hilb}_{I^{[M]}}=
 \operatorname{Hilb}_T-\operatorname{Hilb}_Q$.
At any primitive $M$th root, either $\operatorname{Hilb}_T$
already has a pole, or $\operatorname{Hilb}_{I^{[M]}}$ does.
Therefore $\Phi_M$ divides the least common multiple of
Hilbert denominators of graded ideals of $T$. This proves the
first assertion in (P-F20a) directly from the definition, without
equating a universal denominator to the denominator of the ring.
The same conclusion also follows from Derksen, \emph{Universal
denominators of Hilbert series}, Theorem 1.10(a), printed
pp.~589--590: the exponent of $\Phi_M$ is
$\dim(T/I^{[M]})$.

In the ordinary grading, P-F19 says that a homogeneous invariant
nonzero at $F_n$ has degree divisible by $d_0M$. Repeating the
same ideal, quotient, and discriminant argument in $R$ proves
the second assertion in (P-F20a). Notice that the discriminant's
ordinary degree is $nM=g\,d_0M$, so it again supplies a surviving
positive-degree element.

Now apply Derksen, Theorem 3.6, printed p.~594, to the polynomial
ring $\operatorname{Sym}W$ and the linearly reductive complex group
$SL_n$. That theorem identifies the universal denominator of
its invariant ring with the least common multiple of the
\emph{ordinary} Hilbert denominators of the covariant modules
$(\operatorname{Sym}W\otimes U)^{SL_n}$, over irreducible $U$.
Since $\Phi_{d_0M}$ divides that least common multiple, it divides
the reduced denominator of at least one such module. This is
exactly the pole asserted in (P-F20b).

Finally a sequence with eventual quasi-polynomial period $L$
has a generating function whose poles are $L$th roots of unity:
split its eventual coefficients into residues modulo $L$ and
sum the resulting polynomial sequences. Finite initial terms
add only a polynomial. A primitive $(d_0M)$th pole therefore
forces $d_0M\mid L$. Covariant modules are finitely generated
over $R$, so their Hilbert series are rational with poles at
roots of unity, as follows from a homogeneous finite generating
presentation. Their coefficient sequences are indeed eventually
quasi-polynomial. This completes the proof.
\end{proof}

\paragraph{Why a universal factor alone does not prove an ordinary pole.}
Derksen's Example 4.7, printed pp.~597--598, supplies an exact
counterexample to that general inference. For the permutation action
of $A_{10}$ on ten variables,
\[
 \operatorname{Hilb}_{\mathbb C[x_1,\ldots,x_{10}]^{A_{10}}}(t)
   =\frac{1+t^{45}}{\prod_{i=1}^{10}(1-t^i)}.
\]
The numerator cancels its $\Phi_6$ factor, while the invariant
ring has a homogeneous quotient isomorphic to $\mathbb C[e_6]$
with $\deg e_6=6$. Its universal denominator therefore contains
$\Phi_6$, although its own reduced Hilbert denominator does not.
This is a finite-group invariant domain in characteristic zero,
so normality and rational singularities do not prevent the
cancellation. The primary source is
\url{https://sites.lsa.umich.edu/hderksen/wp-content/uploads/sites/614/2018/05/A.I.a.19.pdf}.

Even preservation of a prime-power factor fails for normal rational
domains. For an integer $M\ge2$, consider
\[
 B_M=\mathbb C[x,y,z,w]/(xz+y^M),\qquad
 \deg(x,y,z,w)=(1,1,M-1,M).
 \tag{P-F20d}
\]
Its underlying surface factor is the cyclic quotient
$\mathbb A^2/\boldsymbol\mu_M$: the invariant generators
$u^M,uv,v^M$ satisfy $u^Mv^M=(uv)^M$, and a sign change gives
the displayed relation. Thus it is normal and has rational
singularities; the latter is also a direct instance of the
reductive-quotient theorem used in GCT VI, Section 4.1. Its
product with the $w$-line retains these properties. The ring is
Cohen--Macaulay as a hypersurface. Its Hilbert series is
\[
 \operatorname{Hilb}_{B_M}(t)
  =\frac{1-t^M}
     {(1-t)^2(1-t^{M-1})(1-t^M)}
  =\frac1{(1-t)^2(1-t^{M-1})},
\]
whereas $B_M/(x,y,z)=\mathbb C[w]$ has Hilbert series
$1/(1-t^M)$. If $M$ is a power of two, the ring's series has
odd period $M-1$, while the quotient has period $M$.
The equality defining the cyclic invariant ring can be checked
without a presentation theorem: an invariant monomial $u^av^b$
has $a-b\equiv0\pmod M$ and factors into the three displayed
generators; reduction by their one relation leaves distinct
monomials in $u,v$, proving injectivity. This example does not
settle a bound comparing the \emph{sizes} of possible pole orders,
since its other pole order $M-1$ is still large.

\paragraph{Factoriality and a connected semisimple group do not restore exact transfer.}
There is a counterexample even with a prime-power factor, a standard
polynomial ambient grading, and a connected semisimple invariant group.
Let
\[
 G=SL_2\times SL_4,\qquad
 V=\operatorname{Mat}_{2\times4}\oplus
       \operatorname{Mat}_{4\times4},\qquad
 C=\mathbb C[V]^G,
\]
where the two factors act by left multiplication on their respective
matrix spaces. Write $p_{ij}$ for the six two-by-two minors of the
first matrix and $d$ for the determinant of the second. Then
\[
 C\cong
 \frac{\mathbb C[p_{12},p_{13},p_{14},p_{23},p_{24},p_{34},d]}
 {\bigl(p_{12}p_{34}-p_{13}p_{24}+p_{14}p_{23}\bigr)},
 \qquad \deg p_{ij}=2,\quad\deg d=4.
 \tag{P-F20f}
\]
For generators, the first fundamental theorem for $SL_n$ gives the
maximal minors: Kraft--Procesi, \emph{Classical invariant theory:
A primer}, Section~8.4, Theorem (FFT for $SL_n$), printed p.~90;
their Section~1.2, Example~2, pp.~3--4, also gives the determinant
case directly. Invariants of the separate actions tensor together:
expand in a basis of one polynomial factor and then take invariants
successively. For completeness, the sole relation in (P-F20f) can
be checked without the second fundamental theorem. The displayed
quadratic relation follows by expanding the minors. It is irreducible,
since a reducible homogeneous quadratic is a product of two linear
forms and has rank at most two, whereas this one has rank six. On
the open set $p_{12}\ne0$, every tuple satisfying it is realized
by the columns
\[
 \binom10,\quad \binom0a,\quad
 \binom{-p_{23}/a}{p_{13}},\quad
 \binom{-p_{24}/a}{p_{14}},\qquad a=p_{12}.
\]
Thus this open set is a dense subset of the irreducible quadratic
hypersurface and there are no further polynomial relations. The
determinant $d$ is independently arbitrary, proving (P-F20f).

The Hilbert series and quotient are consequently
\[
 \operatorname{Hilb}_C(t)
   =\frac{1-t^4}{(1-t^2)^6(1-t^4)}
   =\frac1{(1-t^2)^6},\qquad
 C/(p_{ij}:i<j)=\mathbb C[d],\quad\deg d=4.
 \tag{P-F20g}
\]
The same exact-sequence argument as above forces $\Phi_4$ into
the universal denominator, although the actual Hilbert series is
regular at both primitive fourth roots.

We verify the stronger geometric conditions explicitly. The ring
$C$ is factorial. If a nonzero invariant polynomial is factored
into irreducibles in $\mathbb C[V]$, the connected group $G$
cannot permute its finite set of irreducible-factor lines
nontrivially. Each factor is therefore a semi-invariant, with
an algebraic character of $G$. Every character of $SL_2\times SL_4$
is trivial: the factors are generated by elementary unipotent
one-parameter subgroups, and an algebraic homomorphism from
$\mathbb G_a$ to $\mathbb G_m$ is constant because the only units
in $\mathbb C[t]$ are constants. Every irreducible factor is thus
itself invariant. Existence and uniqueness of factorization in
$C$ follow from those in $\mathbb C[V]$, and $C$ is normal.
As a hypersurface, it is Gorenstein and Cohen--Macaulay; see, for
example, the local-complete-intersection statement in the Stacks
Project, Tag~\texttt{0DWA}.

Rational singularities can also be seen directly in this example.
The quadratic factor is the affine cone over a smooth quadric
$Q^4\subset\mathbb P^5$. Blowing up its vertex gives the smooth
total space of $\mathcal O_{Q^4}(-1)$, with the zero section
contracted to the vertex. For every $k\ge0$, the sequence
\[
 0\longrightarrow\mathcal O_{\mathbb P^5}(k-2)
 \longrightarrow\mathcal O_{\mathbb P^5}(k)
 \longrightarrow\mathcal O_{Q^4}(k)\longrightarrow0
\]
and the projective-space cohomology formula (Stacks Project,
Tag~\texttt{01XT}) give
$H^i(Q^4,\mathcal O_{Q^4}(k))=0$ for every $i>0$.
The projection of this total space to $Q^4$ is affine, with
pushforward $\bigoplus_{k\ge0}\mathcal O_{Q^4}(k)$, so all its
higher structure-sheaf cohomology vanishes. The target cone is
affine; therefore the higher direct images under the resolution
vanish, and its direct image in degree zero is the structure
sheaf by normality. Taking the product of this resolution with
the $d$-line gives the same conclusion for $\operatorname{Spec}C$.

This counterexample rules out the proposed exact transfer even
within factorial Gorenstein rational invariant rings of connected
semisimple groups. It is a representation with several summands
and is not the irreducible cubic representation in P-F19. It does
not rule out a quantitative comparison of maximal denominator
weights; here the positive-form weight is two and the hidden
universal order is only four. The primary first-fundamental-theorem
source is
\url{https://dmi.unibas.ch/fileadmin/user_upload/dmi/Personen/Kraft_Hanspeter/Classical_Invariant_Theory.pdf}.

\paragraph{A precise obstruction to the proposed affine-slice argument.}
Large periods in affine slices do not, from multivariate
quasi-polynomiality alone, force large periods on rays through the
origin. Fix $M=2^a$ with $a\ge2$, and on
$\mathbb Z_{\ge0}^2$ define
\[
 F_M(x,y)=1+\mathbf1_{\{y\ {\rm odd},\ M\mid x\}}.
 \tag{P-F20e}
\]
This is a nonnegative multivariate quasi-polynomial, with
$F_M(0,0)=1$, and it has the rational generating series
\[
 \sum_{x,y\ge0}F_M(x,y)X^xY^y
  =\frac1{(1-X)(1-Y)}+
       \frac{Y}{(1-X^M)(1-Y^2)}.
\]
For every integral direction $(u,v)\in\mathbb Z_{\ge0}^2$,
the function $k\mapsto F_M(ku,kv)$ is either identically one
or is $1+\mathbf1_{\{k\ {\rm odd}\}}$, so its global period
divides two. Indeed, even $k$ makes $kv$ even, and odd $k$ is
invertible modulo $M$; hence for odd $k$ the extra term occurs
exactly when $v$ is odd and $M\mid u$.
But the fixed-target slice is
\[
 F_M(x,1)=1+\mathbf1_{\{M\mid x\}},
\]
whose generating series has a pole at a primitive $M$th root.
Thus cancellation among multivariate Fourier characters can hide
arbitrarily large powers of two on \emph{every} simultaneous ray.

This last example is not asserted to be the Hilbert function of
a branching domain. It pinpoints why a Fourier argument using
only a finite chamber decomposition and global
quasi-polynomiality on each ray cannot finish the transfer.
A successful argument must use additional structure of the actual
branching algebra. P-F20 leaves the weaker positive-form existence
clause and PH2 undecided; its new conclusions are the universal
factor and the actual, but fixed-target, covariant pole.


\subsection{Complete period collapse can hide arbitrary parameter degrees}

\begin{proposition}[P-F21: no quantitative transfer from parameters to positive forms in general]
\label{prop:toricperiodcollapse}
For every integer $D\ge2$, there is a positively graded complex domain
$R_D$ of Krull dimension three with the following properties:
\begin{enumerate}
\item $R_D$ is normal, Cohen--Macaulay, and has rational singularities.
\item Its global Hilbert function and Hilbert series are
\[
 f_D(s)=\frac{D-1}{2}s^2+\frac{D+1}{2}s+1\quad(s\ge0),
 \qquad
 \operatorname{Hilb}_{R_D}(t)=\frac{1+(D-2)t}{(1-t)^3}.
 \tag{P-F21a}
\]
Thus it has a positive numerator and three denominator weights equal to
one, and its positivity and saturation indices are both zero.
\item Every homogeneous system of parameters has at least one degree
divisible by $D$. There is a system of parameter degrees $(1,1,D)$.
Consequently the maximal degree in the lexicographically minimal
nondecreasing parameter-degree tuple is exactly $D$.
\item It has a homogeneous quotient $\mathbb C[z]$ with $\deg z=D$.
In particular its universal Hilbert denominator contains $\Phi_D$,
although its own reduced Hilbert denominator has no nontrivial root
of unity as a zero.
\end{enumerate}
This statement concerns general graded domains; no identification of
$R_D$ with an ordinary irreducible-inner plethysm algebra is asserted.
\end{proposition}

\begin{proof}
Consider the rational triangle
\[
 P_D=\operatorname{conv}\left\{(0,0),(D,0),
           \left(0,\frac{D-1}{D}\right)\right\}
\]
and its integral cone semigroup
\[
 S_D=\{(x,y,s)\in\mathbb Z^3:s\ge0,\ (x,y)\in sP_D\},
 \qquad R_D=\mathbb C[S_D].
\]
The monomial corresponding to $v=(x,y,s)$ is denoted $X^v$ and
has degree $s$. The degree-zero part is $\mathbb C$, and every graded
part is finite-dimensional. The cone has linearly independent
primitive rays
\[
 v_0=(0,0,1),\qquad v_1=(D,0,1),\qquad v_2=(0,D-1,D).
 \tag{P-F21b}
\]
Let $\Pi_D=\{\lambda_0v_0+\lambda_1v_1+\lambda_2v_2:
0\le\lambda_i<1\}$ be the half-open fundamental parallelepiped.
Every element of $S_D$ has a unique decomposition
\[
 v=p+k_0v_0+k_1v_1+k_2v_2,
 \qquad p\in\Pi_D\cap\mathbb Z^3,\quad k_i\in\mathbb Z_{\ge0}:
\]
write its unique nonnegative real ray coordinates and subtract their
integer parts. Conversely all such sums belong to $S_D$.
Therefore $R_D$ is a finite free module, with basis $X^p$, over the
polynomial subring
\[
 A_D=\mathbb C[X^{v_0},X^{v_1},X^{v_2}].
 \tag{P-F21c}
\]
The three displayed monomials are algebraically independent because
the rays are linearly independent. This proves finite generation,
Krull dimension three, and that they form a homogeneous system of
parameters of degrees $(1,1,D)$. It also proves Cohen--Macaulayness:
finite freeness over the polynomial subring gives flat localizations
with zero-dimensional fibers, so the localized rings are
Cohen--Macaulay; in particular the displayed parameter sequence
is regular.

We compute the free-module numerator explicitly. For a lattice point
$p\in\Pi_D$, its first two coordinates satisfy
\[
 x\in\{0,\ldots,D-1\},\quad
 y\in\{0,\ldots,D-2\},\quad
 \lambda_1=x/D,\quad\lambda_2=y/(D-1).
\]
There is a unique $\lambda_0\in[0,1)$ making the third coordinate
integral, and that coordinate is
\[
 s=y+\left\lceil\frac{x}{D}+\frac{y}{D-1}\right\rceil.
\]
For $y=0$, the point $x=0$ has degree zero and the other $D-1$
points have degree one. For $1\le y\le D-2$, the $D-y$ choices
$0\le x\le D-y-1$ have degree $y+1$, while the remaining $y$
choices have degree $y+2$. Indeed $x/D+y/(D-1)\le1$ is equivalent
to $x\le D-y-y/(D-1)$, whose integer upper bound is $D-y-1$.
It follows, including the empty-sum case $D=2$, that
\begin{align*}
 \sum_{p\in\Pi_D\cap\mathbb Z^3}t^{\deg p}
 &=1+(D-1)t+
       \sum_{y=1}^{D-2}\big((D-y)t^{y+1}+yt^{y+2}\big)\\
 &=1+(D-1)(t+\cdots+t^{D-1})+(D-2)t^D\\
 &=(1+(D-2)t)(1+t+\cdots+t^{D-1}).
\end{align*}
Dividing by $(1-t)^2(1-t^D)$, the Hilbert series of $A_D$, proves
the series in (P-F21a). Taking coefficients yields
\[
 f_D(s)=\binom{s+2}{2}+(D-2)\binom{s+1}{2},
\]
which is the stated polynomial for every $s\ge0$. All its ordinary
polynomial coefficients are positive. It is positive at every
nonnegative integer, so both stated indices are zero under the
definitions used above.

For completeness, normality can be checked without a presentation
of $R_D$. The cone is cut out by the integral linear inequalities
\[
 x\ge0,\qquad y\ge0,\qquad
 D(D-1)s-(D-1)x-D^2y\ge0.
\]
For each of these linear forms $\ell$, divide out the gcd of its
coefficients and form
$B_\ell=\mathbb C[\{v\in\mathbb Z^3:\ell(v)\ge0\}]$.
After an integral basis change, $B_\ell$ is a polynomial ring in
one variable over a Laurent polynomial ring in two variables, so
it is integrally closed. Inside the common Laurent polynomial
ring one has
$R_D=\bigcap_\ell B_\ell$, by comparing the supports of Laurent
polynomials. If an element of $\operatorname{Frac}(R_D)$ is integral
over $R_D$, it is integral over every $B_\ell$; since it lies in
the common rational function field, it belongs to each $B_\ell$
and hence to $R_D$. Thus $R_D$ is normal. It is an affine normal
toric ring, so it has rational singularities: Danilov,
\emph{The geometry of toric varieties}, Section~8.5,
Proposition~8.5.1, printed pp.~124--125, proves that a toric
subdivision $\rho:X'\to X$ has
$\rho_*\mathcal O_{X'}=\mathcal O_X$ and
$R^i\rho_*\mathcal O_{X'}=0$ for $i>0$; applying his toric
resolution in Section~8.2 gives the claim. The primary scan is
\url{https://www.math.purdue.edu/~arapura/danilov.pdf}.

Let $J_D$ be the span of all cone monomials not on the extreme ray
$\mathbb R_{\ge0}v_2$. This is an ideal: the nonnegative linear
form $\ell(x,y,s)=(D-1)s-Dy$ vanishes precisely on that ray,
so adding a cone element to an element off the ray stays off it.
Since $v_2$ is primitive, $S_D\cap\mathbb R_{\ge0}v_2
=\mathbb Z_{\ge0}v_2$. Thus
\[
 R_D/J_D\cong\mathbb C[z],\qquad z=X^{v_2},\quad\deg z=D.
 \tag{P-F21d}
\]
For any homogeneous parameter sequence $q_1,q_2,q_3$, the graded
ring $R_D/(q_1,q_2,q_3)$ is finite-dimensional. Some power of
$z$ therefore lies in $(q_1,q_2,q_3)$. In the quotient (P-F21d),
this is impossible if every $q_i$ maps to zero. Some parameter
has nonzero image in $\mathbb C[z]$, and its positive degree is
therefore a multiple of $D$. The tuple $(1,1,D)$ already
constructed is lexicographically minimal among nondecreasing
parameter-degree tuples: all positive degrees are at least one,
and every tuple with first two entries equal to one has last entry
at least $D$. This proves the exact canonical maximum.

Finally the quotient (P-F21d) has Hilbert series $1/(1-t^D)$.
The exact sequence $0\to J_D\to R_D\to\mathbb C[z]\to0$
shows that the least common multiple of the Hilbert denominators
of graded ideals of $R_D$ contains $\Phi_D$. By (P-F21a), the
ring's own reduced denominator is $(1-t)^3$: its numerator
evaluates to $D-1\ne0$ at one. This completes the proof.
\end{proof}

\paragraph{Prior example and precise consequence.}
The period-collapse triangle is classical. McAllister--Woods,
\emph{The minimum period of the Ehrhart quasi-polynomial of a
rational polytope}, arXiv \texttt{math/0310255}, Theorem~2.2 on
printed p.~3 and its proof on pp.~4--5, use
$\operatorname{conv}\{(0,0),(D,0),(1,(D-1)/D)\}$ and obtain
the polynomial in (P-F21a). The integral affine map
$(x,y)\mapsto(D-x-Dy,y)$ carries $P_D$ to that triangle;
on the $s$th dilation use $(x,y)\mapsto(Ds-x-Dy,y)$, which
is a lattice bijection. The source is
\url{https://arxiv.org/pdf/math/0310255}.
The algebraic consequences above are proved here without a claim
of historical priority.

Choosing $D=2^a$ makes the hidden universal factor a prime-power
factor of arbitrarily large order. The Krull dimension remains
three, the positive-form weights remain one, the numerator has
degree at most one and coefficient bit length $O(\log D)$, and
both indices remain zero. Hence normality, Cohen--Macaulayness,
rational singularities, positive Hilbert numerators, and global
quasi-polynomiality together do not give even a qualitative bound
on canonical parameter degrees in terms of the dimension and
positive-form weights. This rules out such a general transfer
as a way of strengthening P-F19. An argument for the weaker
ordinary plethysm PH3 must use additional structure of the actual
plethysm algebra.


\subsection{A fixed Hilbert series with unbounded canonical parameter degrees}

\begin{proposition}[P-F22: the Gorenstein condition does not repair the transfer]
\label{prop:gorensteinperiodcollapse}
Let $q_0=1$, $q_1=2$, and $q_{j+1}=3q_j-q_{j-1}$ for $j\ge1$.
For $j\ge1$ put $b=q_j$, $c=q_{j+1}$ and
\[
 Q_{b,c}=\{(x,y)\in\mathbb R^2:x\ge-1,\ y\ge-1,
                       b^2x+c^2y\le1\}.
\]
Its Ehrhart ring
\[
 A_{b,c}=\mathbb C[\{(x,y,s)\in\mathbb Z^3:
                         s\ge0,\ (x,y)\in sQ_{b,c}\}]
\]
is a positively graded normal Cohen--Macaulay Gorenstein domain
of Krull dimension three with rational singularities. For every
pair in this sequence it has the same Hilbert function and series,
\[
 \dim(A_{b,c})_s=\binom{3s+2}{2}
       =\frac92s^2+\frac92s+1\quad(s\ge0),
 \qquad
 \operatorname{Hilb}_{A_{b,c}}(t)=\frac{1+7t+t^2}{(1-t)^3}.
 \tag{P-F22a}
\]
Its positivity and saturation indices are zero. Nevertheless its
lexicographically minimal nondecreasing homogeneous parameter-degree
tuple is exactly $(1,b,c)$. In particular its canonical maximal
parameter degree is $c$, which tends to infinity with $j$.
No member of this family is identified here with an ordinary
irreducible-inner plethysm multiplicity algebra, and factoriality
is not asserted.
\end{proposition}

\begin{proof}
The recurrence preserves the identity
\[
 1+b^2+c^2=3bc.
 \tag{P-F22b}
\]
Indeed it holds for $(b,c)=(1,2)$, and substitution of
$(c,3c-b)$ into the left minus right side leaves
$1+b^2+c^2-3bc$. Consecutive terms are relatively prime because
$\gcd(c,3c-b)=\gcd(c,b)$, and neither is divisible by three
because reduction of the recurrence modulo three sends
$(b,c)$ to $(c,-b)$. For our pairs $b\ge2$ and
\[
 2b<c<3b.
 \tag{P-F22c}
\]
This holds for $(2,5)$; from $0<b<c$ the next ratio is
$(3c-b)/c=3-b/c$, and induction preserves the displayed bounds.
In particular the sequence is unbounded.

We first count the smaller rational triangle
\[
 P_{b,c}=\operatorname{conv}\{(0,0),(c/b,0),(0,b/c)\}.
\]
We use the following precise published period-collapse criterion.
For $u=q/p$, $v=s/r$ in lowest terms,
Cristofaro-Gardiner--Li--Stanley, \emph{New examples of period
collapse}, Corollary~1.5, printed p.~4 of arXiv
\texttt{1509.01887v1}, states that the triangle with intercepts
$1/u,1/v$ has polynomial Ehrhart function if
\[
 s\mid p,\quad p\mid(rq+1),\quad
 \gcd((rq+1)/p,s)=1,
\]
and
\[
 q\mid r,\quad r\mid(sp+1),\quad
 \gcd((sp+1)/r,q)=1.
\]
This is the source's explicit rational-triangle criterion, not
its separate theorem concerning irrational triangles. Its primary
text is \url{https://arxiv.org/pdf/1509.01887}.

Here take $(p,q,r,s)=(c,b,b,c)$. The divisibilities follow from
\[
 b^2+1=c(3b-c),\qquad c^2+1=b(3c-b).
\]
The two gcd conditions are
$\gcd(3b-c,c)=1$ and $\gcd(3c-b,b)=1$, which follow from
$\gcd(b,c)=\gcd(3,bc)=1$. All hypotheses of the criterion are
therefore satisfied. The Ehrhart function of $P_{b,c}$ is a
quadratic polynomial with leading coefficient its area $1/2$
and constant coefficient one. By (P-F22c), its only lattice
points at dilation one are $(0,0),(1,0),(2,0)$. Its value there
is three, so this polynomial is
\[
 L_{P_{b,c}}(s)=\tfrac12s^2+\tfrac32s+1
                    =\binom{s+2}{2}.
 \tag{P-F22d}
\]
In obtaining the value at zero we use the usual Ehrhart
quasi-polynomial, which counts the single point $0P_{b,c}$;
the criterion says its constituents coincide.

Equation (P-F22b) gives
\[
 Q_{b,c}=3P_{b,c}-(1,1).
\]
Translation by $(-s,-s)$ is integral at every integer dilation,
so (P-F22d) yields $L_{Q_{b,c}}(s)=L_{P_{b,c}}(3s)$.
This proves the Hilbert function in (P-F22a). Summing its
quadratic polynomial, or multiplying its first three differences
by $(1-t)^3$, gives the displayed Hilbert series. Its ordinary
polynomial coefficients are positive and all nonnegative
integer values are positive, so both indices are zero.

The primitive rays of the cone over $Q_{b,c}$ are
\[
 v_0=(-1,-1,1),\quad
 v_1=(3c-b,-b,b),\quad
 v_2=(-c,3b-c,c).
 \tag{P-F22e}
\]
Primitivity of the last two follows from
$\gcd(3c,b)=\gcd(3b,c)=1$. The cone is simplicial and
full-dimensional. As in the full argument of P-F21, subtracting
integer parts of its nonnegative ray coordinates gives a unique
decomposition into a lattice point in the half-open fundamental
parallelepiped plus a nonnegative integral combination of the
three rays. Thus $A_{b,c}$ is finite free over
$\mathbb C[X^{v_0},X^{v_1},X^{v_2}]$, whose generator degrees
are $(1,b,c)$. This proves finite generation, Krull dimension
three, Cohen--Macaulayness, and existence of a parameter sequence
with those degrees. Normality follows by expressing the ring as
the intersection of the three integrally closed halfspace
semigroup rings for
\[
 \ell_1(x,y,s)=x+s,\quad
 \ell_2(x,y,s)=y+s,\quad
 \ell_3(x,y,s)=s-b^2x-c^2y.
 \tag{P-F22f}
\]
These are primitive integral facet forms. Their nonnegative
intersection is precisely the cone: its inequalities imply
$s\ge0$, since $x,y\ge-s$ and
$b^2x+c^2y\le s$ give $(1+b^2+c^2)s\ge0$; at $s=0$
they force $x=y=0$.

The Gorenstein assertion has a direct lattice check. With
$e=(0,0,1)$, each $\ell_i(e)=1$. Hence an integral point is
strictly inside the cone exactly when its translate by $-e$
is an integral point of the cone:
\[
 \operatorname{int}(C_{b,c})\cap\mathbb Z^3
       =e+(C_{b,c}\cap\mathbb Z^3).
 \tag{P-F22g}
\]
Danilov, \emph{The geometry of toric varieties}, Section~4.6,
printed pp.~111--112, identifies the canonical module of a
normal affine toric ring with the ideal spanned by its interior
monomials. Equation (P-F22g) identifies that module with the
principal ideal $X^eA_{b,c}$. Since the ring is Cohen--Macaulay,
it is Gorenstein. Its rational singularities follow from the
toric resolution and vanishing in the same source, Section~8.2
and Proposition~8.5.1, printed pp.~124--125. Both exact locators
are available in the primary scan
\url{https://www.math.purdue.edu/~arapura/danilov.pdf}.

We finally prove the canonical degree assertion; mere existence
of $(1,b,c)$ would not by itself be a lexicographic comparison.
Let $F$ be the two-dimensional cone face spanned by $v_1,v_2$,
equivalently the face $\ell_3=0$. Its degree-$s$ lattice points
correspond, after $u=x+s$, $v=y+s$, to nonnegative integers
satisfying
\[
 b^2u+c^2v=3bc\,s.
\]
Reduction modulo $b$ gives $b\mid v$, and reduction modulo $c$
gives $c\mid u$. Writing $u=cr$, $v=bt$ therefore reduces the
condition exactly to
\[
 br+ct=3s,\qquad r,t\in\mathbb Z_{\ge0}.
 \tag{P-F22h}
\]
There is no solution for $0<s<b$. If $t\ge2$, then
$br+ct\ge2c>4b>3s$. If $t=1$ and $r\ge1$, it is at least
$b+c>3b>3s$. If $t=1,r=0$, then $c=3s$, contrary to
$3\nmid c$. If $t=0$, then $br=3s$, and $\gcd(b,3)=1$
forces $b\mid s$, again impossible.

The monomial face quotient $A_{b,c}\twoheadrightarrow
\mathbb C[F\cap\mathbb Z^3]$ is a graded domain of dimension
two. Every positive-degree element of degree less than $b$ maps
to zero in it. A parameter sequence cannot have two such
elements: its remaining single element would cut a
two-dimensional finitely generated domain to a ring of
dimension at least one, or leave it unchanged, contradicting
finite-dimensionality of the quotient by all three parameters.
Here a nonzero homogeneous positive-degree element is not a
unit, and the dimension assertion is the principal ideal theorem.
Consequently the second entry of every nondecreasing
parameter-degree tuple is at least $b$.

The extreme-ray quotient at $v_2$ is $\mathbb C[z]$ with
$\deg z=c$. If a parameter tuple begins $(1,b)$, both these
parameters map to zero there because $1,b<c$. The third must
map nontrivially, so its degree is a positive multiple of $c$.
The first entry is always at least one. Together with the
existing tuple $(1,b,c)$, these facts prove that this tuple is
lexicographically minimal and its maximal degree is exactly $c$.
\end{proof}

\paragraph{What this additional example excludes.}
The construction uses a standard Markov and Fibonacci period-collapse
family and the published rational-triangle criterion; no historical
priority is claimed. Unlike P-F21, here the entire Hilbert series,
including every numerator coefficient, is independent of the
unbounded parameter degree. Thus even knowing the complete
Hilbert function, fixed Krull dimension, normality,
Cohen--Macaulayness, Gorensteinness, and rational singularities
does not bound the canonical modular index of a general graded
domain. All examples already have a positive form of denominator
weights one. This identifies a further invalid general route
from P-F19 to the weaker positive-form clause; it supplies no
ordinary-plethysm counterexample to that weaker clause.


\subsection{A geometric condition that prevents a cyclotomic cancellation}

\begin{proposition}[P-F23: a unique smooth generic fixed component gives an actual pole]
\label{prop:fixedcomponentpole}
Let $R=\bigoplus_{d\ge0}R_d$ be a finitely generated positively
graded complex domain with $R_0=\mathbb C$. Fix $M\ge2$ and a
primitive $M$th root of unity $\zeta$, and let $g$ act on $R_d$
by multiplication by $\zeta^d$. Let $X=\operatorname{Spec}R$ and
let $X^g$ be the scheme-theoretic fixed locus. Suppose that
$X^g$ has a unique irreducible component $Y$ of maximal dimension
$\delta\ge1$, and that $X$ is smooth at the generic point of $Y$.
Then $\operatorname{Hilb}_R(t)$ has a pole of order exactly
$\delta$ at $t=\zeta$.

More precisely, give $Y$ its reduced induced structure, write
$\mathbb C[Y]=\bigoplus_{k\ge0}\mathbb C[Y]_{Mk}$, and set
\[
 h_Y(z)=\sum_{k\ge0}\dim\mathbb C[Y]_{Mk}\,z^k,
 \qquad
 e_Y=\lim_{z\to1}(1-z)^\delta h_Y(z)>0.
\]
At the generic point of $Y$, the conormal space of $Y$ in $X$
has no $g$-eigenvalue equal to one. If its eigenvalues, with
multiplicity, are $\nu_1,\ldots,\nu_c$, where
$c=\dim X-\delta$, then
\[
 \lim_{t\to\zeta}(1-t^M)^\delta\operatorname{Hilb}_R(t)
       =\frac{e_Y}{\prod_{j=1}^c(1-\nu_j)}\ne0.
 \tag{P-F23a}
\]
In particular every positive rational expression
$H(t)/\prod_i(1-t^{a_i})$ for this Hilbert series has some
denominator weight divisible by $M$.
\end{proposition}

\begin{proof}
Choose positive-degree homogeneous algebra generators of $R$.
List those whose degrees are not divisible by $M$ as
$f_1,\ldots,f_l$, with $\deg f_i=w_i$. Their ideal $I$
is exactly the fixed-locus ideal. Indeed $f_i$ itself is a
nonzero scalar multiple of $g(f_i)-f_i$, and after killing
these generators every remaining algebra generator is fixed
by $g$. Thus
\[
 X^g=\operatorname{Spec}(R/I),
 \qquad (R/I)_d=0\quad\hbox{if }M\nmid d.
 \tag{P-F23b}
\]
The ideal is homogeneous, so every irreducible component of
the fixed locus is graded; equivalently the connected grading
torus cannot permute a finite set of components nontrivially.
In particular $Y$ is graded and its coordinate ring is
concentrated in degrees divisible by $M$.

Let $K_\bullet=K_R(f_1,\ldots,f_l)$ be the graded Koszul
complex, with its $i$th exterior generator given internal degree
$w_i$ and $g$-character $\zeta^{w_i}$. Denote its homology
modules by $N_j$. Multiplication by each $f_i$ on the Koszul
complex is homotopic to zero, so every $N_j$ is annihilated
by $I$ and is a finite graded $R/I$-module. The degreewise
Euler characteristic identity is
\[
 \operatorname{Hilb}_R(t)\prod_{i=1}^l(1-t^{w_i})
       =\sum_{j=0}^l(-1)^j\operatorname{Hilb}_{N_j}(t).
 \tag{P-F23c}
\]
All factors on the left except the Hilbert series are nonzero
at $\zeta$.

We describe carefully the leading contribution of the modules
$N_j$. Every finite graded module over $R/I$ splits into its
internal-degree classes modulo $M$. Each class, shifted by
its residue, is a graded module for the regrading in which
degree $M$ becomes degree one. Its support has dimension at
most $\delta$, so its Hilbert series has pole order at most
$\delta$ at $t=\zeta$. A component of dimension less than
$\delta$ has smaller pole order and contributes zero after
multiplication by $(1-t^M)^\delta$.

For a module $N$ whose localization at the generic point
$\eta$ of $Y$ is a vector space over $\mathbb C(Y)$, the
leading term is
\[
 \lim_{t\to\zeta}(1-t^M)^\delta\operatorname{Hilb}_N(t)
       = e_Y\operatorname{Tr}(g\mid N_\eta).
 \tag{P-F23d}
\]
Here $g$ acts trivially on the field $\mathbb C(Y)$, since it
fixes $Y$ pointwise. One way to verify (P-F23d), including its
grading shifts, is to take a graded prime filtration separately
in each residue class of $N$. All factors of dimension less
than $\delta$ disappear in the displayed limit. Since $Y$ is
the unique maximal-dimensional component, the remaining factors
are shifted copies of $\mathbb C[Y]$. A factor with generator
of internal degree $a$ contributes
$t^a h_Y(t^M)$, so its leading contribution is
$\zeta^a e_Y$. Upon localization, these factors form a
filtration of $N_\eta$ by one-dimensional character spaces of
character $\zeta^a$, giving exactly the trace. This also shows
why lower-dimensional components and torsion cannot cancel
that leading coefficient. The existence and positivity of
$e_Y$ is the ordinary positive leading Hilbert coefficient
of the positively graded domain obtained by regrading
$\mathbb C[Y]$; equivalently take any common Veronese and
use its Hilbert polynomial.

We now compute the generic Koszul homology. Smoothness of
$X$ along a dense open subset of $Y$ implies that the fixed
scheme is smooth there. For clarity, this familiar fact can
be seen directly at a smooth closed fixed point: diagonalize
the finite-order action on the cotangent space and lift each
eigenvector to the maximal ideal by the corresponding averaging
projector. These lifted eigenvectors are regular formal
parameters, and the completed local ring is the power series
ring in them, with the indicated diagonal action. Its fixed
scheme is cut out by exactly the parameters whose eigenvalues
are not one. Thus the fixed scheme is reduced and regularly
embedded at such points, and its conormal representation has
no trivial character. This applies at the generic point of
$Y$ as well.

Put $K=\mathbb C(Y)$. In the regular local ring $R_\eta$,
the localized ideal $I_\eta$ is the maximal ideal, generated
by a regular sequence of length $c=\dim X-\delta$. There
is a natural exact sequence of $K$-representations
\[
 0\longrightarrow L\longrightarrow
 E:=\bigoplus_{i=1}^l K e_i
 \longrightarrow V:=I_\eta/I_\eta^2\longrightarrow0,
 \qquad e_i\longmapsto f_i\bmod I_\eta^2.
 \tag{P-F23e}
\]
The characters on $E$ are $\zeta^{w_i}\ne1$.
The conormal space $V$ has the characters
$\nu_1,\ldots,\nu_c\ne1$, so the kernel $L$ also has
only nontrivial characters.

The localized Koszul homology is naturally
\[
 (N_j)_\eta\cong\bigwedge^j_K L.
 \tag{P-F23f}
\]
For an explicit verification, choose a minimal subcollection of
the $f_i$ that generates $I_\eta$. It is a regular sequence;
subtract its linear combinations from the other generators.
This invertible change of the Koszul presentation gives the
Koszul complex of that regular sequence tensored with the
Koszul complex on $l-c$ zeros. Its homology is an exterior
algebra on $l-c$ copies of $K$. The natural map in homological
degree one sends a Koszul cycle to its coefficient vector modulo
$I_\eta$ and identifies this space with the kernel in
(P-F23e). The exterior-algebra identification is natural and
therefore $g$-equivariant, even though no equivariant choice
of the intermediate change of generators is needed.
In particular the modules satisfy the vector-space hypothesis
used in (P-F23d).

Taking alternating traces gives
\[
 \sum_j(-1)^j\operatorname{Tr}(g\mid (N_j)_\eta)
       =\det(1-g\mid L).
\]
Multiplying (P-F23c) by $(1-t^M)^\delta$, applying
(P-F23d), and using (P-F23e), we obtain
\begin{align*}
 \lim_{t\to\zeta}(1-t^M)^\delta\operatorname{Hilb}_R(t)
 &=e_Y\frac{\det(1-g\mid L)}{
                         \prod_i(1-\zeta^{w_i})}\\
 &=\frac{e_Y}{\det(1-g\mid V)},
\end{align*}
which is (P-F23a). None of its factors is zero. The module
argument already bounded the pole order by $\delta$, and
the nonzero limit makes it exactly $\delta$.

If every $a_i$ in a rational expression were not divisible
by $M$, then every denominator factor $1-t^{a_i}$ would
be nonzero at a primitive $M$th root. Its polynomial numerator
could not create a pole. This proves the final assertion;
numerator positivity is not needed for that step.
\end{proof}

\paragraph{Standard algebra used in the proof.}
The Koszul change of generators, null homotopies, and tensor-product
properties used above are stated in the Stacks Project,
Section~15.29, Lemmas~15.29.3--15.29.6 and 15.29.12
(\url{https://stacks.math.columbia.edu/tag/0621}). The vanishing
of positive Koszul homology for a regular sequence is
Lemma~15.31.2, Tag~062F
(\url{https://stacks.math.columbia.edu/tag/062F}). The
equivariant excess calculation and the leading-character formula
needed here are proved explicitly above, rather than inferred
from an unqualified fixed-point or purity slogan.

\paragraph{Several components and the cubic application.}
The same proof, when every maximal-dimensional fixed component
is generically in the smooth locus of $X$, gives a sum
\[
 \sum_{Y:\,\dim Y=\delta}
       \frac{e_Y}{\det(1-g\mid N^*_{Y/X,\eta_Y})}
 \tag{P-F23g}
\]
for the limit in (P-F23a). Its terms can have different complex
phases. The uniqueness hypothesis avoids that additional
cancellation problem. Generic smoothness is also substantial:
P-F21 and P-F22 show that a universal factor by itself can
disappear completely from the ordinary Hilbert series even
for normal Gorenstein rational domains.

For the chain cubic family, P-F24 below proves generic smoothness
along the one-dimensional chain component, including the full
stabilizer calculation. P-F26 identifies the entire smooth-hypersurface
fixed locus with its orbit and supplies finite two-power lifts.
P-F27 excludes graph-support boundary competitors, even when the
lift order was not assumed to be $2^{n-1}$. The later full support
theorem P-F30 also excludes the remaining singular competitors,
including disconnected graph parts with triple monomials. P-F31
then proves that the fixed scheme has the unique chain curve as
its support and applies the present proposition to obtain the
actual ordinary-series pole. Thus this proposition supplies the
noncancellation step of that application, while its abstract
multiple-component caution remains necessary.


\subsection{The chain contributes a smooth fixed component}

\begin{proposition}[P-F24: full chain stabilizer and local fixed dimension]
\label{prop:chainsmoothcomponent}
Let $n\ge4$, $M=2^{n-1}$, $g=\gcd(n,3)$, and let $F_n$ be the
chain cubic of P-F19. Its full projective linear automorphism group
is the diagonal cyclic group of order $M$. Its stabilizer as an
actual cubic form under $SL_n$ is exactly the scalar subgroup
$\boldsymbol\mu_g$, which acts trivially on all cubic forms.

Let $X=\operatorname{Spec}T^{(n)}$, with the stretch grading of
P-F19, and let $p$ be the image of $F_n$ in this quotient. Then
$X$ is smooth at $p$. The fixed scheme for the grading action of
$\boldsymbol\mu_M$ is smooth of dimension one at $p$.
The closure of the grading orbit of $p$ is consequently an
irreducible one-dimensional component of this fixed scheme,
and its nonzero grading orbit consists of smooth points of $X$.
If $3\nmid n$, this component has coordinate ring
$\mathbb C[z]$ with $\deg z=M$; its regraded leading Hilbert
coefficient $e_Y$ in P-F23 is therefore one.
This does not classify the other components of the fixed scheme.
\end{proposition}

\begin{proof}
We first determine the full projective automorphism group; computing
only its diagonal subgroup would not suffice. Half the Hessian matrix
of $F_n$ is the symmetric tridiagonal matrix with diagonal
\[
 (x_2,x_3,\ldots,x_n,3x_n)
\]
and successive off-diagonal entries $x_1,\ldots,x_{n-1}$.
We claim that its projective rank-at-most-two locus consists exactly of
\[
 [e_1],\qquad
 [a e_{n-1}+b e_n]\quad
 \text{with }a=0\text{ or }a^2=3b^2.
 \tag{P-F24a}
\]
The latter are three distinct points on a line not containing $[e_1]$.

Here is a full rank check. A tridiagonal block of length $l$ with
all its $l-1$ off-diagonal entries nonzero has kernel dimension at
most one: its first equation determines the second coordinate from
the first, and the subsequent equations determine the rest in order.
The zero off-diagonal entries divide the whole matrix into such
blocks. Its rank is therefore at least the number of nonzero
entries among $x_1,\ldots,x_{n-1}$. At rank at most two, at most
two of these coordinates are nonzero.

If precisely two are nonzero, say $x_i,x_j$ with $i<j$, first
suppose $j\ge i+2$. The block on indices $i,i+1$ has determinant
$-x_i^2$: its first diagonal entry $x_{i+1}$ is zero. It has
rank two. The separate block containing the other nonzero
off-diagonal has rank at least one, a contradiction. If
$j=i+1$ and $i\ge2$, the singleton on index $i-1$ has nonzero
diagonal $x_i$ and is separate from the length-three block on
$i,i+1,i+2$. Their ranks are at least one and two. If $i=1$,
then for $n\ge5$ the block on $1,2,3$ has determinant
$-x_2^3\ne0$. For $n=4$ this remains true when $x_4=0$;
when $x_4\ne0$, the singleton on index four has rank one and
the length-three block has rank at least two. This excludes
every case with two nonzero off-diagonal entries.

If the only nonzero one is $x_j$ with $2\le j\le n-2$,
the block on $j,j+1$ has determinant $-x_j^2$, and the
singleton on $j-1$ has diagonal $x_j$, again giving rank
at least three. If $j=1$, its length-two block already has
rank two; $x_n\ne0$ would contribute nonzero diagonal entries
on the disjoint last two indices. Thus only $[e_1]$ remains.
If $j=n-1$, put $a=x_{n-1}$ and $b=x_n$. The only nonzero
part is the singleton $a$ on index $n-2$ and the block
\[
 \begin{pmatrix}b&a\\a&3b\end{pmatrix}.
\]
Its total rank is two exactly when $a^2=3b^2$.
Finally, if there are no nonzero off-diagonal entries, the
projective point is $[e_n]$ and the rank is two. This proves
(P-F24a).

Under a linear change of coordinates the Hessian transforms by
congruence, so any projective linear automorphism of $F_n$
preserves (P-F24a). The three collinear points form its unique
collinear triple. Hence its remaining point $[e_1]$ is fixed.
Write a representing linear transformation in coordinate forms
$\ell_1,\ldots,\ell_n$, with
\[
 F_n(\ell_1,\ldots,\ell_n)=cF_n(x_1,\ldots,x_n),\qquad c\ne0.
\]
Fixing $[e_1]$ means $\ell_1=ax_1+L(x_2,\ldots,x_n)$,
$a\ne0$, while the other $\ell_i$ do not involve $x_1$.
Comparing the coefficients of $x_1^2$ gives
$a^2\ell_2=cx_2$, so $\ell_2=bx_2$ with $b\ne0$.
The coefficient of $x_1$ is $2aL\ell_2$ and must vanish.
Thus $L=0$. Restriction to $x_1=0$ leaves an automorphism
of the chain cubic in $n-1$ variables, whose first coordinate
form is a scalar multiple of its first variable.

This proves by induction that all the coordinate forms are
diagonal once the base case $n=4$ is checked. For that case
write the remaining ternary cubic as
\[
 G(y,z,w)=y^2z+z^2w+w^3.
\]
Its transformation has coordinate forms $by$, $uy+A(z,w)$,
and $vy+B(z,w)$, where $A,B$ are independent linear forms.
The coefficient of $y$ in its transformed cubic is
\[
 vA^2+2uAB+3vB^2.
\]
It is zero identically, so independence of $A,B$ implies
$u=v=0$. Comparison of the coefficient of $y^2$ now gives
$b^2A=cz$, so $A$ is a nonzero multiple of $z$.
Write $B=\alpha z+\beta w$, with $\beta\ne0$. The
coefficient of $zw^2$ is $3\alpha\beta^2$; hence
$\alpha=0$. This is the required diagonal form and completes
the induction.

For diagonal projective transformations choose a scalar
representative with last entry one. The last monomial then
forces the common multiplier to be one, and the other
monomials give $a_{i+1}=a_i^{-2}$. The last-entry condition is
$a_1^{(-2)^{n-1}}=1$. Thus the full projective group is
cyclic of order $M$. For linear transformations preserving
$F_n$ exactly, rather than its projective class, the same
recurrence instead ends with $a_n^3=1$. They form the
cyclic group of order $3M$, with entries
\[
 a_i=a_1^{(-2)^{i-1}},\qquad a_1^{3M}=1.
\]
Their determinant is $a_1^S$, where
\[
 S=1-2+4-\cdots+(-2)^{n-1}
   =\frac{1-(-2)^n}{3}.
\]
The integer $S$ is odd and satisfies $S\equiv n\pmod3$.
Consequently $\gcd(S,3M)=g$. The determinant-one kernel
has order $g$; when $g=3$, its entries are all the same
cube root of unity, because $(-2)^i\equiv1\pmod3$.
This proves the asserted $SL_n$ stabilizer.

By P-F19 the chain hypersurface is smooth. The standard
stability theorem for smooth hypersurfaces of degree at
least three implies that its nonzero affine $SL_n$ orbit
is closed. An exact source is Huybrechts, \emph{The geometry
of cubic hypersurfaces}, Chapter~3, Corollary~1.6, printed
pp.~149--150, together with Exercise~1.4 on p.~149 for the
affine formulation. The cited degree and dimension conditions
hold here, since $n\ge4$. Passing to $SL_n/\boldsymbol\mu_g$
does not change the quotient and makes the stabilizer trivial.
Luna's \emph{Slices \'etales}, Section~III.1, the unnumbered
slice theorem on p.~97 and Remark~1$^\circ$ on pp.~97--98, identify an
\'etale neighborhood of the quotient point with a neighborhood
in the quotient of a smooth normal slice by its stabilizer.
The stabilizer here is trivial, so the slice and the quotient
are smooth. Thus $X$ is smooth at $p$.

It remains to determine the local fixed dimension; smoothness
alone does not give it. Let $\zeta$ be primitive of order $M$,
and put
\[
 A=\operatorname{diag}
       (\zeta,\zeta^{-2},\ldots,\zeta^{(-2)^{n-1}}).
\]
This transformation fixes $F_n$. Its determinant is
$\zeta^S$ with $S$ odd. On degree-$s$ elements of the
canonical algebra $T^{(n)}$, it acts by the character
$\det(A)^{-3s/g}$, up to replacing $A$ by its inverse
according to the coordinate-action convention. Since
$3/g$ and $S$ are odd, this character generates exactly
the grading $\boldsymbol\mu_M$ action.

The only degree-three monomials fixed by $A$ are the $n$
monomials of $F_n$. To see this, the nonzero residue weights
$(-2)^i\pmod M$, $0\le i\le n-2$, have distinct
2-adic valuations, and the remaining weight is zero.
In a sum of three such weights that is zero modulo $M$,
the least valuation cannot occur only once. Its index
therefore occurs twice. If it is $i<n-1$, the third
weight must be $-2(-2)^i=(-2)^{i+1}$, giving exactly
a chain monomial. At the last nonzero weight, twice that
weight is zero and the third is the zero weight. Three
zero weights give $x_n^3$; three equal nonzero weights
never sum to zero because three is odd. This exhausts
the possibilities. Thus the fixed subspace of the cubic
representation has dimension $n$.

All eigenvalues of $A$ on the underlying $n$-space are
distinct, so its fixed subspace on $\mathfrak{sl}_n$
under conjugation consists of the traceless diagonal
matrices and has dimension $n-1$. The map
$\mathfrak{sl}_n\to T_{F_n}(SL_nF_n)$ is an isomorphism,
because the stabilizer is finite. Luna's smooth quotient
description, or the corresponding differential of the
local principal bundle, gives the exact tangent sequence
\[
 0\longrightarrow\mathfrak{sl}_n F_n
 \longrightarrow W^*
 \longrightarrow T_pX\longrightarrow0.
\]
It is $A$-equivariant. Taking fixed vectors is exact for
this finite-order action in characteristic zero. Hence
\[
 \dim(T_pX)^{\boldsymbol\mu_M}=n-(n-1)=1.
 \tag{P-F24b}
\]
The fixed scheme of a finite-order automorphism of a
smooth complex variety is smooth near a fixed point:
averaging local parameters into eigenspaces and passing
to the completed local ring identifies its fixed ideal
with the parameters of nontrivial character. This is
also the local argument used in P-F23. It follows from
(P-F24b) that the fixed scheme is smooth of dimension one
at $p$.

The positive-degree discriminant is nonzero at $p$, so
the grading orbit of $p$ has dimension one. All of it
is fixed by $\boldsymbol\mu_M$, since the grading group
is abelian. Its irreducible closure is contained in a
component of the fixed scheme through $p$. That component
has dimension one locally at its smooth point $p$, and
therefore has dimension one globally. The two irreducible
closed sets agree. Grading translations preserve smoothness
of $X$, so the entire nonzero orbit consists of smooth
points.

For the last assertion assume $3\nmid n$, so $g=1$.
Restriction to the grading orbit is the homomorphism
\[
 T^{(n)}\longrightarrow\mathbb C[u,u^{-1}],\qquad
 f\in T_s\longmapsto f(p)u^s.
\]
All exponents are nonnegative, and P-F19 shows that a
nonzero image has $M\mid s$. The discriminant has stretch
degree exactly $M$ and is nonzero at $p$. The image therefore
contains a nonzero multiple of $u^M$ and is exactly
$\mathbb C[u^M]$. This image is the coordinate ring of
the orbit closure, proving the claim. Its Hilbert series is
$1/(1-t^M)$, so $e_Y=1$. This completes the proof.
\end{proof}

\paragraph{The geometric gate isolated at this step.}
The diagonal subgroup itself is a known computation:
Zheng, \emph{On abelian automorphism groups of hypersurfaces},
arXiv:\texttt{2004.09008v3}, Lemma~3.2, printed p.~5,
computes the type-$T_n$ diagonal group. The Hessian argument
above verifies the full-group assertion needed here.
No historical priority is claimed. The local result establishes
the generic-smoothness hypothesis of P-F23 for the chain component.
What remains is to control all the other components of
$X^{\boldsymbol\mu_M}$, including those represented by
singular polystable cubics, sufficiently to prevent their
leading contributions from cancelling the chain contribution.
Uniqueness of the maximal-dimensional component would suffice.
This local proposition alone does not assert it; P-F30--P-F31
below establish it by a complete support classification and
thereby close the global noncancellation gate.

Primary sources for the cited geometric theorems are
\url{https://www.math.uni-bonn.de/people/huybrech/Notes.pdf}
and \url{https://www.numdam.org/item/10.24033/msmf.110.pdf}.


\subsection{Binary determinant rays: a bounded-hole subfamily}

The complete zero classification for invariants of a single binary form
already rules out one possible source of unbounded realized holes.
The classification used here is known: Brouwer--Draisma--Popoviciu,
\emph{The degrees of a system of parameters of the ring of invariants
of a binary form}, \emph{Transformation Groups} 20 (2015),
953--967, Proposition~4.2, printed p.~957. The following is its
specialization to the canonical ordinary-plethysm grading, rather
than a new classification.

\begin{proposition}[P-F25: the known binary-form zero classification on stretching rays]
\label{prop:binarydeterminantzeros}
Let $d\ge2$, set $g=\gcd(d,2)$ and $q=2/g$, and take the ordinary
plethysm data
\[
 W=\operatorname{Sym}^d\mathbb C^2,\qquad
 \lambda=(q),\qquad \pi=(d/g,d/g).
\]
Its canonical stretching function is
\[
 f_d(s)=\dim\operatorname{Hom}_{GL_2}
   \bigl((\det)^{ds/g},\operatorname{Sym}^{qs}W\bigr).
 \tag{B-H1}
\]
For $d=2,3$, it is nonzero exactly when $s$ is even.
For $d\ge4$, its positive-integer zeros are exactly those in the table:
\[
\begin{array}{c|l}
 d&\{s\ge1:f_d(s)=0\}\\ \hline
 5&\{1,3,5,7\}\\
 6&\{1,3,5,7,9,11,13\}\\
 7&\{1,3,5\}\\
 10&\{1,3,5,7\}\\
 14&\{1,3,5\}\\
 9,11,13&\{1,3\}\\
 d\equiv2\pmod4,\ d\ge18&\{1,3\}\\
 \text{all other }d\ge4&\{1\}.
\end{array}
 \tag{B-H2}
\]
In particular, after excluding identically zero congruence classes,
every realized stretching hole has index at most thirteen. This
bound is attained by binary sextics, where $f_6(13)=0$ and
$f_6(s)>0$ for every $s\ge14$.
\end{proposition}

\begin{proof}
Write $R_d=(\operatorname{Sym}W)^{SL_2}$ and denote its ordinary
coefficient degree by $D$. The scalar $-I_2$ forces $dD$ to be
even whenever $(R_d)_D\ne0$, equivalently $q\mid D$.
At degree $D=qs$, a general linear transformation acts on this
invariant subspace by $(\det)^{dD/2}=(\det)^{ds/g}$.
Indeed it is a scalar times a special linear matrix, and the scalar
acts with weight $dD$. Evaluation at compatible powers of a volume
vector therefore identifies the canonical algebra in (B-H1) with
the whole ring $R_d$, regraded by $\deg (R_d)_{qs}=s$.
As in P-F19, this evaluation respects Cartan multiplication.

Consequently $f_d(s)$ is the invariant dimension denoted
$h^d_{qs}$ in the cited Proposition~4.2. Substituting $D=qs$ in
that proposition gives (B-H2). In particular, for odd $d$ all odd
ordinary degrees are absent before regrading, while their finite
even-degree exceptions are divided by two. For even $d$ the
ordinary and stretch gradings coincide. The two exceptional rings
of orders two and three have invariant degrees divisible by two
and four, respectively, giving the stated even-stretch support.
The maximum finite zero in the table is thirteen and occurs when
$d=6$. This proves all the asserted consequences of the known
classification.
\end{proof}

\paragraph{What this does not say about the source's strict SH.}
GCT VI, Definition~1.2.4, tests a nonzero constituent polynomial
$f_i(X)$ at \emph{every} integer $X\ge1$, not only those integers
in the residue class represented by $i$. Thus (B-H2), which
classifies the actual values $f_d(s)$, does not by itself prove
$s(f_d)\le13$, nor a bound for $p(f_d)$. Off-residue constituent
values are additional data. The bound above concerns realized
holes and excludes this family as a source of arbitrarily late
actual zero coefficients; a strict-SH or PH2 conclusion would need
an extra constituent-positivity argument.

\paragraph{Polynomial generator bounds and the remaining parameter gate.}
For this same binary family, precise polynomial bounds are already
known for different invariants. Derksen,
\emph{Polynomial bounds for rings of invariants},
\emph{Proceedings of the American Mathematical Society} 129 (2001),
955--963, Theorem~1.1, Proposition~1.2, and Example~2.1,
printed p.~956, gives
\[
 \sigma(\operatorname{Sym}^d\mathbb C^2)\le2d^3,
 \qquad
 \beta(\operatorname{Sym}^d\mathbb C^2)
       \le\frac32(d+1)d^6.
 \tag{B-H3}
\]
Here $\sigma$ bounds the degrees of a collection of homogeneous
invariants defining the null cone, and $\beta$ bounds algebra
generator degrees. The coordinate-ring convention in that source
agrees with $R_d$ because $\operatorname{Sym}^d\mathbb C^2$ is
self-dual as an $SL_2$-module. The source obtains the first inequality from
$SL_2\subset\mathbb A^4$, its one defining equation of degree two,
group dimension three, and matrix coefficients of the degree-$d$
representation of degree at most $d$. The representation has
finite kernel, as required by Proposition~1.2.

These statements do not supply a polynomial bound on the canonical
modular index. The same paper, Section~3, equation~(3.1), printed
p.~957, constructs a genuine homogeneous system of parameters by
raising null-cone equations to a common least-common-multiple
degree. Its Section~4 explicitly avoids this step when proving
the improved \emph{generator} bound: Proposition~4.2 concerns a
finite module over a possibly dependent collection of invariants,
and its proof places the parameters in a polynomial extension of
the original ring. It does not place a small-degree parameter
system in the original invariant ring.

For clarity, a verified explicit, though not polynomial, consequence
is available. Put
\[
 B_d=\left\lfloor\frac{2d^3}{q}\right\rfloor,
 \qquad L_d=\operatorname{lcm}(1,2,\ldots,B_d).
\]
In the stretch grading the null-cone equations can be chosen with
degrees at most $B_d$. Their powers of degree $L_d$ still define
the irrelevant maximal ideal set-theoretically. Homogeneous
normalization gives a parameter system with every degree $L_d$;
the prime-avoidance extension argument in P-F17 also bounds the
source's lexicographically minimal canonical index by $L_d$.
For $d\ge3$ the invariant ring has dimension $d-2$
(Brouwer--Draisma--Popoviciu, Section~1, p.~953, with the
$d\ge3$ scope made explicit in Proposition~2.1, p.~954).
The Cohen--Macaulay theorem and the global quasi-polynomial
input for the actual canonical ray, together with P-F9, yield
\[
 \chi(f_d)\le L_d,\qquad
 s(f_d)\le p(f_d)\le(d-3)L_d.
 \tag{B-H4}
\]
The case $d=2$ is the separately explicit polynomial ring with
stretch generator degree two and both indices zero. Formula
(B-H4) is a computable envelope, not a proof of polynomial PH2 or
PH3. Converting the already polynomial null-cone bounds into a
polynomial bound for actual parameter degrees remains a separate
family-specific gate.

The primary sources are
\url{https://pure.tue.nl/ws/files/17749340/Degrees2015.pdf}
and
\url{https://sites.lsa.umich.edu/hderksen/wp-content/uploads/sites/614/2018/05/A.I.a.7.pdf}.


\subsection{Finite lifts and the entire smooth fixed locus}

\begin{proposition}[P-F26: finite two-power lifts and the smooth chain locus]
\label{prop:finitecubiclift}
Let $n\ge4$, $M=2^{n-1}$, $g=\gcd(n,3)$, and
$X=\operatorname{Spec}T^{(n)}$ with the canonical stretch grading
of P-F19. Write $\Delta$ for the discriminant, viewed as a
homogeneous function on $X$, and let $p_n$ be the quotient point
of the chain cubic $F_n$. Then the following statements hold.

\begin{enumerate}
\item Every nonzero closed point $p\in X^{\boldsymbol\mu_M}$
has a representative cubic $F$ with closed $SL_n$ orbit and an
exact-form stabilizer element $A\in GL_n$ of finite two-power
order such that $\det A$ has order exactly $M$.
\item If $\Delta(p)\ne0$, the lift $A$ can be chosen of order
exactly $M$, and $F$ is $GL_n$-equivalent to $F_n$.
\item Scheme-theoretically, the part of the fixed scheme in
the smooth-hypersurface locus is the single grading orbit:
\[
 X^{\boldsymbol\mu_M}\cap D(\Delta)
      =\mathbb G_m\cdot p_n.
 \tag{P-F26a}
\]
This is a smooth reduced curve. Every irreducible component of
$X^{\boldsymbol\mu_M}$ other than the chain component of P-F24
is contained in $V(\Delta)$.
\end{enumerate}
The first statement does not assert that the lift has order $M$
when $F$ is singular.
\end{proposition}

\begin{proof}
Put $\gamma=3/g$, an odd integer. Recall the actual ring
identification from P-F19:
\[
 T^{(n)}_s=\bigl(\operatorname{Sym}^{(n/g)s}
                         (\operatorname{Sym}^3\mathbb C^n)\bigr)^{SL_n}.
\]
The $GL_n$-character on this subspace is $(\det)^{\gamma s}$.
On points, with the convention $(B\cdot F)(x)=F(B^{-1}x)$,
one consequently has
\[
 P(B\cdot F)=(\det B)^{-\gamma s}P(F)
       \quad(P\in T^{(n)}_s).
 \tag{P-F26b}
\]
Fix a primitive $M$th root $\zeta$. Since $\gamma$ is odd,
there is a unique $\eta\in\boldsymbol\mu_M$ with
$\eta^{-\gamma}=\zeta$, and $\eta$ is primitive. Choose
$B\in GL_n$ with determinant $\eta$. Thus the action of $B$
on the quotient is the grading action with parameter $\zeta$.
This choice obtains a determinant of order $M$ directly; a
scalar-only lift could instead introduce an unnecessary odd
factor into its determinant order.

Every fiber of the affine reductive quotient contains exactly
one closed orbit. A precise primary source for this fact is
Luna, \emph{Slices \'etales}, Section~I.1, printed p.~84.
Choose $F$ in that closed orbit over $p$. The form $B\cdot F$
has a closed $SL_n$ orbit, since $B$ normalizes $SL_n$.
It has the same quotient point by (P-F26b) and the fixed-point
hypothesis. Hence $B\cdot F=S\cdot F$ for some $S\in SL_n$.
Replacing $B$ by $S^{-1}B$ produces a matrix fixing $F$
exactly and still of determinant $\eta$. Denote this matrix
temporarily by $C$.

We next replace $C$ by a finite two-power element without
changing its determinant. Let $C=C_sC_u$ be its commuting
semisimple and unipotent parts. Their actions on cubic forms
are again semisimple and unipotent. Decompose $F$ into the
eigenspaces of $C_s$ in that representation. On an eigenspace
of eigenvalue $a\ne1$, the operator $aC_u-1$ is invertible,
because every eigenvalue of $C_u$ is one. Since $C F=F$,
the components in those eigenspaces are zero. It follows that
$C_sF=F$; also $\det C_s=\det C=\eta$.

Diagonalize $C_s$, with eigenvalues $\alpha_1,\ldots,\alpha_n$
on the underlying vector space. Let
$\Gamma=\langle\alpha_1,\ldots,\alpha_n\rangle\subset\mathbb C^*$.
This is a finitely generated abelian group. Its torsion
subgroup is finite cyclic, and a choice of a free complement
gives a group homomorphism
\[
 \rho:\Gamma\longrightarrow\Gamma_{\mathrm{tors},2}
\]
that is the identity on its two-primary torsion subgroup,
and kills its odd torsion and the chosen free complement.
Here $\eta=\prod_i\alpha_i$ belongs to that two-primary
torsion subgroup, so $\rho(\eta)=\eta$.

Every monomial $x_1^{e_1}\cdots x_n^{e_n}$ with nonzero
coefficient in $F$ satisfies $\prod_i\alpha_i^{e_i}=1$;
inverting all the eigenvalues, as in the action on forms,
gives the same equality. Applying $\rho$ preserves this
relation. Therefore
\[
 A=\operatorname{diag}(\rho(\alpha_1),\ldots,
                       \rho(\alpha_n))
 \tag{P-F26c}
\]
fixes $F$ exactly, has finite two-power order, and has
determinant $\eta$. This proves the first statement.

Now assume $\Delta(p)\ne0$. The representative $F$ defines
a smooth hypersurface. Write the order of $A$ as $2^\ell$.
Because its determinant has order $M$, one has $\ell\ge n-1$.
At least one eigenvalue of $A$ has order $2^\ell$.
We give a direct smoothness argument for the reverse bound
and its equality case.

Suppose the eigenvalue on variable $x_i$ has nontrivial
two-power order. The coefficient of $x_i^3$ in an invariant
cubic must be zero. Thus the coordinate point $[e_i]$ lies
on the hypersurface. Smoothness there implies the presence
of a monomial $x_i^2x_j$ with nonzero coefficient for some
$j\ne i$. Invariance forces
\[
 \alpha_j=\alpha_i^{-2},
 \tag{P-F26d}
\]
where in this paragraph $\alpha_i$ denotes the eigenvalue
of the finite-order matrix $A$ itself. Its order is half
the order of $\alpha_i$. Starting with an eigenvalue of
order $2^\ell$ and iterating gives eigenvalues of orders
\[
 2^\ell,2^{\ell-1},\ldots,2,1.
\]
They occupy $\ell+1$ distinct variables. Hence $\ell+1\le n$.
Consequently $\ell=n-1$, and the chosen lift has order
exactly $M$.

In the equality case these are all the variables. After
permutation their eigenvalues are
\[
 \xi,\xi^{-2},\ldots,\xi^{(-2)^{n-1}}=1,
 \qquad \operatorname{ord}(\xi)=M.
 \tag{P-F26e}
\]
For completeness, the invariant degree-three monomials for
this spectrum are exactly those of the chain. The nonzero
residue weights $1,-2,4,\ldots,(-2)^{n-2}$ have distinct
two-adic valuations. In a vanishing sum of three weights,
a least nonzero valuation cannot occur only once. If it
occurs three times, the sum is not zero, since three is odd.
Thus it occurs twice, and the third weight must be minus
twice that weight, which is precisely the next chain weight.
At the last nonzero weight the next is zero. The all-zero
case gives the final cube. Therefore
\[
 F=c_1x_1^2x_2+\cdots+c_{n-1}x_{n-1}^2x_n+c_nx_n^3.
 \tag{P-F26f}
\]
Every $c_i$ is nonzero. For $i<n$ this follows again from
smoothness at $[e_i]$; if $c_n=0$, the final coordinate
point is singular as well. Starting with $b_n^3c_n=1$
and successively solving $b_i^2b_{i+1}c_i=1$ gives nonzero
complex numbers $b_i$. The diagonal substitution
$x_i\mapsto b_ix_i$ turns (P-F26f) into $F_n$. This proves
the second statement, including the equality classification.

The $GL_n$-action on the quotient is precisely the grading
action by (P-F26b): the map $\mathbb C^*\to\mathbb C^*$,
$z\mapsto z^{-\gamma}$, is surjective. Thus every fixed
closed point in $D(\Delta)$ belongs to the grading orbit
of $p_n$. Conversely that entire orbit is fixed by
$\boldsymbol\mu_M$ and lies in $D(\Delta)$, by P-F24 and
discriminant nonvanishing.

P-F24 also shows that the fixed scheme is smooth of dimension
one at every point of this orbit. Its intersection with
$D(\Delta)$ consequently is reduced and has exactly the
closed points of that orbit. Both carry their reduced
locally closed scheme structures, so they agree, proving
(P-F26a). More explicitly, any nonreduced or additional
finite-type locus would have a closed complex point,
contradicting the preceding local smoothness and point
classification. If a different irreducible fixed component
met $D(\Delta)$, it would pass through a smooth point of
the fixed scheme on this orbit. That point has a unique
component, namely the chain component of P-F24. This is a
contradiction. Thus all other components lie in $V(\Delta)$.
\end{proof}

\paragraph{Known order bound, and the precise remaining lift issue.}
For smooth hypersurfaces the two-power order bound also follows
from Zheng, \emph{On abelian automorphism groups of hypersurfaces},
arXiv:\texttt{2004.09008v3}, Theorem~4.8, printed pp.~11--12.
An exact-form stabilizer of two-power order has the same
projective order, since its scalar kernel is contained in
$\boldsymbol\mu_3$.
At degree three its factors $|1-(-2)^a|$ are odd, and the only
two-power factors in the theorem are $2^{b-1}$ with $b\le n$.
The elementary coordinate-point argument above proves the needed
bound and makes its equality case explicit. Theorem~4.3 alone
states an abstract subgroup embedding into the diagonal group
of a simple hypersurface; it should not be cited by itself as
conjugacy of the original representation. No historical priority
is claimed for the order bound or its chain realization.

For a singular closed-orbit representative $F$, put
$H=\operatorname{Stab}_{GL_n}(F)$. Once the determinant image
contains $\boldsymbol\mu_M$, an exact-order-$M$ lift exists
if and only if the sequence of complex groups
\[
 1\longrightarrow H\cap SL_n\longrightarrow
   \det^{-1}(\boldsymbol\mu_M)\cap H
   \xrightarrow{\det}\boldsymbol\mu_M\longrightarrow1
 \tag{P-F26g}
\]
has a group-theoretic section. Indeed an order-$M$ element
mapping to a generator defines a section, and the image
of a generator under a section has order exactly $M$.
Such splitting is not automatic for abstract group extensions:
the squaring map from a cyclic group of order $2M$ onto a
cyclic group of order $M$ has no such section when $M$ is a
power of two. This abstract example is not asserted to occur
as the stabilizer extension of a cubic. The proof above
establishes splitting in the smooth case. At this step it
confines all possible competing fixed components to the
discriminant boundary without assuming a splitting there.
P-F30 subsequently excludes every nonzero singular competitor
at this extremal determinant order, even for larger-order lifts;
P-F31 then proves uniqueness and noncancellation. The abstract
extension warning explains why that additional work is needed.

The primary texts are
\url{https://www.numdam.org/item/10.24033/msmf.110.pdf}
and \url{https://arxiv.org/pdf/2004.09008}.


\subsection{The singular support gate requires three distinct variables}

\begin{proposition}[P-F27: extremal determinant order for graph supports]
\label{prop:cubicgraphsupport}
Let $n\ge2$, and let $F\in\mathbb C[x_1,\ldots,x_n]$ be a
cubic whose nonzero monomials all have the forms $x_i^3$ or
$x_i^2x_j$ with $i\ne j$. Let $B$ be the integer matrix whose
columns are their exponent vectors. Assume that the vector
$\mathbf1=(1,\ldots,1)^t$ belongs to the rational column span
of $B$. This hypothesis holds, in particular, if $F$ is
$SL_n$-semistable.

Suppose a diagonal matrix $A$ of finite two-power order fixes
$F$ exactly. Then the order of $\det A$ is at most $2^{n-1}$.
If that order is $2^{n-1}$, the support is, after a permutation
of variables, exactly
\[
 x_1^2x_2,\ x_2^2x_3,\ldots,\ x_{n-1}^2x_n,\ x_n^3.
 \tag{P-F27a}
\]
All their coefficients are nonzero, the matrix $A$ itself has
order $2^{n-1}$, and a diagonal change of variables turns $F$
into the smooth chain cubic $F_n$.

Consequently, for $n\ge4$, every nonzero fixed quotient point
in the discriminant boundary of P-F26 has the following
property: in any eigenbasis for any finite two-power
exact-form lift with determinant of order $2^{n-1}$, its
closed-orbit representative contains a nonzero monomial
$x_ix_jx_k$ with $i,j,k$ pairwise distinct.
\end{proposition}

\begin{proof}
We first explain the rational-span assumption. If a
positive-degree $SL_n$-invariant polynomial in the
coefficients of $F$ does not vanish at $F$, at least one of
its coefficient monomials has nonzero value. Invariance
under the diagonal determinant-one torus means that the
sum of the cubic exponent vectors, counted with that
coefficient monomial's nonnegative integer multiplicities,
has all coordinates equal. The common coordinate is
positive, because the coefficient degree is positive.
Dividing by it expresses $\mathbf1$ as a nonnegative
rational combination of columns of $B$. Only membership
in their rational span will be used below.

Write the eigenvalues of $A$ as $\alpha_1,\ldots,\alpha_n$.
They lie in one finite cyclic two-group. The map
\[
 \psi:\mathbb Z^n\longrightarrow\mathbb C^*,\qquad
 (u_i)\longmapsto\prod_i\alpha_i^{u_i}
 \tag{P-F27b}
\]
annihilates the lattice $L$ generated by the columns of $B$,
since every supported monomial is fixed. Thus it factors
through $\mathbb Z^n/L$, and
$\psi(\mathbf1)=\det A$. This is the precise lattice
interpretation of the determinant character; the order
of the determinant need not equal the order of $A$.

Make a graph on the variables. A supported monomial
$x_i^2x_j$ gives an edge with coefficient two at vertex
$i$ and coefficient one at vertex $j$; a cube gives a loop
with coefficient three. Distinct oriented monomials may
give two edges on the same unordered pair. Decompose the
graph into connected components, ignoring loops for
connectivity. The exponent lattice and rational span
split over these components. In a component with $v$
vertices, therefore, its all-ones vector $\mathbf1_v$
belongs to the rational span of its support columns.
The determinant is the product of the images of these
component vectors under $\psi$.

We prove that each component contributes a determinant
character of order at most $2^{v-1}$, with a stronger
bound $2^{v-2}$ in the rank-deficient case when $v\ge2$.
An isolated vertex must have a cube, since otherwise
$\mathbf1_v$ would not belong to the column span. Its
relation is $3e_i=0$, and every character into a two-group
is trivial on $e_i$. This handles $v=1$.

For $v\ge2$, choose a spanning tree of the underlying
connected graph, and let $T$ be the $v\times(v-1)$ matrix
of its support columns. Its columns are independent:
removing a leaf removes a row with exactly one nonzero
entry, equal to one or two, and induction proves rank
$v-1$. Hence the full component matrix has rank either
$v-1$ or $v$.

\emph{The rank-deficient case.}
Suppose its rank is $v-1$. Then its rational column span
is that of $T$, so $Tc=\mathbf1_v$ has a unique rational
solution $c$. For every vertex $r$, let $T^{(r)}$ be the
square matrix obtained by deleting row $r$. Repeatedly
remove leaves other than $r$. This proves
\[
 D_r:=\det T^{(r)}=\pm2^{e_r},\qquad 0\le e_r\le v-1.
 \tag{P-F27c}
\]
Cramer's rule applied after deleting row $r$ shows that
the least positive integer $q$ with $qc$ integral divides
$|D_r|$. It therefore divides the greatest common divisor
of all these maximal minors. Also
$q\mathbf1_v=T(qc)$ lies in the tree lattice. Equivalently,
the torsion class of $\mathbf1_v$ in the tree cokernel has
order $q$, and its image in the full support cokernel has
order dividing $q$. This proves the needed maximal-minor
assertion directly, without assuming full rank.

The signed cofactor vector of $T$ is a nonzero vector
in the kernel of $T^t$. Along an edge with coefficients
two and one, its two endpoint coordinates consequently
satisfy $2u_i+u_j=0$. Their absolute values differ by a
factor of two. Since their absolute values are the
$|D_r|$, two adjacent exponents in (P-F27c) differ by one.
As all these exponents are at most $v-1$, their minimum
is at most $v-2$. Thus
\[
 q\mid 2^{\min_r e_r},\qquad \min_r e_r\le v-2.
 \tag{P-F27d}
\]
The component determinant character has order at most
$2^{v-2}$, as asserted. Notice that free directions in
the cokernel may allow the matrix eigenvalues themselves
to have much larger order; the span hypothesis makes
$\mathbf1_v$ a torsion class and is what controls its
character.

\emph{The full-rank case.}
Suppose instead that the component matrix has rank $v$.
One support column outside the span of $T$ extends its
columns to a square nonsingular matrix $C$. Its support
graph is connected with $v$ edges and $v$ vertices, and
has exactly one cycle, possibly a loop or a double edge.
The finite cokernel of $C$ has order $|\det C|$; one may
also use the adjugate identity to see directly that
$\det C$ kills every cokernel class. Since its image under
$\psi$ is a two-group, the component determinant order
is at most $2^{v_2(\det C)}$.

Let the unique cycle have length $l$. Expanding the
matrix successively along leaf rows contributes one
factor of either one or two for each of its $v-l$
tree edges. If $l=1$, the remaining loop contributes
three, and hence
\[
 v_2(\det C)\le v-1.
 \tag{P-F27e}
\]
If $l\ge2$, the cycle determinant has exactly two
permutation terms. For some $0\le a\le l$ they give,
up to an overall sign,
\[
 2^a+(-1)^{l-1}2^{l-a}.
 \tag{P-F27f}
\]
If $a=l-a$, then $l$ is even and the determinant is zero,
which is excluded. Otherwise its two-adic valuation is
$\min(a,l-a)\le\lfloor(l-1)/2\rfloor$. Therefore
\[
 v_2(\det C)
 \le v-l+\left\lfloor\frac{l-1}{2}\right\rfloor
 \le v-2\qquad(l\ge2).
 \tag{P-F27g}
\]
This proves the component bound in every case, including
multiple edges and rank deficiency.

Because all component characters have two-power order,
their product has order at most their largest order.
Every component has at most $n$ vertices. This proves
$\operatorname{ord}(\det A)\le2^{n-1}$.

Now suppose equality holds. At least one component must
have $n$ vertices, since every smaller component has
bound at most $2^{n-2}$. Thus the support graph is
connected. It has full rank by (P-F27d), and the chosen
unicyclic basis must have a loop by (P-F27g). Moreover,
equality in (P-F27e) requires every leaf expansion factor
to be two. Root its remaining tree at the loop vertex.
Each basis edge therefore has coefficient two at its
child and one at its parent. The eigenvalue relations are
\[
 \alpha_{\mathrm{root}}^3=1,\qquad
 \alpha_{\mathrm{parent}}=\alpha_{\mathrm{child}}^{-2}.
 \tag{P-F27h}
\]
Since all eigenvalues have two-power order, the root
eigenvalue is one. A vertex at depth $d$ then has
eigenvalue order at most $2^d$. If the maximum tree
depth is $h$, the order of every eigenvalue, and therefore
of their product, is at most $2^h$. Equality for the
determinant forces $h\ge n-1$. A rooted tree with $n$
vertices has $h\le n-1$, with equality only for a path
rooted at an endpoint. Thus this basis is the chain
support in (P-F27a).

The deepest eigenvalue must have order $2^{n-1}$: all
other vertices have smaller depth, and a product of
elements of smaller two-power order could not have the
required determinant order. Write this eigenvalue as
$\xi$, and order the path from its deepest vertex to the
root. Then the entire eigenvalue list is
\[
 \xi,\ \xi^{-2},\ \xi^4,\ldots,
 \xi^{(-2)^{n-1}}=1,
 \qquad \operatorname{ord}(\xi)=2^{n-1}.
 \tag{P-F27i}
\]
In particular, $A$ has this same order, although this
was not assumed.

These eigenvalues permit no additional cubic support.
Indeed their nonzero residue weights modulo $2^{n-1}$
have distinct two-adic valuations. In a zero sum of
three weights, the least nonzero valuation cannot occur
once, and cannot occur three times because three is
odd. It occurs twice at the same vertex; the third
weight must be its negative double, the next chain
weight. The all-zero case gives the root cube. Thus the
fixed cubic monomials are exactly (P-F27a). They are all
present, since they were the columns of the chosen basis.
Their nonzero coefficients can be made one by first
rescaling the root variable using a cube root and then
successively rescaling its predecessors using square
roots. Hence $F$ is the smooth chain cubic up to a
diagonal substitution, by the smoothness proof in P-F19.

Finally, P-F26 supplies a finite two-power lift with
primitive determinant for every nonzero fixed quotient
point. Its closed-orbit representative is semistable,
so it satisfies the span hypothesis in every coordinate
basis. If that representative is singular and any such
lift eigenbasis had only graph-support monomials, the
result just proved would identify it with a smooth
chain. This contradiction proves the last assertion.
\end{proof}

\paragraph{Connected graph part, even with additional triple monomials.}
The same equality classification holds for an arbitrary cubic satisfying
the rational-span hypothesis if its graph made from its square-linear
monomials is connected; cubes are retained as loops. Thus the presence
of a three-distinct-variable monomial by itself is not enough to produce
a competitor. In the discriminant boundary of P-F26, this graph must be
disconnected in every finite-lift eigenbasis.

Here is the additional proof. Let $G$ be the matrix consisting only
of the cube and square-linear support columns. Since its graph is
connected, its rank is $n$ or $n-1$. If it has rank $n$, its graph
subform satisfies the span hypothesis by itself. The proposition
therefore gives the chain eigenvalue list (P-F27i), which permits no
additional triple monomials. The full form is the chain.

Suppose instead that $G$ has rank $n-1$, and take the spanning-tree
matrix $T$ used above. Its signed cofactor vector has the form
\[
 u_i=\varepsilon 2^e(-2)^{h_i},\qquad
 \varepsilon\in\{1,-1\},\quad e\ge0,\quad
 h_i\ge0,\quad \min_i h_i=0.
 \tag{P-F27j}
\]
Indeed its coordinates are signed powers of two and the relation
$2u_i+u_j=0$ along each tree edge fixes both the change in exponent
and the change in sign. Connectivity makes the sign formula in
(P-F27j) consistent. Put $H=\max_i h_i$. The tree-minor calculation
gives $e+H\le n-1$.

If every triple support column lies in the rational span of $T$,
then so does $\mathbf1$, and the rank-deficient argument already
bounds the determinant order by $2^{n-2}$. Otherwise choose a triple
column $\beta=e_i+e_j+e_k$ outside that span, with $i,j,k$ distinct.
Cofactor expansion gives a nonsingular square matrix with
\[
 \det(T,\beta)=\pm 2^e
       \bigl((-2)^{h_i}+(-2)^{h_j}+(-2)^{h_k}\bigr).
 \tag{P-F27k}
\]
A nonzero sum of three powers of $-2$ has two-adic valuation strictly
less than the largest exponent occurring, unless all three exponents
are equal. To verify this, a unique smallest exponent gives exactly
that valuation. If two smallest exponents are $r$ and the third is
$t>r$, the sum is $(-2)^r(2+(-2)^{t-r})$: it vanishes when $t=r+1$,
and has valuation $r+1<t$ when $t\ge r+2$. Three equal exponents
$r$ give valuation $r$. Consequently the valuation in (P-F27k) is
at most $e+H$, with equality possible only if all three chosen
vertices have height $H$.

If $e+H\le n-2$, this already contradicts determinant order
$2^{n-1}$, since the determinant of the square support submatrix
annihilates the determinant character. In the only remaining case,
$e+H=n-1$, choose a vertex $r$ of height $H$. Its deleted-row tree
minor has absolute value $2^{n-1}$. Root the tree at $r$ and perform
the leaf expansion for this minor. Every one of its $n-1$ factors
must be two. Thus each tree edge has coefficient two at its child
and one at its parent, and the cofactor relation says
$h_{\rm parent}=h_{\rm child}+1$. The root is the unique vertex of
height $H$. Three distinct vertices cannot all have that height.
Therefore the valuation in (P-F27k) is again at most $n-2$, the
same contradiction. This completes the connected-graph extension.

\paragraph{Scope of the reduction.}
The proof treats graph supports and, by the extension just proved,
arbitrary supports with connected graph part. It does not classify
all singular semistable cubics. It does not infer
that the determinant order bounds an arbitrary finite lift
before imposing the extremal graph-support case. Three
distinct-variable monomials are exactly the remaining
cubic support type, and the unresolved case has a disconnected
graph part. Thus P-F27 removes the rank-deficiency and
larger-order-lift obstructions for the stated classes, while
leaving the disconnected hypergraph support problem unresolved by
this proposition alone. P-F30 subsequently resolves that remaining
case by component contraction and a determinant bound, and P-F31
uses it to control the global fixed support and ordinary pole.
No historical priority is claimed for this elementary
lattice and graph reduction.


\subsection{The exact local coefficient of the chain component}

\begin{proposition}[P-F28: chain multiplicity and cyclotomic valuation]
\label{prop:chainlocalcoefficient}
Let $n\ge4$, $M=2^{n-1}$, $g=\gcd(n,3)$, and
$m=\binom{n+2}{3}$. Let $Y\subset X=\operatorname{Spec}T^{(n)}$
be the reduced chain orbit closure of P-F24. Its grading
stabilizer at every nonzero point is exactly
$\boldsymbol\mu_M$. In the notation of P-F23 its leading
graded multiplicity is
\[
 e_Y=\lim_{z\to1}(1-z)
           \sum_{s\ge0}\dim\mathbb C[Y]_{Ms}z^s=1.
 \tag{P-F28a}
\]

Fix a primitive $M$th root $\zeta$, and let $g_\zeta$
denote its grading action. Normalize the cyclotomic
valuation above $2$ by $v(1-\zeta)=1$. The generic
conormal representation of $Y$ in $X$ has dimension
$m-n^2$, and
\[
 v\!\left(\det(1-g_\zeta\mid N^*_{Y/X,\eta_Y})\right)
       =3M-\binom{n+2}{2}.
 \tag{P-F28b}
\]
Consequently the local coefficient attached to this
component by the calculation of P-F23,
\[
 \mathcal C_n(\zeta):=
   \frac{e_Y}{\det(1-g_\zeta\mid N^*_{Y/X,\eta_Y})},
\]
is nonzero and satisfies
\[
 v(\mathcal C_n(\zeta))=\binom{n+2}{2}-3M.
 \tag{P-F28c}
\]
This is a local coefficient. No claim is made that it
equals a Laurent coefficient of the entire Hilbert series,
or dominates the contributions of other possible fixed components
by this proposition alone. P-F30--P-F31 below prove that it is
the actual normalized leading coefficient.
\end{proposition}

\begin{proof}
Put $\gamma=3/g$. By P-F24, the full group of linear
transformations fixing the chain form exactly is cyclic
of order $3M$, and its determinant image has order
$3M/g$. On the quotient, a matrix $B\in GL_n$ acts by
the grading parameter $(\det B)^{-\gamma}$, with the
point-action convention of (P-F26b). Thus the image
of this exact-form stabilizer in the grading torus has
order
\[
 \frac{3M/g}{\gcd(3M/g,\,3/g)}=M.
 \tag{P-F28d}
\]
This image is the entire grading stabilizer of the
quotient point. To see the converse inclusion, let a
grading parameter $u$ fix that point and choose a
matrix $B$ such that $(\det B)^{-\gamma}=u$.
The chain form and $B$ applied to it have closed
$SL_n$ orbits and the same quotient point. The
unique-closed-orbit property of the affine reductive
quotient therefore gives $S\in SL_n$ such that
$B\cdot F_n=S\cdot F_n$. The matrix $S^{-1}B$
fixes the form exactly and induces the same parameter
$u$. This proves the converse and hence the exact
stabilizer assertion. Grading translations give it
at every point of the same nonzero orbit.

Restrict homogeneous functions to the map
$\mathbb G_m\to Y$, $u\mapsto u\cdot p_n$.
A function of degree $d$ restricts to $P(p_n)u^d$.
The image of the coordinate ring is consequently
\[
 \mathbb C[Y]=\mathbb C[u^d:d\in S],
 \tag{P-F28e}
\]
where $S\subset\mathbb Z_{\ge0}$ is a finitely generated
additive semigroup containing a positive integer.
Finite generation follows by restricting a finite
set of homogeneous algebra generators; positivity
follows, for example, from the nonzero discriminant.
The common stabilizer of these monomials is
$\boldsymbol\mu_{\gcd S}$, so (P-F28d) gives
$\gcd S=M$.
This description also shows that there are no other
nonzero points of $Y$: the map $\mathbb A^1\to Y$
from (P-F28e) is finite and surjective, since $u$
is integral over the ring containing some $u^a$,
and its points with $u\ne0$ are exactly the grading
orbit.

After dividing its degrees by $M$, this semigroup
has greatest common divisor one and contains every
sufficiently large integer. Here is the elementary
argument for the latter assertion. Choose a positive
element $a$ of the semigroup. The residues of the
semigroup elements generate all of the finite group
$\mathbb Z/a\mathbb Z$, because their greatest common
divisor with $a$ is one. A subsemigroup of a finite
group is a subgroup: the inverse of an element of
order $b$ is its $(b-1)$-fold sum. Hence there is a
semigroup element in every residue class modulo $a$.
Adding nonnegative multiples of $a$ to these finitely
many representatives proves the assertion. Each
nonzero graded piece in (P-F28e) is one-dimensional.
The regraded Hilbert series therefore differs from
$1/(1-z)$ by a polynomial. This proves $e_Y=1$.

We next compute the conormal valuation. Write
$r=n-1$, so $M=2^r$. For any nonzero residue
$b\bmod M$, one has
\[
 v(1-\zeta^b)=2^{v_2(b)},
 \qquad 0\le v_2(b)\le r-1.
 \tag{P-F28f}
\]
For completeness, the polynomial
$\Phi_{2^r}(1+Z)=(1+Z)^{2^{r-1}}+1$
is Eisenstein at $2$: its constant term is $2$,
and reduction modulo $2$ leaves only the leading
term. Thus $1-\zeta$ is a uniformizer of the
cyclotomic extension at $2$.
If $b=2^j c$ with $c$ odd, the quotient
$(1-\zeta^{2^j c})/(1-\zeta^{2^j})$ is a unit,
since its geometric sum is congruent to the odd
integer $c$ modulo the uniformizer; one may choose
$c$ positive in its residue class. Moreover
\[
 1-\zeta^{2^j}
   =\prod_{k=0}^{2^j-1}
       \left(1-\zeta^{\,1+(M/2^j)k}\right).
 \tag{P-F28g}
\]
Every exponent on the right is odd, and every factor
has valuation one. This gives (P-F28f).
In particular this valuation is unchanged on
replacing $\zeta$ by any primitive $M$th root,
or on replacing a character by its inverse.

Choose a primitive $M$th root $\xi$ so that the
diagonal chain automorphism
\[
 A=\operatorname{diag}
    (\xi,\xi^{-2},\xi^4,\ldots,\xi^{(-2)^{n-1}})
\]
induces $g_\zeta$ on the quotient. This is possible:
its determinant exponent
$S_0=(1-(-2)^n)/3$ and the integer $\gamma$ are
both odd, so the quotient parameter
$\xi^{-\gamma S_0}$ runs through all primitive
$M$th roots. Put $w_i=(-2)^i\bmod M$ for
$0\le i\le n-1$. The weights $w_0,\ldots,w_{n-2}$
are nonzero and have respective two-adic valuations
$0,\ldots,n-2$, while $w_{n-1}=0$.

By P-F24, the tangent space of $X$ at the chain point
is the quotient of the cubic representation by the
tangent space of the $SL_n$ orbit. Its fixed part is
one-dimensional, namely the tangent direction of
the grading orbit. Taking moving parts and then
duals therefore gives the conormal representation
of $Y$. The cubic space has dimension $m$, the
$SL_n$ orbit has dimension $n^2-1$, and the fixed
quotient direction has dimension one. The conormal
dimension is consequently $m-n^2$.
The same character data holds at the generic point
of $Y$, since its nonzero grading orbit is dense
and the grading torus commutes with $g_\zeta$.

It suffices to subtract the determinant valuation
of the moving adjoint representation from that
of the moving cubic representation. Duals do not
change these valuations by (P-F28f).
The off-diagonal adjoint weights are $w_i-w_j$
for $i\ne j$. If $i<j$, their two-adic valuation
is $i$, including the case $j=n-1$ of the zero
weight. The adjoint determinant valuation is thus
\[
 V_{\mathrm{ad}}
   =2\sum_{i=0}^{n-2}(n-1-i)2^i
   =4M-2n-2.
 \tag{P-F28h}
\]

For cubic monomials, fix the smallest index $i$
of a nonzero weight that occurs. There are three
cases according to its multiplicity. If it occurs
once, the other two indices are greater than $i$,
and there are $\binom{n-i}{2}$ unordered choices
with repetition. The weight sum has two-adic
valuation $i$. If it occurs three times, there is
one monomial, also with valuation $i$.
If it occurs twice, the remaining index is some
$j>i$. The choice $j=i+1$ is a fixed chain
monomial and supplies no moving factor. Each
$j\ge i+2$ instead gives weight $2w_i+w_j$ of
valuation $i+1$; there are $n-i-2$ such choices
for $0\le i\le n-3$. The final all-zero-weight
cube is fixed. These cases exhaust the cubic
monomials, so
\begin{align*}
 V_{\mathrm{cub}}
 &=\sum_{i=0}^{n-2}
        \left[\binom{n-i}{2}+1\right]2^i
      +\sum_{i=0}^{n-3}(n-i-2)2^{i+1}\\
 &=7M-\frac{n^2+7n+6}{2}.
 \tag{P-F28i}
\end{align*}
For the finite sums used here, the identities are
\[
 \sum_{i=0}^{n-2}(n-1-i)2^i=2^n-n-1,\qquad
 \sum_{i=0}^{n-2}\binom{n-i}{2}2^i
    =2^{n+1}-\frac{n^2+3n+4}{2}.
 \tag{P-F28j}
\]
They follow by induction on $n$: the left sides
are twice their values at $n-1$, plus respectively
$n-1$ and $\binom n2$, and both identities begin
at $n=2$. The second sum in (P-F28i) is twice the
first sum of (P-F28j) with $n$ replaced by $n-1$.

Subtracting (P-F28h) from (P-F28i) yields
\[
 V_{\mathrm{cub}}-V_{\mathrm{ad}}
     =3M-\frac{n^2+3n+2}{2}
     =3M-\binom{n+2}{2}.
 \tag{P-F28k}
\]
This is (P-F28b). Each conormal eigenvalue is
nontrivial by P-F24, so its determinant factor
is nonzero. Combining it with $e_Y=1$ proves
(P-F28c).
\end{proof}

\paragraph{Role in the subsequent global argument.}
The local calculation does not compare possible boundary
components or establish their absence. P-F30 supplies
that separate support classification, and P-F31 then
verifies the global hypotheses of P-F23. Only at that
stage does the local coefficient computed here become
the actual normalized leading coefficient of the whole
Hilbert series. Thus the proof of the global pole uses
classification of the fixed locus, without an assumed
inequality between cyclotomic valuations of different
components.


\subsection{Pure cubes and squarefree triples cannot contribute an extremal boundary point}

\begin{proposition}[P-F29: the complementary squarefree-support bound]
\label{prop:cubicsquarefreesupport}
Let $n\ge2$, and let a cubic $F\in\mathbb C[x_1,\ldots,x_n]$
have only monomials of the forms $x_i^3$ and $x_ix_jx_k$
with $i,j,k$ pairwise distinct. Let $B$ be its support
exponent matrix, of rank $r$ over $\mathbb Q$, and assume
$\mathbf1\in\operatorname{im}_{\mathbb Q}B$. Let $l$ be the
number of variables whose cubes occur in the support.
If a diagonal matrix $A$ of finite two-power order fixes
$F$ exactly, then
\[
 \operatorname{ord}(\det A)\le
 \begin{cases}
  3^{(r-l)/2},&r<n\text{ or }l>0,\\
  3^{n/2-1},&r=n\text{ and }l=0.
 \end{cases}
 \tag{P-F29a}
\]
In both cases this is strictly less than $2^{n-1}$.

Consequently, every nonzero competing fixed quotient point
in the discriminant boundary of P-F26 has, in every
eigenbasis of every finite two-power exact-form lift with
determinant order $2^{n-1}$, both a monomial $x_i^2x_j$
with $i\ne j$ and a monomial in three pairwise distinct
variables.
\end{proposition}

\begin{proof}
As in P-F27, the eigenvalue character
$\psi:\mathbb Z^n\to\mathbb C^*$ of $A$ annihilates
the lattice generated by the support columns, and
$\psi(\mathbf1)=\det A$. Its image is a two-group.
The $l$ cube columns $3e_i$ are linearly independent.
Include all of them in a selection of $r$ support columns
forming a basis of the rational column span, and write
$C$ for the resulting $n\times r$ matrix. There is a
unique rational vector $c$ with
\[
 Cc=\mathbf1.
 \tag{P-F29b}
\]

Choose a nonzero $r\times r$ minor $C_I$ by selecting
rows. Those rows must contain all $l$ cube-variable
indices: otherwise a cube column would be zero in the
minor. After reordering rows and columns, expansion in
the cube columns gives
\[
 \det C_I=\pm3^l\det E,
 \tag{P-F29c}
\]
where $E$ is an $(r-l)\times(r-l)$ matrix with entries
zero or one and at most three ones in each column.
For $r=l$ the empty determinant is interpreted as one.
Hadamard's inequality gives
\[
 0<|\det E|\le3^{(r-l)/2}.
 \tag{P-F29d}
\]
For clarity, this use of Hadamard is elementary: the
absolute determinant is the product of the successive
orthogonalized column lengths, each at most the original
Euclidean column length, which here is at most $\sqrt3$.

Let $q$ be the least positive integer such that $qc$ is
integral. Cramer's rule in $C_Ic=\mathbf1_I$ shows that
$q$ divides $|\det C_I|$. Also
$q\mathbf1=C(qc)$ belongs to the support lattice, so
the order of $\psi(\mathbf1)$ divides $q$.
Because that order is a power of two, the odd factor
$3^l$ in (P-F29c) does not contribute, and
\[
 \operatorname{ord}(\det A)
       \le 2^{v_2(\det E)}
       \le |\det E|\le3^{(r-l)/2}.
 \tag{P-F29e}
\]
If $r<n$ or $l>0$, then $r-l\le n-1$.
Since $\sqrt3<2$ and $n-1>0$, this is strictly less
than $2^{n-1}$, proving the first case. This reasoning
allows a rank-deficient support matrix and an arbitrary
two-power order for the matrix $A$ itself.

In the remaining case $r=n$ and $l=0$, the selected
matrix $C$ is square and every column has exactly three
ones. Thus
$\mathbf1^tC=3\mathbf1^t$, so reduction modulo three
makes its rows dependent and $3\mid\det C$.
The two-part of a nonzero integer divisible by three
is at most one third of its absolute value. Therefore
\[
 \operatorname{ord}(\det A)
    \le 2^{v_2(\det C)}
    \le\frac{|\det C|}{3}
    \le3^{n/2-1}<2^{n-1}.
 \tag{P-F29f}
\]
The last strict inequality follows immediately for
$n=2$, and for $n>2$ from
$3^{(n-2)/2}<2^{n-2}<2^{n-1}$.
This proves (P-F29a).

For the quotient consequence, P-F26 supplies the finite
two-power lifts and P-F27 supplies the rational-span
hypothesis from semistability in every coordinate basis.
Absence of a three-distinct-variable monomial would put
the support in P-F27's graph class and force a smooth
chain, contrary to discriminant vanishing. Absence of
an $x_i^2x_j$ monomial would put it in the support class
of this proposition and contradict the primitive
determinant order. Hence both types must occur.
\end{proof}

\paragraph{Remaining support class.}
This is a complementary exclusion, not a classification
of arbitrary singular semistable cubics. The remaining
representatives must mix square-linear monomials with
three-distinct-variable monomials. By the connected-graph
extension of P-F27, their square-linear graph must also
be disconnected. The interaction of triple monomials
with these different graph components is the remaining
exponent-lattice problem at this step. P-F30 below resolves it
by contracting the graph components and bounding the resulting
triple-column minors. P-F31 then obtains the actual pole.
No historical priority is
claimed for the elementary determinant estimate.


\subsection{The full cubic support theorem}

\begin{theorem}[P-F30: extremal two-power determinant character forces the chain]
\label{thm:fullcubicsupport}
Let $n\ge2$, let $F\in\mathbb C[x_1,\ldots,x_n]$ be
a cubic, and let $B$ be the matrix of exponent vectors
of its nonzero monomials. Suppose
$\mathbf1=(1,\ldots,1)^t\in\operatorname{im}_{\mathbb Q}B$.
If a diagonal matrix $A$ of finite two-power order
fixes $F$ exactly, then
\[
 \operatorname{ord}(\det A)\le2^{n-1}.
 \tag{P-F30a}
\]
If equality holds, then, after a permutation of variables,
the support of $F$ is exactly
\[
 x_1^2x_2,\ x_2^2x_3,\ldots,\ x_{n-1}^2x_n,\ x_n^3.
 \tag{P-F30b}
\]
Its coefficients can be made one by a diagonal substitution,
and $A$ itself has order $2^{n-1}$. In particular, the
form is smooth. The bound is attained by the chain.

The rational-span hypothesis holds for every
$SL_n$-semistable cubic. Thus this theorem applies
to every eigenbasis of the finite lifts in P-F26;
it imposes no prior bound on the order of the lift.
\end{theorem}

\begin{proof}
The exponent character
\[
 \psi:\mathbb Z^n\longrightarrow\mathbb C^*,
 \qquad (a_i)\longmapsto\prod_i\alpha_i^{a_i},
\]
where $\alpha_i$ are the eigenvalues of $A$, annihilates
the support lattice and has two-power image. Its value
on $\mathbf1$ is $\det A$.
We will use the following elementary denominator fact
throughout. If $K$ is a selection of support columns
forming a basis of the rational support span, then
$Kz=\mathbf1$ has a unique rational solution $z$.
For any nonzero maximal square minor $K_I$, Cramer's
rule shows that the least positive integer $d$ for
which $dz$ is integral divides $|\det K_I|$.
Since $d\mathbf1=K(dz)$ is in the support lattice,
the order of $\det A$ divides $d$. Consequently
\[
 \operatorname{ord}(\det A)\le2^{v_2(\det K_I)}.
 \tag{P-F30c}
\]
This argument requires neither full rank of $B$
nor a bound on the orders of the individual $\alpha_i$.

Form the graph whose edges are the supported monomials
$x_i^2x_j$, with coefficient two at $i$ and one at $j$,
and whose loops are the supported cubes. Include all
$n$ vertices, even those isolated in this graph.
Let $G$ be the matrix of these graph-support columns.
Every other support column has the form
$e_i+e_j+e_k$ for three distinct vertices.

In a connected graph component on $v$ vertices,
a spanning-tree matrix has rank $v-1$. Thus the
graph-component matrix has rank $v$ or $v-1$.
This includes an isolated vertex: it has rank one
if its cube occurs, and rank zero otherwise.
Write $f$ for the number of full-rank components,
$c$ for the number of rank-deficient components,
and $v_i$ for their vertex counts. Then
\[
 \operatorname{rank}G=n-c.
 \tag{P-F30d}
\]

Choose a square support basis $C_i$ in each full-rank
component by adjoining one independent graph column
to a spanning tree. The unicyclic determinant
calculation in P-F27 gives
\[
 v_2(\det C_i)\le v_i-1.
 \tag{P-F30e}
\]
For precision, leaf expansion contributes a factor
one or two for each edge off the unique cycle.
A loop contributes the odd factor three.
A cycle of length $l\ge2$ contributes, up to sign,
$2^a+(-1)^{l-1}2^{l-a}$; if nonzero, its valuation
is at most $\lfloor(l-1)/2\rfloor$. These alternatives
give (P-F30e), including loops and double edges.

Choose a spanning-tree matrix $T_i$ in each deficient
component. Its signed cofactor vector has the form
\[
 2^{e_i}u_i,\qquad
 (u_i)_v=\varepsilon_i(-2)^{h_{iv}},
 \quad \varepsilon_i\in\{1,-1\},\quad
 \min_v h_{iv}=0.
 \tag{P-F30f}
\]
All exponents here are nonnegative integers.
Indeed every maximal tree minor is, up to sign,
a power of two, by successive leaf expansion.
The cofactor relation along an edge with
coefficients two and one is $2u_v+u_w=0$.
Thus adjacent heights differ by one and their
signs are opposite; connectivity gives the
single sign convention in (P-F30f).
For a single isolated deficient vertex, the
empty minor is one and we use $e_i=h_{iv}=0$.

Put
\[
 b_i=v_i-1-e_i,\qquad H_i=\max_v h_{iv}.
\]
Since a tree minor is the product of $v_i-1$
factors, each one or two, one has
\[
 0\le H_i\le b_i.
 \tag{P-F30g}
\]
We need one additional equality property: if
$H_i=b_i$, the vertex of height $H_i$ is unique.
To prove it, choose such a vertex as root.
Its deleted-row minor has absolute value
$2^{e_i+H_i}=2^{v_i-1}$. Expand this minor by
removing leaves toward that root. Every factor
must be two. Thus each edge has coefficient
two at its child and one at its parent.
The cofactor relation gives
$h_{\rm parent}=h_{\rm child}+1$.
Heights strictly decrease away from the root,
so no other vertex has its height.
The assertion is also immediate for an isolated vertex.

Normalize the entries in each of these cofactor
vectors by $2^{b_i}$. For sums over distinct
vertices of one deficient component, the following
bounds hold:
\[
 \left|\sum_{v\in E}\frac{(u_i)_v}{2^{b_i}}\right|
 \le
 \begin{cases}
  1,&|E|=1,\\
  5/4,&|E|=2,\\
  3/2,&|E|=3.
 \end{cases}
 \tag{P-F30h}
\]
Here a case is used only when the component has
enough distinct vertices. If $H_i<b_i$, every
entry has absolute value at most $1/2$, which
proves all three bounds. If $H_i=b_i$, there is
only one entry of absolute value one. Change
the overall sign to make this entry positive.
Every entry of absolute value $1/2$ is negative,
because its height is $H_i-1$; every remaining
entry has absolute value at most $1/4$.
A pair avoiding the top entry has absolute sum
at most one; a pair containing it has absolute
sum at most $5/4$. A triple avoiding it has
absolute sum at most $3/2$. For a triple containing
it, either the other two entries are both $-1/2$,
giving zero, exactly one is $-1/2$, giving absolute
sum at most $3/4$, or neither is, giving at most
$1+1/4+1/4=3/2$. This proves (P-F30h).

The chosen full-component bases and deficient-component
trees together form a rational basis for the columns
of $G$. Extend it to a basis for all support columns
by adjoining $q$ independent triple columns
$\beta_1,\ldots,\beta_q$. Then
\[
 0\le q\le c,\qquad
 r:=\operatorname{rank}B=n-c+q.
 \tag{P-F30i}
\]
If $q=0$, the graph subform itself satisfies the
rational-span hypothesis. P-F27 applies directly:
it gives (P-F30a), and equality gives the chain
eigenvalue list and the chain graph support.
That eigenvalue list admits no other cubic
monomials, so there can be no additional triple
monomials in $F$. This proves all the required
claims in the case $q=0$.
We therefore assume $q>0$ for the rest of the argument.

Define the $c\times q$ compressed matrix
\[
 D_{ij}=u_i^t(\beta_j|_i),
 \tag{P-F30j}
\]
where $\beta_j|_i$ is the restriction to the vertices
of deficient component $i$. The map on the entire
vertex space that sends a vector to these $c$
cofactor pairings is surjective, and its kernel
is precisely the rational graph-column span:
on a full component the graph columns span all
coordinates, while on a deficient component the
kernel of $u_i^t$ is the tree-column span.
Consequently $D$ has rank $q$. Choose $q$
deficient-component indices $J$ for which the
square matrix $D_J$ has nonzero determinant.

Let $K$ be the $n\times r$ original support-basis
matrix just selected. We now specify an actual
nonzero $r\times r$ minor of it. Keep every row
in each full component and in each deficient
component indexed by $J$. In every other deficient
component, delete one row at a vertex of height zero,
and keep the remaining rows. The resulting row
set $I$ has size $n-(c-q)=r$. Its determinant is
\[
 \det K_I
   =\pm\left(\prod_{\text{full }i}\det C_i\right)
          2^{\sum_{\text{deficient }i}e_i}
          \det D_J.
 \tag{P-F30k}
\]
We justify this factorization explicitly.
In a full component the graph basis is square;
eliminating the triple-column entries there
using its graph columns and expanding factors
out $\det C_i$. In a deficient component not
in $J$, the retained tree block is square.
Its deleted row had height zero, so its determinant
is $\pm2^{e_i}$; eliminate the triple entries
there and factor out that determinant as well.
These column eliminations do not alter other
components, since every graph column is supported
in its own component.
The remaining matrix consists of the tall
tree blocks for $i\in J$ and $q$ triple columns.
Expansion by the maximal minors of each tree
block leaves one row pairing for each such
component. Its signed cofactor vector is
$2^{e_i}u_i$. The remaining $q\times q$
determinant is exactly the pairings (P-F30j),
giving the remaining factors
$2^{\sum_{i\in J}e_i}\det D_J$.
This proves (P-F30k), and proves that the
specified original minor is nonzero.

Divide row $i$ of $D_J$ by $2^{b_i}$.
Each resulting column is formed from at most
three distinct vertices of the original graph.
If they belong to three selected components,
its squared Euclidean norm is at most three.
If two lie in one selected component and one
in another, (P-F30h) gives squared norm at most
\[
 (5/4)^2+1=41/16<3.
\]
If all three lie in one component, the squared
norm is at most $(3/2)^2<3$. Ignoring vertices
in unselected or full components only leaves
cases with fewer than three contributing vertices,
which satisfy the same bound. Every normalized
column therefore has Euclidean norm at most
$\sqrt3$. Hadamard's determinant inequality yields
\[
 0<|\det D_J|
       \le 2^{\sum_{i\in J}b_i}3^{q/2}.
 \tag{P-F30l}
\]

Suppose first that $q<c$, or that $q=c$ and $f>0$.
Using (P-F30e), (P-F30k), and (P-F30l) gives
\begin{align*}
 v_2(\det K_I)
 &\le \sum_{\text{full }i}(v_i-1)
      +\sum_{\text{deficient }i}e_i
      +\sum_{i\in J}b_i+\frac q2\log_2 3\\
 &=n-f-c-\sum_{\substack{\text{deficient }i\\i\notin J}}b_i
              +\frac q2\log_2 3\\
 &<n-f-c+q.
 \tag{P-F30m}
\end{align*}
The last inequality is strict because $q>0$
and $\sqrt3<2$; the omitted $b_i$ are nonnegative.
If $q<c$, its right side is at most $n-1$.
If $q=c$ and $f>0$, it is likewise at most $n-1$.
The valuation is an integer, so in both cases
\[
 v_2(\det K_I)\le n-2.
 \tag{P-F30n}
\]

The only case remaining is $q=c$ and $f=0$.
Here $K$ is a square $n\times n$ support matrix,
and $I$ is its full row set. Every original
support column has sum three. Hence
$\mathbf1^tK=3\mathbf1^t$, and reduction modulo
three shows that $3$ divides its nonzero determinant.
By (P-F30k), which now has no full-component
factors, $3$ also divides $\det D$.
The two-part of this determinant is consequently
at most $|\det D|/3$. Since
$\sum_i(e_i+b_i)=\sum_i(v_i-1)=n-c$,
we obtain
\begin{align*}
 2^{v_2(\det K)}
 &=2^{\sum_i e_i}\,2^{v_2(\det D)}\\
 &\le 2^{\sum_i e_i}\frac{|\det D|}{3}\\
 &\le 2^{n-c}3^{c/2-1}
  <\frac{2^n}{3}
  <2^{n-1}.
 \tag{P-F30o}
\end{align*}
Here $c=q>0$ and $3^{c/2}<2^c$ justify the
first strict inequality. Thus (P-F30n) holds
in this case as well.

Combining (P-F30c) and (P-F30n) proves the
stronger bound
$\operatorname{ord}(\det A)\le2^{n-2}$
whenever $q>0$. All extremal cases therefore
occur in the already treated case $q=0$,
where P-F27 proves precisely (P-F30b), the
coefficient normalization, and the exact order
of $A$. Conversely the chain is fixed by
the diagonal matrix with eigenvalues
$\xi,\xi^{-2},\ldots,\xi^{(-2)^{n-1}}$
for a primitive $2^{n-1}$th root $\xi$.
Its determinant exponent
$1-2+4-\cdots+(-2)^{n-1}$ is odd, so its
determinant has order $2^{n-1}$. The bound
is therefore attained.

Finally, P-F27's coefficient-monomial argument
shows that any positive-degree $SL_n$ invariant
nonzero at $F$ expresses $\mathbf1$ as a
nonnegative rational combination of the
support exponents. This supplies the
rational-span hypothesis for every semistable
cubic, in every coordinate basis, and completes
the proof.
\end{proof}

\paragraph{The structural conclusion.}
The theorem closes the mixed-support case left by
P-F27 and P-F29. It uses the actual determinant
character, rather than identifying that character's
order with an arbitrary lift order. Its proof
allows rank-deficient exponent matrices and uses
an explicitly constructed maximal minor in those
cases. Combining it with P-F26 rules out
nonzero singular representatives in the
$\boldsymbol\mu_{2^{n-1}}$ fixed quotient locus.
The ensuing global fixed-locus and Hilbert-pole
conclusions are stated separately.


\subsection{Actual poles at roots of exponential order in ordinary plethysm}

\begin{theorem}[P-F31: the chain pole refutes the positive-form clause of PH3]
\label{thm:actualcubicpole}
Let $n\ge4$, put
\[
 g=\gcd(n,3),\qquad d_0=n/g,\qquad
 M=2^{n-1},\qquad m=\binom{n+2}{3},
\]
and use the ordinary plethysm ray of P-F19,
\[
 \mu=(3),\qquad \lambda=(d_0),\qquad
 \pi=((3/g)^n).
\]
Write
\[
 f_n(s)=\dim\operatorname{Hom}_{GL_n}
 \left((\det)^{3s/g},\operatorname{Sym}^{d_0s}
                    (\operatorname{Sym}^3\mathbb C^n)\right),
 \qquad
 H_n(t)=\sum_{s\ge0}f_n(s)t^s.
 \tag{P-F31a}
\]
Let $T^{(n)}$ be its canonical Cartan multiplicity algebra,
$X=\operatorname{Spec}T^{(n)}$, and let $Y\subset X$ be the
reduced closure of the chain grading orbit. Then:

\begin{enumerate}
\item The support of the fixed scheme $X^{\boldsymbol\mu_M}$
is exactly $Y$. Away from the graded origin this fixed scheme
is the smooth chain orbit, and $X$ is smooth there. Possible
nilpotent or embedded structure at the origin is not excluded.
\item The ordinary Hilbert series $H_n(t)$ has a simple pole
at every primitive $M$th root of unity $\zeta$. More precisely,
with $g_\zeta$ the grading action,
\[
 \lim_{t\to\zeta}(1-t^M)H_n(t)
   =\frac{1}{\det(1-g_\zeta\mid N^*_{Y/X,\eta_Y})}
   \ne0.
 \tag{P-F31b}
\]
\item In every expression
\[
 H_n(t)=\frac{Q(t)}{\prod_{j=1}^{r}(1-t^{a_j})^{b_j}},
 \qquad Q(t)\in\mathbb C[t],\quad a_j,b_j\in\mathbb Z_{>0},
 \tag{P-F31c}
\]
some denominator weight $a_j$ is divisible by $M$. Thus every
such expression, in particular every positive rational form,
has maximum denominator weight at least $2^{n-1}$.
\end{enumerate}
Since the source parameter is
\[
 h=\dim GL_m+\operatorname{ht}(\lambda)+\operatorname{ht}(\pi)
   =\binom{n+2}{3}^{\!2}+1+n=O(n^6),
 \tag{P-F31d}
\]
the first, positive-form existence clause of GCT VI,
Hypothesis~1.6.7 is false for this ordinary irreducible-inner
plethysm family. The stronger canonical parameter clause was
already refuted by P-F19. No superpolynomial positivity or
saturation index is asserted here.
\end{theorem}

\begin{proof}
By P-F19 the actual canonical algebra, including multiplication,
is the entire cubic invariant ring
\[
 T^{(n)}\simeq
   \bigoplus_{s\ge0}
      \mathbb C[(\operatorname{Sym}^3\mathbb C^n)^*]^{SL_n}_{d_0s},
 \tag{P-F31e}
\]
with the coefficient grading divided by $d_0$. Scalar central
characters force every nonzero invariant degree to be a multiple
of $d_0$, so this is a regrading of the whole ring, rather than
a smaller subring. In particular $T^{(n)}$ is a finitely generated
positively graded complex domain with degree-zero part $\mathbb C$.
Its Hilbert series is exactly (P-F31a).

Consider a nonzero closed point $p$ of the fixed scheme.
P-F26 gives a representative cubic $F$ with closed $SL_n$ orbit
and an exact-form stabilizer element $A$ of finite two-power order
whose determinant has order $M$. Since $p$ is not the graded
origin, some positive-degree invariant is nonzero at $F$.
The supported cubic exponent vectors therefore span
$\mathbf1$ rationally, by the coefficient-monomial argument
in P-F27. The complete support theorem P-F30 applies to this
$F$ and $A$, without assuming that the lift has order $M$
or that the support is a graph. Its equality case identifies
$F$, after a permutation and diagonal rescaling in an eigenbasis,
with the smooth chain cubic $F_n$.

Consequently $p$ lies on the $GL_n$ orbit of the chain quotient
point. On the quotient this is exactly the grading orbit, by
(P-F26b) and the surjectivity of the character
$z\mapsto z^{-3/g}$ of $\mathbb G_m$. Conversely P-F24 shows
that the entire nonzero chain grading orbit is fixed by
$\boldsymbol\mu_M$.

The graded origin is also fixed and is in the closure of this
orbit. Indeed the map $u\mapsto u\cdot p_n$ extends from
$\mathbb G_m$ to $\mathbb A^1$: each positive-degree homogeneous
function restricts to a constant times a positive power of $u$,
and thus takes value zero at $u=0$. It follows that the closed
points of the fixed scheme are exactly those of $Y$. A
finite-type closed subset over $\mathbb C$ is determined by
its closed points, so the support of $X^{\boldsymbol\mu_M}$
is $Y$. This is an irreducible curve.

By P-F24 the ambient quotient $X$ and the fixed scheme are
smooth at every point of the nonzero chain orbit, and the fixed
scheme there has dimension one. Since these are all its
nonzero closed points, any nonreduced structure or embedded
associated point is confined to the origin. It is unnecessary
to assert reducedness at that origin: the fixed scheme has a
unique maximal-dimensional component $Y$, and $X$ is smooth
at its generic point. These are precisely the hypotheses
needed in P-F23.

Apply P-F23 with dimension $\delta=1$. It gives
\[
 \lim_{t\to\zeta}(1-t^M)H_n(t)
   =\frac{e_Y}{\det(1-g_\zeta\mid N^*_{Y/X,\eta_Y})}.
 \tag{P-F31f}
\]
All conormal eigenvalues are nontrivial. The exact grading
stabilizer and semigroup calculation in P-F28 give $e_Y=1$.
Thus the right hand side is nonzero, proving (P-F31b).
The factor $1-t^M$ has a simple zero at $\zeta$, so this is
an actual pole of order one of $H_n$.

This application includes the possible structure at the
origin. In the Koszul proof of P-F23, terms supported at
that zero-dimensional point have finite-dimensional graded
homology and hence polynomial Hilbert series. They vanish
after multiplication by $1-t^M$ in the limit. They cannot
cancel the one-dimensional generic contribution. No global
smoothness or reducedness of the fixed scheme was used.

For additional precision, P-F28 now computes the actual
normalized leading coefficient in (P-F31b). If the cyclotomic
valuation above two is normalized by $v(1-\zeta)=1$, it has
valuation
\[
 \binom{n+2}{2}-3\,2^{n-1}.
 \tag{P-F31g}
\]
The ordinary residue of $H_n(t)\,dt$ at $t=\zeta$ is
$-\zeta/M$ times the right hand side of (P-F31b), because
$(1-t^M)'|_{t=\zeta}=-M\zeta^{-1}$. The nonvanishing
in (P-F31b), rather than any comparison of phases, is
what is needed below.

Finally suppose (P-F31c) is any rational expression of
the indicated form. Its numerator is regular at $\zeta$.
If none of its denominator factors vanished there, the
expression would also be regular there, contradicting
the proved pole. Hence $\zeta^{a_j}=1$ for some $j$.
Primitivity of $\zeta$ is exactly the assertion that
$M\mid a_j$. This argument does not require a nonnegative
numerator, a prescribed number of denominator factors,
or a presentation arising from an h.s.o.p.

The exact source definition of a positive form is
GCT VI, Definition~1.2.7, printed p.~15, equation~(1.4);
its modular index is the maximum of the weights $a_j$.
Hypothesis~1.6.7, printed p.~25, first asks for the
existence of such a form with modular index polynomial
in $\dim G+\operatorname{ht}(\lambda)+\operatorname{ht}(\pi)$,
and then makes the stronger assertion for its canonical
minimal parameter form. The parameter in this family is
(P-F31d), whereas every allowed form has a weight at
least $2^{n-1}$. For any fixed polynomial in $h$, its
value at (P-F31d) is bounded by a fixed polynomial in
$n$, which is eventually smaller than $2^{n-1}$.
Therefore even the first existence clause fails.
\end{proof}

\paragraph{Consequences and remaining scope.}
The result closes the ordinary-pole gate that P-F19 alone
left open. It uses the full cubic support theorem and the
unique generically smooth fixed component; it does not
infer an ordinary pole merely from an orbit degree period,
a large parameter degree, a universal denominator, or
cohomological purity. The cancellation examples P-F21--P-F22
remain valid and explain why those shorter arguments
would not have sufficed.

Every quasipolynomial period of $f_n$ is divisible by $M$.
Indeed a quasipolynomial of period $L$ has a generating
function with denominator a power of $1-t^L$: sum its
constituent polynomials over each arithmetic progression,
using the elementary generating functions for powers.
The pole in (P-F31b) forces $M\mid L$. This does not say
that $M$ is the entire minimal period or classify its
other prime-power factors.

The implication from bounded positive-form weights to
PH2 in P-F12 remains correct, but failure of its premise
does not refute PH2 or SH. Nor do large denominator
weights alone establish a counting-complexity lower bound,
an obstruction search algorithm, or a failure of every
possible revised geometric formulation. Here the inner
shape is the single irreducible cubic shape $(3)$, and
only the number of variables grows; reducible torus
embeddings and determinant twists are not responsible
for the obstruction. No historical priority claim is
made for the support classification or the resulting
Hilbert-pole consequence.

The exact primary source for the conjecture and its
positive-form convention is
\url{https://arxiv.org/pdf/0704.0229},
Definition~1.2.7, p.~15, and Hypothesis~1.6.7, p.~25.


\subsection{An exterior-power reduction for the surviving index question}

\begin{proposition}[P-F32: rectangular complement and the first nonzero stretch]
\label{prop:cubicexteriorduality}
For $n\ge1$, put $g=\gcd(n,3)$ and
\[
 f_n(s)=\dim\left(\operatorname{Sym}^{ns/g}
                         (\operatorname{Sym}^3\mathbb C^n)\right)^{SL_n}.
\]
For $s\ge1$, write $k=3s/g$, $d=ns/g=nk/3$, and
$N(k)=(k-1)(k-2)/2$. Then
\[
 f_n(s)=\operatorname{mult}_{(\det)^n}
       \left(\bigwedge^d\left(\bigwedge^3\mathbb C^k\right)\right).
 \tag{P-F32a}
\]
In particular $f_n(s)=0$ if $n>N(k)$ and $f_n(s)=1$ if $n=N(k)$.
If $1\le n<N(k)$, put $n'=N(k)-n$, $g'=\gcd(n',3)$
and $s'=kg'/3$. Then $s'$ is an integer and
\[
                         f_n(s)=f_{n'}(s').
 \tag{P-F32b}
\]
Consequently, for each $k\ge5$ and $n=\binom{k-1}{2}$,
the first positive stretch is exactly
\[
 s_* = \frac{\gcd(n,3)k}{3},\qquad
 f_n(s)=0\ (1\le s<s_*),\qquad f_n(s_*)=1.
 \tag{P-F32c}
\]
For these $n\ge6$, the strict saturation and positivity
indices in GCT VI satisfy
\[
                     p(f_n)\ge s(f_n)\ge s_*-1.
 \tag{P-F32d}
\]
This is a polynomial lower bound, of order $\sqrt n$.
It does not refute the narrower polynomial-index PH2 or SH(a).
\end{proposition}

\begin{proof}
Every summand of $\operatorname{Sym}^d(\operatorname{Sym}^3
\mathbb C^n)$ has total weight $3d=nk$. Its only possible
polynomial $GL_n$ summand restricting trivially to $SL_n$
is $(\det)^k$, of highest weight $(k^n)$. Thus $f_n(s)$
is the coefficient of $s_{(k^n)}$ in $h_d[h_3]$.

Let $\omega$ be the involution with
$\omega(s_\lambda)=s_{\lambda'}$. The conjugation rule gives
\[
                         \omega(h_d[h_3])=e_d[e_3].
 \tag{P-F32e}
\]
An exact primary locator is De Boeck--Paget--Wildon,
\emph{Plethysms of symmetric functions and highest weight
representations}, Lemma~2.2, printed p.~11; Proposition~3.3, p.~21,
identifies plethysms with the corresponding Schur-functor
characters. Here the needed identity also has a direct
power-sum proof. Write $p_\alpha$ for the power-sum product,
$\ell(\alpha)$ for partition length and $z_\alpha$ for the
centralizer factor. The identities
\[
 h_j=\sum_{\alpha\vdash j}\frac{p_\alpha}{z_\alpha},
 \quad
 e_j=\sum_{\alpha\vdash j}
          \frac{(-1)^{j-\ell(\alpha)}p_\alpha}{z_\alpha},
 \quad \omega(p_a)=(-1)^{a-1}p_a
\]
imply $\omega(p_a[h_3])=(-1)^{a-1}p_a[e_3]$:
for the term indexed by $\alpha\vdash3$, the respective
signs differ by $(-1)^{3a-3}=(-1)^{a-1}$.
Expanding $h_d$ then proves (P-F32e).
The conjugate rectangle is $(n^k)$. Specialization to
$k$ variables preserves its coefficient, since Schur
functions of length exceeding $k$ vanish and the others
remain linearly independent. The character $e_d[e_3]$
belongs to $\bigwedge^d(\bigwedge^3\mathbb C^k)$, and
$s_{(n^k)}$ belongs to $(\det)^n$. This proves (P-F32a),
also when $k=1,2$ and the exterior cube is zero.

For $k\ge3$, set $E=\mathbb C^k$ and $W=\bigwedge^3E$.
Its dimension is $L=\binom{k}{3}=kN(k)/3$. Thus $d>L$
exactly when $n>N(k)$, proving the vanishing assertion.
Moreover
\[
                       \det W=(\det E)^{N(k)}:
 \tag{P-F32f}
\]
each diagonal variable occurs in precisely $\binom{k-1}{2}$
exterior-cube basis weights, and one-dimensional $GL(E)$
characters are determined on the diagonal torus.
When $n=N(k)$, (P-F32a) is the multiplicity of this
character in itself, hence one.

When $d<L$, the perfect exterior pairing gives
\[
 \bigwedge^dW\simeq
        \left(\bigwedge^{L-d}W\right)^*\otimes\det W.
 \tag{P-F32g}
\]
The multiplicity of $(\det E)^n$ therefore equals that
of $(\det E)^{N(k)-n}$ in $\bigwedge^{L-d}W$.
Here $L-d=kn'/3$ is integral. If $3\nmid n'$ this
forces $3\mid k$, and otherwise $g'=3$. Thus
$s'=kg'/3$ is integral and $n's'/g'=L-d$.
Applying (P-F32a) with $n',s'$ proves (P-F32b).

For $n=N(k)$, the identity $kn=3\binom k3$ similarly
makes $s_*=gk/3$ integral. The equality case just proved
gives $f_n(s_*)=1$. For $1\le s<s_*$ the integer
$k_s=3s/g$ is smaller than $k$, and $N(k_s)<N(k)$,
including $k_s=1,2$. The vanishing criterion proves
(P-F32c).

To obtain the strict-index bound, we must exclude an
identically zero constituent on the residue of the
displayed realized zero. B\"urgisser--Ikenmeyer,
\emph{Fundamental invariants of orbit closures},
arXiv:\texttt{1511.02927v2}, Theorem~2.3, printed p.~5,
gives reduced generic stabilizer period one for cubic
forms in $n>3$ variables. Together with generic polystability
(Proposition~2.10, p.~7), their Theorem~3.4, p.~8, and
the generic degree-monoid discussion on p.~9 show that
the group generated by nonzero ordinary invariant
degrees is $(n/g)\mathbb Z$. The greatest common divisor
of the corresponding stretch degrees is therefore one.

The invariant ring is a finitely generated domain.
The degrees of a finite nonzero homogeneous generating
set generate exactly its support semigroup: products
of such generators are nonzero, and every homogeneous
element is a polynomial in them. A finitely generated
semigroup of nonnegative integers with greatest common
divisor one contains all sufficiently large integers.
Indeed, for any positive member $a$, its residues form
a subsemigroup, hence a subgroup, of $\mathbb Z/a\mathbb Z$.
The greatest-common-divisor condition makes this the
whole group. Choose one semigroup representative per
residue and add multiples of $a$.
It follows that $f_n(s)>0$ for all sufficiently large
$s$, so every constituent polynomial is nonzero.

The actual zero at $s_*-1\ge1$ prevents strict saturation
after any shift $u<s_*-1$: the appropriate shifted
constituent is nonzero and vanishes at the positive
integer $s_*-1-u$. Hence $s(f_n)\ge s_*-1$.
A nonzero polynomial with nonnegative coefficients
has positive value at every positive integer, so
$p(f_n)\ge s(f_n)$. This proves (P-F32d).
\end{proof}

\paragraph{Scope and next use.}
The complement identity reduces coefficients near the
exterior-dimension boundary to cubic invariants in
$N(k)-n$ variables. It supplies no general polynomial
upper bound for either index. The conjugation and
exterior-duality identities are standard, and no
historical priority is claimed for these deductions.
Primary source URLs are
\url{https://kar.kent.ac.uk/87807/1/PlethysmByWeightsKAR.pdf}
and \url{https://arxiv.org/pdf/1511.02927}.



\subsection{Large actual poles and fixed reciprocity do not force large indices}

\begin{proposition}[P-F33: actual period with zero strict indices]
\label{prop:largeperiodzeroindices}
For every integer $M\ge3$, give
$\mathbb C[x,y,z,w]$ the weights $(1,1,1,M)$ and set
\[
 R_M=\frac{\mathbb C[x,y,z,w]}
 {\,w(xy-z^2)+x^{M+2}+y^{M+2}\,}.
 \tag{P-F33a}
\]
This is a normal Cohen--Macaulay Gorenstein domain of
dimension three, with rational singularities and
$a$-invariant $-1$. Its Hilbert series is
\[
 H_M(t)=\frac{1-t^{M+2}}{(1-t)^3(1-t^M)},
 \qquad H_M(t^{-1})=-tH_M(t).
 \tag{P-F33b}
\]
It has a simple pole at every primitive $M$th root.
Nevertheless, if $f_M(s)=\dim(R_M)_s$, its global
quasipolynomial constituents of period $M$ are
\[
 P_r(X)=\left(\frac12+\frac1M\right)X(X+1)
                +\frac{(r+1)(M-r)}M,
 \qquad 0\le r<M.
 \tag{P-F33c}
\]
Their coefficients are all positive. In the strict
definitions of GCT VI,
\[
 p(f_M)=s(f_M)=0,
 \tag{P-F33d}
\]
while the minimum quasipolynomial period is exactly $M$.
Thus even actual poles at roots of arbitrarily large
multiplicative order, together with fixed
Krull dimension, normality, rational singularities,
Gorenstein symmetry, and a fixed negative $a$-invariant,
does not force an increasing positivity or saturation index.
No identification with an ordinary plethysm algebra is asserted.
\end{proposition}

\begin{proof}
Write $Q=xy-z^2$ and $F=x^{M+2}+y^{M+2}$.
The polynomial $Q$ is an irreducible quadratic in
$\mathbb C[x,y,z]$, and it does not divide $F$:
the point $(x,y,z)=(1,0,0)$ has $Q=0$ and $F=1$.
Hence $Q$ and $F$ are relatively prime. The polynomial
$wQ+F$ is primitive and linear in $w$ over this
unique factorization domain. It is irreducible over
its fraction field, so Gauss's lemma makes it
irreducible in the polynomial ring. It is therefore
prime, proving that $R_M$ is a domain of dimension three.
As a hypersurface it is Cohen--Macaulay and Gorenstein.

We check its singular locus rather than infer normality.
The partial derivatives of its equation are
\[
 wy+(M+2)x^{M+1},\quad
 wx+(M+2)y^{M+1},\quad -2wz,\quad xy-z^2.
 \tag{P-F33e}
\]
If $w\ne0$, their vanishing first gives $z=0$ and
$xy=0$. If $x=0$, the first derivative gives $y=0$;
if $y=0$, the second gives $x=0$.
If $w=0$, the first two derivatives give $x=y=0$,
and the fourth then gives $z=0$.
Conversely all four derivatives vanish on the entire
$w$-axis. The singular locus is exactly that axis.
It has codimension two, so the hypersurface is regular
in codimension one. Its Cohen--Macaulay property gives
$S_2$, and Serre's criterion proves normality.

The equation has weighted degree $M+2$ and the sum
of the four variable weights is $M+3$. The graded
hypersurface resolution, or equivalently its graded
$\operatorname{Ext}^1$ into the canonical module
of the polynomial ring, gives
\[
 \omega_{R_M}\simeq R_M((M+2)-(M+3))=R_M(-1).
 \tag{P-F33f}
\]
Thus $a(R_M)=-1$. The same one-step resolution
gives (P-F33b); direct substitution $t\mapsto t^{-1}$
verifies its displayed reciprocity identity.

Here are details verifying rational singularities.
Away from the origin, the only singular points
have $w\ne0$. On this open set, blow up the ideal
$(x,y,z)$. The $x$-chart of the strict transform,
with $y=xY$ and $z=xZ$, has equation
\[
 w(Y-Z^2)+x^M(1+Y^{M+2})=0.
 \tag{P-F33g}
\]
At its exceptional divisor $x=0$, the derivative
with respect to $Y$ is the unit $w$.
The $y$-chart is symmetric. In the $z$-chart,
$x=zX,\ y=zY$, the equation is
\[
 w(XY-1)+z^M(X^{M+2}+Y^{M+2})=0.
 \tag{P-F33h}
\]
At $z=0$, one has $XY=1$, so a derivative with
respect to $X$ or $Y$ is nonzero. Thus these charts
are smooth along the exceptional divisor, and
away from it the blowup is the identity on the
already smooth complement of the axis.
They give a resolution of $\operatorname{Spec}(R_M)_w$.

The exceptional divisor is the smooth conic
$xy-z^2=0$ in $\mathbb P^2\times\mathbb G_m$,
isomorphic to $\mathbb P^1\times\mathbb G_m$.
Its conormal line bundle is the restriction of
$\mathcal O_{\mathbb P^2}(1)$, hence is
the pullback of $\mathcal O_{\mathbb P^1}(2)$.
Writing $E$ for this divisor, the successive
quotients of its ideal powers are therefore
the pullbacks of $\mathcal O_{\mathbb P^1}(2j)$
for $j\ge0$.
Each has zero cohomology in positive degrees:
the base $\mathbb G_m$ is affine and
$H^i(\mathbb P^1,\mathcal O(2j))=0$ for $i>0$.
The exact sequences for successive thickenings
show the same vanishing for every infinitesimal
thickening of $E$.

By the theorem on formal functions, the completions
along the axis of all higher direct-image modules
of the structure sheaf under this resolution are
zero. These modules are coherent and supported on
the axis, since the resolution is an isomorphism
off it. A power of the axis ideal therefore
annihilates each module, so completion does not
change it. The modules themselves are zero.
This proves rationality on the punctured spectrum.

The graded rational-singularity criterion now applies:
Kurano--Sato--Singh--Watanabe,
\emph{Multigraded rings, diagonal subalgebras, and
rational singularities}, Theorem~5.6, printed p.~3262,
states that a normal $\mathbb N$-graded algebra of
finite type over a characteristic-zero field is
rational when it is Cohen--Macaulay, its punctured
spectrum is rational, and its $a$-invariant is negative.
All three hypotheses have just been checked.
Its formulation allows a local degree-zero ring;
here that ring is the field $\mathbb C$.
This proves the claimed rational singularities,
including at the graded origin.

We next compute every constituent and its strict index.
Write $s=Mk+r$ with $k\ge0$ and $0\le r<M$.
Using the power-series expansion of (P-F33b),
pair each summand after the first with the
corresponding shifted numerator summand. This gives
\begin{align*}
 f_M(s)
 &=\binom{s+2}{2}
      +\sum_{j=1}^{k}
           \left(\binom{s-jM+2}{2}
                        -\binom{s-jM}{2}\right)\\
 &=\binom{s+2}{2}
      +\sum_{j=1}^{k}\bigl(2(s-jM)+1\bigr)\\
 &=\binom{s+2}{2}+k(2s+1)-Mk(k+1).
 \tag{P-F33i}
\end{align*}
The pairing also holds for $r=0$ and $r=1$:
the formal expressions $\binom02$ and $\binom12$
are both zero, exactly matching the absent
negative-degree terms. Thus no initial exceptional
values have been omitted.
Substituting $k=(s-r)/M$ into (P-F33i) gives
exactly $P_r(s)$ in (P-F33c).
This proves global quasipolynomiality, including $s=0$.

The quadratic and linear coefficients of every
$P_r$ are $1/2+1/M>0$, and its constant coefficient
is $(r+1)(M-r)/M>0$. Hence all constituent coefficients
are nonnegative without any shift, which is the
definition of strict positivity in GCT VI,
Definition~1.2.2, printed p.~13. Moreover every
$P_r(X)$ is positive for every real $X\ge0$,
and in particular for every integer $X\ge1$.
This verifies strict saturation in Definition~1.2.4,
printed p.~14, on all constituent values, not only
on the residue class on which that constituent
represents $f_M$. Both indices are therefore zero.

Finally, if $\zeta$ is primitive of order $M$,
the numerator of (P-F33b) at $\zeta$ is
$1-\zeta^2\ne0$, since $M\ge3$.
The factor $1-t^M$ has a simple zero there
and $1-\zeta\ne0$. This is an actual simple pole.
Every quasipolynomial period must consequently
be divisible by $M$, by the elementary
generating-function argument used in P-F31.
Formula (P-F33c) already supplies period $M$,
so this is exactly the minimum period.
\end{proof}

\paragraph{Primary algebraic references and scope.}
Serre's criterion is Stacks Project
Lemma~10.157.4, Tag~031S; the complete-intersection
Gorenstein property is Lemma~47.21.8, Tag~0DWA.
The cohomology used for the exceptional conic is
Lemma~30.8.1, Tag~01XT, and the formal-functions
step is Theorem~30.20.5, Tag~02OC.
The rational-singularity criterion is available at
\url{https://www.math.utah.edu/~singh/publications/kurano_sato_singh_watanabe.pdf}.
The strict-index definitions are in
\url{https://arxiv.org/pdf/0704.0229}.

This proposition settles a side question about a
proposed inference from poles and reciprocity.
It neither bounds nor disproves PH2 or SH for the
ordinary cubic plethysm sequence in P-F31.
Any such conclusion needs additional
representation-specific information about its
constituent polynomials.

One further obstruction to an elementary shortcut is
visible in the same formula. The actual coefficients
satisfy $f_M(s)=\binom{s+2}{2}$ for every $0\le s<M$,
but the average of the constituent constant terms is
\[
 \frac1M\sum_{r=0}^{M-1}P_r(0)
       =\frac{(M+1)(M+2)}{6M}.
 \tag{P-F33j}
\]
To check the sum, expand
$(r+1)(M-r)=M(r+1)-r(r+1)$ and use
the finite sums of $r$ and $r^2$.
It grows linearly with $M$, despite the unchanged
initial coefficients. Thus those initial values alone
do not bound the off-residue constituent average
that would enter a Fourier comparison with a pole
coefficient.


\subsection{A direct Fourier gate for the surviving index hypotheses}

\begin{proposition}[P-F34: a large cubic Fourier amplitude and the remaining average-polynomial gate]
\label{prop:cubicfourierindexgate}
For the ordinary cubic plethysm function $f_n$ of P-F31, let
$L$ be any quasipolynomial period and write
\[
 f_n(s)=q_r(s)\quad\hbox{when }s\equiv r\pmod L,
 \qquad
 B_n(x)=\frac1L\sum_{r=0}^{L-1}q_r(x).
 \tag{P-F34a}
\]
The polynomial $B_n$ is independent of the chosen period.
It is the contribution from the pole of $H_n(t)$ at $t=1$.
Its degree is
\[
 D=\binom{n+2}{3}-n^2=\binom n3.
\]
Put $M=2^{n-1}$ and
\[
 \mathcal C_n(\zeta)=\lim_{t\to\zeta}(1-t^M)H_n(t),
 \qquad A_n(\zeta)=\mathcal C_n(\zeta)/M
\]
for a primitive $M$th root $\zeta$. Then the following are
necessary conditions on the source's literal indices:
\[
 B_n\bigl(p(f_n)\bigr)\ge |A_n(\zeta)|,
 \qquad
 B_n\bigl(s(f_n)+1\bigr)\ge |A_n(\zeta)|.
 \tag{P-F34b}
\]
In the standard complex embedding
$\zeta=\exp(2\pi i/M)$, the amplitude satisfies, for every
$n\ge250$,
\[
                         |A_n(\zeta)|\ge2^{n^4/192}.
 \tag{P-F34c}
\]

Write $B_n(x)=\sum_{j=0}^{D}b_{n,j}x^j$ and
$U_n=\sum_j|b_{n,j}|>0$. Consequently
\[
 \max\{1,p(f_n)\},\quad s(f_n)+1
      \ \ge\ \left(\frac{2^{n^4/192}}{U_n}\right)^{1/D}
       \qquad(n\ge250).
 \tag{P-F34d}
\]
In particular, an additional estimate
\[
 U_n\le2^{C n^3\log_2(n+2)}
 \tag{P-F34e}
\]
with an absolute constant $C$ would imply, for all sufficiently
large $n$,
\[
 p(f_n)\ge\frac{2^{n/32}}{(n+2)^{6C}},\qquad
 s(f_n)+1\ge\frac{2^{n/32}}{(n+2)^{6C}}.
 \tag{P-F34f}
\]
Estimate (P-F34e) is not proved or assumed here. Thus the
proposition establishes a quantitative necessary test and
a sufficient remaining condition for an index refutation;
it does not itself refute narrow PH2 or SH.
\end{proposition}

\begin{proof}
We first give the general Fourier calculation, keeping the
normalization of a simple pole explicit. If $f$ is a
quasipolynomial of period $L$, with constituent polynomials
$q_r(x)$, define for each $L$th root of unity $\rho$
\[
 P_\rho(x)=\frac1L\sum_{r=0}^{L-1}\rho^r q_r(x).
 \tag{P-F34g}
\]
Finite Fourier inversion gives
$q_r(x)=\sum_{\rho^L=1}\rho^{-r}P_\rho(x)$, and hence
\[
 \sum_{s\ge0}f(s)t^s
   =\sum_{\rho^L=1}\sum_{s\ge0}P_\rho(s)(t/\rho)^s.
 \tag{P-F34h}
\]
For a nonzero polynomial $P$ of degree $a$, the generating
function $\sum_{s\ge0}P(s)z^s$ has its only pole at $z=1$,
of order $a+1$. One proof is to expand $P(s)$ in the
basis $\binom{s+j}{j}$, whose generating functions are
$(1-z)^{-j-1}$; the coefficient of the top basis member
is nonzero. Thus a simple pole of the left side at
$\rho=\zeta$ means precisely that $P_\zeta$ is a nonzero
constant. If $\zeta$ is primitive of order $M\mid L$,
its summand is $P_\zeta/(1-t/\zeta)$, and
\[
 \lim_{t\to\zeta}
       \frac{1-t^M}{1-t/\zeta}=M.
 \tag{P-F34i}
\]
For the family at hand, P-F31 therefore proves
\[
                   P_\zeta(x)=A_n(\zeta)
                         \quad\hbox{for every }x.
 \tag{P-F34j}
\]
There is no dependence on $x$ in this Fourier coefficient.

The coefficient for $\rho=1$ is $B_n$. Enlarging the
period only repeats the constituent polynomials, so this
average is intrinsic. Formula (P-F34h) also shows directly
that its generating function is the principal part of
$H_n$ at $t=1$. More explicitly, if that principal part is
\[
 \sum_{j=1}^{D+1}\frac{c_{n,j}}{(1-t)^j},
 \quad\hbox{then}\quad
 B_n(x)=\sum_{j=1}^{D+1}c_{n,j}
                        \binom{x+j-1}{j-1}.
 \tag{P-F34k}
\]
No exponential-period enumeration is needed to specify
this one polynomial once the Laurent coefficients at one
are known. No efficient procedure for finding them is
asserted.
The ring in P-F31 has dimension
$m-n^2+1$, with $m=\binom{n+2}{3}$, by the smooth tangent
calculation in P-F24. Its Hilbert series has a pole of
that order at one. This is the usual Hilbert dimension
statement for a finitely generated positively graded
domain, obtainable by taking a common Veronese and its
Hilbert polynomial. Consequently $B_n$ has degree
$D=m-n^2=\binom n3$.

For any real $x$ at which all $q_r(x)$ are nonnegative,
(P-F34g) and the triangle inequality imply
\[
 |A_n(\zeta)|
   =\left|\frac1L\sum_r\zeta^r q_r(x)\right|
   \le\frac1L\sum_r q_r(x)=B_n(x).
 \tag{P-F34l}
\]
After a positivity shift $p=p(f_n)$, every polynomial
$q_r(x+p)$ has nonnegative coefficients. Its constant
coefficient $q_r(p)$ is therefore nonnegative.
After a saturation shift $s=s(f_n)$, every nonzero
polynomial $q_r(x+s)$ is positive at every positive
integer $x$; at $x=1$ it is positive, while an identically
zero constituent has value zero. Applying (P-F34l) at
$p$ and $s+1$ proves (P-F34b).
These are the literal conventions of GCT VI,
Definition~1.2.2, printed p.~13, and Definition~1.2.4,
printed p.~14. In particular the saturation definition
requires positivity of each nonzero constituent at every
positive integer, not only at integers in its original
residue class.

We now prove the absolute-value estimate, rather than
trying to infer it from the cyclotomic valuation alone.
Put $g=\gcd(n,3)$, $\gamma=3/g$, and
$S_0=(1-(-2)^n)/3$. Choose
\[
 \zeta=e^{2\pi i/M},\qquad \xi=\zeta^{-g},\qquad
 A=\operatorname{diag}
        (\xi,\xi^{-2},\ldots,\xi^{(-2)^{n-1}}).
 \tag{P-F34m}
\]
The induced grading parameter is exactly $\zeta$:
\[
 \xi^{-\gamma S_0}=\zeta^{g\gamma S_0}
                 =\zeta^{1-(-2)^n}=\zeta.
\]
Here $g$ is odd and at most three, so $\xi$ is still
primitive of order $M$.
Let $w_i=(-2)^i$ for $0\le i\le n-2$, and let
$w_{n-1}=0$; these are integer representatives for the
chain weights modulo $M$.

The exact moving tangent sequence in P-F24 and P-F28
gives the following absolute-value identity:
\[
 |\mathcal C_n(\zeta)|
   =\frac{\displaystyle\prod_{i\ne j}
                     |1-\xi^{w_i-w_j}|}
          {\displaystyle\prod_{\substack{0\le i\le j\le k<n\\
                    w_i+w_j+w_k\not\equiv0\pmod M}}
                     |1-\xi^{w_i+w_j+w_k}|}.
 \tag{P-F34n}
\]
Indeed the moving cubic representation minus the moving
adjoint representation is the moving quotient tangent
representation, whose dual is the conormal space.
Determinants multiply in the resulting exact sequences.
Replacing a root-of-unity eigenvalue by its inverse
leaves $|1-\alpha|$ unchanged, so possible dual or
point-action conventions do not affect (P-F34n).
There are $n(n-1)$ numerator factors and $m-n$
denominator factors, since precisely the $n$ chain
cubic monomials are fixed.

For every nontrivial $M$th root $\alpha$,
\[
                         |1-\alpha|\ge4/M.
 \tag{P-F34o}
\]
To verify this, choose its angle with distance
$2\pi a/M$ from an integer multiple of $2\pi$,
where $1\le a\le M/2$. Then
$|1-\alpha|=2\sin(\pi a/M)\ge4a/M\ge4/M$,
using concavity of sine on $[0,\pi/2]$.
All numerator factors in (P-F34n) are nontrivial,
because the chain eigenvalues are distinct.

Let $q=\lfloor n/2\rfloor$, and use only cubic monomials
in the first $q$ variables, indexed $0,\ldots,q-1$.
There are
\[
 L_q=\binom{q+2}{3}-(q-1)
 \tag{P-F34p}
\]
moving monomials among them: exactly the $q-1$ consecutive
chain monomials are fixed. For each of these moving
monomials, its integer weight $b$ satisfies
$|b|\le3\cdot2^{q-1}$. Therefore
\[
 |1-\xi^b|=|1-e^{-2\pi igb/M}|
    \le \frac{2\pi g|b|}{M}
    \le \frac{9\pi\,2^q}{M}.
 \tag{P-F34q}
\]
Every other denominator factor is at most two.
Combining (P-F34n)--(P-F34q), and dividing by $M$
to pass from $\mathcal C_n$ to $A_n$, gives
\begin{align*}
 \log_2|A_n(\zeta)|
 &\ge L_q\bigl(n-q-\log_2(9\pi)\bigr)
       -n(n-1)(n-3)-m+1.
 \tag{P-F34r}
\end{align*}
This estimate is entirely in the specified complex
embedding.

For completeness, explicit elementary constants yield
(P-F34c). For $n\ge10$ one has
\[
 L_q\ge\frac{q^3}{6}\ge\frac{(n-1)^3}{48},\qquad
 n-q-\log_2(9\pi)\ge n/2-5\ge0,
\]
since $9\pi<32$. Substitution in (P-F34r) and expansion
of $m=n(n+1)(n+2)/6$ give
\[
 \log_2|A_n(\zeta)|
 \ge
 \frac{n^4-125n^3+369n^2-351n+106}{96}.
 \tag{P-F34s}
\]
For $n\ge250$, the first two terms in the numerator
are at least $n^4/2$, and
$369n^2-351n+106>0$. Thus the right side is at least
$n^4/192$. This proves (P-F34c).

Finally, for $x\ge0$,
\[
 |B_n(x)|\le U_n\max\{1,x\}^{D}.
\]
Together with (P-F34b)--(P-F34c), this proves
(P-F34d). If the additional bound (P-F34e) holds,
then for sufficiently large $n$ the quantity
$n^4/192-Cn^3\log_2(n+2)$ is positive. Since
$D=\binom n3\le n^3/6$, (P-F34d) gives
\[
 \log_2\max\{1,p(f_n)\},\quad\log_2(s(f_n)+1)
       \ge n/32-6C\log_2(n+2).
\]
The right hand side tends to infinity, so the maximum
with one can be removed for all sufficiently large
$n$. This proves (P-F34f), as a conditional consequence
only.
\end{proof}

\paragraph{The exact remaining estimate.}
The pole-at-one polynomial in (P-F34k) is a smaller
piece of data than the entire family of constituent
polynomials over its exponentially large period.
An estimate specific to the cubic invariant ring,
such as (P-F34e), would combine with the proved large
Fourier amplitude to refute polynomial index bounds.
No such estimate is supplied by P-F34. Its necessary
inequality is not an equivalence to PH2 or SH.

In particular $f_n(0)=1$, or bounds on finitely many
small-stretch coefficient values, do not bound $U_n$:
those values sample different constituents and do not
control their values away from their residue classes.
P-F33 supplies a concrete warning: it has actual large
primitive-root amplitudes and a growing average constant,
yet both strict indices are zero, even with normality,
rational singularities, Gorensteinness and negative
$a$-invariant. Thus no generic purity, normality or
Gorenstein argument can replace the remaining
ring-specific estimate. Large cyclotomic valuation
also cannot replace the absolute-value calculation
above, because the complex conjugate amplitudes may
have very different sizes.

The exact source for the index conventions is
\url{https://arxiv.org/pdf/0704.0229}, Definitions~1.2.2
and 1.2.4, printed pp.~13--14. The Fourier identity
and the quantitative bounds required here are proved
in full above; no historical priority is claimed.


\paragraph{A sharper asymptotic amplitude and a weaker sufficient estimate.}
The same product gives the more precise asymptotic formula
\[
 \log_2|A_n(e^{2\pi i/M})|=\frac{n^4}{24}+O(n^3).
 \tag{P-F34t}
\]
Here is a full proof, including control of the factors that are
not close to one. For $n\ge6$, set $K=n-6$ and first use only
monomials in variables with nonzero chain weights. For a monomial
whose largest index is $k$, the number of moving monomials is
\[
 N_0=1,\qquad N_k=\binom{k+2}{2}-1=\frac{k^2+3k}{2}
                      \quad(k\ge1).
\]
The subtraction removes precisely $x_{k-1}^2x_k$.
For any remaining monomial with this largest index, its integer
weight $b$ satisfies
\[
                       2^{k-2}\le |b|\le3\cdot2^k.
\]
To check the lower bound, if the largest index occurs two or
three times, the largest term dominates the other terms by at
least $3\cdot2^{k-1}$. If it occurs once and both smaller indices
are at most $k-2$, the triangle inequality gives $2^{k-1}$.
If one smaller index is $k-1$, the sum of these two terms is
$(-2)^k/2$; the other smaller term has magnitude at most
$2^{k-2}$, unless both smaller indices are $k-1$, which is
exactly the removed fixed monomial. The case $k=0$ is the
single cube and satisfies the displayed bounds directly.

For $k\le K$, $g|b|\le9\cdot2^{n-6}\le M/2$.
The same elementary sine bounds as above therefore give
\[
 \frac{2^k}{M}\le |1-\xi^b|
       \le\frac{18\pi\,2^k}{M}<\frac{64\,2^k}{M}.
\]
In particular the base-two logarithm of each such factor is
$k-n+\epsilon$ with $1\le\epsilon\le7$.
The main contribution to the negative logarithm of the cubic
moving determinant is consequently
\begin{align*}
 S_n&=\sum_{k=0}^{n-6}(n-k)N_k\\
    &=\frac{n^4+6n^3-181n^2+798n-720}{24}
      =\frac{n^4}{24}+O(n^3).
\end{align*}
The exact equality follows from the standard sums of the first
three powers of $k$ after substituting
$N_k=(k^2+3k)/2$ for $k\ge1$ and $N_0=1$.
The total logarithmic error in these bulk factors is at most
$7m=O(n^3)$.

There are at most $3n^2$ remaining cubic factors. Indeed a
monomial either contains the last variable of weight zero,
giving at most $\binom{n+1}{2}$ choices, or its largest index
is one of $n-5,n-4,n-3,n-2$, giving at most
$4\binom n2$ choices. Their sum is less than $3n^2$.
Each moving factor, regardless of angle, lies between $4/M$
and two, so its logarithm has absolute value at most $n$.
Thus all these exceptional factors together contribute
$O(n^3)$. The $n(n-1)$ adjoint factors obey the same bound
and also contribute $O(n^3)$. Finally division by $M$ changes
the logarithm by only $n-1$. This proves (P-F34t) with an
absolute implied constant; for example the estimates just
given imply absolute error at most $12n^3$ for $n\ge6$.

It follows that the sufficient remaining condition can be
weakened from (P-F34e) to
\[
 \limsup_{n\to\infty}\frac{\log_2 U_n}{n^4}<\frac1{24}.
 \tag{P-F34u}
\]
To verify the implication with its quantifiers, choose
$\delta>0$ such that
$\log_2 U_n\le(1/24-2\delta)n^4$ for every sufficiently
large $n$. Formula (P-F34t) then gives
$\log_2(|A_n|/U_n)\ge\delta n^4$ for all sufficiently
large $n$. The already proved index inequalities and
$D\le n^3/6$ imply
\[
                   p(f_n),\quad s(f_n)+1\ge2^{6\delta n}
\]
for all sufficiently large $n$; the maximum with one in
(P-F34d) disappears because this lower bound exceeds one.
Condition (P-F34u), like (P-F34e), remains unproved for the
actual cubic ring. P-F38 excludes deriving it from bounded
constraint entries, cubic-scale dimension, and a bounded
containing box alone. The sharper amplitude therefore
strengthens the precise remaining condition without settling
that condition or either surviving index hypothesis.


\subsection{Exact indices for the four-variable cubic calibration}

\begin{proposition}[P-F35: both quaternary-cubic indices are twenty-three]
\label{prop:quaternaryindices}
For the ordinary cubic-invariant stretching function of P-F31 with
$n=4$,
\[
 f_4(s)=\dim\left(\operatorname{Sym}^{4s}
                  (\operatorname{Sym}^3\mathbb C^4)\right)^{SL_4},
\]
the strict positivity and saturation indices of GCT VI are exactly
\[
                          p(f_4)=s(f_4)=23.
 \tag{P-F35a}
\]
Its minimal quasipolynomial period is $120$. In particular the
primitive eighth-root poles of P-F31 occur in the classical
explicit series and coexist with these finite indices. This
single dimension gives no asymptotic index bound.
\end{proposition}

\begin{proof}
Salmon's invariant-ring theorem gives algebraically independent
invariants of ordinary degrees $8,16,24,32,40$, together with
one invariant of degree $100$ whose square lies in the polynomial
ring they generate. An exact inspected locator is Huybrechts,
\emph{The geometry of cubic hypersurfaces}, Chapter~4,
Theorem~4.1, printed p.~196; the following proof of Corollary~4.2
also states the resulting graded decomposition. The ordinary
degrees in the polynomial subring are divisible by eight, whereas
those in its degree-$100$ translate are congruent to four modulo
eight. Thus the sum of these two modules is direct; it spans
by the monic quadratic relation. Dividing ordinary degrees by
four gives
\[
 H_4(t)=\sum_{s\ge0}f_4(s)t^s
       =\frac{1+t^{25}}{(1-t^2)(1-t^4)(1-t^6)(1-t^8)(1-t^{10})}.
 \tag{P-F35b}
\]
This also directly checks the nonzero primitive eighth-root
pole: the factor $1-t^8$ vanishes simply there, the other four
factors do not vanish, and $1+\zeta^{25}=1+\zeta\ne0$.

Let $P(y)$ be the global quasipolynomial with generating series
\[
                  \sum_{j\ge0}P(j)z^j
                       =\frac1{\prod_{a=1}^5(1-z^a)}.
\]
We will use its constituent polynomials explicitly. Put
\[
 Q(y)=\frac{y^4}{2880}+\frac{y^3}{96}
            +\frac{31y^2}{288}+\frac{85y}{192}
            +\frac{50651}{86400}.
\]
For an integer residue $r$, define the periodic numbers
\[
 a_r=\begin{cases}1&r\equiv0,1\pmod4,\\-1&r\equiv2,3\pmod4,\end{cases}
 \quad
 b_r=\begin{cases}2&3\mid r,\\-1&3\nmid r,\end{cases}
 \quad
 c_r=\begin{cases}4&5\mid r,\\-1&5\nmid r.\end{cases}
\]
Then the period-$60$ constituent on $r$ is
\[
 P_r(y)=Q(y)+\frac{(-1)^r(2y+15)}{128}
                   +\frac{a_r}{16}+\frac{b_r}{27}+\frac{c_r}{25}.
 \tag{P-F35c}
\]
For completeness, this formula follows by coefficient extraction
from the following exact partial-fraction identity, which can also
be checked by multiplying through by $\prod_{a=1}^5(1-z^a)$:
\begin{align*}
 \frac1{\prod_{a=1}^5(1-z^a)}
 &=\frac1{120(1-z)^5}+\frac1{24(1-z)^4}
   +\frac{31}{288(1-z)^3}+\frac3{16(1-z)^2}
   +\frac{20831}{86400(1-z)}\\
 &\quad+\frac1{64(1+z)^2}+\frac{13}{128(1+z)}
   +\frac{z+1}{16(z^2+1)}
   +\frac{z+2}{27(z^2+z+1)}\\
 &\quad+\frac{z^3+2z^2+3z+4}
                  {25(z^4+z^3+z^2+z+1)}.
 \tag{P-F35d}
\end{align*}
The first line gives $Q$ after using the binomial-series formula.
The two $1+z$ terms give $(-1)^r(2y+15)/128$.
The last three periodic terms give respectively $a_r/16$,
$b_r/27$ and $c_r/25$: multiplication by $1-z$ reduces
the last two numerators to $2-z-z^2$ and
$4-z-z^2-z^3-z^4$. This proves (P-F35c) for every
nonnegative integer argument and hence identifies the whole
constituent polynomials.

Every polynomial $P_r(y-1)$ has nonnegative coefficients.
Indeed its coefficients of degrees four, three and two are
respectively
\[
              \frac1{2880},\qquad\frac{13}{1440},
                       \qquad\frac{113}{1440}.
\]
Its linear coefficient is
\[
             \frac{247}{960}+\frac{(-1)^r}{64}
                     \ \ge\ \frac{29}{120}>0.
\]
Its constant term is
\[
 \frac{20831}{86400}+\frac{13(-1)^r}{128}
                        +\frac{a_r}{16}+\frac{b_r}{27}
                        +\frac{c_r}{25}.
\]
The sum $13(-1)^r/128+a_r/16$ takes the four values
$21/128,-5/128,5/128,-21/128$ on residues modulo four.
Since $b_r\ge-1$ and $c_r\ge-1$, the constant term is at least
\[
 \frac{20831}{86400}-\frac{21}{128}-\frac1{27}-\frac1{25}=0.
 \tag{P-F35e}
\]
All coefficient inequalities have therefore been verified
for every constituent without an interpolation assumption.

Equation (P-F35b) says that the even constituent polynomials
of $f_4$ are of the form $P_r(s/2)$ and the odd ones are
of the form $P_r((s-25)/2)$. To justify the latter globally,
for each odd residue modulo $120$ the equality holds for
all sufficiently large nonnegative $s$; thus it is the
polynomial identity on that residue. The series is a global
quasipolynomial, either by the canonical-ring theorem already
used or directly because its numerator over $(1-t^{120})^5$
has degree $595<600$.

After shifting by $23$, the odd polynomials become
$P_r(s/2-1)$ and therefore have nonnegative coefficients
by (P-F35e) and its preceding calculations. The even
polynomials become $P_r((s+23)/2)$, obtained from the
coefficientwise nonnegative polynomial $P_r(y-1)$ by
substituting $y=(s+25)/2$. A polynomial with nonnegative
coefficients remains such after substitution of an affine
polynomial with nonnegative coefficients. Thus $p(f_4)\le23$.

Every constituent has positive degree-four coefficient
$1/46080$, so none is the zero polynomial. On the other
hand (P-F35b) gives the actual value $f_4(23)=0$:
the denominator contains only even degrees, and the odd
numerator translate first starts at $25$. For any shift
$u<23$, the corresponding nonzero constituent still
vanishes at the positive integer $23-u$. Hence strict
saturation fails for all such shifts, giving
$s(f_4)\ge23$. Coefficientwise positivity implies strict
saturation for every nonzero constituent on positive
integers, so
\[
                   23\le s(f_4)\le p(f_4)\le23.
\]
This proves (P-F35a).

Finally $120$ is a period by the denominator in (P-F35b).
Its actual primitive eighth-root, third-root and fifth-root
poles force every quasipolynomial period to be divisible
by $8,3,5$, respectively. At a primitive third or fifth root,
only the factors of the corresponding divisible degrees
can vanish and the numerator is nonzero ($1+t^{25}$ equals
$1+t$ at a third root and equals two at a fifth root).
Thus their least common multiple $120$ is the minimal period.
\end{proof}

The exploratory exact rational computation is preserved in
\texttt{quaternary\_index\_calibration.py} and its output file.
The proof above supplies explicit constituent formulas and
uniform inequalities, independently of the finite search.
The exact source URL is
\url{https://www.math.uni-bonn.de/people/huybrech/Notes.pdf}.


\subsection{A concrete descent obstruction in the nonstandard cubic algebra}

The general nonstandard setting is not an ordinary Laurent-polynomial
version of the standard canonical-basis problem. The following obstruction
uses the actual cubic example in GCT VII--VIII, rather than an unrelated
positive-semidefinite countermodel.

\paragraph{Exact known example.}
In GCT VII, Section~7.1.1--7.1.2, printed pp.~18--20,
equations~(19)--(24), let $B_0$ be the nonstandard algebra for
$H=GL_2$, $X=\operatorname{Sym}^3\mathbb C^2$, and tensor degree three,
over $F=\mathbb C(q)$. It has dimension $21$ and a basis of the
alternating monomials $Q_\sigma$ of length at most ten, including
the empty monomial. The displayed splitting extension is
\[
 K=F(x),\qquad x^2=D(q),\qquad
 D(q)=(q^8+1)^2q^{24}A(q),
 \tag{P-F36a}
\]
where
\[
 A(q)=5q^{16}+8q^{12}-4q^{10}+18q^8-4q^6+8q^4+5.
 \tag{P-F36b}
\]
Over $K$, the algebra is a direct sum of one one-dimensional
simple block and five four-dimensional simple blocks.
The blocks $U_1,U_2$ are interchanged by $x\mapsto-x$;
see equation~(22), the definitions following it, and equation~(23).
These are explicit assertions for the worked example, not an
assumption that the conjectural algorithms work in every degree.

GCT VIII repeats the discriminant in Section~5.3, equation~(18),
printed p.~31. In its Figure~6, p.~38, the coefficient of
$Q_{1212121}$ in the displayed element $u_1^1$ is
\[
 c(q)=-\frac{P(q)+x}{2q^{20}},
 \tag{P-F36c}
\]
where
\[
 \begin{split}
 P(q)={}&2q^{30}+3q^{28}+2q^{26}+6q^{24}-2q^{22}
          +22q^{20}-2q^{18}\\
       &+6q^{16}+2q^{14}+3q^{12}+2q^{10}.
 \end{split}
\]
The corresponding element $u_2^1$ is obtained by replacing $x$ by
$-x$. This coefficient is recorded in the monomial basis just
specified; it is not asserted here to be a coefficient in a
different ambient canonical basis.

\begin{proposition}[P-F36: power regrading and rational projectors do not remove the quadratic datum]
\label{prop:nonstandarddescent}
For the displayed nonstandard cubic algebra:
\begin{enumerate}
\item For every positive integer $k$, the rational function $D(t^k)$
is not a square in $\mathbb C(t)$. Consequently, after the power
regrading $q=t^k$, the coefficient $c(t^k)$ is not in
$\mathbb C(t)$. In particular $c(q)$ is not the character of any
finite rationally graded vector space, even allowing arbitrary
complex coefficients or virtual signs. After any power regrading,
no nonzero rational-function scalar normalization makes the
coefficient a finite Laurent character.
\item There is no $F$-linear projector on $B_0$ whose scalar
extension has image exactly $U_1$ or exactly $U_2$. More generally,
if a graded $B_0$-bimodule or the cohomology of a complex of
$B_0$-bimodules is defined over $F$, then after extension to $K$
the multiplicities of $U_1$ and $U_2$ as simple bimodules are equal
in each degree. Thus neither block alone is a descended graded
object of that type.
\end{enumerate}
The same descent conclusions hold after any power regrading.
\end{proposition}

\begin{proof}
We first prove that $A(q)$ is not a square even in $\mathbb C(q)$.
If the square of a rational function is a polynomial, that rational
function has no finite pole and is itself a polynomial. Suppose
therefore that $A(q)=R(q)^2$ with $R\in\mathbb C[q]$.
Since $A$ is even, the integral-domain identity
\[
 (R(-q)-R(q))(R(-q)+R(q))=0
\]
makes $R$ either even or odd. Its constant term is nonzero because
$A(0)=5$, so it is even. Its degree is eight. Write
\[
 R(q)=a q^8+b q^6+c_0 q^4+d q^2+e.
\]
The coefficients at degrees $16,14,12,10$ respectively give
\[
 a^2=5,\qquad b=0,\qquad c_0=4/a,\qquad d=-2/a.
\]
The coefficient at degree eight then gives
\[
 2ae+c_0^2=18,
 \qquad ae=\frac{37}{5}.
\]
But the constant term gives $e^2=5$, and hence
$(ae)^2=a^2e^2=25$, contrary to $ae=37/5$.
Thus $A$ is not a square.

Over $\mathbb C$, a polynomial is a square up to a nonzero scalar
exactly when all its roots have even multiplicity. The scalar
itself has a square root. Since $A$ is not a square and $A(0)\ne0$,
there is a nonzero root $\alpha$ of odd multiplicity. For every
$k\ge1$, each root of $t^k=\alpha$ is nonzero and the derivative
$kt^{k-1}$ is nonzero there. The corresponding root of $A(t^k)$
therefore has the same odd multiplicity. It follows that
$A(t^k)$ is not a square in $\mathbb C(t)$. Multiplication by the
explicit nonzero square factor in (P-F36a) does not change its
square class, so $D(t^k)$ is not a square either.

In the quadratic extension obtained by adjoining a square root
$x_k$ of $D(t^k)$, formula~(P-F36c) reads
\[
 c(t^k)=-\frac{P(t^k)+x_k}{2t^{20k}}.
\]
Its rational coefficient of $x_k$ is nonzero. If this expression
belonged to $\mathbb C(t)$, solving for $x_k$ would put the
square root in that field, a contradiction. Multiplying it by
a nonzero element of $\mathbb C(t)$ cannot change this conclusion.

A finite rationally graded vector space has a character that
becomes a finite Laurent polynomial in $t$ after clearing the
denominators of its grading. A virtual difference gives a signed
Laurent polynomial; allowing arbitrary scalar coefficients still
gives a Laurent polynomial. All lie in $\mathbb C(t)$.
This proves the first assertion. It concerns finite Laurent
characters after power regrading, not arbitrary local Puiseux
series: algebraic functions can have such local expansions.

For the second assertion, let $\tau$ be the nontrivial automorphism
of $K/F$. Scalar extension of any $F$-linear map commutes with
the natural semilinear action of $\tau$. Its image is therefore
$\tau$-stable. But $\tau(U_1)=U_2$ and the two blocks are distinct
summands in the displayed direct decomposition. Hence that image
cannot be either block alone. This already proves the projector
claim without imposing bimodule linearity on the projector.

For the more general claim, scalar extension of an $F$-defined
$B_0$-bimodule carries this semilinear action compatibly with
the algebra. The split semisimple algebra $B_0\otimes_F K$
has distinct simple bimodules $U_i$. Conjugating a decomposition
interchanges the $U_1$ and $U_2$ multiplicities. Since the module
itself is unchanged up to the natural semilinear identification,
uniqueness of the simple decomposition makes these multiplicities
equal. Apply the same reasoning separately in each preserved
grading degree.

Extension by a field is exact. For any complex of $F$-defined
bimodules its cohomology after extension is the extension of its
cohomology before extension. Thus the same argument applies to
every cohomological degree, and also to signed Euler classes.
Finally, the first part proves that the quadratic extension
remains nontrivial after every power regrading, so the entire
descent argument repeats over $\mathbb C(t)$.
\end{proof}

\paragraph{The sharper transfer requirement.}
A direct purity or determinant construction whose output is an
ordinary finite graded multiplicity character cannot identify
the displayed coefficient, even after a fractional grading or
rational scalar normalization. Likewise a construction that stays
over the rational function field and yields an actual bimodule
projector cannot select the two conjugate blocks separately.
It must retain the quadratic extension and its descent datum,
or work with a conjugate sum and then provide an additional
comparison that recovers the desired separate objects.

This does not refute arbitrary theta-lift constructions or
algebraic coefficient systems: those can carry nontrivial field
extensions. It also does not refute GCT VIII's relaxed conjectures,
which explicitly concern coefficients of minimal polynomials,
with permitted signs and factors, rather than identifying every
raw algebraic coefficient with a positive Laurent polynomial;
see Conjectures~2.8--2.9, pp.~16--17, and 3.5--3.6, pp.~24--25.
The supplied Alp\"oge \emph{Remark}, pp.~3--4, steps~4, 8--12,
does not specify a comparison of its automorphic projectors with
these two nonstandard blocks. P-F36 states a necessary condition
on such a comparison; it makes no assertion that the cited
theta lifts lack the required extension or local systems.

The exact primary sources are
\url{https://arxiv.org/pdf/0709.0749}, Section~7.1.1--7.1.2,
pp.~18--20, equations~(19)--(24), and
\url{https://arxiv.org/pdf/0709.0751}, Section~5.3, p.~31,
equation~(18), and Figure~6, p.~38.


\subsection{An actual unitary automorphic realization of one cubic plethysm ring}

There is a precise automorphic dictionary in a fixed dimension that is
stronger than an analogy between two sign arguments. We use the following
known theorem with its exact grading, and then transfer the pole already
proved here.

\paragraph{Known dictionary.}
Looijenga--Swierstra, \emph{The period map for cubic threefolds},
arXiv:\texttt{math/0608279v1}, Theorem~3.1, printed p.~12, identifies
$\mathbb C[\operatorname{Sym}^3(\mathbb C^5)^*]^{SL_5}$ with the
ring of meromorphic automorphic forms on their arithmetic complex
$10$-ball quotient, allowing poles along the specified hyperplane
arrangement. Write $\mathcal M_k$ for its weight-$k$ space.
The isomorphism multiplies automorphic weight by $5/3$.
The central-character explanation on p.~11 gives
$\mathcal M_k=0$ unless $3\mid k$; negative weights vanish
by the theorem. The Hermitian lattice and automorphic line
bundle are defined on pp.~10--11. This is a known ring
isomorphism, including multiplication, rather than a
new construction of an automorphic model.

\begin{corollary}[P-F37: exact automorphic regrading and inherited poles]
\label{cor:automorphiccubic}
With those meromorphic spaces and their multiplication,
\[
 \bigoplus_{s\ge0}\mathcal M_{3s}\simeq T^{(5)},\qquad
 \dim\mathcal M_{3s}=f_5(s).
 \tag{P-F37a}
\]
The left hand side is graded by $s$ in this formula.
In the original automorphic weight, put
\[
 G(u)=\sum_{k\ge0}\dim\mathcal M_k\,u^k.
\]
Then
\[
                         G(u)=H_5(u^3),
 \tag{P-F37b}
\]
and $G$ has an actual simple pole at every primitive
$48$th root $\eta$. More precisely,
\[
 \lim_{u\to\eta}(1-u^{48})G(u)
                          =\mathcal C_5(\eta^3)\ne0.
 \tag{P-F37c}
\]
Consequently every expression for $G$ with polynomial numerator
and denominator factors $1-u^a$ has some $a$ divisible by $48$.
Every homogeneous system of parameters in this automorphic
ring has some weight at least $48$.
\end{corollary}

\begin{proof}
The cited graded isomorphism sends $\mathcal M_{3s}$
to the ordinary coefficient-degree-$5s$ part of the cubic
invariant ring. P-F19 identifies the latter, with its actual
multiplication and degree $s$, with $T^{(5)}_s$.
Composing the two graded ring isomorphisms proves (P-F37a).
The central-character vanishing outside weights divisible
by three then gives (P-F37b) coefficient by coefficient.

For $n=5$, P-F31 has $M=2^4=16$. If $\eta$ is primitive
of order $48$, then $\eta^3$ is primitive of order $16$.
Substitution $t=u^3$ in P-F31b gives (P-F37c) directly,
since $1-t^{16}=1-u^{48}$. The map $u\mapsto u^3$
has nonzero derivative at $\eta$, so a simple pole
remains simple. Equivalently, the nonzero finite limit
in (P-F37c) and the simple zero of $1-u^{48}$ give
this assertion without any coordinate convention.
If none of the denominator weights were divisible by
$48$, every denominator factor would be nonzero at
$\eta$, contradicting this pole. This proves the
assertion about arbitrary polynomial numerators.

Finally an h.s.o.p. in the automorphic ring, with
its original weights divided by three, is an h.s.o.p.
in $T^{(5)}$ under (P-F37a). All its nonzero weights
are divisible by three. P-F19 forces one of the
regraded parameter degrees to be at least $16$,
so its original automorphic weight is at least $48$.
This proves the last claim independently of choosing
a rational Hilbert presentation.
\end{proof}

\paragraph{What this supplies for the theta-lift question.}
This gives an identified unitary automorphic object carrying
one actual ordinary plethysm algebra and all its stretching
coefficients. It also shows that an automorphic realization
preserves, rather than removes, the grading obstructions proved
in the algebraic model. Its dimension is fixed here; no
uniform family in growing $n$ or index estimate is obtained.

A transfer from the supplied note's theta lifts must still
identify their output with these particular weight spaces,
including the permitted arrangement poles, the change of
weight and multiplication. The ball has unitary signature
$(10,1)$, whereas the supplied note's iterated lifts use
its specified totally definite unitary data. No theorem
identifying these constructions is stated in either the
cited period-map theorem or the supplied \emph{Remark}.
This observation does not preclude an additional comparison,
and does not assert that theta lifts can never have the
needed analytic type. It makes the remaining comparison
concrete in a case where an automorphic dictionary is
already available.

The primary source is
\url{https://arxiv.org/pdf/math/0608279}, Theorem~3.1,
p.~12, with the grading discussion on pp.~10--11.


\subsection{Why small constraint entries do not bound the average polynomial}

\begin{proposition}[P-F38: an exponential obstruction for bounded-entry vector partitions]
\label{prop:boundedentryaverage}
For every integer $n\ge24$, put
\[
 M=2^n,\qquad r=2\lfloor n^3/2\rfloor.
 \tag{P-F38a}
\]
There is a rational polytope $P_n$ of dimension $r$ in the
unit box $[0,1]^{r+n+1}$, defined by nonnegativity and $n+1$
independent equations with integer coefficients of absolute
value at most two and right-hand sides zero or one, such that
its Ehrhart quasipolynomial $F_n(s)=\#(sP_n\cap\mathbb Z^{r+n+1})$
has average constituent polynomial
\[
 B_n^{\mathrm{ex}}(X)=\frac1L\sum_{a=0}^{L-1}q_a(X)
      =\sum_{j=0}^{r}b_jX^j
\]
whose coefficient norm satisfies
\[
 \sum_j|b_j|\ \ge\ |b_0|\ \ge\ 2^{n^4/2}.
 \tag{P-F38b}
\]
Here $L$ is any quasipolynomial period. The polytope, its
integral lattice, and the coefficient lower bound are explicit.
The same counting function has an equivalent unit-box
presentation with $r+2n+1$ variables and $2n+1$ independent
equations in which every matrix entry is in $\{0,1,2\}$,
every right-hand side is one, and each column sum is at most
three.
Consequently, a bound of the form
$2^{C n^3\log_2(n+2)}$ for all such average-polynomial norms
does not follow from the dimension, the number of equations,
small constraint entries, or containment in a unit box.
This is not an assertion about the full cubic weight system
or its Weyl alternating sum in P-F34.
\end{proposition}

\begin{proof}
Use coordinates
$x_1,\ldots,x_r,y_0,\ldots,y_n$. Define $P_n$ by
\[
 x_j\ge0,\quad y_i\ge0,\qquad
 y_{i+1}=2y_i\ (0\le i<n),\qquad
 \sum_{j=1}^{r}x_j+y_n=1.
 \tag{P-F38c}
\]
The $n$ chain equations are independent, and the last equation
is independent of them because it has a nonzero $x_1$
coefficient. Positive solutions exist, so the dimension is
$(r+n+1)-(n+1)=r$. Every coordinate is at most one:
$x_j\le1$ and $0\le y_i\le y_n\le1$.

For every nonnegative integer $s$, the integer solutions of
the dilated equations are in bijection with
\[
 x_1,\ldots,x_r,c\in\mathbb Z_{\ge0},
 \qquad \sum_{j=1}^{r}x_j+Mc=s.
 \tag{P-F38d}
\]
The forward map takes $c=y_0$; the inverse sets $y_i=2^ic$.
Both maps preserve the integer lattice exactly. Thus this is
not merely a rational equivalence of polytopes.

For the nonnegative-matrix presentation, introduce
$z_0,\ldots,z_{n-1}\ge0$ and replace the $n$ chain equations by
\[
 y_{i+1}+z_i=1,\qquad 2y_i+z_i=1\quad(0\le i<n),
 \qquad \sum_{j=1}^{r}x_j+y_n=1.
 \tag{P-F38d'}
\]
Subtracting the two equations in each pair recovers the
chain. Conversely the chain uniquely gives
$z_i=1-y_{i+1}\in[0,1]$; after dilation this becomes
$z_i=s-y_{i+1}$. Thus adjoining or removing the $z_i$
preserves integer points exactly for every $s$.
These $2n+1$ equations are independent. A linear relation
among their rows has zero coefficient on the last row
by inspecting an $x_j$ column. Each $z_i$ column makes
the two row coefficients in its pair opposite.
The $y_0$ column then kills the first pair, and the
$y_1,y_2,\ldots,y_{n-1}$ columns kill all successive
pairs inductively. The dimension is still $r$.
Every entry is $0$, $1$, or $2$. An internal $y_i$
column has sum three, and all other column sums are
at most two. All coordinates remain in the unit box.

Summing over all these integers gives
\[
 \sum_{s\ge0}F_n(s)t^s
          =\frac1{(1-t)^r(1-t^M)}.
 \tag{P-F38e}
\]
The rational function is proper and all its poles are roots
of unity. Its partial fractions therefore give a global
quasipolynomial, with degree $r$ and a period dividing $M$.
The Fourier argument of P-F34 identifies the average
polynomial with the part belonging to the pole at one.

We prove a quantitative coefficient estimate for (P-F38e)
for every even $r\ge4$ and $M\ge128r$. Write
\[
 T(z)=\frac{z}{1-e^{-z}},\qquad
 T(z)^r=\sum_{j\ge0}a_jz^j,\qquad a_0=1.
 \tag{P-F38f}
\]
The constant coefficient $b_0$ of the average polynomial
is the coefficient of $z^{-1}$ in
$((1-e^{-z})^r(1-e^{-Mz}))^{-1}$.
Indeed
\[
 \sum_{s\ge0}s^je^{-zs}
       =\left(-\frac{d}{dz}\right)^j\frac1{1-e^{-z}}
       =j!z^{-j-1}+\text{a function holomorphic at }z=0.
 \tag{P-F38g}
\]
For $j=0$ the same formula holds with the holomorphic
constant term $1/2$. Terms belonging to poles different
from one are holomorphic near $z=0$, so (P-F38g)
extracts exactly $b_0$.

Let $B_k$ denote the Bernoulli numbers with $B_1=-1/2$.
Their generating series gives
\[
 \frac1{1-e^{-Mz}}
 =\frac1{Mz}+\frac12+
      \sum_{\substack{k\ge2\\ k\ \mathrm{even}}}
                     \frac{B_k}{k!}M^{k-1}z^{k-1}.
\]
Consequently
\[
 b_0=\frac{a_r}{M}+\frac{a_{r-1}}2
       +\sum_{\substack{2\le k\le r\\ k\ \mathrm{even}}}
                 \frac{B_k}{k!}M^{k-1}a_{r-k}.
 \tag{P-F38h}
\]
The leading term as a Laurent polynomial in $M$ is
\[
             L_r=\frac{B_r}{r!}M^{r-1}.
 \tag{P-F38i}
\]
We now bound every other term relative to it.

For $|z|\le1/2$, the exponential series gives
\[
 |(1-e^{-z})-z|
       \le\sum_{k\ge2}\frac{|z|^k}{k!}\le |z|^2.
\]
For example the last bound follows from
$\sum_{k\ge2}u^k/k!\le u^2/(2(1-u))\le u^2$
when $0\le u\le1/2$.
On $|z|=1/r$ this yields
\[
 |T(z)|^r\le(1-1/r)^{-r}\le e^2<8.
\]
The middle inequality follows from
$-\log(1-u)\le2u$ for $0\le u\le1/2$.
Cauchy's coefficient estimate therefore proves
\[
                         |a_j|\le8r^j
                         \quad\hbox{for every }j\ge0.
 \tag{P-F38j}
\]

The even Bernoulli-number formula is
\[
 \frac{|B_k|}{k!}=\frac{2\zeta(k)}{(2\pi)^k}
           \qquad(k\ge2\ \hbox{even}).
 \tag{P-F38k}
\]
It is NIST DLMF equation~25.6.2, equivalently equation~24.8.1
at $x=0$. This exact identity also follows by integrating the
Fourier coefficients of the Bernoulli polynomials by parts.
In particular $\zeta(k)<2$ and $\zeta(r)>1$, so
\[
 \frac{|B_k|/k!}{|B_r|/r!}<2(2\pi)^{r-k}
        \quad(2\le k\le r-2,\ k\ \hbox{even}),\qquad
 |L_r|\ge\frac{2M^{r-1}}{(2\pi)^r}.
\]
The bound $\zeta(k)<2$ needs no asymptotic estimate:
$\sum_{\ell\ge2}\ell^{-k}\le\sum_{\ell\ge2}\ell^{-2}<1$.

Set $\eta=2\pi r/M$. The assumption $M\ge128r$ and
$2\pi<8$ give $\eta<1/16$. For the nonleading Bernoulli
terms in (P-F38h), put $j=r-k$. Equations (P-F38j)--(P-F38k)
give the total relative bound
\[
 \sum_{\substack{2\le k\le r-2\\ k\ \mathrm{even}}}
 \frac{|B_k|M^{k-1}|a_{r-k}|}{k!\,|L_r|}
       \le16\sum_{\substack{j\ge2\\ j\ \mathrm{even}}}\eta^j
       =\frac{16\eta^2}{1-\eta^2}
       <\frac{16}{255}.
 \tag{P-F38l}
\]
The two remaining terms satisfy
\[
 \frac{|a_{r-1}|/2}{|L_r|}
        \le4\pi\eta^{r-1}<16\eta^{r-1},
 \qquad
 \frac{|a_r|/M}{|L_r|}\le4\eta^r.
 \tag{P-F38m}
\]
Since $r\ge4$, the sum of the three relative bounds is
less than
\[
        \frac{16}{255}+\frac{16}{16^3}
                          +\frac4{16^4}<\frac12.
\]
There is thus no cancellation of even half the leading
Bernoulli term. We have proved
\[
              |b_0|\ge\frac12|L_r|
                         \ge\frac{M^{r-1}}{(2\pi)^r}.
 \tag{P-F38n}
\]

For the values (P-F38a), $M\ge128r$ for every $n\ge24$.
At $n=24$, $2^{24}>128\cdot24^3$; thereafter the ratio
$2^n/n^3$ increases, because
$2(n/(n+1))^3>1$ for $n\ge4$.
Also $r\ge n^3-1$ and $\log_2(2\pi)<3$, whence
\[
 \begin{split}
 \log_2|b_0|
 &\ge (r-1)n-r\log_2(2\pi)\\
 &>r(n-3)-n\\
 &\ge n^4-3n^3-2n+3
 >n^4/2\qquad(n\ge24).
 \end{split}
 \tag{P-F38o}
\]
This proves the stated coefficient-norm lower bound.
For each fixed $C$, $n^4/2>Cn^3\log_2(n+2)$ eventually,
so the proposed universal smaller bound is impossible
on this class of bounded-entry systems.
\end{proof}

\paragraph{The two-dimensional warning and the remaining cubic gate.}
Already the triangle
$\operatorname{conv}\{(0,0),(1,0),(0,1/M)\}$ lies inside
the unit square and has Ehrhart series
$1/((1-t)^2(1-t^M))$. Writing $s=Mk+a$, $0\le a<M$,
its exact count is
\[
 \sum_{c=0}^{k}(s-Mc+1)
   =\frac{s^2}{2M}+
       \left(\frac12+\frac1M\right)s+
       1+\frac a2-\frac aM-\frac{a^2}{2M}.
 \tag{P-F38p}
\]
Thus its average constituent constant is
$(M^2+6M+5)/(12M)$, unbounded while the containing box
is fixed. Formula (P-F38c) shows that encoding a large
denominator by small constraint entries does not repair
the issue, and (P-F38b) gives the quantitative
$2^{\Omega(n^4)}$ obstruction at cubic-scale dimension.
If the nonnegative presentation is parameterized instead by
its row rank $N=2n+1$, the stronger displayed bound in
(P-F38o) gives
\[
 \liminf_{n\to\infty}
   \frac{\log_2\bigl(\sum_j|b_j|\bigr)}{(2n+1)^4}
       \ge\frac1{16}>\frac1{24}.
 \tag{P-F38q}
\]
This follows by dividing $n^4-3n^3-2n+3$ by $(2n+1)^4$
and taking its limit. Thus the sharper constant threshold
in (P-F34u) also cannot be justified from these general
features alone.

This example is an ordinary vector-partition problem,
even in the nonnegative-matrix presentation (P-F38d').
It does not have the full list of degree-three
monomial weights and the Weyl alternating numerator
of the actual cubic invariant ring. Those specific
features would have to supply an additional cancellation
or estimate if one seeks the smaller $U_n$ bound in P-F34.
No upper or lower bound on that actual $U_n$ is inferred.
The primary identity used above is available at
\url{https://dlmf.nist.gov/25.6.E2}.


\subsection{The exact minimal-polynomial target in the cubic example}

P-F36 distinguishes the algebraic coefficient from a finite Laurent
character. The worked example itself permits a more precise positive
comparison after retaining its quadratic extension. The positivity
recorded below for this one coefficient was displayed computationally
in GCT VIII; the elementary identities make its scope explicit.

\paragraph{The source's normalization and target.}
Use $A,D,P,x,c$ from P-F36, with the exact displayed discriminant
$D=(q^8+1)^2q^{24}A$. GCT VIII, Section~5.3, pp.~31--32,
defines $v_1^1$ by replacing $Q_\sigma$ with $P_\sigma$ in
$u_1^1$, without changing the scalar coefficients. Thus the coefficient
$c$ of $P_{1212121}$ in $v_1^1$ is precisely (P-F36c).
Section~5.4.1, p.~48, explains that these convenient coefficients must
be rescaled by suitable powers of the positive Laurent polynomial
$\widehat f_p$ displayed in GCT VII, Section~7.1.7, p.~28.
It studies their traces and norms, equivalently their quadratic
minimal-polynomial coefficients. Figure~21, p.~55, and
Figures~22--23, pp.~56--57, display the resulting coefficient vectors.

Set $F=q^{11}\widehat f_p$. The monomial factor merely centers its
exponents. To specify this normalization without a convention about
the ordering of the displayed source terms, write
\[
 \begin{gathered}
 F(q)=f_0+\sum_{j=1}^{16}f_j(q^j+q^{-j}),\\
 (f_0,\ldots,f_{16})=
 (16,16,16,16,17,18,18,18,16,14,12,10,8,6,4,2,1).
 \end{gathered}
 \tag{P-F39a}
\]
In particular $F(1)=400$ and $F$ has positive coefficients.
No unimodality of $F$ itself is needed or asserted.
For a finite vector $v=(v_0,\ldots,v_d)$, abbreviate
\[
 [v](q)=v_0+\sum_{j=1}^d v_j(q^j+q^{-j}).
\]

\begin{proposition}[P-F39: positive minimal-polynomial coefficients and the remaining comparison]
\label{prop:nonstandardtracenorm}
Let $a=Fc\in\mathbb C(q,x)$, with $x^2=D$.
\begin{enumerate}
\item The minimal polynomial of $a$ over $\mathbb C(q)$ is
\[
                   Z^2+T_1(q)Z+N_1(q),
 \qquad T_1=FT,\quad N_1=F^2N,
 \tag{P-F39b}
\]
where
\[
 \begin{split}
 T(q)&=22-2(q^2+q^{-2})+6(q^4+q^{-4})
       +2(q^6+q^{-6})\\
     &\hspace{15mm}+3(q^8+q^{-8})+2(q^{10}+q^{-10}),\\
 N(q)&=138-14(q^2+q^{-2})+70(q^4+q^{-4})
      +20(q^6+q^{-6})\\
     &\quad+31(q^8+q^{-8})+26(q^{10}+q^{-10})
      +6(q^{12}+q^{-12})\\
     &\quad+9(q^{14}+q^{-14})+3(q^{16}+q^{-16})
      +3(q^{18}+q^{-18})+(q^{20}+q^{-20}).
 \end{split}
 \tag{P-F39c}
\]
Both $T_1$ and $N_1$ are symmetric Laurent polynomials with positive
integer coefficients, nonincreasing from the middle towards either
end. Their values at one are respectively $17600$ and $71680000$.
Each is the Tate-degree character of a finite direct sum of
Tate-twisted cohomologies of projective spaces, with their usual
Lefschetz operators.
\item For every positive real $q$, the rational matrix
\[
 M_0(q)=\frac1{2q^{20}}
          \begin{pmatrix}P(q)&D(q)\\1&P(q)\end{pmatrix}
 \tag{P-F39k}
\]
is positive definite and self-adjoint for the positive definite
Gram matrix $\operatorname{diag}(1,D(q))$. Its eigenvalues
are $-c$ and $-c'$. Nevertheless
\[
       \det(I+tM_0)=1+T(q)t+N(q)t^2
 \tag{P-F39l}
\]
has the negative Laurent coefficients displayed in (P-F39c).
The scaled matrix $M=FM_0$ is positive for the same form and has
\[
       \det(I+tM)=1+T_1(q)t+N_1(q)t^2,
 \tag{P-F39m}
\]
whose coefficients are positive Laurent polynomials.
Thus this actual quadratic coefficient admits a positive-operator
and determinant realization, while pointwise operator positivity
alone does not imply Laurent coefficient positivity even in this
example.
\item There is an explicit finite \emph{\'{e}tale} quadratic cover
of a nonempty open subset of the $q$-line whose rank-two pushforward
of the structure sheaf has a multiplication endomorphism with
characteristic polynomial (P-F39b). Nevertheless this endomorphism
cannot be identified with a horizontal endomorphism of a finite
complex local system on that open subset, or after any connected
finite dominant cover. Its trace is nonconstant.
\end{enumerate}
The cohomological character realizations in the first assertion
identify the two rational coefficients separately. They do not
identify a nonstandard block, multiplication map, or theta-lift
projector with those cohomology objects.
\end{proposition}

\begin{proof}
The conjugate of $c$ is $c'=-(P-x)/(2q^{20})$. Therefore
\[
       -(c+c')=P/q^{20}=T,
 \qquad cc'=(P^2-D)/(4q^{40})=N.
 \tag{P-F39d}
\]
Expanding gives exactly (P-F39c). Since $F$ is a nonzero element
of the rational-function field, P-F36 implies that $a=Fc$ is
not rational. It has degree two, so its characteristic polynomial
$(Z-Fc)(Z-Fc')$ is also its minimal polynomial. This proves
(P-F39b).

For completeness, the entire finite coefficient calculation is
given here. Multiplication of (P-F39a) by (P-F39c) gives
$T_1=[\alpha]$ and $N_1=[\beta]$, where
\[
\begin{split}
\alpha={}&(708,708,694,680,673,666,642,618,565,512,464,416,365,314,\\
&262,210,170,130,106,82,62,42,30,18,11,4,2),
\end{split}
 \tag{P-F39e}
\]
and
\[
\begin{split}
\beta={}&(2212002,2205476,2186350,2155076,2112378,2058980,1995675,\\
&1923256,1842975,1756084,1663849,1567536,1468102,1366504,\\
&1263855,1161268,1059834,960644,864459,772040,684050,601152,\\
&523831,452572,387494,328716,276117,229576,188839,153652,\\
&123592,98236,77108,59732,45640,34364,25488,18596,13346,9396,\\
&6493,4384,2894,1848,1142,672,377,196,95,40,15,4,1).
\end{split}
 \tag{P-F39f}
\]
These identities can be checked coefficient by coefficient using
the ordinary convolution rule. More explicitly, put $f_{-j}=f_j$
for $|j|\le16$, and $f_j=0$ otherwise. Put
$t_{\pm 2j}=t_{2j}$ with $(t_0,t_2,\ldots,t_{10})=(22,-2,6,2,3,2)$,
and set the other $t_j$ equal to zero. Define $n_j$ in the same
way from the coefficients in (P-F39c). Then the displayed entries
are exactly
\[
       \alpha_j=\sum_i f_i t_{j-i},\qquad
       \beta_j=\sum_{i,k} f_i f_k n_{j-i-k}.
 \tag{P-F39g}
\]
The support bounds make every sum finite. These formulas,
together with the explicit lists, verify the polynomial identities
without interpolation or numerical approximation. Every entry
of both lists is positive and each list is nonincreasing. Their
values at one can either be summed from the lists or obtained
from $F(1)=400$, $T(1)=44$, and $N(1)=448$.

Here is an actual cohomological realization of the resulting
characters. For any nonincreasing nonnegative integral vector
$v=(v_0,\ldots,v_d)$, put $v_{d+1}=0$ and
$m_j=v_j-v_{j+1}$. Telescoping gives
\[
 [v](q)=\sum_{j=0}^d m_j q^{-j}(1+q+\cdots+q^{2j}).
 \tag{P-F39h}
\]
The coefficient of $q^k$ on the right is
$\sum_{j\ge|k|}(v_j-v_{j+1})=v_{|k|}$, which proves
the identity for every exponent. The cohomology of
$\mathbb P^{2j}$ consists of one Tate class in each degree
$0,2,\ldots,4j$. Its twist by $(j)$ has Tate-degree
character $q^{-j}(1+q+\cdots+q^{2j})$, where
$\mathbb Q_\ell(-i)$ is recorded as $q^i$.
This is the projective-space case of Milne's projective-bundle
formula, \emph{Lectures on \'{E}tale Cohomology}, Theorem~23.2,
printed pp.~139--140. The hyperplane class gives the usual
Lefschetz chain. Thus
\[
       \bigoplus_{j=0}^d
       H^*(\mathbb P^{2j},\mathbb Q_\ell)(j)^{\oplus m_j}
 \tag{P-F39i}
\]
realizes $[v]$ with nonnegative multiplicities and even Tate
cohomology. Apply this to $v=\alpha$ and $v=\beta$.
This assertion concerns the specified direct sums and their
Tate grading, and does not require all summands to have one
common cohomological degree or weight.

We next prove the positive-operator assertions for all $q>0$.
Put $u=q^2+q^{-2}\ge2$ and $v=u-2=(q-q^{-1})^2\ge0$.
Direct expansion gives
\[
 q^{-8}A(q)=5u^4-12u^2-4u+12
          =36+108v+108v^2+40v^3+5v^4>0.
 \tag{P-F39n}
\]
Consequently $D(q)>0$. In $T$, group the constant and the
$q^{\pm2},q^{\pm4}$ terms as
\[
             22-2u+6(u^2-2)=6u^2-2u+10>0.
\]
All remaining terms are positive for $q>0$, so $T>0$ and
$P=q^{20}T>0$. The same grouping in $N$ gives
\[
             138-14u+70(u^2-2)=70u^2-14u-2>0
                         \quad(u\ge2).
\]
For example $70u^2-14u-2=u(70u-14)-2\ge250$.
Again all remaining terms are positive, so $N>0$.
Equation (P-F39d) now yields $P^2-D=4q^{40}N>0$.

Let $G=\operatorname{diag}(1,D)$. Then
\[
       GM_0=\frac1{2q^{20}}
          \begin{pmatrix}P&D\\D&DP\end{pmatrix}
\]
is symmetric, proving self-adjointness for $G$.
Its first principal minor is positive, and its determinant
is $D(P^2-D)/(4q^{40})>0$. The quadratic form is therefore
positive definite: explicitly, completing the square gives
\[
 P z_1^2+2D z_1z_2+DPz_2^2
   =P(z_1+D z_2/P)^2+\frac{D(P^2-D)}Pz_2^2.
\]
This proves the assertion without relying on an abstract
spectral criterion. The eigenvalues are
$(P\pm\sqrt D)/(2q^{20})=-c,-c'$, and both are positive.
Since $F$ has positive coefficients, $F(q)>0$ for $q>0$;
the same claims hold for $M=FM_0$. The two determinant
identities follow by the trace and determinant formulas
already proved. The negative coefficients in (P-F39c)
prove the asserted failure of the stronger inference.

To retain the quadratic datum geometrically, let
\[
 R=\mathbb C[q,q^{-1},D^{-1},F^{-1}],\qquad
 S=R[x]/(x^2-D),\qquad U=\operatorname{Spec}R.
\]
The extension is free of rank two with basis $1,x$.
Because $D$ is invertible, so is $x$, and the derivative
$2x$ of its defining monic polynomial is invertible.
Hence $\operatorname{Spec}S\to U$ is finite \'{e}tale.
It is connected: P-F36 proves irreducibility over the
fraction field, and $S$ embeds in that quadratic field
since it is free over $R$.
Multiplication by $a=-F(P+x)/(2q^{20})$ on $S$ is represented by
\[
 -\frac{F}{2q^{20}}
   \begin{pmatrix}P&D\\1&P\end{pmatrix}.
 \tag{P-F39j}
\]
Taking its determinant polynomial gives exactly (P-F39b).
This is an endomorphism of an algebraic vector bundle,
with the quadratic cover and its involution retained.

Finally an endomorphism of a complex local system has,
in local flat frames, locally constant matrices. Its trace
is therefore locally constant, and is constant on a connected
base. Equivalently a horizontal endomorphism of the associated
flat bundle has constant characteristic-polynomial coefficients.
Here the trace is $-T_1(q)$, which is nonconstant: for example
its coefficient at $q^{26}$ is $2$, by (P-F39e).
On any connected finite dominant cover of $U$, at least one
irreducible component dominates $U$. The inclusion of function
fields on that component is injective and fixes the constants,
so a nonconstant rational function on $U$ remains nonconstant
there. In particular its pullback is not constant on the cover.
Consequently the multiplication endomorphism cannot become
a horizontal local-system endomorphism after such a cover.
This proves the third assertion.
\end{proof}

\paragraph{What a sign or purity argument would still have to identify.}
Before the prescribed rescaling, the rational minimal-polynomial
coefficients $T$ and $N$ already satisfy the source's positive-at-one
condition, but have negative Laurent coefficients. After one power
of the exact scaling factor, the two coefficients in (P-F39b)
have positive unimodal characters, and (P-F39i) is a concrete
Tate--Lefschetz model for each. Their construction here is a
coefficient-level realization; it supplies no correspondence
with the nonstandard multiplication, cell modules, or automorphic
projectors. Likewise the finite cover in (P-F39j) supplies the
algebraic eigenvalue datum, while its varying multiplication
operator is not a morphism of the cover's local system of
constant coefficients.

This locates a specific missing comparison rather than ruling
out local systems, algebraic coefficient fields, or theta lifts
in general. A construction may treat $q$ as a formal grading
variable rather than as a coordinate on a base; the horizontal
obstruction applies only to the latter literal interpretation.
No asymptotic bound in tensor degree, or conjecture for all
nonstandard coefficients, follows from this fixed calculation.
Even the explicit positive matrix family in (P-F39k) does not
supply a bimodule intertwiner or identify an automorphic
determinant with the nonstandard multiplication. The negative
coefficients of its unscaled determinant explain precisely why
a sign-only statement about such a family is insufficient to
recover coefficientwise positivity; the source's specific
normalization is substantive data.
The formal saturation and almost-positivity targets are
GCT VIII Conjectures~3.5--3.6, pp.~24--25; the worked
monomial-basis calculation is not silently identified with a
coefficient in a different ambient canonical basis.

Primary sources are
\url{https://arxiv.org/pdf/0709.0749}, Section~7.1.7, p.~28;
\url{https://arxiv.org/pdf/0709.0751}, Section~5.3, pp.~31--32,
Figure~6, p.~38, Section~5.4.1, p.~48, and Figures~21--23,
pp.~55--57; and
\url{https://www.jmilne.org/math/CourseNotes/LEC.pdf},
Theorem~23.2, pp.~139--140.


\subsection{A large local Klein character and an exact cancellation at one}

\begin{proposition}[P-F40: the Klein normal representation does not supply an unsigned global Laurent bound]
\label{prop:kleinphase}
Let $n\ge4$ with $3\nmid n$, and put
\[
 Q=2^n-(-1)^n,\qquad P=Q/3,\qquad
 K_n=\sum_{i=0}^{n-1}x_i^2x_{i+1},\qquad x_n=x_0.
 \tag{P-F40a}
\]
The cubic $K_n$ is smooth. Set $w_i=(-2)^i$ and
$\rho=e^{2\pi i/P}$. The matrix
\[
 \beta=\operatorname{diag}(\rho^{w_0},\ldots,\rho^{w_{n-1}})
 \tag{P-F40b}
\]
fixes $K_n$ exactly, belongs to $SL_n$, and has order $P$.
Let $W$ be the space of cubic forms with this substitution
convention and let
\[
 S=W/(\mathfrak{sl}_n\cdot K_n),\qquad D=\binom n3.
 \tag{P-F40c}
\]
This is the normal representation to the $SL_n$ orbit, restricted
to the actual cyclic stabilizer subgroup $\Gamma=\langle\beta\rangle$.
It satisfies
\[
 \dim S=D+1,\qquad\dim S^\beta=1,\qquad
 \det(\beta\mid S)=1.
 \tag{P-F40d}
\]
For each $a$ coprime to $P$, define
\[
 C_a=\frac1{\det(1-\beta^a\mid S^*_{\rm mov})}.
\]
Then
\[
                 C_{-a}=\overline{C_a}=(-1)^D C_a.
 \tag{P-F40e}
\]
For the standard root in (P-F40b), the absolute value nevertheless
has the same large asymptotic scale as the chain calculation:
\[
                         \log_2|C_1|=\frac{n^4}{24}+O(n^3).
 \tag{P-F40f}
\]
The implied constant is absolute as $n$ ranges over $3\nmid n$.

For the cyclic normal-representation model, with ordinary
polynomial degree on $\operatorname{Sym}(S^*)$, Molien averaging is
\[
 \operatorname{Hilb}_{\operatorname{Sym}(S^*)^\Gamma}(t)
       =\frac1P\sum_{a=0}^{P-1}
                       \frac1{\det(1-t\beta^a\mid S^*)}.
 \tag{P-F40g}
\]
Every summand with $(a,P)=1$ has at most a simple pole at one,
with normalized coefficient $C_a/P$.
If $n\equiv3\pmod4$, the pair of summands indexed by $a$ and
$-a$ is holomorphic at one: its possible simple poles cancel
exactly. Thus the entire primitive-order part of (P-F40g)
is holomorphic at one in this parity. If $n\not\equiv3\pmod4$,
the two normalized coefficients add to $2C_a/P$, which is real;
its sign is specified below and is not assumed positive.
Already for $n=8$ the signs are mixed: $C_1<0<C_3$.

These statements concern an actual stabilizer representation
and its cyclic finite-quotient model. No identification of
(P-F40g), or any of its summands, with the global pole-at-one
polynomial $B_n$ of P-F34 is asserted. In particular (P-F40f)
is neither a lower nor an upper bound on the actual $U_n$.
\end{proposition}

\begin{proof}
The partial derivatives of $K_n$ are
\[
                 2x_i x_{i+1}+x_{i-1}^2
\]
with cyclic indices. If one coordinate vanishes at a common
zero of these derivatives, the preceding coordinate also
vanishes, and induction around the cycle gives the origin.
If no coordinate vanishes, multiplying all $n$ equations gives
$2^n(\prod_i x_i)^2=(-1)^n(\prod_i x_i)^2$, which is impossible
in characteristic zero. This proves smoothness.

For comparison with the known family, $K_n$ is the type-$K_n$
polynomial of Zheng, Section~2.2. Lemma~3.1, printed p.~5,
identifies its diagonal projective automorphism group as
cyclic of order $Q/3$. We prove the particular lift and
representation properties needed here directly.
Let $r=-2$ and $S_0=\sum_{i=0}^{n-1}r^i=(1-r^n)/3$.
Then $S_0=\pm P$. The relations
$2w_i+w_{i+1}=0$ for $i<n-1$ and
$2w_{n-1}+w_0=1-r^n=3S_0$ show that $\beta$ fixes
all $n$ terms of $K_n$. Also
$\det\beta=\rho^{S_0}=1$, and its first eigenvalue is $\rho$,
so its order is $P$.
One may obtain this lift by cubing the exact-form diagonal
matrix with a primitive $Q$th root; that original matrix has
determinant of order three. This distinction between linear
and projective orders is essential here.

The hypothesis $3\nmid n$ implies $3\nmid P$. Indeed
$(-2)^n=(1-3)^n\equiv1-3n\pmod9$, so
$(1-(-2)^n)/3\equiv n\pmod3$.
Every $w_i$ is congruent to one modulo three.
Consequently, if $P$ divides a difference $w_i-w_j$, then
$Q=3P$ divides it. But for distinct indices its absolute
value is at most $2^{n-1}+2^{n-2}<Q$.
The eigenvalues of $\beta$ are therefore distinct.

We next determine all fixed cubic monomials, rather than
assuming that the displayed support is the whole fixed space.
A cubic monomial with indices $i,j,k$ is fixed exactly when
$P\mid b$, where $b=r^i+r^j+r^k$. Since $3\mid b$ and
$(3,P)=1$, this is equivalent to $Q\mid b$.
The bound $|b|\le3\cdot2^{n-1}<2Q$ leaves only
$b=0,Q,-Q$.
If $b=0$, the least index occurring cannot have multiplicity
one or three, as is seen by dividing by its power of two.
It has multiplicity two, and
$2r^i+r^j=0$ forces $j=i+1$.
These are the $n-1$ nonwrapping terms of $K_n$.

If $b=\pm Q$, it is odd. Thus the index zero occurs an odd
number of times. It cannot occur three times, since
$Q\ge15$, so it occurs once and the other two indices
$j,k$ are positive. For the target $b=1-r^n$, one has
$r^j+r^k=-r^n$. If $j\ne k$, the two-adic valuation of
the left side is $\min(j,k)<n$; if $j=k$, the equation
forces $j=k=n-1$. This is the wrapping monomial
$x_{n-1}^2x_0$. For the other target $b=r^n-1$,
one would have $r^j+r^k=r^n-2$, whose two-adic valuation
is one. Equal indices would give valuation at least two.
For unequal indices the smaller must be one, and subtracting
$r^1=-2$ would force the larger to be $n$, which is not
an available index. There are no other fixed monomials.
Thus $\dim W^\beta=n$.

We justify the orbit dimension without assuming the full
finite stabilizer group has been classified. Let
$R=\mathbb C[x_0,\ldots,x_{n-1}]$ and
$\mathfrak m=(x_0,\ldots,x_{n-1})$.
The $n$ quadratic partial derivatives have common zero only
at the origin. In the regular local ring $R_{\mathfrak m}$
they generate an ideal of colength finite, so they are a
regular sequence. The exact Koszul complex says that any
relation among them has each coefficient in their generated
ideal. In particular a relation with linear coefficients
has every coefficient in $\mathfrak m^2R_{\mathfrak m}$,
and hence every such linear coefficient is zero.
A Lie algebra element fixing $K_n$ gives precisely such a
relation. Thus the map
$\mathfrak{sl}_n\to W$, $X\mapsto X\cdot K_n$,
is injective. The regular-sequence and Koszul facts used
here are Stacks Project Lemmas~10.106.3, 10.104.2 and
15.31.2, Tags~00NQ, 02JN and 062F, respectively.

The quotient sequence in (P-F40c) is $\Gamma$-equivariant.
Finite-group invariants over $\mathbb C$ are exact by
averaging. As $\beta$ has distinct eigenvalues,
$\dim(\mathfrak{sl}_n)^\beta=n-1$. Therefore
\[
 \dim S=\binom{n+2}{3}-(n^2-1)=D+1,
 \qquad \dim S^\beta=n-(n-1)=1.
\]
The determinant of the cubic representation at any matrix
$B$ equals $(\det B)^{(n+1)(n+2)/2}$: by symmetry each
coordinate eigenvalue occurs with total exponent
$3\binom{n+2}{3}/n$ in the product of all cubic monomial
eigenvalues. In particular it is one at $\beta$.
The adjoint determinant is also one, since off-diagonal
weights occur in inverse pairs and diagonal weights are one.
Taking determinants in the exact sequence gives
$\det(\beta\mid S)=1$. Its fixed part has determinant one,
so both the moving part and its dual have determinant one.
All these conclusions hold for $\beta^a$ when $(a,P)=1$;
passing to such a power does not change any fixed eigenspace.

Let $\alpha_1,\ldots,\alpha_D$ be the nontrivial eigenvalues
of $\beta^a$ on $S^*$. Their product is one, and
\[
 \prod_{j=1}^D(1-\alpha_j^{-1})
     =(-1)^D\prod_{j=1}^D(1-\alpha_j).
\]
They are roots of unity, so inverse also means complex
conjugate. Taking reciprocals proves (P-F40e).
For an exact sign description, write
$\alpha_j=e^{2\pi i\theta_j}$ with $0<\theta_j<1$ and put
$h_a=\sum_j\theta_j\in\mathbb Z$; integrality follows
from their product being one. Then
\[
 C_a=(-1)^{h_a}i^D |C_a|.
 \tag{P-F40h}
\]
Indeed $1-e^{2\pi i\theta}=-2i\,e^{\pi i\theta}
\sin(\pi\theta)$ and the sine is positive in this interval.
The integer $h_a$ is computable directly from the cubic
weights:
\[
 h_a=\sum_{0\le i\le j\le k<n}
       \left\{-\frac{a(w_i+w_j+w_k)}P\right\}
                   -\frac{n(n-1)}2,
 \tag{P-F40i}
\]
where braces denote fractional part. Fixed cubic weights
have fractional part zero; each off-diagonal adjoint inverse
pair contributes one to its age. Additivity in the exact
sequence proves (P-F40i).
For even $D$ this gives the exact real sign
$(-1)^{D/2+h_a}$. It supplies no positivity assertion for
sums over different $a$.

The large absolute value has an independent elementary
proof. The moving determinant quotient gives
\[
 |C_1|=
 \frac{\prod_{i\ne j}|1-\rho^{w_i-w_j}|}
 {\prod_{i\le j\le k,\ P\nmid w_i+w_j+w_k}
                    |1-\rho^{w_i+w_j+w_k}|}.
 \tag{P-F40j}
\]
Dualizing does not change these absolute values.
For largest index $k\le n-6$, the moving cubic counts
are exactly $N_0=1$ and $N_k=(k^2+3k)/2$ for $k\ge1$,
as in P-F34. Every corresponding nonzero integer weight
satisfies $2^{k-2}\le|b|\le3\cdot2^k$.
Since $P\ge2^{n-2}$, these weights satisfy $|b|<P/2$.
The sine inequalities give
\[
       \frac{2^k}{P}\le|1-\rho^b|
                   \le\frac{6\pi\,2^k}{P}.
\]
Because $\log_2P=n+O(1)$, the logarithm of each bulk
factor is $k-n+O(1)$ with an absolute constant.
Their total main contribution to $\log_2|C_1|$ is
\[
 \sum_{k=0}^{n-6}(n-k)N_k=\frac{n^4}{24}+O(n^3),
\]
with the exact sum already displayed in P-F34.
The only remaining cubic factors have their largest
index among the last five indices. There are $O(n^2)$
of these. All nontrivial $P$th-root factors lie between
$4/P$ and two, so each of these logarithms is $O(n)$.
The $n(n-1)$ adjoint factors satisfy the same estimate.
The total error is therefore $O(n^3)$, proving (P-F40f).
Dividing by $P$ would change this logarithm by only $O(n)$.

Finally, averaging the trace projector over $\Gamma$ on
every symmetric degree gives (P-F40g). For a diagonal
operator, the generating trace on its symmetric algebra
is the product of the geometric series
$(1-t\alpha)^{-1}$, which proves the formula directly.
When $(a,P)=1$, its one-dimensional fixed subspace gives
\[
 \frac1{\det(1-t\beta^a\mid S^*)}
       =\frac{\phi_a(t)}{1-t},\qquad
 \phi_a\text{ holomorphic at }1,\quad\phi_a(1)=C_a.
\]
If $D$ is odd, (P-F40e) makes
$\phi_a(1)+\phi_{-a}(1)=0$; the numerator of the paired
term is divisible by $1-t$ as a germ. The pair is
therefore holomorphic at one. If $D$ is even, its
normalized coefficient is $2C_a$ instead.
The integer $D=n(n-1)(n-2)/6$ is odd exactly when
$n\equiv3\pmod4$: among three consecutive factors,
its numerator has two-adic valuation one exactly
in that case, while the denominator has valuation one.
This proves every cancellation assertion.

For completeness, mixed signs in even codimension can be checked
exactly in the small case $n=8$, rather than left as a hypothetical
possibility. Here $P=85$ and $D=56$. For a residue $b$, write
$\mathcal O_b=\{b(-2)^j\bmod85:0\le j<8\}$, with residues in
$\{0,\ldots,84\}$. The moving dual cubic and adjoint representations
decompose into the following orbits. The last two columns give the
sum of the orbit's residues divided by $85$, before and after
multiplication by three.
\[
\begin{array}{c|rrrr}
 b&\text{cubic multiplicity}&\text{adjoint multiplicity}
       &\sum_{u\in\mathcal O_b}u/85
       &\sum_{u\in\mathcal O_b}(3u\bmod85)/85\\\hline
 1&1&0&4&4\\
 3&3&2&4&4\\
 6&3&2&4&4\\
 7&1&0&3&4\\
 9&2&1&4&3\\
15&2&1&4&4\\
18&2&1&4&5
\end{array}
\]
Here is a direct way to verify the whole finite table. Rotation
of the eight variable indices acts freely on cubic monomials,
because a nontrivial rotational stabilizer has even order,
which cannot divide the total monomial degree three.
The fifteen monomial orbits have representatives: the cube
$(0,0,0)$; the seven triples $(0,0,d)$ for $1\le d\le7$;
and the seven squarefree triples
\[
 (0,1,2),(0,1,3),(0,1,4),(0,1,5),(0,1,6),(0,2,4),(0,2,5).
\]
The last list follows by writing the three positive cyclic gaps
with sum eight, up to cyclic rotation. Reducing the negative
weight sum of each representative modulo $85$ gives exactly
the displayed cubic multiplicities; $(0,0,1)$ is the fixed
orbit and is omitted. The seven off-diagonal adjoint orbits
have representatives $(-2)^d-1$, $1\le d\le7$, giving the
second multiplicity column. Each displayed residue orbit has
eight distinct elements. Its residues, or their multiples by
three, give the two elementary integer sums in the table.
Subtracting the adjoint multiplicities leaves each of the seven
displayed orbits once in $S^*_{\rm mov}$. Thus
\[
 h_1=4+4+4+3+4+4+4=27,\qquad
 h_3=4+4+4+4+3+4+5=28.
\]
Formula (P-F40h), with $D/2=28$, proves
\[
                              C_1<0<C_3.
 \tag{P-F40k}
\]
This establishes mixed local signs even when inverse pairing
does not cancel a simple pole.
\end{proof}

\paragraph{Three fixed variables remove the odd-codimension pair cancellation.}
There is also a precise limit to the parity shortcut. Let
$n\ge11$ with $3\nmid n$, put $m=n-3$, and consider the smooth
cubic in $n$ variables
\[
             F=K_m(x_0,\ldots,x_{m-1})+y_1^3+y_2^3+y_3^3.
 \tag{P-F40l}
\]
Extend the preceding Klein automorphism on the $x$ variables by
the identity on the three $y$ variables, and denote it by
$\widetilde\beta$. It is in $SL_n$ and has order
$P_m=(2^m-(-1)^m)/3$. Its normal representation
$\widetilde S=W_n/(\mathfrak{sl}_n\cdot F)$ satisfies
\[
 \dim\widetilde S=D+1,\qquad
 \dim\widetilde S^{\widetilde\beta}=2,
 \qquad \det(\widetilde\beta\mid\widetilde S)=1.
 \tag{P-F40m}
\]
Here and below $D=\binom n3$ uses the total number of variables.
To prove the fixed-space assertion, first note that no two
Klein eigenvalues have product one when $m\ge5$ and $3\nmid m$.
Equal indices would require the odd integer $P_m>1$ to divide
a power of two. For distinct indices, after removing a power
of $-2$, such a pair would require
$P_m\mid1+(-2)^d$ with $1\le d\le m-1$.
For $d\le m-2$ its nonzero absolute value is smaller than
$P_m$, since $2^{m-2}+1<P_m$ for $m\ge5$.
For $d=m-1$ its absolute value lies strictly between
$P_m$ and $2P_m$; equality with $P_m$ would give
$2^{m-1}=2$ when $m$ is even and $2^{m-1}=-2$ when $m$
is odd. Neither is possible. Thus there are no such pairs.

No Klein eigenvalue itself is one. Consequently no mixed
cubic monomial involving both the $x$ and $y$ variables is
fixed: its weight is either one Klein weight or the sum of
two. The fixed cubic subspace has dimension
$m+\binom{3+2}{3}=m+10$.
The centralizer in $\mathfrak{sl}_n$ has dimension
$m+3^2-1=m+8$, because the Klein eigenvalues are distinct
and different from the eigenvalue one on the three fixed
variables. Their difference is two. Smoothness follows
from the independent gradient blocks, and the same regular
partials argument gives Lie-orbit injectivity. This proves
the dimension assertions in (P-F40m); its determinant
assertion follows from the same cubic/adjoint determinant
calculation as before.

Put
\[
 \widetilde C_a=
 \frac1{\det(1-\widetilde\beta^a\mid
                         \widetilde S^*_{\rm mov})}
                \quad((a,P_m)=1).
\]
The moving codimension is now $D-1$, so
$\widetilde C_{-a}=(-1)^{D-1}\widetilde C_a$.
In particular, if $n\equiv3\pmod4$, the inverse pair in the
cyclic normal-model Molien sum has normalized double-pole
coefficient
\[
                    \frac{2\widetilde C_a}{P_m}\ne0
                    \quad\hbox{at }t=1.
 \tag{P-F40n}
\]
Indeed each term is $(1-t)^{-2}$ times a holomorphic function
with value $\widetilde C_a$ at one. The two values are equal
and nonzero in this parity, so this particular pair has an
actual pole of order two. This does not exclude cancellation
against other pairs in the complete average.

In the standard embedding one again has
\[
                  \log_2|\widetilde C_1|
                             =\frac{n^4}{24}+O(n^3).
 \tag{P-F40o}
\]
To check this without a new estimate, use the bulk calculation
above on the first $m$ variables: its contribution is
$m^4/24+O(m^3)$. Cubic monomials involving one of the three
fixed variables number only $O(n^2)$, as do all adjoint
factors and all exceptional top-index factors. Each nontrivial
$P_m$th-root factor has logarithm $O(n)$, so these change the
answer by $O(n^3)$. Finally $m=n-3$ gives
$m^4/24=n^4/24+O(n^3)$.
Thus an odd value of $D$ in the full $n$-variable calculation does not remove every
large local contribution at one: a higher-dimensional fixed
normal space changes the relevant parity and Laurent order.

\paragraph{Scope of the local calculation and the later global comparison.}
The representation $S$ comes from a genuine stabilizer subgroup
of a smooth cubic. The full stabilizer may be larger than
$\Gamma$, and the cyclic model (P-F40g) is not claimed to be
the entire local quotient. Even a correctly identified local
quotient would not identify its ordinary tangent-coordinate
Hilbert grading with the global stretching grading of P-F31.
A global equivariant Riemann--Roch or Molien--Weyl comparison,
with all strata and their signs controlled, would be required
to express the actual $B_n$ through such local data.
P-F44 below supplies an exact global comparison for the
primitive-Klein collection by the Weyl--Brion expansion;
its normalization is $1/(nP)$ and it retains the entire
complementary signed contribution. The cyclic model itself,
and in particular the three-fixed-variable extension, is
not identified with the global Hilbert series. The exact
odd-codimension cancellation shows why summing absolute
local amplitudes cannot replace that comparison.

The classical family and its diagonal projective group are in
Zheng, \emph{On abelian automorphism groups of hypersurfaces},
Section~2.2 and Lemma~3.1, printed pp.~3 and~5,
\url{https://arxiv.org/pdf/2004.09008}.
The commutative-algebra premises are directly available at
\url{https://stacks.math.columbia.edu/tag/00NQ},
\url{https://stacks.math.columbia.edu/tag/02JN}, and
\url{https://stacks.math.columbia.edu/tag/062F}.
No historical priority is claimed for this family or for
Molien averaging; the phase and scope computations are
proved explicitly above.


\subsection{A specific Borcherds theta construction is the cubic discriminant}

The automorphic ring in P-F37 has a known theta construction for a
particular nonzero invariant. This is substantially more specific than
asking for an unidentified theta realization of an entire multiplicity
space. The construction and its divisor are known; the following
corollary records their exact compatibility with the plethysm grading.
No historical priority is claimed for the construction or this
identification.

\paragraph{Known geometric input.}
Let $\mathbb B$ be the cubic-threefold ball of dimension ten, with
nodal and chordal arrangements $\mathcal H_\Delta$ and $\mathcal H_c$.
Casalaina-Martin--Laza, Proposition~7.5, printed p.~25,
construct a weight-$48$ form $\phi$ from the Borcherds form
$\Phi_{12}$ by an $E_6$ quasi-pullback and restriction to this ball.
Its orders along these two arrangements are respectively $3$ and $84$.
Their Lemma~7.7 and Corollary~7.8, p.~26, give ramification
orders $3,6$ and the quotient divisor $H_\Delta+14H_c$.
The source is \url{https://arxiv.org/pdf/0710.5329}.

Write $U=\mathbb C^5$, $\mathscr S=\operatorname{Sym}^3U^*$,
and let $\Delta_{\rm cub}$ be the usual discriminant polynomial
on $\mathscr S$, with an arbitrary fixed nonzero normalization.
The period-map isomorphism of P-F37 will be denoted by
\[
 \mathcal P:\bigoplus_k\mathcal M_k
       \longrightarrow \mathbb C[\mathscr S]^{SL(U)};
 \qquad\deg\mathcal P(g)=\frac53\operatorname{wt}(g).
 \tag{P-F41a}
\]

\begin{corollary}[P-F41: an identified theta image in the actual cubic plethysm algebra]
\label{cor:thetadiscriminant}
The form $\phi$ has trivial multiplier on the group used in
(P-F41a), and
\[
             \mathcal P(\phi)=c\,\Delta_{\rm cub}
             \quad\text{for some }c\in\mathbb C^*.
 \tag{P-F41b}
\]
Consequently this particular multiplicative Borcherds theta
construction identifies a nonzero vector in the actual graded piece
\[
 T^{(5)}_{16}
    =\left(\operatorname{Sym}^{80}
                    (\operatorname{Sym}^3\mathbb C^5)\right)^{SL_5},
 \tag{P-F41c}
\]
with automorphic weight $48$, ordinary coefficient degree $80$,
and stretching degree $16$. Under the full $GL_5$ action its
character is $\det^{48}$. It identifies the polynomial subalgebra
$\mathbb C[\phi]$ with $\mathbb C[\Delta_{\rm cub}]$,
including their products and these gradings.

Here ``Borcherds theta construction'' has a precise meaning:
$\Phi_{12}$ is the multiplicative lift associated with the
regularized orthogonal theta integral for the input
$\delta(\tau)^{-1}=q^{-1}+24+\cdots$ on $II_{2,26}$,
where $\delta$ is the elliptic modular discriminant.
The integral gives the logarithm of its Petersson norm;
$\Phi_{12}$ is not asserted to be the value of an ordinary
linear convergent theta integral. The subsequent operations
are quasi-pullback and restriction.
\end{corollary}

\begin{proof}
We first check the multiplier instead of discarding it implicitly.
Let $\Lambda$ be the Eisenstein primitive cohomology lattice of
a cyclic cubic fourfold and $\Gamma_{\rm mon}$ its monodromy group.
Allcock--Carlson--Toledo, Lemma~2.3, printed p.~11, sends a
nodal meridian to an exact order-three complex reflection fixing
its mirror pointwise. Their Theorem~8.1 and proof, p.~74,
identify $\Gamma_{\rm mon}$ with the subgroup of
$\operatorname{Aut}_E(\Lambda)$ having determinant in $\mu_3$;
the proof also records that these meridians generate the
monodromy image. These are primary statements in
\url{https://arxiv.org/pdf/math/0608287}.

Suppose initially that $\phi$ has a scalar multiplier $\chi$.
This allows for the character in a Borcherds automorphy statement.
View it as a homogeneous function of degree $-48$ on the
negative cone over the ball. For a nodal reflection $r$, choose
a generic point of its mirror and a linear slice of that cone
defined by a nonzero linear functional invariant under $r$.
Such a slice exists because $r$ fixes the mirror pointwise.
On this slice choose normal coordinate $z$ with
$r(z)=\omega z$, where $\omega^3=1$ and $\omega\ne1$.
Its order-three zero gives a local expansion
\[
               \phi(z,y)=z^3 u(z,y),\qquad u(0,y_0)\ne0.
\]
The automorphy equation $\phi(rZ)=\chi(r)\phi(Z)$,
compared in its first nonzero normal coefficient at $y_0$,
gives $\chi(r)=\omega^3=1$. It follows that
$\chi$ is trivial on $\Gamma_{\rm mon}$.
If the automorphy group is taken to be the full unitary group,
the same conclusion holds there: a unitary lattice determinant
is an Eisenstein unit, hence belongs to $\mu_6$; the subgroup
just described has index two and the other coset is represented
by $-I$, because $\operatorname{rank}_E\Lambda=11$.
Homogeneity gives $\phi(-Z)=(-1)^{-48}\phi(Z)=\phi(Z)$.

For completeness, this also matches precisely the group in
Looijenga--Swierstra rather than merely its projectivization.
Their group consists of the $\mu_3$-equivariant isometries
of full middle cohomology fixing the polarization class
(Section~3, printed p.~10). It contains the monodromy group.
The proof of Allcock--Carlson--Toledo Theorem~8.1 explicitly
notes that $-I$ on primitive cohomology does not extend to
full cohomology while fixing that class. Since the full
unitary group has just the two cosets above, its extendable
subgroup is exactly $\Gamma_{\rm mon}$. Thus the form is
invariant for the group of (P-F41a), with no need to take a power.

Apply the actual graded ring isomorphism of
Looijenga--Swierstra Theorem~3.1, p.~12, and its grading
calculation on pp.~10--11, as in P-F37. The proof on p.~13
identifies the automorphic and invariant line bundles over
the stable locus, so the identification also preserves
vanishing at a stable nodal cubic.
It sends the nonzero form $\phi$ to a nonzero invariant
$I\in\mathbb C[\mathscr S]^{SL_5}$ of degree
$(5/3)48=80$. A holomorphic form on the ball is in particular
allowed in the arrangement-meromorphic ring of that theorem.
At a nodal cubic, the extended period map meets the nodal
mirror; this discriminant identification is also explicit
in Allcock--Carlson--Toledo Theorem~1.1, printed p.~3.
Since $\phi$ vanishes there, its pullback $I$
vanishes on the nodal discriminant locus.

We specify why this locus is dense in the whole discriminant.
The incidence variety
\[
 \mathscr I=\{(F,p)\in\mathscr S\times\mathbb P^4:
                         dF(p)=0\}
\]
is a rank-$30$ vector bundle over $\mathbb P^4$: after sending
$p$ to $[1:0:0:0:0]$, the five conditions set the coefficients
of $x_0^3,x_0^2x_1,\ldots,x_0^2x_4$ to zero.
It is irreducible of dimension $34$ and projects properly
onto the discriminant. Requiring singularities at two
distinct points imposes ten independent conditions. Indeed
move the ordered pair to $[1:0:\cdots:0]$ and
$[0:1:0:\cdots:0]$; the two lists of five cubic monomials
are disjoint. The ordered-pair incidence therefore has
dimension $8+(35-10)=33$.
The locus in $\mathscr I$ where the quadratic part at $p$
is degenerate is also proper closed. To see properness,
\[
        F_0=x_0\sum_{i=1}^4x_i^2+\sum_{i=1}^4x_i^3
 \tag{P-F41d}
\]
has its only singular point at $[1:0:\cdots:0]$, where its
affine quadratic part is nondegenerate. In fact a critical
point with $x_0=0$ forces all $x_i=0$. Otherwise set
$x_0=1$; every $x_i$ is $0$ or $-2/3$, and
$\partial F_0/\partial x_0=\sum_i x_i^2=0$ forces
all four to vanish. Thus the exceptional loci just described
have images, or image closures, of dimension at most $33$.
Away from them a singular cubic has a unique ordinary node.
The nodal locus is therefore dense in the dimension-$34$
discriminant hypersurface.

The exact algebraic discriminant facts are
Bus\'e--Jouanolou, Proposition~4.7(i), printed p.~40,
and Theorem~4.26, p.~53, in
\url{https://arxiv.org/pdf/1210.4697}.
Over $\mathbb C$ their generic discriminant is a prime
polynomial and its coefficient degree is
$5(3-1)^4=80$. It cuts out the singular-hypersurface locus.
Consequently $I$, vanishing on the dense nodal locus,
vanishes on its whole irreducible closure and is divisible
by $\Delta_{\rm cub}$. Both polynomials have degree $80$,
so $I=c\Delta_{\rm cub}$ with $c\ne0$.
This proves (P-F41b). The change-of-variables formula in
the same source, Proposition~4.13, pp.~42--43, gives
$\Delta_{\rm cub}(F\circ g)=\det(g)^{3\cdot2^4}
\Delta_{\rm cub}(F)$, proving the stated $GL_5$ character.
The regrading by five is the actual Cartan/plethysm ring
identification already made in P-F19 and P-F37, so this
proves (P-F41c), not merely an equality of dimensions.
Since the discriminant is a nonconstant polynomial,
the homomorphism $\mathbb C[t]\to\mathbb C[\Delta_{\rm cub}]$
sending $t$ to $\Delta_{\rm cub}$ is injective:
if $p(\Delta_{\rm cub})=0$ for a nonzero $p\in\mathbb C[t]$,
factor $p$ over $\mathbb C$ and use that the polynomial
ring is a domain to obtain the contradiction that
$\Delta_{\rm cub}$ is constant.
The ring isomorphism then gives the assertion about
$\mathbb C[\phi]$ and all its products.

We finish by verifying the precise meaning of the theta
construction. For $M=II_{2,26}$ the weakly holomorphic
input $\delta^{-1}$ has the required weight $-12=1-26/2$
and has integral Fourier coefficients at nonpositive
exponents. Borcherds, \emph{Automorphic forms with
singularities on Grassmannians}, Section~6, printed p.~22,
defines the theta integral by truncation and meromorphic
regularization, taking a constant term at the regularizing
parameter zero. Theorem~13.3, pp.~48--49, gives its
multiplicative lift of weight $24/2=12$; item~(3)
relates the logarithm of the lift's norm to that integral.
The input's only negative term is $q^{-1}$, giving the
order-one root divisors of $\Phi_{12}$.
The identification of this product with $\Phi_{12}$ is
explicit in Gritsenko, \emph{24 faces of the Borcherds
modular form $\Phi_{12}$}, Section~1, equation~(1),
p.~4, and Section~3.2, pp.~11--12.

The primitive cubic-fourfold lattice embeds into
$II_{26,2}$ with orthogonal complement $E_6$.
The operation in Casalaina-Martin--Laza Proposition~7.5
first removes the $36$ root-hyperplane factors containing
the smaller domain and then restricts to it; restriction
to the cubic-threefold ball follows. Removing a linear
factor increases automorphic weight by one, so the resulting
weight is $12+36=48$. These are exactly the geometric
operations producing the $\phi$ used above. Thus the
theta origin and the invariant identification refer to
the same form. This completes the proof.
\end{proof}

\paragraph{Exact remaining comparison with the supplied method.}
The supplied \emph{Remark}, printed p.~3, already mentions
$O(2,n)$, so the preceding orthogonal construction is not
excluded merely by its group name. Its item~8 on p.~4
specifies iterated theta lifts between totally definite
unitary groups; items~9--10 and~12 concern their norms,
divisibility and periods. The concrete construction above
instead uses a weakly holomorphic scalar input on $SL_2$,
a regularized multiplicative lift to $O(2,26)$, an $E_6$
quasi-pullback, and restriction to $U(10,1)$.
No identification of these operations with the supplied
iterated lift or its determinant has been established.

The next comparison can now be stated for a named output:
it must recover $\phi$ (equivalently $\Delta_{\rm cub}$),
or an identified larger graded collection containing it,
with the three gradings in (P-F41c) respected.
The existing construction supplies neither a basis for
every plethysm multiplicity space nor a uniform positive
counting formula, efficient oracle, or PH2/SH estimate.
It concerns the fixed inner representation
$\operatorname{Sym}^3\mathbb C^5$.

The regularized-theta primary sources are
\url{https://arxiv.org/pdf/alg-geom/9609022}, Section~6 p.~22,
Theorem~13.3 pp.~48--49, and
\url{https://arxiv.org/pdf/1203.6503}, equation~(1) p.~4
and Section~3.2 pp.~11--12.


\subsection{An unconditional bound for the actual cubic average polynomial}

\begin{proposition}[P-F42: a Weyl--cone bound on the actual coefficient norm]
\label{prop:cubicaverageupper}
Use the ordinary cubic stretching function $f_n$, average polynomial
$B_n$, and coefficient norm $U_n$ of P-F34. For every $n\ge4$, put
\[
 m=\binom{n+2}{3},\qquad a=m-n.
\]
Then
\[
 U_n\le
 8n!\,m^{3n}\bigl(288n^3\,3^n\bigr)^{m-n}.
 \tag{P-F42a}
\]
In particular,
\[
 \log_2 U_n\le
 \frac{\log_2 3}{6}\,n^4+O(n^3\log n),\qquad
 \limsup_{n\longrightarrow\infty}
       \frac{\log_2 U_n}{n^4}
       \le\frac{\log_2 3}{6}.
 \tag{P-F42b}
\]
The implied constant is absolute. This is an unconditional estimate
for the actual cubic invariant ring, including its Weyl alternation.
Its leading constant is greater than $1/24$, so it does not establish
the sufficient remaining condition in P-F34 or refute PH2 or SH.
It is a coefficient bound, not a claim of a polynomial-time procedure
for computing $B_n$.
\end{proposition}

\begin{proof}
\emph{1. The exact vector-partition formula.}
Let
\[
 g=\gcd(n,3),\quad c=3/g,\quad d_0=n/g,\quad
 \delta=(n-1,n-2,\ldots,0).
\]
Index the columns of an $n$ by $m$ matrix $A$ by all degree-three
monomials in $n$ variables, and let a column $\alpha_i$ be that
monomial's exponent vector. Each column is nonnegative with
coordinate sum three and Euclidean norm at most three.
The columns $3e_j$ show that $A$ has rank $n$. Define
\[
 p(b)=\#\{\beta\in\mathbb Z_{\ge0}^{m}:A\beta=b\}.
\]
Then
\[
 f_n(s)=
 \sum_{\sigma\in S_n}\operatorname{sgn}(\sigma)\,
       p(cs\mathbf1+b_\sigma),\qquad
 b_\sigma=\delta-\sigma\delta.
 \tag{P-F42c}
\]
In particular $|(b_\sigma)_j|\le n-1$ and
$\sum_j(b_\sigma)_j=0$.

Here is the coefficient extraction proving (P-F42c).
The monomial character of
$\operatorname{Sym}^{d_0s}(\operatorname{Sym}^3\mathbb C^n)$
counts the vectors $\beta$ with $\sum_i\beta_i=d_0s$.
If $A\beta=cs\mathbf1+b_\sigma$, the last equation is automatic:
summing coordinates gives $3\sum_i\beta_i=ncs=3d_0s$.
Multiply this character by the alternant
$\sum_{\sigma}\operatorname{sgn}(\sigma)x^{\sigma\delta}$
and extract the coefficient of $x^{cs\mathbf1+\delta}$.
The Weyl character formula says that an irreducible summand with
highest weight $\lambda$ contributes its alternant with exponents
the permutations of $\lambda+\delta$.
The vector $cs\mathbf1+\delta$ is strictly decreasing, so it occurs
only when $\lambda=cs\mathbf1$, and then with coefficient one.
This is precisely the $SL_n$-invariant summand.
The Weyl formula used here is also displayed in GCT VI \cite{PGCTVI},
Section~5.1, equation~(5.1). No cancellation or estimate has yet
been discarded.

\emph{2. One perturbation works for every integer right-hand side.}
We shall bound the average constituent polynomial of each function
$p(cs\mathbf1+b_\sigma)$ separately. Choose a rational vector
$\epsilon\in(0,1)^m$ as follows.
For each rational hyperplane spanned by a rank-$(n-1)$ subset
of columns of $A$, choose a nonzero primitive integer normal
$\ell_H$. There are finitely many such hyperplanes. Take
$\epsilon$ sufficiently small and outside their inverse-image
hyperplanes so that
\[
             0<|\ell_H A\epsilon|<1
                    \quad\hbox{for every }H.
 \tag{P-F42d}
\]
This is possible: $A$ has full rank, so each linear functional
$\ell_H A$ is nonzero, and a finite union of proper hyperplanes
cannot cover an open positive box. Rational points outside that
union exist in a sufficiently small such box.
For every integer vector $b$, the vector $b+A\epsilon$ lies on
none of the specified hyperplanes, since $\ell_H b$ is an integer
and the nonzero number $\ell_H A\epsilon$ has absolute value less
than one.

For $b=cs\mathbf1+b_\sigma$, consider
\[
 P_s^\epsilon=
 \{\beta\in\mathbb R^m:A\beta=b,\ \beta_i\ge-\epsilon_i\}.
 \tag{P-F42e}
\]
Its integer points are exactly the nonnegative integer solutions
of $A\beta=b$. It is bounded, because the column sums give
$\sum_i\beta_i=d_0s$ and every coordinate has a lower bound.
For all sufficiently large $s$, it is nonempty of dimension
$a=m-n$: the translated right-hand side $b+A\epsilon$ has
strictly positive coordinates, and the columns $3e_j$ allow a
solution with every translated variable $u_i=\beta_i+\epsilon_i$
strictly positive.
At a vertex, the positive coordinates of $u$ must correspond
to linearly independent columns. There must be exactly $n$ of
them, since otherwise $b+A\epsilon$ would lie in one of the
hyperplanes excluded above. Thus the polytope is simple, with
vertices indexed by feasible $n$-column bases.
For a fixed basis, every basic coordinate is affine in $s$.
Its strict positivity is therefore eventually constant.
There is consequently one threshold beyond which the set of
feasible bases is constant, simultaneously for all the finitely
many shifts $b_\sigma$. No bound on that threshold is needed.

\emph{3. Bounds on the bases and one common generic direction.}
For an $n$-column basis $J$, write
\[
 C=A_J,\qquad d=|\det C|,\qquad I=\{1,\ldots,m\}\setminus J.
\]
Hadamard's inequality applied to columns, and then to every
cofactor, gives
\[
 1\le d\le3^n,\qquad
 |(C^{-1})_{ij}|\le 3^{n-1}/d.
 \tag{P-F42f}
\]
There are at most $\binom mn\le m^n$ bases.
We choose a single real vector
$\ell=(\ell_1,\ldots,\ell_m)\in[-1,1]^m$.
For each basis and each $i\in I$, put
\[
 w_{J,i}=\ell_i-\ell_J C^{-1}\alpha_i,\qquad
 Z_J=\prod_{i\in I}\max\{1,|w_{J,i}|^{-1}\}.
 \tag{P-F42g}
\]
There exists such a vector for which every $w_{J,i}$ is nonzero
and
\[
                 Z_J\le4m^{2n}9^a
                         \quad\hbox{for every basis }J.
 \tag{P-F42h}
\]
To prove this existence assertion, take the coordinates of
$\ell$ independently and uniformly in $[-1,1]$.
For a uniform random variable $X$ on this interval and any real
number $v$,
\[
 \mathbb E\bigl[\max\{1,|X-v|^{-1}\}^{1/2}\bigr]
 \le1+\frac12\int_{-1}^{1}|u|^{-1/2}\,du=3.
 \tag{P-F42i}
\]
Conditional on $\ell_J$, the variables $w_{J,i}$ for $i\in I$
are independent translates of uniform variables. Hence
$\mathbb E[Z_J^{1/2}]\le3^a$.
Markov's inequality bounds the probability that (P-F42h) fails
for a given basis by $1/(2m^n)$. The union bound over all bases
is at most $1/2$. The probability that any $w_{J,i}$ is zero is
zero, so a simultaneous choice exists. This argument is used
only to prove a bound; no randomized algorithm is asserted.

For this chosen direction, and a shift $b_0=b_\sigma$, set
\[
 \lambda=c\ell_J C^{-1}\mathbf1,\qquad
 v=\ell_J C^{-1}b_0.
\]
The cofactor bounds, the column sum three, and $c\le3$ imply
\[
 |w_{J,i}|\le\frac{(n+1)3^n}{d},\qquad
 |\lambda|\le\frac{n^2 3^n}{d},\qquad
 |v|\le\frac{n^3 3^n}{d}.
 \tag{P-F42j}
\]
For the first inequality use
$1+n3^n/d\le(n+1)3^n/d$, since $d\le3^n$.
Keeping this factor $d$ in the denominator will be useful.

\emph{4. Exact finite-character formula for each tangent cone.}
Let $G_J=\mathbb Z^n/C\mathbb Z^n$, a finite abelian group of
order $d$, and let $\widehat G_J$ be its character group.
At a vertex with basis $J$, the tangent cone retains precisely
the inequalities $\beta_i\ge-\epsilon_i$ for $i\in I$.
Its integer points therefore have $\beta_i\ge0$ for $i\in I$,
and
\[
 \beta_J=C^{-1}\left(cs\mathbf1+b_0-
                            \sum_{i\in I}\alpha_i\beta_i\right).
 \tag{P-F42k}
\]
The remaining coordinates are integers exactly when the vector
in parentheses belongs to $C\mathbb Z^n$.
Finite character orthogonality imposes this congruence.
Consequently the exponential generating function of the cone,
under the specialization $\beta\mapsto e^{z\ell\cdot\beta}$, is
the following meromorphic function near $z=0$:
\[
 \frac1d\sum_{\chi\in\widehat G_J}
 \chi(b_0)^{-1}\chi(c\mathbf1)^{-s}
 \frac{\exp(z(\lambda s+v))}
      {\prod_{i\in I}(1-\chi(\alpha_i)e^{z w_{J,i}})}.
 \tag{P-F42l}
\]
One may first sum the geometric series in a multivariable
domain of convergence for this pointed cone, apply the finite
character filter, and then specialize the resulting rational
identity. Equivalently, before specialization use independent
exponential variables and the rational expressions for
$\beta_J$ in (P-F42k). The formula then follows by analytic
continuation near a generic line. We do not require all
$w_{J,i}$ to have the same sign or claim a common univariate
domain of absolute convergence for all cones.

Brion's identity says that the sum of these rational cone
generating functions over the feasible vertices is the
generating function of the polytope. The rational-polyhedron
version used here is Barvinok \cite{CBarvinokGF}, \emph{The Complexity of
Generating Functions for Integer Points in Polyhedra and
Beyond}, Theorem~2.5, printed p.~9. In particular, the
coefficient of $z^0$ in that sum is $p(cs\mathbf1+b_0)$.
The same identity applies to these lower-dimensional
polyhedra in their ambient lattice; alternatively one can
restrict to their rational affine hull before applying it.

\emph{5. A uniform analytic bound on every character polynomial.}
Set
\[
                      R=\frac1{16n^3 3^n}.
 \tag{P-F42m}
\]
Fix a basis and a character, and let $k$ be the number of
indices $i\in I$ for which $\chi(\alpha_i)=1$.
There are $a$ denominator factors in total.
For $|z|\le R$, (P-F42j) gives
\[
 |zw_{J,i}|\le\frac1{8d},\qquad
 R(|v|+|\lambda|)\le\frac18.
 \tag{P-F42n}
\]
If $\chi(\alpha_i)\ne1$, it is a root of unity of order at
most $d$, so
$|1-\chi(\alpha_i)|\ge4/d$.
The elementary estimate $|e^u-1|\le2|u|$ for $|u|\le1/2$
then gives, throughout this disk,
\[
 \left|\frac1{1-\chi(\alpha_i)e^{z w_{J,i}}}\right|
          \le d/2.
 \tag{P-F42o}
\]
There are no such factors when $d=1$.
For the factors with $\chi(\alpha_i)=1$, the removable
function $z/(1-e^{z w_{J,i}})$ is analytic on this disk and
satisfies
\[
 \left|\frac{z}{1-e^{z w_{J,i}}}\right|
            \le2/|w_{J,i}|.
 \tag{P-F42p}
\]
For example, the latter inequality follows by bounding the
Taylor remainder in $e^u-1=u+O(u^2)$ when $|u|\le1/2$.

Remove the root-of-unity factor $\chi(c\mathbf1)^{-s}$ from
one summand of (P-F42l), and write the rest as
$d^{-1}\chi(b_0)^{-1}z^{-k}e^{\lambda sz}F(z)$, where
\[
 F(z)=e^{vz}
 \prod_{\chi(\alpha_i)=1}
             \frac{z}{1-e^{z w_{J,i}}}
 \prod_{\chi(\alpha_i)\ne1}
             \frac1{1-\chi(\alpha_i)e^{z w_{J,i}}}.
\]
This is analytic on $|z|\le R$.
The constant term of the character summand is
$\chi(c\mathbf1)^{-s}P_{J,\chi,b_0}(s)$, where
$P_{J,\chi,b_0}$ is a polynomial of degree at most $k$.
Writing $\|P\|_1$ for the sum of the absolute values of
the ordinary monomial coefficients, Cauchy's estimate gives
\begin{align*}
 \|P_{J,\chi,b_0}\|_1
 &\le \frac1d
       \sum_{j=0}^k\frac{|\lambda|^j}{j!}
          R^{-(k-j)}\max_{|z|=R}|F(z)|\\
 &\le \frac{2Z_J}{d}\,2^a d^{a-k}R^{-k}\\
 &\le \frac{2Z_J}{d}\,(2/R)^a.
 \tag{P-F42q}
\end{align*}
For the second line, the product of the reciprocal
$|w_{J,i}|$ over any subset is bounded by $Z_J$,
(P-F42o)--(P-F42p) give
$2^k(d/2)^{a-k}\le2^a d^{a-k}$, and
$e^{R(|v|+|\lambda|)}\le2$.
For the last line use $dR\le1$, a consequence of
$d\le3^n$ and (P-F42m).
This absorbs the nontrivial-character denominators into
the same radius bound; multiplying by an additional
factor $d^a$ is unnecessary.

\emph{6. Averaging constituents and completing the estimate.}
For every sufficiently large integer $s$, the same finite
set of feasible bases is used. Formula (P-F42l) and Brion's
identity therefore express $p(cs\mathbf1+b_0)$ as a finite
sum of terms
\[
              \chi(c\mathbf1)^{-s}P_{J,\chi,b_0}(s).
 \tag{P-F42r}
\]
Take a common period for all these roots of unity.
On each residue class, agreement for all sufficiently large
integers identifies the entire eventual constituent
polynomial. Its average over the residue classes retains
exactly the terms with $\chi(c\mathbf1)=1$, by finite
root-of-unity orthogonality. Thus we are averaging
constituent polynomials, not values at a selected set of
integers.
The average polynomial contributed by one basis has norm
at most $2Z_J(2/R)^a$, since there are at most $d$ characters
and each bound in (P-F42q) includes the factor $1/d$.
There are at most $m^n$ bases and $n!$ Weyl shifts.
The triangle inequality in (P-F42c), together with
(P-F42h), now gives
\[
 \begin{split}
 U_n
 &\le n!\,m^n\cdot
          2\bigl(4m^{2n}9^a\bigr)(2/R)^a\\
 &=8n!\,m^{3n}\bigl(288n^3 3^n\bigr)^a.
 \end{split}
 \tag{P-F42s}
\]
The actual function $f_n$ is globally quasipolynomial by
P-F31. Hence the average of the eventual constituents
just computed is exactly its intrinsic polynomial $B_n$
from P-F34: equality on infinitely many integers of
each residue class identifies those polynomials.
Possible cancellations lowering the degree from at most
$a$ to $\binom n3$ are harmless for this upper bound.

Finally, $a=n^3/6+O(n^2)$ and
$\log_2(n!m^{3n})=O(n\log n)$.
Taking logarithms of (P-F42s) proves (P-F42b).
All constants in the proof are uniform in the Weyl shift,
the basis, the finite character, and the stretch parameter.
\end{proof}

\paragraph{Exact remaining quantitative gate.}
P-F34 proves
$\log_2|A_n(e^{2\pi i/2^{n-1}})|=n^4/24+O(n^3)$.
The unconditional coefficient-norm bound above has leading
constant $(\log_2 3)/6>1/24$.
It therefore places the actual norm on an explicit
$2^{O(n^4)}$ scale but leaves the required comparison open.
The proof discards all Weyl signs at its final norm estimate.
Any improvement based on cancellations among the full cubic
weight list and the Weyl numerator requires an additional
argument; neither Brion's identity nor the generic
small-entry bounds alone provide it.

\paragraph{A coefficient-sensitive refinement of the remaining gate.}
Eventual feasibility supplies a stronger estimate for the coefficient
of the stretch parameter than the cofactor estimate used above.
Write \(B_n(X)=\sum_{j=0}^{D}b_{n,j}X^j\), where
\(D=\binom n3\). Then the same construction proves, for every
\(0\le j\le D\),
\[
 |b_{n,j}|
 \le 8n!\,m^{3n}\,18^a
       (16n^3\,3^n)^{a-j}\frac{n^j}{j!},
 \qquad a=m-n.
 \tag{P-F42t}
\]
This estimate concerns the actual coefficients after Weyl
alternation; it does not assume their signs.

Here is the additional proof. A basis that is feasible for all
sufficiently large \(s\) has
\(C^{-1}(c\mathbf1)\ge0\) entrywise. Indeed its basic
coordinates are affine functions of \(s\), with this vector
of slopes, so a negative slope would make that coordinate
negative for all sufficiently large \(s\).
The column sums of \(C\) are three, whence
\(\mathbf1^{\mathsf T}C^{-1}=\mathbf1^{\mathsf T}/3\).
Consequently
\[
 \sum_i (C^{-1}\mathbf1)_i=n/3,
 \qquad
 |\lambda|=|c\ell_JC^{-1}\mathbf1|
              \le cn/3\le n.
 \tag{P-F42u}
\]
This also applies when a limiting basic slope is zero; strict
positivity of the slope is not required.

Use the same radius \(R=(16n^3\,3^n)^{-1}\), direction,
and analytic function \(F\) as in steps~3--5.
For one character summand with \(k\) fixed denominator
factors, the coefficient of \(s^j\) is zero when \(j>k\).
When \(0\le j\le k\), its absolute value is at most
\[
 \frac1d\,\frac{n^j}{j!}\,
       R^{-(k-j)}
       \max_{|z|=R}|F(z)|
 \le
 \frac{2Z_J}{d}\,2^a d^{a-k}
       R^{-(k-j)}\frac{n^j}{j!}
 \le
 \frac{2Z_J}{d}\,2^a
       R^{-(a-j)}\frac{n^j}{j!}.
 \tag{P-F42v}
\]
For the middle inequality, the bound for \(F\) used in
(P-F42q) remains valid, with \(e^{R|v|}\le2\);
no sum over \(j\) is taken here. The last inequality is
exactly \((dR)^{a-k}\le1\).
Averaging retains a subset of at most \(d\) characters.
Summing over at most \(m^n\) bases and \(n!\) Weyl
shifts, and substituting \(Z_J\le4m^{2n}9^a\), gives
(P-F42t). Every basis in that sum is eventually feasible,
so (P-F42u) holds for each of them.

This proves that a definite upper-degree portion of the
average polynomial already lies below the Fourier-amplitude
threshold. More precisely, fix a real constant
\[
 1-\frac1{4\log_2 3}<\theta<1,
 \qquad J_n=\lceil\theta D\rceil,
 \qquad V_{n,\theta}=\sum_{0\le j<J_n}|b_{n,j}|.
 \tag{P-F42w}
\]
Then
\[
 \limsup_{n\to\infty}
 \frac{\log_2\left(\sum_{j=J_n}^{D}|b_{n,j}|\right)}{n^4}
 \le \frac{(1-\theta)\log_2 3}{6}<\frac1{24}.
 \tag{P-F42x}
\]
We use \(\log_2 0=-\infty\) if the indicated tail vanishes.
To verify the estimate, the ratio of the right-hand side
of (P-F42t) at \(j+1\) to that at \(j\) is
\(nR/(j+1)<1\).
The whole tail is therefore at most \(D+1\) times
the bound at \(J_n\). Taking logarithms,
all terms except \(n(a-J_n)\log_2 3\) are
\(O(n^3\log n)\) from above, uniformly in this fixed
\(\theta\); the negative term \(-\log_2(J_n!)\)
can be omitted for this upper estimate.
Since \(a=n^3/6+O(n^2)\) and
\(J_n=\theta n^3/6+O(n^2)\), (P-F42x) follows.

It follows, with the same strict threshold and with no
assumption on a sign or a cancellation, that
\[
 \limsup_{n\to\infty}\frac{\log_2 U_n}{n^4}<\frac1{24}
 \quad\Longleftrightarrow\quad
 \limsup_{n\to\infty}
        \frac{\log_2 V_{n,\theta}}{n^4}<\frac1{24}.
 \tag{P-F42y}
\]
The forward implication uses \(V_{n,\theta}\le U_n\).
For the reverse implication use
\(U_n=V_{n,\theta}+\sum_{j=J_n}^D|b_{n,j}|\),
(P-F42x), and
\(\log_2(x+y)\le1+\max\{\log_2x,\log_2y\}\)
for nonnegative \(x,y\) not both zero.
For example \(\theta=7/8\) is valid: the highest eighth
of the coefficient indices has upper leading constant
\((\log_2 3)/48<1/24\). The unresolved quantitative
gate is thereby confined to the complementary lower
coefficients. Their bound is not supplied here, so the
strict positivity and saturation hypotheses remain open.

\paragraph{Removing pure-cube factors improves the constant and the tail.}
The full determinant of a basis overestimates the orders of
its finite characters when it contains pure-cube columns.
Taking account of this gives the stronger bounds
\[
 \begin{split}
 U_n&\le8n!\,m^{3n}
        \bigl(1152n^4\,5^{n/2}\bigr)^a,\\
 |b_{n,j}|&\le8n!\,m^{3n}\,18^a
        \bigl(64n^4\,5^{n/2}\bigr)^{a-j}
                       \frac{n^j}{j!}
                 \quad(0\le j\le D).
 \end{split}
 \tag{P-F42z}
\]
In particular
\[
 \limsup_{n\to\infty}\frac{\log_2 U_n}{n^4}
             \le\frac{\log_2 5}{12}.
 \tag{P-F42aa}
\]
This is still greater than \(1/24\).

We give all of the changes to the proof, including the distinction
between the order of a finite group and an integer that annihilates
it. Let \(q\) of the selected basis columns be pure cubes.
They are distinct columns \(3e_i\). After permuting rows and
columns, the basis has the block form
\[
 C=\begin{pmatrix}3I_q&B\\0&D_0\end{pmatrix},
 \qquad r=n-q,\qquad d'=|\det D_0|.
 \tag{P-F42ab}
\]
If \(r=0\), use the empty determinant \(d'=1\).
Otherwise \(D_0\) is a nonsingular integer \(r\)-by-\(r\)
matrix. Each of its columns is the restriction of a cubic
column that is not a pure cube, hence has Euclidean norm
at most \(\sqrt5\). Put \(T=5^{n/2}\).
Hadamard's inequality gives \(1\le d'\le5^{r/2}\le T\),
and every cofactor of \(D_0\) has absolute value at most
\(5^{(r-1)/2}\le T\). The latter assertion includes \(r=1\),
whose cofactor is the empty determinant one.
The entries of \(B\) are at most two, so every row sum of
\(B\) is at most \(2r\).
The block inverse
\[
 C^{-1}=
 \begin{pmatrix}
  I_q/3&-BD_0^{-1}/3\\
  0&D_0^{-1}
 \end{pmatrix}
\]
therefore satisfies
\[
       \|C^{-1}\|_{\max}\le \frac{nT}{d'}.
 \tag{P-F42ac}
\]
For the off-diagonal block use
\((2r/3)T/d'\le nT/d'\);
the diagonal \(1/3\) entries obey the same bound because
\(d'\le T\). If \(r=0\), (P-F42ac) follows directly from
\(C^{-1}=I_n/3\).

Set \(E=3d'\). The displayed block inverse shows that
\(EC^{-1}\) is an integer matrix: the off-diagonal block
is \(-B(d'D_0^{-1})\), and the lower block is
\(3d'D_0^{-1}\). Thus \(E\mathbb Z^n\subseteq C\mathbb Z^n\),
so \(E\) annihilates \(G_J=\mathbb Z^n/C\mathbb Z^n\).
Every value of every character of \(G_J\) is consequently
an \(E\)th root of unity.
The number of characters is still
\(d=|\det C|=3^q d'\); it must not be replaced by \(E\)
in the finite filter. The filter retains its factor \(1/d\)
and its sum over \(d\) characters.

Keep the same common generic direction and the same
\(Z_J\) bound. From (P-F42ac), the column sum three,
and \(|(b_0)_i|\le n-1\), we obtain
\[
 |w_{J,i}|\le\frac{4n^2T}{d'},\qquad
 |v|\le\frac{n^4T}{d'},\qquad |\lambda|\le n.
 \tag{P-F42ad}
\]
For the first inequality,
\(|\ell_JC^{-1}\alpha_i|\le3n\|C^{-1}\|_{\max}\);
the additional one is at most \(T/d'\).
For the second, use
\(\sum_{i,j}|\ell_i(C^{-1})_{ij}(b_0)_j|
 \le n^2(n-1)\|C^{-1}\|_{\max}\).
The last inequality is the feasible-slope result (P-F42u).
Take
\[
                  R'=(64n^4T)^{-1}.
 \tag{P-F42ae}
\]
For \(n\ge4\), these estimates give
\[
 R'|w_{J,i}|\le\frac{3}{16n^2E}\le\frac1{8E},
 \qquad R'(|v|+|\lambda|)\le\frac1{32},
 \qquad ER'\le\frac{3}{64n^4}\le1.
 \tag{P-F42af}
\]
Every nontrivial character value \(\alpha\) satisfies
\(|1-\alpha|\ge4/E\).
The analytic proof in step~5 therefore applies on
\(|z|\le R'\) with the nontrivial denominator bound \(E/2\),
while a fixed denominator still has bound
\(2/|w_{J,i}|\) after multiplication by \(z\).
For the coefficient of \(s^j\), the same Cauchy calculation
now reads
\[
 \bigl|[s^j]P_{J,\chi,b_0}\bigr|
 \le
 \frac{2Z_J}{d}\,2^a E^{a-k}
           (R')^{-(k-j)}\frac{n^j}{j!}
 \le
 \frac{2Z_J}{d}\,2^a
           (R')^{-(a-j)}\frac{n^j}{j!},
 \tag{P-F42ag}
\]
with coefficient zero for \(j>k\).
The last inequality uses \((ER')^{a-k}\le1\).
When bounding the full coefficient norm before summing,
the Taylor coefficients of \(e^{\lambda sz}\) contribute
at most \(e^{R'|\lambda|}\); together with the numerator
factor \(e^{R'|v|}\) this is at most two by (P-F42af).
Hence the full norm per character is at most
\((2Z_J/d)(2/R')^a\), just as in the original proof.
There are at most \(d\) retained characters, at most \(m^n\)
bases, and \(n!\) shifts. Substituting
\(Z_J\le4m^{2n}9^a\) proves both estimates in (P-F42z),
with \(18/R'=1152n^4T\).
Taking logarithms proves (P-F42aa).

The tail and equivalence arguments (P-F42w)--(P-F42y)
now hold with the weaker requirement
\[
       1-\frac1{2\log_2 5}<\theta<1,
 \qquad
 \limsup_{n\to\infty}
 \frac{\log_2\left(\sum_{j=\lceil\theta D\rceil}^{D}
                            |b_{n,j}|\right)}{n^4}
       \le\frac{(1-\theta)\log_2 5}{12}<\frac1{24}.
 \tag{P-F42ah}
\]
Indeed the ratio of the coefficient bounds at successive
indices is \(nR'/(j+1)<1\), and the only term of order
\(n^4\) in their logarithms is
\((a-j)(n/2)\log_2 5\).
For the simple choice \(\theta=4/5\), the highest fifth
has upper leading constant \((\log_2 5)/60<1/24\).
Thus the strict \(1/24\) bound for the actual full norm
\(U_n\) is equivalent to that bound for
\(\sum_{j<\lceil4D/5\rceil}|b_{n,j}|\).
This is a smaller unresolved set of coefficients;
no bound of the required size is proved for it here.

\paragraph{The corresponding signed test at polynomial arguments.}
The controlled tail can also be removed from the exact
pointwise inequality in P-F34, without replacing the
remaining polynomial by its coefficient norm.
Put
\[
 J_n=\lceil4D/5\rceil,\qquad
 L_n(X)=\sum_{j<J_n}b_{n,j}X^j,\qquad
 \eta=\frac1{24}-\frac{\log_2 5}{60}>0.
 \tag{P-F42ai}
\]
For every fixed real \(K\ge0\), all sufficiently large \(n\)
satisfy the uniform estimate
\[
 \sup_{0\le x\le n^K}|B_n(x)-L_n(x)|
       \le 2^{-\eta n^4/2}\,
                 |A_n(e^{2\pi i/2^{n-1}})|.
 \tag{P-F42aj}
\]
Indeed the left side is at most
\(n^{KD}\sum_{j=J_n}^D|b_{n,j}|\).
The proof of (P-F42ah) bounds its base-two logarithm by
\((\log_2 5)n^4/60+O_K(n^3\log n)\).
The amplitude formula (P-F34t) has logarithm
\(n^4/24+O(n^3)\), so their ratio has logarithm at most
\(-\eta n^4+O_K(n^3\log n)\), which is at most
\(-\eta n^4/2\) for all sufficiently large \(n\).

In particular, if either \(x=p(f_n)\) or \(x=s(f_n)+1\)
is at most \(n^K\), its necessary Fourier inequality gives
\[
 L_n(x)\ge
       (1-2^{-\eta n^4/2})
                 |A_n(e^{2\pi i/2^{n-1}})|.
 \tag{P-F42ak}
\]
Thus a bound of the form
\(\sup_{0\le x\le n^K}L_n(x)\le(1-\epsilon_K)|A_n|\),
with some fixed \(\epsilon_K>0\) for every fixed \(K\)
and all sufficiently large \(n\), would also exclude
polynomially bounded strict indices. This is a signed
value comparison and can be useful even if the coefficient
norm has limiting quartic exponent \(1/24\).
For example it would suffice to have
\(b_{n,0}\le(1-2\epsilon_K)|A_n|\) and
\(\sum_{1\le j<J_n}|b_{n,j}|n^{Kj}\le\epsilon_K|A_n|\)
eventually for each fixed \(K\).
These remain sufficient tests for excluding small indices;
the necessary Fourier inequality is not an equivalence
characterizing either index.
No such comparison, or corresponding upper bound on
the constant term \(b_{n,0}\), is established here.


\subsection{A binary-degree-twelve parameter construction correcting a divisibility claim}
\label{sec:binary12-hsop}

\paragraph{P-F43 (a parameter system with no degree divisible by ten).}
Let
\[
 R=\mathbb C[\operatorname{Sym}^{12}\mathbb C^2]^{SL_2},
 \qquad X=\operatorname{Spec}R,
 \qquad A=5544=2^3\,3^2\,7\,11,
 \qquad B=3465=3^2\,5\,7\,11.
\]
For every sufficiently large integer $k\equiv1\pmod {10}$,
there is a homogeneous system of parameters of $R$ with degrees
\[
 \underbrace{Ak,Ak,Ak,Ak,Ak}_{5\text{ entries}},\quad
 \underbrace{Bk,Bk,Bk,Bk,Bk}_{5\text{ entries}}.
 \tag{P-F43a}
\]
One explicit valid choice is $k=10^{12}+1$, giving five copies
of each of the degrees
\[
 a_0=5\,544\,000\,000\,005\,544,\qquad
 b_0=3\,465\,000\,000\,003\,465.
 \tag{P-F43a'}
\]
None of these ten degrees is divisible by ten. This contradicts
the literal assertion of Brouwer--Draisma--Popoviciu,
\emph{The degrees of a system of parameters of the ring of
invariants of a binary form}, Theorem~3.4(ii), printed p.~955,
at binary order $12$ and $t=10$. It also contradicts the
degree-ten necessity stated in the Note on Brouwer's
Poincar\'e-series page. The result does not decide whether
that page's much smaller tuple
$(2,4,5,5,6,6,7,8,9,11)$ is the degree sequence of a parameter
system. The construction supplies explicit valid degrees,
but not explicit polynomial parameters or an efficient
procedure for constructing them.

\paragraph{Proof: uniform eventual nonvanishing from finite generation.}
The grading defines an action of $\mathbb C^*$ on $X$, whose
vertex is denoted $0$. For $x\ne0$, set
\[
 S_x=\{d\ge0:\text{some }f\in R_d\text{ satisfies }f(x)\ne0\},
 \qquad g_x=\gcd(S_x\setminus\{0\}).
 \tag{P-F43b}
\]
Choose finitely many homogeneous positive-degree algebra
generators $h_1,\ldots,h_N$ of $R$. There is at least one
$h_i(x)\ne0$ for every $x\ne0$. The semigroup $S_x$ is exactly
the additive semigroup generated by the degrees of the
$h_i$ which are nonzero at $x$. One inclusion follows by
multiplying these nonzero values. For the other, express a
homogeneous $f$ as a weighted-homogeneous polynomial in the
$h_i$; a nonzero evaluation has at least one nonzero monomial
term, and that term has the same degree as $f$.

Only finitely many semigroups $S_x$ therefore occur. Each
$S_x/g_x$ is a finitely generated subsemigroup of
$\mathbb Z_{\ge0}$ with gcd one and contains all sufficiently
large integers. For completeness, choose a positive element
$a$ of such a semigroup. Its residues modulo $a$ form a
subsemigroup of the finite group $\mathbb Z/a$, hence a
subgroup, and gcd one makes it the whole group. Choose one
semigroup representative of each residue; adding multiples
of $a$ proves eventual membership. Taking a maximum over
the finitely many semigroups gives an integer $C$ such that
\[
 d\ge C\quad\Longrightarrow\quad
 \bigl(d\in S_x\ \Longleftrightarrow\ g_x\mid d\bigr)
 \qquad\text{for every }x\in X\setminus\{0\}.
 \tag{P-F43c}
\]
Moreover the stabilizer of $x$ for the grading action is
exactly $\mu_{g_x}$: a scalar $\lambda$ fixes $x$ if and only
if $\lambda^d=1$ for every degree with a nonzero value there.
In particular every nonvertex grading orbit has dimension one.

\paragraph{Scalar orders on the binary quotient.}
The quotient map
\[
 \pi:\operatorname{Sym}^{12}\mathbb C^2\longrightarrow X
\]
has one closed $SL_2$-orbit in each fiber. The nonzero closed
orbits are represented either by a stable binary form, whose
root multiplicities are all less than six, or by
$c x^6y^6$ with $c\ne0$. These are the standard binary-form
stability and closed-orbit criteria; see Alper--Isaev,
\emph{Associated forms and hypersurface singularities: the
binary case}, Section~5.2, printed p.~13, paragraph preceding
Proposition~5.1. In the strictly semistable case the affine
description also follows directly: put a multiplicity-six
root at $x=0$ and apply the diagonal one-parameter subgroup
which keeps the nonzero $x^6y^6$ term and sends all the other
terms to zero. A closed orbit must contain this limit.

Suppose first that a stable form $F$ represents $x$. It has
at least three distinct roots, so its projective automorphism
group is finite: it permutes that finite set, and a projective
linear transformation fixing three distinct points is the
identity. Choose a primitive $g_x$th root $\lambda$.
The closed orbits of $\lambda F$ and $F$ have the same
quotient point, so some $S\in SL_2$ satisfies
$S\cdot F=\lambda F$. Its projective class has finite order.
Consequently $S$ is diagonalizable, and after an $SL_2$
change of coordinates it has the form
$\operatorname{diag}(z,z^{-1})$. Write
\[
 t=\operatorname{ord}(z^2),\qquad
 F=\sum_{i=0}^{12}c_i x^{12-i}y^i.
\]
All nonzero coefficients have indices in a single residue
class $j$ modulo $t$. Since $F$ is not a monomial, it has
two distinct such indices, and therefore $t\le12$.
In the convention where the indicated diagonal matrix acts
on the displayed monomials with weight $z^{12-2i}$, one has
$\lambda=(z^2)^{6-j}$; the inverse action convention gives
the same order. Thus
\[
               g_x=\frac{t}{\gcd(t,6-j)},\qquad 1\le t\le12.
 \tag{P-F43d}
\]
If $10\mid g_x$, then $g_x=t=10$. Two supported indices
modulo ten force $j\in\{0,1,2\}$. Formula (P-F43d) gives
orders $5,2,5$, respectively, a contradiction.

For $F=cx^6y^6$, every projective transformation taking $F$
to a scalar multiple preserves its unordered pair of roots.
An $SL_2$ lift is diagonal or antidiagonal. A diagonal lift
fixes $F$, and an antidiagonal lift multiplies it by
$(-1)^6=1$. Hence $g_x=1$ in this case. We have proved
\[
              1\le g_x\le12,\qquad 10\nmid g_x
                    \quad(x\in X\setminus\{0\}).
 \tag{P-F43e}
\]

\paragraph{Dimensions of the two exceptional loci.}
Let
\[
 E_2=X^{\mu_2}=\{0\}\cup\{x\ne0:2\mid g_x\},
 \qquad
 E_5=X^{\mu_5}=\{0\}\cup\{x\ne0:5\mid g_x\}.
\]
These are closed conical subsets, and (P-F43e) gives
$E_2\cap E_5=\{0\}$. We claim
\[
                       \dim E_2\le5,\qquad\dim E_5\le2.
 \tag{P-F43f}
\]
For $1\le t\le12$ and $0\le j<t$, let $U_{t,j}$ be the
linear space of binary forms whose supported indices are
congruent to $j$ modulo $t$. The preceding diagonalization
places every point in question in the image of the stable
part of one of finitely many such spaces. The diagonal
$SL_2$ torus preserves $U_{t,j}$ and has one-dimensional
orbits on its stable part, because a stable form has finite
projective stabilizer. Consequently
\[
 \dim\pi(U_{t,j}\cap\mathrm{stable})
       \le \dim U_{t,j}-1.
 \tag{P-F43g}
\]
If $g_x$ is even, $t$ is even. For $t=2$, (P-F43d) forces
$j=1$; there are six odd indices between zero and twelve,
so (P-F43g) bounds the image dimension by five. For $t\ge3$
there are at most $\lfloor12/t\rfloor+1\le5$ indices in any
residue class, giving image dimension at most four. If
$5\mid g_x$, then $t=5$ or $10$ and each residue class has
at most three indices, so its image dimension is at most
two. Strictly semistable closed representatives contribute
to neither locus. Taking the finite union and its closure
proves (P-F43f).

\paragraph{The entire degree spaces have complementary base loci.}
Choose $k\equiv1\pmod {10}$ sufficiently large that
$a=Ak$ and $b=Bk$ are at least $C$. For each integer
$g\le12$, one has
\[
 g\mid a\ \Longleftrightarrow\ 5\nmid g,
 \qquad
 g\mid b\ \Longleftrightarrow\ 2\nmid g.
\]
Indeed $A$ contains every prime power at most twelve except
five, and $k$ is not divisible by five; $B$ contains every
odd prime power at most twelve, and $k$ is odd. By
(P-F43c), the common zero locus in $X$ of the entire space
$R_a$ is $E_5$, and that of $R_b$ is $E_2$.

Since $\dim X=12-2=10$, a nonempty subsequence of at most
five $a$-degrees has base locus of codimension at least
eight, a subsequence of at most five $b$-degrees has base
locus of codimension at least five, and any mixed
subsequence has base locus $\{0\}$ of codimension ten.
These are precisely the hypotheses of Dixmier's criterion,
as stated in Brouwer--Draisma--Popoviciu, Proposition~5.1,
printed p.~958 (preprint p.~6). They already imply (P-F43a).

Here is an independent incidence proof of this last
selection step. Let $\mathcal V=R_a^5\times R_b^5$ be the
space of parameter choices. Stratify $X\setminus\{0\}$
into $E_5\setminus\{0\}$, $E_2\setminus\{0\}$, and their
complement. For a fixed point of the first stratum, the
condition that all ten choices vanish imposes five
independent linear conditions on $\mathcal V$; the other
five vanish automatically. The associated incidence
variety has dimension at most $\dim\mathcal V+2-5$.
For the second stratum the corresponding bound is
$\dim\mathcal V+5-5$, and for the complement it is
$\dim\mathcal V+10-10$. These statements follow equally
by viewing the incidence as the kernel of the evaluation
map of constant rank over each stratum.

Each of these incidence varieties is invariant under
the grading $\mathbb C^*$ action on the point coordinate,
with the parameter choice held fixed: all ten equations
are homogeneous. Its irreducible components are preserved
by this connected group. Every nonempty fiber of an
irreducible component over $\mathcal V$ contains a
one-dimensional grading orbit. The dimension of its image
is therefore at most the dimension of the component minus
one. All the bad images have dimension strictly less than
$\dim\mathcal V$. There are only finitely many components,
so their closures cannot cover $\mathcal V$. A generic
choice has the vertex as its entire common zero set.

For such a choice $f_1,\ldots,f_{10}$, the graded algebra
$R/(f_1,\ldots,f_{10})$ is finite-dimensional by the
Nullstellensatz. Lifting a finite homogeneous vector-space
basis and inducting on degree shows that $R$ is finite as
a module over $\mathbb C[f_1,\ldots,f_{10}]$. This subring
has dimension ten because a finite integral extension
preserves dimension. It is generated by ten elements and
is a domain, so these elements are algebraically
independent: a nonzero prime kernel in the polynomial
ring on ten variables would lower its dimension. They
are therefore a homogeneous system of parameters.
Finally $5\nmid a$ and $2\nmid b$, proving the assertion
about divisibility by ten.\hfill$\square$

\paragraph{An explicit uniform threshold and degree tuple.}
To verify the advertised value $k=10^{12}+1$, use Derksen,
\emph{Polynomial bounds for rings of invariants},
Theorem~1.1, Proposition~1.2, and Example~2.1, printed
p.~956. In characteristic zero these give for binary
forms of degree $d$ an algebra-generator degree bound
$\frac32(d+1)d^6$. The representation of $SL_2$ has finite
kernel, as required by the cited proposition, and its
matrix coefficients have degree at most $d$ in the
matrix entries. At $d=12$ we may therefore choose the
generators above with degrees at most
\[
              \mathcal B=\frac32\cdot13\cdot12^6
                        =58\,226\,688.
 \tag{P-F43h}
\]

Here is the needed elementary uniform conductor estimate.
If an additive semigroup is generated by positive integers
at most $\mathcal B$ and has gcd $g$, divide its generators
by $g$ and let $h$ be the smallest normalized generator.
In the directed graph of residues modulo $h$, with an
edge for addition of each normalized generator, every
residue is reachable from zero. A shortest such path
repeats no vertex, so has at most $h-1$ edges. Hence each
residue has a semigroup representative at most
\[
            (h-1)\mathcal B/g\le(\mathcal B/g)^2.
\]
Every larger integer in the same residue belongs to the
normalized semigroup, by adding copies of $h$.
If $N\ge\mathcal B^2$ and $g\mid N$, then
$N/g\ge\mathcal B^2/g\ge(\mathcal B/g)^2$.
Thus the original semigroup contains every $g$-multiple
at least $\mathcal B^2$. This proof applies uniformly to
all the subsets of generators nonzero at a point, so
$C=\mathcal B^2$ is valid in (P-F43c).

Exact integer arithmetic gives
\[
 \begin{split}
 \mathcal B^2&=3\,390\,347\,195\,449\,344,\\
 3465(10^{12}+1)&=3\,465\,000\,000\,003\,465
                        >\mathcal B^2.
 \end{split}
\]
The other degree is larger still, and $10^{12}+1\equiv1$
modulo ten. The proof just given therefore applies to
the concrete tuple (P-F43a'). The only arithmetic
computation here checks displayed integer values; the
existence assertion follows from the full argument above.

\paragraph{Exact source discrepancy and scope.}
Theorem~3.4(ii) asserts at least $\lfloor(n-1)/t\rfloor$
degrees divisible by $t$ for every integer $t>1$ and every
even binary order $n$. Its preceding Lemma~3.3(ii) gives
instead $\lfloor(n-j)/t\rfloor$, where $j$ is the least
nonnegative integer with $\gcd(n/2-j,t)=1$.
For $(n,t)=(12,10)$ this least $j$ is three, so the lemma
only gives zero. The printed improvement writes $n=at+b$
and replaces a condition on $\gcd(n-2i,2t)$ by one on
$\gcd(b-2i,2t)$. This replacement is invalid when $a$ is
odd: at $n=12,t=10,a=1,b=2,i=0$ the two gcds are four
and two. The construction above proves that this is a
defect in the stated necessity, not merely an unfilled
proof step. It makes no priority claim and does not
invalidate other independently proved results in that paper.

The exact primary sources used here are
\begin{itemize}
\item Brouwer--Draisma--Popoviciu, Theorem~3.4(ii), p.~955,
and Proposition~5.1, p.~958:
\url{https://pure.tue.nl/ws/files/17749340/Degrees2015.pdf};
the latter is Proposition~5.1, p.~6, in
\url{https://arxiv.org/pdf/1404.5722}.
\item Alper--Isaev, Section~5.2, p.~13, paragraph preceding
Proposition~5.1 (binary stability and the strictly semistable
closed orbit):
\url{https://maths-people.anu.edu.au/~alperj/papers/binary.pdf}.
\item Brouwer's Poincar\'e-series page, Note paragraph
following the two displayed expressions for binary order
twelve:
\url{https://aeb.win.tue.nl/math/poincare.html}.
\item Derksen, \emph{Polynomial bounds for rings of
invariants}, Theorem~1.1, Proposition~1.2, and Example~2.1,
p.~956:
\url{https://sites.lsa.umich.edu/hderksen/wp-content/uploads/sites/614/2018/05/A.I.a.7.pdf}.
\end{itemize}


\subsection{The primitive Klein contribution inside the actual average polynomial}

\begin{proposition}[P-F44: an exact signed global cone contribution]
\label{prop:kleinglobalcontribution}
Let $n\ge4$ with $3\nmid n$, and retain $B_n$ and the exact
Weyl--Brion expansion of P-F42. Set
\[
 P=\frac{2^n-(-1)^n}{3},\qquad D=\binom n3,
\]
and let $C_a$ be the moving dual-normal determinant coefficient
of the Klein cubic in P-F40, for $a\in(\mathbb Z/P\mathbb Z)^\times$.
There is the following exact decomposition of the \emph{actual}
average constituent polynomial:
\[
 B_n(X)=\mathcal K_n+\mathcal R_n(X),\qquad
 \mathcal K_n=
   \frac1{nP}\sum_{\substack{a\bmod P\\(a,P)=1}}C_a
   =\frac1{nP}\operatorname{Tr}_{\mathbb Q(\zeta_P)/\mathbb Q}(C_1).
 \tag{P-F44a}
\]
Here $\mathcal R_n$ is given explicitly below by all the other
balanced character terms in the global Weyl--Brion expansion.
The removed terms are exactly the primitive Klein characters
on directed Hamiltonian-cycle bases, with all three central
cube-root lifts included.
They have degree zero as polynomials in $X$, and their total
is independent of the perturbation and specialization direction
used in P-F42.

If $n\equiv3\pmod4$, then
\[
                              \mathcal K_n=0.
 \tag{P-F44b}
\]
Thus the primitive Klein terms cancel in this particular
part of the actual global average polynomial, despite the
large individual amplitudes in P-F40.
For even $D$, formula (P-F44a) is a signed trace, and no
nonvanishing, sign, or dominance over $\mathcal R_n$ is claimed.
The normalization is $1/(nP)$, rather than the $1/P$ appearing
in the cyclic normal-model Molien sum.
\end{proposition}

\begin{proof}
We first identify a general collection of degree-zero cone terms.
Use the matrix $A$, the Weyl shifts $b_\sigma=\delta-\sigma\delta$,
and the character-polynomials $P_{J,\chi,b_\sigma}$ of P-F42.
Write $z_i=\chi(e_i)$ for a character of
$\mathbb Z^n/A_J\mathbb Z^n$.
Suppose the following hold:
\[
 \begin{gathered}
 A_J^{-1}\mathbf1\ \hbox{has all coordinates strictly positive},
 \qquad \chi(3\mathbf1)=1,\\
 z_1,\ldots,z_n\ \hbox{are distinct},\qquad
 \{\alpha:z^\alpha=1,\ |\alpha|=3\}=J.
 \end{gathered}
 \tag{P-F44c}
\]
The character is balanced because $c=3$ when $3\nmid n$.
The first condition means that $J$ is eventually feasible
for every Weyl shift and every fixed perturbation allowed
in P-F42. The last condition means that no nonbasic
denominator factor has character one, so the number $k$
in that proof is zero.
Consequently (P-F42l) has a constant term independent of
the stretch variable and the specialization direction:
\[
 P_{J,\chi,b_\sigma}(X)=
 \frac{z^{-b_\sigma}}
 {d_J\prod_{\alpha\notin J}(1-z^\alpha)},\qquad
 d_J=|\det A_J|.
 \tag{P-F44d}
\]
It is retained by averaging, since the character is balanced.

Put
\[
 \Delta(z)=z^{-\delta}
                  \sum_{\sigma\in S_n}
                         \operatorname{sgn}(\sigma)z^{\sigma\delta}
          =\prod_{i<j}(1-z_j/z_i).
 \tag{P-F44e}
\]
Summing (P-F44d) over the signed Weyl shifts therefore gives
\[
 T(J,z)=
 \frac{\Delta(z)}
 {d_J\prod_{\alpha\notin J}(1-z^\alpha)}.
 \tag{P-F44f}
\]
Permuting the variables permutes pairs $(J,z)$, preserves
all the conditions in (P-F44c), and preserves the denominator
product and $d_J$. This action on pairs is free because the
$z_i$ are distinct.
The elementary alternant identity
\[
 \sum_{\tau\in S_n}\Delta(\tau z)
       =\prod_{i\ne j}(1-z_i/z_j)
 \tag{P-F44g}
\]
then shows that a complete variable-permutation orbit of
such pairs contributes exactly
\[
 \frac1{d_J}\,
 \frac{\prod_{i\ne j}(1-z_i/z_j)}
      {\prod_{\alpha\notin J}(1-z^\alpha)}.
 \tag{P-F44h}
\]
For clarity, if
$A_\delta(z)=\sum_\sigma\operatorname{sgn}(\sigma)
z^{\sigma\delta}$, the sum in (P-F44g) is
$A_\delta(z)A_\delta(z^{-1})$.
The Vandermonde product gives precisely the right-hand
side of (P-F44g), including its sign.

We now apply this calculation to the Klein bases.
A directed Hamiltonian cycle on the $n$ labelled variables
gives the $n$ monomials $x_i^2x_{\nu(i)}$, where $\nu$ is
an $n$-cycle. There are $(n-1)!$ distinct bases of this
form. For the standard cycle, with indices $0,\ldots,n-1$,
the basis matrix is $2I$ plus a cyclic permutation matrix.
Thus
\[
 d_J=2^n-(-1)^n=3P,\qquad
                  A_J^{-1}\mathbf1=\tfrac13\mathbf1.
 \tag{P-F44i}
\]
The determinant formula follows directly from the two
nonzero terms in its determinant expansion.

Let $\rho=e^{2\pi i/P}$ and $w_i=(-2)^i$.
For each $\omega\in\mu_3$ and each unit $a\bmod P$, take
\[
                   z_i=\omega\rho^{a w_i}.
 \tag{P-F44j}
\]
These are $3\varphi(P)$ distinct characters for the
standard basis. Distinctness follows from
$\mu_3\cap\mu_P=\{1\}$, since $3\nmid P$.
Indeed $(-2)^n=(1-3)^n\equiv1-3n\pmod9$, so
$(1-(-2)^n)/3\equiv n\pmod3$; this integer is $P$
up to sign, and $3\nmid n$.
They annihilate every basis monomial and hence are
characters of the finite group of order $3P$.
Their determinant is $\omega^n$, because the Klein
matrix in P-F40 belongs to $SL_n$. In particular
$\chi(3\mathbf1)=1$.
P-F40 proves that their eigenvalues are distinct and
their fixed cubic monomials are exactly the displayed
$n$ cycle monomials. Multiplication by $\omega$ has no
effect on a cubic monomial character. Thus every
condition in (P-F44c) holds.
Transport these characters by variable permutations to
every directed-cycle basis. Denote the resulting set
of pairs by $\mathscr K_n$.

We count its variable-permutation orbits precisely.
The subgroup preserving the standard directed-cycle
basis is the cyclic group of $n$ rotations. It acts
on $a$ by multiplication by powers of $-2$.
The order of $-2$ modulo $P$ is exactly $n$: its
$n$th power is one, and the distinct Klein eigenvalues
proved in P-F40 exclude a smaller positive exponent.
Thus each orbit on units has size $n$.
The central index $\omega$ cannot be mixed by a variable
permutation: it is determined by the determinant
$\omega^n$, since $3\nmid n$.
It follows that the orbits of $\mathscr K_n$ under all
of $S_n$ are indexed by
\[
 \omega\in\mu_3,\qquad
 a\in(\mathbb Z/P\mathbb Z)^\times/\langle-2\rangle.
 \tag{P-F44k}
\]
Equivalently,
$|\mathscr K_n|=3(n-1)!\varphi(P)$ and every full
permutation orbit has size $n!$, giving the same count.

The determinant quotient in (P-F44h) is independent of
$\omega$. For the Klein cubic, exactness of the tangent
sequence identifies it with the reciprocal moving
\emph{primal} normal determinant. P-F40 defines $C_a$
using the dual normal representation, so the quotient
is $C_{-a}$, rather than $C_a$:
\[
 \frac{\prod_{i\ne j}(1-z_i/z_j)}
      {\prod_{\alpha\notin J}(1-z^\alpha)}
       =\frac1{\det(1-\beta^a\mid S_{\rm mov})}
       =C_{-a}.
 \tag{P-F44l}
\]
This explicitly accounts for dual conventions and
possible signs.
The quantity $C_a$ is constant on a class modulo
multiplication by $-2$, because rotation conjugates
the corresponding Klein automorphisms.
Combining (P-F44h)--(P-F44l), the total contribution
of $\mathscr K_n$ is
\[
 \sum_{\omega\in\mu_3}
 \sum_{[a]\bmod\langle-2\rangle}\frac{C_{-a}}{3P}
   =\frac1{nP}\sum_{(a,P)=1}C_{-a}
   =\frac1{nP}\sum_{(a,P)=1}C_a.
 \tag{P-F44m}
\]
The last equality reindexes units by negation.
Each $C_a$ is the Galois conjugate of $C_1$ under
$\rho\mapsto\rho^a$, which proves the trace formula
and rationality in (P-F44a).
P-F40 proves $C_{-a}=(-1)^D C_a$.
Since $D$ is odd exactly when $n\equiv3\pmod4$,
pairing units under negation proves (P-F44b).
Here $P>2$, so no unit is its own negative.

To specify the remainder without making any local-to-global
identification, choose any perturbation and common direction
as in P-F42. For a Weyl shift $\sigma$, let
$\mathscr B_\sigma$ be its eventually feasible bases.
Then
\[
 \mathcal R_n(X)=
 \sum_{\sigma\in S_n}\operatorname{sgn}(\sigma)
 \sum_{J\in\mathscr B_\sigma}
 \sum_{\substack{\chi\in\widehat{\mathbb Z^n/A_J\mathbb Z^n}\\
                 \chi(3\mathbf1)=1,\ (J,\chi)\notin\mathscr K_n}}
              P_{J,\chi,b_\sigma}(X).
 \tag{P-F44n}
\]
All sets and polynomial terms in this expression are
exactly those already constructed in P-F42.
Every Klein basis occurs in every $\mathscr B_\sigma$
by (P-F44i), so removing precisely $\mathscr K_n$
from that global expansion removes the constant
(P-F44m) and nothing else.
This proves the asserted decomposition of $B_n$.
Although the displayed representation of the remainder
uses the chosen perturbation and direction, its value
is the intrinsic polynomial $B_n-\mathcal K_n$.
\end{proof}

\paragraph{What this does and does not settle.}
Formula (P-F44a) supplies a global comparison for this
specified collection of terms, derived from the exact
lattice-point expansion. It does not identify the cyclic
normal-model Hilbert series with the global invariant
Hilbert series.
The large value of an individual $C_a$ does not bound
its Galois trace, and the trace can still cancel against
the constant term of $\mathcal R_n$.
More explicitly,
\[
 U_n=
   \bigl|\mathcal K_n+[X^0]\mathcal R_n\bigr|
       +\sum_{j\ge1}\bigl|[X^j]\mathcal R_n\bigr|.
 \tag{P-F44o}
\]
Thus the remaining quantitative gate is an explicitly
signed coefficient estimate for (P-F44n), together with
its possible constant-term cancellation with the trace.
No unsigned dominance or index lower bound is inferred.

\paragraph{A squarefree eigenbasis for the moving Klein normal module.}
There is a useful exact simplification of the coefficient in
(P-F44a). As a representation of the specified cyclic group
$\Gamma=\langle\beta\rangle$,
\[
                  S_{\rm mov}\cong\bigwedge\nolimits^3\mathbb C^n,
 \qquad
 C_a=
 \prod_{0\le i<j<k<n}
       \left(1-\rho^{-a(w_i+w_j+w_k)}\right)^{-1}.
 \tag{P-F44p}
\]
The exterior power here carries the diagonal action with
eigenvalues $\rho^{w_i}$ on $\mathbb C^n$.
The isomorphism is a statement about this cyclic stabilizer
representation, not an assertion of an $SL_n$-module structure
on a normal space at a nonfixed point.

Here is a direct proof, including the cyclic boundary cases.
Let $T=\mathbb C[x_0,\ldots,x_{n-1}]$ and let $J_K$ be the
Jacobian ideal of $K_n$. Its generators are
\[
                 x_i^2+2x_{i+1}x_{i+2},
                         \qquad i\bmod n.
 \tag{P-F44q}
\]
The smoothness proof in P-F40 shows that the partial
derivatives have common zero only at the origin.
They form a regular sequence of quadrics. There are
therefore no relations with linear coefficients among
them: equivalently, the degree-three part of the first
Koszul syzygy module is zero, as each basic Koszul
syzygy already has degree four. Hence
\[
                    \dim(T/J_K)_3
                       =\binom{n+2}{3}-n^2=D.
 \tag{P-F44r}
\]

Every degree-three monomial reduces modulo $J_K$ to a
linear combination of squarefree monomials, in fact to
a scalar multiple of one or to zero. Cubes satisfy
$x_i^3=-2x_ix_{i+1}x_{i+2}$.
For $x_i^2x_j$ with $j\ne i$, use
$x_i^2x_j=-2x_{i+1}x_{i+2}x_j$.
If $j$ is different from $i+1,i+2$, this is squarefree.
When $j=i+1$, put $u_i=x_i^2x_{i+1}$; then
$u_i=-2u_{i+1}$, and going once around the cycle gives
$u_i=(-2)^n u_i$. Thus $u_i=0$ in characteristic zero.
When $j=i+2$, a second reduction gives
\[
                  x_i^2x_{i+2}
                     =4x_{i+1}x_{i+3}x_{i+4}.
 \tag{P-F44s}
\]
These three indices are distinct modulo every $n\ge4$:
their pairwise positive differences are one, two, and
three. In particular, at $n=4$ the final index is $i$,
and the three remaining indices are still distinct.
There are exactly $D$ squarefree cubic monomials, so
the spanning argument and (P-F44r) make their classes
a basis.

The degree-three Jacobian subspace is the image
$\mathfrak{gl}_n\cdot K_n$.
The absence of linear syzygies makes this Lie-orbit
map injective. Its subspace
$\mathfrak{sl}_n\cdot K_n$ has codimension one, with
the complementary invariant Euler line represented
by $K_n$ itself.
It follows that the quotient normal module $S$ in
P-F40 is an extension of $(T/J_K)_3$ by a trivial
one-dimensional representation.
Representations of the finite group $\Gamma$ over
$\mathbb C$ split by averaging.
Moreover P-F40 proves $\dim S^\Gamma=1$, so
$(T/J_K)_3$ has no fixed vectors and is exactly the
moving module.
Each squarefree basis vector has eigenvalue
$\rho^{w_i+w_j+w_k}$. Sending
$e_i\wedge e_j\wedge e_k$ to its squarefree class is
therefore a $\Gamma$-equivariant vector-space
isomorphism. This proves (P-F44p), including the
negative exponent required for the dual coefficient.

Thus the specified global constant contribution is
the completely explicit cyclotomic trace
\[
 \mathcal K_n=\frac1{nP}
 \operatorname{Tr}_{\mathbb Q(\rho)/\mathbb Q}
 \left[
   \prod_{0\le i<j<k<n}
          (1-\rho^{-(w_i+w_j+w_k)})^{-1}
 \right].
 \tag{P-F44t}
\]
This expression does not remove the signed trace or the
remaining polynomial in (P-F44n), and it is not used to
claim a lower bound for $U_n$.

\paragraph{Three exact finite trace values, including a negative contribution.}
The specified global contributions in three small cases are
\[
              \mathcal K_4=\frac1{25},\qquad
              \mathcal K_5=\frac2{121},\qquad
              \mathcal K_8=-\frac{624}{85}.
 \tag{P-F44u}
\]
For $n=4$, $P=5$ and the four dual squarefree weights
are all the nonzero residues modulo five.
For $n=5$, $P=11$ and the variable weights are
$1,9,4,3,5$ modulo eleven. Their sum is zero, and their
ten pair sums are the ten nonzero residues. Thus the
ten dual squarefree triple weights are also all the
nonzero residues. In both cases
$C_a=1/P$, since
$\prod_{j=1}^{P-1}(1-\rho^j)=P$.
Equation (P-F44a) gives the first two displayed values.

Here is a complete finite polynomial certificate for the
third value. Put
\[
 \Phi(x)=\Phi_{85}(x)
       =\frac{(x^{85}-1)(x-1)}{(x^5-1)(x^{17}-1)},\qquad
 D_8(x)=\prod_{0\le i<j<k<8}(1-x^{r_{ijk}}),
\]
where $r_{ijk}$ is the representative in $\{1,\ldots,84\}$
of $-((-2)^i+(-2)^j+(-2)^k)$ modulo $85$.
Define $N(x)=\sum_{j=0}^{63}N_jx^j$ by the following
table; a row labelled $j$ gives $N_j,\ldots,N_{j+7}$:
\[
\begin{array}{r|rrrrrrrr}
 j&N_j&N_{j+1}&N_{j+2}&N_{j+3}&N_{j+4}&N_{j+5}&N_{j+6}&N_{j+7}\\\hline
 0&-1167&0&1048&-808&0&708&0&454\\
 8&0&-708&254&0&0&-708&100&-340\\
16&0&708&-708&340&-100&0&708&100\\
24&-354&-340&-708&708&100&100&-694&100\\
32&708&-708&354&-694&340&-100&-708&808\\
40&354&-354&708&-708&-100&354&100&0\\
48&-708&0&-694&354&0&340&0&-694\\
56&808&808&-254&0&-1048&-100&808&-708
\end{array}
\]
Exact integer polynomial multiplication and monic division give
\[
                         D_8(x)N(x)\equiv17\pmod{\Phi(x)}.
 \tag{P-F44v}
\]
For a reproducible check requiring no symbolic inversion,
first obtain $\Phi$ by dividing
$1+x^5+\cdots+x^{80}$ by $1+x+\cdots+x^{16}$.
Starting at one, multiply the $56$ factors defining $D_8$,
reducing by this monic polynomial after each multiplication.
Multiply the result by the displayed $N$ and reduce once
more; the remainder is the constant $17$.
The independent integer-arithmetic verifier
\texttt{verify\_klein85\_certificate.py} implements exactly
these divisions and was executed successfully.
All the inputs to this finite identity are displayed above.
It follows that $C_1=N(\rho)/17$ for a primitive
$85$th root $\rho$.

The trace of $\rho^j$ is $64$ for $j=0$, is $-4$ for a
positive multiple of five less than $64$, is $-16$ for
a positive multiple of seventeen less than $64$, and
is one for the remaining positive indices below $64$.
Indeed, summing over residues coprime to $85$ and using
inclusion--exclusion on multiples of five and seventeen
gives
\[
 \sum_{(a,85)=1}\rho^{aj}
   =85\,\mathbf1_{85\mid j}
          -17\,\mathbf1_{17\mid j}
          -5\,\mathbf1_{5\mid j}+1.
\]
The table gives respectively
$N_0=-1167$, the three other coefficient-group sums
$-2934$, $1416$, and $744$. Hence
\[
 \operatorname{Tr}(C_1)
  =\frac{64(-1167)-4(-2934)-16(1416)+744}{17}
  =-4992.
 \tag{P-F44w}
\]
Dividing by $nP=8\cdot85$ proves the last value in
(P-F44u).
It is an exact negative contribution to the specified
part of the actual average polynomial.
It is not asserted to equal its full constant coefficient,
and gives no value or bound for $U_8$.


\subsection{The exact ball range of the full cubic-cover eigenspace}
\label{sec:cyclic-cover-range}

\paragraph{P-F45 (the natural full-eigenspace construction has a fixed ball range).}
Let $n\ge4$ and let $F(x_1,\ldots,x_n)$ define a smooth cubic
hypersurface in $\mathbb P^{n-1}$. Put
\[
 Y_F=\{w^3+F(x_1,\ldots,x_n)=0\}\subset\mathbb P^n,
 \qquad d=\dim Y_F=n-1,
 \qquad \sigma(w,x)=(\zeta w,x),
\]
where $\zeta$ is a primitive cube root of unity. Let $V_\zeta$
be the full $\zeta$-eigenspace in the primitive middle cohomology
$H^d(Y_F,\mathbb C)_{\rm prim}$, and let $\mathcal D_n$ be
the connected full classifying domain of polarized Hodge
filtrations on $V_\zeta\oplus\overline{V_\zeta}$, with this
deck action and these Hodge numbers. Then
\[
 h^{n-1-p,p}(V_\zeta)=\binom n{3p+2-n},
 \qquad
 h^{n-1-p,p}(V_{\zeta^2})=\binom n{3p+1-n},
 \tag{P-F45a}
\]
where a binomial coefficient outside $0,\ldots,n$ is zero.
For $n\ge4$, the domain $\mathcal D_n$ is a complex ball
if and only if $n=4$ or $n=5$. In these cases its dimensions
are four and ten, respectively. The first signatures of the
Hermitian polarization on $V_\zeta$, with the convention
that the pieces with even $p$ are positive, are
\[
\begin{array}{c|c|c}
n&\text{nonzero pairs }(p,h^{n-1-p,p})& (h_+,h_-)\\ \hline
4&(1,4),(2,1)&(1,4)\\
5&(1,1),(2,10)&(10,1)\\
6&(2,15),(3,6)&(15,6)\\
7&(2,7),(3,35),(4,1)&(8,35)\\
8&(2,1),(3,56),(4,28)&(29,56).
\end{array}
\tag{P-F45b}
\]
For $n=6$, $\mathcal D_6$ is the matrix ball of type
$I_{15,6}$, of complex dimension $90$, and is not a complex
ball. For every $n\ge7$, $\mathcal D_n$ contains a nonconstant
compact holomorphic curve isomorphic to $\mathbb P^1$;
in particular it cannot be biholomorphic to a bounded domain
in any complex affine space.

\paragraph{Proof of the Hodge formula and signatures.}
The eigenspace residue formula is van Geemen--Izadi,
\emph{Half twists and the cohomology of hypersurfaces},
Lemma~2.4, printed p.~285, with the cyclic-cover convention
of Section~2.1, pp.~283--284. It says that the $i$th
eigenspace in $H^{k-p,p}(Y_k)_{\rm prim}$ is the degree
$3(p+1)-k-1-i$ part of the Jacobian algebra of the branch
cubic in $k+1$ variables. Substituting $k=n-1$ and
$i=1,2$ gives the two degrees in (P-F45a).
For a smooth cubic the $n$ quadratic partial derivatives
have no common projective zero. They form a homogeneous
regular sequence in the polynomial ring on $n$ variables:
their ideal has height $n$, and a polynomial ring is
Cohen--Macaulay. Successive quotient by a homogeneous
nonzerodivisor of degree two multiplies its Hilbert series
by $1+t$ after canceling one factor $1-t$. Thus its
Jacobian algebra has Hilbert series
\[
                      \frac{(1-t^2)^n}{(1-t)^n}
                            =(1+t)^n.
\]
This proves (P-F45a). The same coefficient count can be
read directly from the squarefree monomial basis of the
Fermat cubic, as in the cited lemma's proof.

For clarity about the eigenspace convention, the residue
volume form acquires a factor $\zeta$ under $\sigma$.
In the Jacobian algebra of $w^3+F$ the remaining powers
of $w$ are $1,w$; these therefore contribute eigenvalues
$\zeta,\zeta^2$, respectively. In particular the first
formula does not interchange the two characters silently.

Let $Q$ denote the real bilinear polarization on primitive
middle cohomology. Since $\sigma$ preserves $Q$, the two
eigenspaces pair with one another and each is isotropic
for the bilinear pairing. The form
\[
                        h(v,u)=i^d Q(v,\overline u)
                        \quad(v,u\in V_\zeta)
\]
is Hermitian: the symmetry $Q(u,v)=(-1)^dQ(v,u)$ verifies
the conjugate symmetry directly. The Hodge--Riemann
relations give orthogonality of different Hodge pieces,
with sign $(-1)^p$ on $H^{d-p,p}(V_\zeta)$ for this
normalization. An exact primary statement is Griffiths,
\emph{Periods of integrals on algebraic manifolds: summary
of main results and discussion of open problems},
equation~(9.9), printed p.~254. Its sign convention is
$(-i)^d(-1)^{d-p}Q(v,\overline v)>0$, which is exactly
$(-1)^p h(v,v)>0$. Evaluating (P-F45a) gives (P-F45b).

\paragraph{The domains with two Hodge pieces.}
When precisely two adjacent Hodge pieces are nonzero,
the full filtration is specified by a maximal positive
or negative subspace for $h$. Choose a fixed orthogonal
splitting $\mathbb C^{a+b}=P\oplus N$ with positive
dimension $a$ and negative dimension $b$. Every negative
$b$-plane projects isomorphically onto $N$: a nonzero
vector in the kernel would be both positive and negative.
It is therefore the graph of a unique map $Z:N\to P$,
and negativity is equivalent to $I-Z^*Z>0$. This gives
holomorphic coordinates when the higher Hodge piece is
negative. If the higher piece is positive, use instead
its graph $P\to N$; transposing its matrix gives the
same displayed matrix domain holomorphically. Thus in
either case the domain is
\[
              \{Z\in\operatorname{Mat}_{a\times b}(\mathbb C):
                        \|Z\|_{\rm op}<1\}.
 \tag{P-F45c}
\]
For $n=4,5$ one of $a,b$ is one, so the operator norm is
the Euclidean norm and (P-F45c) is the four- or ten-ball.
For $n=6$ the two dimensions are $15,6$; its dimension
is $15\cdot6=90$.

Here is a direct proof that this last matrix domain is
not biholomorphic to a ball, without relying on a
classification of symmetric domains. For a domain $D$
containing the origin, define the intrinsic infinitesimal
seminorm
\[
 c_D(0;v)=\sup\{|df_0(v)|:f:D\to\mathbb D
                  \text{ holomorphic},\ f(0)=0\}.
\]
For the domain (P-F45c), the Schwarz lemma on the disc
$z\mapsto z v/\|v\|_{\rm op}$ gives
$c_D(0;v)\le\|v\|_{\rm op}$. Conversely, choose unit
vectors $u,w$ with $|u^*vw|=\|v\|_{\rm op}$; the function
$Z\mapsto u^*Zw$ maps $D$ to $\mathbb D$ and attains
this derivative. Thus $c_D(0;v)=\|v\|_{\rm op}$.
The identical argument for a Euclidean ball gives its
Euclidean norm. The seminorm is preserved by the
derivative of a biholomorphism, directly from its definition.

A putative biholomorphism to a ball can be composed with
a ball automorphism to send the origin to the origin.
For completeness, when $a\ne0$ the usual map
\[
 z\longmapsto
 \frac{a-P_a z-\sqrt{1-\|a\|^2}(I-P_a)z}
      {1-\langle z,a\rangle},
 \qquad P_a z=\frac{\langle z,a\rangle}{\|a\|^2}a,
\]
interchanges $0$ and $a$; expanding the squared norm
gives
$1-\|\phi_a(z)\|^2=(1-\|a\|^2)(1-\|z\|^2)
/|1-\langle z,a\rangle|^2$, and direct substitution makes
it an involution. The case $a=0$ is immediate.
The derivative would now identify the operator norm
linearly with a Hermitian norm. This is impossible when
$a,b\ge2$: for $X=E_{11}$ and $Y=E_{22}$, the operator
norms of $X,Y,X+Y,X-Y$ are all one, so the parallelogram
identity has sides two and four. Every Hermitian norm
satisfies that identity. This proves the assertion for
$n=6$.

\paragraph{Compact curves for every higher dimension.}
For $n\ge7$, let $j_0\in\{0,1,2\}$ be congruent to
$2-n$ modulo three. The three integers
$j_0,j_0+3,j_0+6$ all lie between zero and $n$.
For $n=7,8$ this is immediate, and for $n\ge9$ their
maximum is at most eight. They correspond in (P-F45a)
to three consecutive nonzero Hodge pieces, at indices
$p_0,p_0+1,p_0+2$. The first and third have the same
definite sign for $h$ and are orthogonal. Choose one
nonzero line in each and let their sum be a definite
two-dimensional subspace $S$. Keep every other Hodge
direction fixed. For every line $\ell\subset S$, place
$\ell$ in the first of these Hodge pieces and its
$h$-orthogonal complement in $S$ in the third. The
middle piece stays fixed. This gives polarized Hodge
decompositions with the same dimensions and deck action.

Their filtrations vary holomorphically with
$\ell\in\mathbb P(S)\simeq\mathbb P^1$: at every
filtration level the moving contribution is either
$0$, the line $\ell$, or the fixed space $S$. Although
the orthogonal-complement description of the lower
summand uses conjugation, the filtration itself does
not; it is the holomorphic tautological-line construction.
The conjugate eigenspace filtration is equivalently
the annihilator filtration under $Q$, which is also
holomorphic. The curve is injective, since a filtration
level between the two chosen pieces recovers $\ell$,
and lies in the connected component of the starting
Hodge structure. It is therefore a nonconstant compact
holomorphic $\mathbb P^1$ in $\mathcal D_n$.
Every holomorphic map from $\mathbb P^1$ to a domain
in $\mathbb C^N$ is constant, by the maximum principle
applied to its coordinate functions. Hence
$\mathcal D_n$ is not biholomorphic to a bounded domain,
and in particular not to a ball. Together with the
two-piece cases this proves the exact claimed range.
\hfill$\square$

\paragraph{Scope for the theta comparison.}
The domain in this proposition is the \emph{full}
polarized eigenspace classifying domain of this particular
cyclic triple cover. It is not the horizontal period image
of a restricted family. The result explains why the actual
four- and ten-ball constructions used for cubic surfaces
and cubic threefolds do not, with the same full-eigenspace
definition, give a sequence of balls as $n$ grows.
It excludes neither special ball subdomains, selected
sub-Hodge structures, half twists, different covers, nor
other orthogonal or unitary theta constructions. It supplies
no obstruction to those alternatives and no positivity or
complexity bound. The fixed-$n=5$ identification in P-F41
is unaffected.

The primary sources used for the Hodge premises are
\begin{itemize}
\item van Geemen--Izadi, Section~2.1, pp.~283--284,
Lemma~2.4, p.~285 (the exact cyclic-cover eigenspace formula):
\url{https://mathweb.ucsd.edu/~eizadi/papers/MathZ2002.pdf}.
\item Griffiths, Section~2, pp.~234--235 (the full classifying
domain and its compact flag fibers), and equation~(9.9),
p.~254 (the polarization relations):
\url{https://publications.ias.edu/sites/default/files/periodsofintegral.pdf}.
\end{itemize}


\subsection{The exact positive theta input for the cubic discriminant}
\label{sec:theta-input}

\paragraph{P-F46 (theta contraction, with its gluing and coefficient spaces).}
Use the cubic discriminant form $\phi$ and the period-ring map
$\mathcal P$ from P-F41. Let $L=II_{2,26}$, and let
$M\subset L$ be the primitive cubic-fourfold lattice with
signature $(2,20)$; this is the sign reversal of the
$(20,2)$ convention in Casalaina-Martin--Laza. Its orthogonal
complement is $K(-1)$, where $K=E_6$ is positive definite.
For an even lattice write $A_M=M^\vee/M$ and use the
discriminant quadratic form $(x,x)/2$ modulo $\mathbb Z$.
The unimodular gluing gives an isometry of finite quadratic
modules
\[
                 \iota:A_M\longrightarrow A_K.
\]
Both groups have order three. Write their classes as
$0,+,-$, with $q(\pm)=2/3$ modulo $\mathbb Z$, matching
the two nonzero classes by $\iota$. Define
\[
 \begin{split}
 \delta(\tau)&=q\prod_{j\ge1}(1-q^j)^{24},
                  \qquad q=e^{2\pi i\tau},\\
 \theta_{K+\lambda}(\tau)
       &=\sum_{x\in K+\lambda}q^{(x,x)/2},\\
 g(\tau)&=\frac1{\delta(\tau)}
       \sum_{\mu\in A_M}\theta_{K+\iota(\mu)}(\tau)e_\mu.
 \end{split}
 \tag{P-F46a}
\]
Here the Weil representation convention is
$\rho_M(T)e_\mu=e^{2\pi i q(\mu)}e_\mu$.
Then $g$ is a weakly holomorphic modular form of weight
$-9$ and type $\rho_M$, and its Borcherds lift satisfies
\[
       \Psi_M(g)|_{\mathbb B}\doteq\phi,
       \qquad
       \mathcal P\bigl(\Psi_M(g)|_{\mathbb B}\bigr)
                   \doteq\Delta_{\rm cub}.
 \tag{P-F46b}
\]
The notation $\doteq$ means equality up to a nonzero
constant; the scalar normalization of a Borcherds product
is not being fixed implicitly. Thus the output has
automorphic weight $48$, coefficient degree $80$, and
stretching degree $16$, exactly as in P-F41.

Every Fourier coefficient of every component of $g$ is
a nonnegative integer. More precisely, let $p_{24}(r)$
count partitions of $r$ with each part assigned one of
twenty-four colors, and put $p_{24}(r)=0$ for $r<0$.
Then
\[
 [q^t]g_\mu
   =\sum_{r\ge0}p_{24}(r)
       \#\{x\in K+\iota(\mu):(x,x)/2=t+1-r\}.
 \tag{P-F46c}
\]
All these sums are finite. The complete terms of exponent
at most zero are
\[
        (q^{-1}+96)e_0
                   +27q^{-1/3}(e_++e_-).
 \tag{P-F46d}
\]
In particular the constant coefficient giving the
Borcherds weight is $96$, so that $96/2=48$.

The input $g$ is not an ordinary square-integrable modular
input: its unregularized weight-$-9$ Petersson norm
diverges at the cusp. This assertion concerns the displayed
input itself, not all possible theta comparisons.

There is also a precise obstruction to a direct Hilbert
coefficient dictionary. Set $\nu_0=0$ and
$\nu_\pm=2/3$, and let
\[
                  c_\mu(N)=[q^{N+\nu_\mu}]g_\mu.
\]
For every fixed coset $\mu$, every integer $a\ge1$,
and every integer $b$, the sequence $c_\mu(as+b)$
cannot equal the actual cubic stretching multiplicity
$f_5(s)=\dim T^{(5)}_s$ for all sufficiently large $s$.
In fact it eventually exceeds $f_5(s)$. The same conclusion
holds for a positive rational affine reindexing on any
infinite arithmetic progression on which the exponents
belong to the component's Fourier lattice.

\paragraph{Proof of the exact contraction and the overlattice normalization.}
The primitive embedding with complement $E_6$ is the one
used in Casalaina-Martin--Laza, Proposition~7.5, printed
p.~25, in the construction already checked in P-F41.
Set $L'=M\oplus K(-1)$. Because $L$ is unimodular and
the two orthogonal factors are primitive, its gluing
subgroup
$I=L/L'\subset A_M\oplus A_{K(-1)}$
is the graph of the isometry $\iota:A_M\to A_K$.
One can see the sign explicitly: isotropy says
$q_M(\mu)-q_K(\iota\mu)=0$. There is therefore no
dual-Weil-representation sign left over in the final
$M$-valued theta series.

More explicitly, the projections of $I$ to the two
discriminant groups are injective by primitivity.
Unimodularity gives $|I|^2=|A_M||A_K|$, while
injectivity gives $|I|\le\min(|A_M|,|A_K|)$.
These two statements force $|I|=|A_M|=|A_K|$.
The determinant-three calculation for $K$ below
therefore makes all three orders equal to three,
and both projections are isomorphisms.
In the pullback convention of Ma, equation~(2.1),
pp.~5--6, the scalar vector $1\in\mathbb C[A_L]$
pulls back to
\[
                  \sum_{\mu\in A_M}
                      e_\mu\otimes\overline e_{\iota\mu}.
 \tag{P-F46e}
\]
Here $\overline e_\lambda$ is the basis vector for
$K(-1)$, identified with the dual of $e_\lambda$ for
$K$. Contracting this sum with $\Theta_K$ gives
exactly the sum in (P-F46a), with coefficient one
on each coset and no division by three or extra
overlattice sum. This is also Ma's Example~3.13,
printed p.~18, the unimodular specialization of
his Theorem~1.1 and equation~(1.1), p.~2.

That theorem applies: $M$ has rank $22$, and its
Witt index is at most two, strictly smaller than
$\operatorname{rank}M-2=20$. Its scalar input on
$L$ is $\delta^{-1}$, of weight $-12$, with
integral principal part and constant coefficient
$24$. The positive-definite rank-six theta series
has weight three and type $\rho_K\simeq\rho_M$;
Ma's theta-contraction formula consequently gives
the weight-$-9$ form $g$. This is the required
input weight $1-20/2$ for $M$.
The theorem identifies $\Psi_M(g)$, up to a
nonzero scalar, with the $E_6$ quasi-pullback of
$\Phi_{12}$. Restricting to the same ball as in
Casalaina-Martin--Laza gives $\phi$; applying
P-F41 gives (P-F46b). This proves an actual
source-to-output identity, not merely a similarity
between two theta constructions.

\paragraph{The coefficient counts and the complete principal part.}
The product identity
\[
             \delta^{-1}=q^{-1}
                    \sum_{r\ge0}p_{24}(r)q^r
\]
follows by expanding each factor $(1-q^j)^{-24}$:
it records an arbitrary number of parts of size $j$
in each of twenty-four colors. Multiplication by
the lattice sum proves (P-F46c). For a fixed exponent
$t$ only $0\le r\le t+1$ can contribute, and
positive definiteness makes each norm shell finite.
This proves both finiteness and nonnegative integrality.

Here is an elementary count of all the lattice terms
needed for (P-F46d), including the nonzero cosets.
Take
\[
 A_2=\{(a,b,c)\in\mathbb Z^3:a+b+c=0\},
 \qquad v=(2,-1,-1)/3,
 \qquad
 K=A_2^3+\mathbb Z(v,v,v).
 \tag{P-F46f}
\]
The gluing vector has norm two and integral pairings
with $A_2^3$, so this is an even integral lattice.
Its index over $A_2^3$ is three, giving determinant
$3^3/3^2=3$. To identify it with $E_6$ rather than
assume the identification, let $\alpha_i,\beta_i$
be simple roots in its three $A_2$ factors, chosen
so that $v_i=(2\alpha_i+\beta_i)/3$, and put
$u=v_1+v_2+v_3$.
The six roots
\[
             -u,\ \alpha_1,\beta_1,\alpha_2,\beta_2,\alpha_3
\]
have the $E_6$ Cartan matrix: the central root $-u$
has branches of lengths two, two, and one.
They generate $K$, since
$\beta_3=3u-2\alpha_1-\beta_1-2\alpha_2-\beta_2-2\alpha_3$.
This verifies the stated lattice and determinant convention.

The $A_2$ lattice has six vectors of norm two.
Each nonzero coset $\pm v+A_2$ has minimum norm
$2/3$, attained by exactly three vectors.
Indeed the vectors in $v+A_2$ have coordinates
$k_i-1/3$ with $k_i\in\mathbb Z$ and $\sum k_i=1$;
their squared norm is $\sum k_i^2-1/3$.
Its minimum is $2/3$, exactly when the $k_i$ are
a permutation of $(1,0,0)$. Negation proves the
other coset case.

The zero coset of $K$ is the union of $A_2^3$ and
the two classes $(\pm v+A_2)^3$. Its norm-two
vectors therefore number
\[
                         3\cdot6+2\cdot3^3=72.
\]
There are no additional possibilities: an $A_2$
factor of norm two already accounts for the full
norm in the first class, and each of the three
factors must have minimum norm in the other two.

For the dual quotient, identify
$(A_2^\vee/A_2)^3$ with $(\mathbb Z/3)^3$.
The gluing subgroup is the diagonal line
$H=\langle(1,1,1)\rangle$, and its orthogonal
subgroup consists of $(a,b,c)$ with $a+b+c=0$.
Thus $A_K=H^\perp/H$ has three elements.
Represent a nonzero class by $(1,-1,0)$.
Adding $H$ gives the three cyclic permutations
of this pattern. Each has minimum norm $4/3$,
with $3^2=9$ vectors at that norm; hence the
coset has exactly $27$ minimum vectors.
The negative class has the same count. This proves
\[
 \theta_K=1+72q+O(q^2),\qquad
 \theta_{K+\pm}=27q^{2/3}+O(q^{5/3}).
 \tag{P-F46g}
\]
Since $p_{24}(0)=1$ and $p_{24}(1)=24$,
(P-F46g) gives (P-F46d) and the weight $48$.
The norm and coset calculation also verifies
$q(\pm)=2/3$ in the convention used above.

\paragraph{A proved failure of direct coefficient reindexing.}
The zero vector in the zero coset, and one minimum vector
in either nonzero coset, show from (P-F46c) that for
every integer $N\ge0$,
\[
                         c_\mu(N)\ge p_{24}(N+1)
                                            \ge p(N+1),
 \tag{P-F46i}
\]
where $p$ is the ordinary partition function. The second
inequality assigns a fixed color to every part of an
ordinary partition; the first uses a lattice vector
of half-norm exactly $\nu_\mu$.

For an elementary lower bound on $p(r)$, put
$k=\lfloor\sqrt r\rfloor$ when $r\ge1$.
Every subset of $\{2,\ldots,k\}$ has sum at most
$k(k+1)/2-1\le r$. Turn it into a partition of $r$
by adding parts equal to one. Distinct subsets give
distinct partitions, since their parts greater than
one recover the subset. Hence
\[
                         p(r)\ge2^{\lfloor\sqrt r\rfloor-1}.
 \tag{P-F46j}
\]
On the actual plethysm side, $\dim\operatorname{Sym}^3
\mathbb C^5=35$, so the inclusion of invariants gives
\[
       f_5(s)\le
       \dim\operatorname{Sym}^{5s}(\mathbb C^{35})
                       =\binom{5s+34}{34}.
 \tag{P-F46k}
\]
For fixed $a>0$ and $b$, (P-F46i)--(P-F46j) grow at
least exponentially in $\sqrt s$, whereas (P-F46k)
is a polynomial of degree $34$. More explicitly,
the logarithm of the lower bound divided by $\sqrt s$
has positive lower limit, while the logarithm of the
upper bound divided by $\sqrt s$ tends to zero.
This proves the strict inequality for all sufficiently
large $s$. A compatible rational affine reindexing
becomes an integral affine one upon restricting $s$
to its fixed arithmetic progression; the right side
of (P-F46k) remains polynomial in the new parameter.
This proves the claimed obstruction without using
an asymptotic formula for partitions or assuming
an exact Hilbert dimension for the invariant ring.

\paragraph{The precise cusp obstruction to unchanged ordinary input.}
Uniformly for $x$ in a bounded real interval,
\[
                g_0(x+iy)=e^{-2\pi i x}e^{2\pi y}
                                      +96+O(e^{-2\pi y})
                         \quad(y\longrightarrow\infty).
\]
The expansion is uniform because the remaining power
series converges absolutely on every sufficiently
small closed $q$-disc. Thus for all large $y$,
$|g_0(x+iy)|\ge\tfrac12e^{2\pi y}$.
With the standard invariant norm on the Weil
representation, $\|g\|^2\ge|g_0|^2$.
Consequently the unregularized Petersson integral
over the cusp satisfies
\[
 \int_Y^\infty\int_{-1/2}^{1/2}
            \|g(x+iy)\|^2 y^{-9}\,\frac{dx\,dy}{y^2}
       \ \ge\ \frac14\int_Y^\infty e^{4\pi y}y^{-11}\,dy
       =\infty.
 \tag{P-F46h}
\]
In particular $g$ has exponential, not polynomial,
cusp growth. It cannot simply be substituted for an
ordinary square-integrable automorphic input while
retaining that input hypothesis. Borcherds's
regularization handles precisely a weakly holomorphic
input of this kind; the statement does not assert
that every theta lift of it diverges or that no
regularized comparison is possible.
\hfill$\square$

\paragraph{The same form is a unitary Borcherds lift.}
There is also an exact unitary interpretation, not merely
an analogy with an orthogonal construction. Equip $M$ with
a generator $\sigma$ of its cubic-cover deck group,
chosen to match the complex embedding of the period ball,
so that
$\sigma^2+\sigma+1=0$ and $\sigma$ preserves its integral
bilinear form $B$. Put $F=\mathbb Q(\zeta)$,
$\mathcal E=\mathbb Z[\zeta]$, where
$\zeta=e^{2\pi i/3}$ acts by $\sigma$. On this
$\mathcal E$-lattice define the form, linear in the first
variable,
\[
 h(x,y)=\frac{1-\zeta^2}{3}B(x,y)
           +\frac{\zeta-\zeta^2}{3}B(x,\sigma y).
 \tag{P-F46l}
\]
The identities $B(\sigma x,y)=-B(x,y)-B(x,\sigma y)$
and $B(\sigma x,\sigma y)=B(x,y)$ give
$h(\sigma x,y)=\zeta h(x,y)$ and
$h(y,x)=\overline{h(x,y)}$. Taking the field trace gives
$\operatorname{Tr}_{F/\mathbb Q}h(x,y)=B(x,y)$.
Also $B(x,\sigma x)=-B(x,x)/2$, so
$h(x,x)=B(x,x)/2\in\mathbb Z$.
Thus the Hermitian signature is $(1,10)$, and the
underlying quadratic lattice is exactly $M$, with no
rescaling of its quadratic form.

Write $\Lambda$ for this Hermitian lattice. For
$x,y\in\Lambda$ and $a\in\mathcal E$, the trace
of $\overline a h(x,y)$ is $B(x,ay)\in\mathbb Z$.
By trace duality this says
$h(x,y)\in\mathfrak D_F^{-1}$. It proves integrality
in Hofmann's convention, and the diagonal calculation
proves evenness. The same trace-duality argument proves
that its Hermitian dual $\Lambda'$ is exactly $M^\vee$:
if $v\in M^\vee$, apply integrality to $ax$ for every
$a\in\mathcal E$ and $x\in M$; the converse follows
by taking the trace. Consequently the finite quadratic
module is precisely $A_M$, with
$q(v)=h(v,v)=B(v,v)/2$, and Hofmann's Weil
representation in equation~(20), p.~19, is precisely
$\rho_M$, not its dual.

The isotropic hypotheses of Hofmann's Theorem~8.1
also hold. The displayed lattice in the proof of
Casalaina-Martin--Laza, Proposition~7.5, p.~25,
contains two hyperbolic planes; after sign reversal
this supplies a nonzero isotropic vector of $M$.
Since $\mathcal E$ is a principal ideal domain,
saturating its $F$-line in $\Lambda$ gives a primitive
isotropic vector $\ell\in\Lambda$.
Choose $y\in M\otimes\mathbb Q$ with
$a=h(\ell,y)\ne0$, and put
\[
              y'=y-\frac{h(y,y)}{2a}\ell.
\]
Direct expansion gives $h(y',y')=0$ and
$h(\ell,y')=a$. Clearing denominators gives an
isotropic $\ell'\in\Lambda\subseteq\Lambda'$ with
$h(\ell,\ell')\ne0$. These are exactly Assumption~2.1,
p.~6, and the hypotheses in Theorem~8.1, p.~24;
no unimodularity hypothesis from the special introductory
version of that theorem is being imposed.

Now $g$ has weight $1-10=-9$, the required Weil type,
and integral negative Fourier coefficients. Hofmann's
Theorem~8.1 therefore produces a unitary Borcherds
product $\Xi_g$, of weight $96/2=48$, initially for
the unitary discriminant kernel
$\Gamma^U_\Lambda\subset U(\Lambda)$, with a
finite-order multiplier. His definition in
equation~(28), p.~25, is
\[
        \Xi_g=\alpha^*\Psi_M(g)\doteq\phi,
        \qquad \mathcal P(\Xi_g)\doteq\Delta_{\rm cub}.
 \tag{P-F46m}
\]
Here $\alpha$ is the holomorphic embedding obtained
by viewing a positive Hermitian line as an oriented
positive quadratic two-plane, as in his Section~4,
pp.~13--15. This is the actual cubic-cover ball
embedding already used in P-F41 and (P-F46b).
Thus the source supplies the exact unitary lift name
for the same automorphic form. Its equality with
$\phi$, together with the group and multiplier check
in P-F41, extends this particular form to the exact
cubic period group with its stated trivial multiplier;
we do not identify that group with the discriminant
kernel. Theorem~8.1(4) further gives, after compatible
scalar normalization,
$\Xi_{r g}=\Xi_g^r$ for every positive integer $r$.
These are exactly the discriminant powers of output
degrees $(48r,80r,16r)$, rather than an assertion about
all graded pieces of the invariant ring.

\paragraph{What passes, and the remaining supplied-method comparison.}
The positive input spaces in (P-F46c) have a concrete
basis indexed by a lattice vector and a twenty-four-colored
partition. The established map sends their generating
series, through the regularized multiplicative lift,
to one algebra element, the cubic discriminant. The
coefficient $96$ determines its automorphic weight;
it is not identified with the dimension of a graded
plethysm multiplicity space. The powers of that one
element give the polynomial subalgebra from P-F41,
with degrees $(48r,80r,16r)$ for its $r$th power.

This sharpens the comparison with the supplied
\emph{Remark}, printed p.~4, item~8: the known
construction already contains a specific positive-definite
theta factor, but also the explicit inverse modular
discriminant, the coset contraction, and the regularized
multiplicative unitary Borcherds lift (P-F46m). Items~9--10 and~12 of that note
concern the norms, divisibility, and periods of its
iterated unitary lifts. No map identifying those
operations with (P-F46a)--(P-F46b) is provided by the
present argument. Such a comparison must account for
the displayed principal part and the divergent ordinary
input norm, as well as the multiplication and three
output gradings. This is a concrete remaining operation,
not an exclusion based on a group name.

In particular, the growth argument proves that the
positive Fourier coefficient spaces in (P-F46c) cannot
themselves be all the graded pieces of the invariant
ring under a positive affine reindexing. It does not
exclude subquotients, a different choice of spaces,
or a nonlinear transform. No basis
of those pieces, Hilbert-series identity with this
input, uniform positivity algorithm, or PH2/SH bound
is obtained. The theorem is an exact specialization of
a known quasi-pullback formula, with explicit coefficient
and normalization checks; no historical novelty is claimed.

The additional primary source is Shouhei Ma,
\emph{Quasi-pullback of Borcherds products},
Theorem~1.1 and equation~(1.1), p.~2; Weil conventions
and pullback, Section~2.1 and equation~(2.1), pp.~4--6;
theta contraction, equation~(3.1), p.~12; Example~3.13,
p.~18:
\url{https://arxiv.org/pdf/1801.08333}.

The unitary interpretation uses Eric Hofmann,
\emph{Borcherds Products on Unitary Groups},
Hermitian lattice conventions and Assumption~2.1, p.~6;
discriminant kernel, p.~7; underlying quadratic lattice
and holomorphic embedding, Section~4, pp.~13--15;
Weil representation, equation~(20), p.~19;
Theorem~8.1, p.~24; and its defining pullback,
equation~(28), p.~25:
\url{https://arxiv.org/pdf/1210.2542}.


\subsection{What the natural exterior-volume powers actually detect}
\label{sec:exterior-volume-fermat}

\paragraph{P-F47 (an explicit coefficient map to the Fermat degree monoid).}
Let $q,n\ge1$, put $k=3q$, $d=nq$, and let
$E=\mathbb C^k$ with ordered basis $e_1,\ldots,e_k$.
Write $W=\bigwedge^3E$ and $e_J=e_i\wedge e_j\wedge e_l$
for a triple $J=\{i<j<l\}$. Define
\[
 \theta_k=\sum_{\mathcal M}\epsilon(\mathcal M)
                    e_{J_1}\wedge\cdots\wedge e_{J_q}
       \ \in\ \bigwedge^qW,                         \tag{P-F47a}
\]
where the sum is over the partitions $\mathcal M$ of $[k]$
into triples, their blocks $J_1,\ldots,J_q$ are listed
lexicographically, and $\epsilon(\mathcal M)$ is the sign
of their concatenated increasing lists as a permutation
of $[k]$. Then $\theta_k$ spans a copy of $\det E$, and
\[
 \theta_k^{\wedge n}\ne0
 \quad\Longleftrightarrow\quad
 \exists I\in\mathbb C[\operatorname{Sym}^3\mathbb C^n]^{SL_n}_d
             \text{ with }I(x_1^3+\cdots+x_n^3)\ne0.
                                                        \tag{P-F47b}
\]
More precisely, every coefficient of the exterior power
is, up to an explicitly fixed sign, the evaluation of a
determinant-contraction invariant on the Fermat cubic.
Conversely these coefficients span all evaluations of
degree-$d$ invariants on that cubic.

This is a statement about nonvanishing at one specified
form. It is not an equivalence between this particular
exterior power and nonvanishing of the whole degree-$d$
invariant space. In particular, failure of (P-F47b) does
not imply $f_n(gq)=0$, where $g=\gcd(n,3)$.

\paragraph{The canonical volume inclusion.}
The map $\det E\longrightarrow\bigwedge^q(\bigwedge^3E)$
giving (P-F47a) is explicitly
\[
 v_1\wedge\cdots\wedge v_k\longmapsto
 \frac{1}{q!(3!)^q}\sum_{\sigma\in\mathfrak S_k}
 \operatorname{sgn}(\sigma)
 \bigwedge_{a=0}^{q-1}
 (v_{\sigma(3a+1)}\wedge v_{\sigma(3a+2)}
                          \wedge v_{\sigma(3a+3)}).
                                                        \tag{P-F47c}
\]
It is alternating and multilinear in the $v_i$, so is
well-defined and $GL(E)$-equivariant. Permuting two
blocks contributes a sign $(-1)^9$ in their flattened
lists and a sign $-1$ in the outer wedge; these cancel.
Permuting entries of a triple similarly contributes
two equal signs. Thus each term of (P-F47a) occurs
exactly $q!(3!)^q$ times before division. Its coefficients
are $\pm1$, so the map is nonzero. This proves the
claimed determinant character without choosing a
metric or an unproved splitting of a Schur functor.

\paragraph{Determinant contractions and their values.}
Consider a tableau $T$ with $n$ rows and $k$ columns,
whose entries are labels $1,\ldots,d$, each occurring
three times, with no repeated label in one column.
For each label $a$, let $J_a$ be its three columns.
For a symmetric cubic tensor $F=(F_{ijk})$, define
\[
 I_T(F)=\sum_{\sigma_1,\ldots,\sigma_k\in\mathfrak S_n}
       \left(\prod_{j=1}^k\operatorname{sgn}(\sigma_j)\right)
       \prod_{a=1}^d
       F_{\sigma_{j_1}(r_1),\sigma_{j_2}(r_2),
                                      \sigma_{j_3}(r_3)},
                                                        \tag{P-F47d}
\]
where $(r_b,j_b)$ are the three positions of label $a$.
The order of those positions does not matter because
$F$ is symmetric. Contracting the $n$ indices in each
column with the alternating determinant tensor proves
\[
                    I_T(hF)=(\det h)^k I_T(F).
                                                        \tag{P-F47e}
\]
Thus it is a degree-$d$ polynomial $SL_n$-invariant.
An exact primary construction and transformation proof
are B\"urgisser--Ikenmeyer, \emph{Fundamental invariants
of orbit closures}, equation~(3.7) and Theorem~3.11,
printed pp.~11--12. Their Remark~3.13, p.~13, states
that the semistandard tableaux of this shape and content
span the entire invariant space. Equivalently, one can
polarize a degree-$d$ invariant, use the determinant
contractions spanning the $SL_n$-invariants in the
$3d=nk$ tensor positions, and then symmetrize within
each cubic and among the $d$ cubic factors. A repeated
label in a column gives zero after this symmetrization.

In the tensor convention of (P-F47d), the Fermat cubic
$F_n=\sum_i x_i^3$ has $F_{iii}=1$ and all its other
tensor entries zero. Its nonzero summands are exactly
the colorings $c:\{1,\ldots,d\}\to[n]$ such that in
each column of $T$ every color occurs once. A coloring
determines its $k$ permutations uniquely. Consequently
\[
 I_T(F_n)=\sum_c\prod_{j=1}^k
                  \operatorname{sgn}(c\text{ in column }j).
                                                        \tag{P-F47f}
\]
If $J_a=J_b$ for two distinct labels, exchange their
colors. The colors are different in each of their
three common columns. This is a fixed-point-free
involution on the admissible colorings, and it changes
the product of column signs by $(-1)^3=-1$.
Therefore $I_T(F_n)=0$ whenever two triples coincide.

\paragraph{The exact coefficient sign.}
It remains to consider a set $H=\{J_1,\ldots,J_d\}$
of distinct triples, listed in a fixed order, with
each vertex in exactly $n$ triples. Choose the tableau
by placing the incident labels in any fixed order in
each column. A coloring in (P-F47f) is precisely an
ordered factorization of $H$ into $n$ perfect matchings
of $[k]$. Expansion of $\theta_k^{\wedge n}$ shows that
its coefficient at $e_{J_1}\wedge\cdots\wedge e_{J_d}$
is the signed sum of these same ordered factorizations.
There is no omitted factorial: each matching occurs
once in (P-F47a), and the $n$ factors are ordered.

For an explicit comparison of the signs, distinguish
the $3d=nk$ incidences $(a,j)$ with $j\in J_a$.
Let $A$ be their order by label $a$, then by increasing
vertex $j$ within that triple. Let $B$ be their order
by vertex $j$, then by the chosen tableau row. Write
$\varepsilon_T$ for the sign of this fixed reordering
$A\to B$. Write $\tau_{n,k}$ for the sign of the fixed
transpose from the grid order $(j,i)$ to $(i,j)$.
For a coloring $c$, the sign of reordering $A$ to
the order by color and then vertex can be computed
in two ways. Grouping triples by color has the same
sign as grouping the outer wedge factors, because
interchanging two triples interchanges nine incidences.
Ordering the vertices inside each color gives the
matching signs in (P-F47a). This is exactly its exterior
coefficient contribution. Alternatively, first use
$A\to B$, then sort each column by color, and finally
transpose the grid. This gives
\[
 [e_{J_1}\wedge\cdots\wedge e_{J_d}]
                \theta_k^{\wedge n}
       =\varepsilon_T\tau_{n,k}\,I_T(F_n).
                                                        \tag{P-F47g}
\]
Both fixed signs are defined by specified permutations,
so (P-F47g) is an actual coefficient identity.
Sets of triples that are not $n$-regular have coefficient
zero, since a factorization into perfect matchings
would give degree $n$ at every vertex.

If the exterior power is nonzero, a nonzero coefficient
and (P-F47g) give the required invariant. Conversely,
the spanning statement for $I_T$ shows that a nonzero
invariant evaluation has a nonzero tableau evaluation.
The repeated-triple involution eliminates nonsimple
tableaux, and (P-F47g) then gives a nonzero exterior
coefficient. This proves (P-F47b) completely.

\paragraph{Exact nonvanishing ranges that follow immediately.}
For every $q$, $\theta_{3q}\ne0$. If $q$ is odd and
$n\ge2$, graded anticommutativity gives
$\theta_{3q}\wedge\theta_{3q}=0$, so its $n$th power
vanishes. This is consistent directly with Fermat:
a transposition of two variables fixes $F_n$ and has
determinant $-1$, whereas $k=3q$ is odd in (P-F47e).

For $q=2$, $k=6$, the twenty triples pair into ten
complementary pairs. Hence (P-F47a) has the form
\[
          \theta_6=\sum_{a=1}^{10}\epsilon_a u_a\wedge v_a,
             \qquad \epsilon_a\in\{1,-1\},              \tag{P-F47h}
\]
where the $u_a,v_a$ are a basis of $\bigwedge^3\mathbb C^6$.
Its $n$th power is nonzero exactly for $0\le n\le10$:
for $n\le10$ a product of any $n$ distinct pairs has
coefficient $n!\prod\epsilon_a\ne0$, and for $n>10$
the exterior degree exceeds twenty. Thus degree $2n$
invariants nonzero at Fermat exist exactly for $n\le10$.
This recovers the cubic case of B\"urgisser--Ikenmeyer,
Proposition~3.24, pp.~15--16; it is not a priority claim.

More generally the dimension bound gives
\[
      \theta_{3q}^{\wedge n}=0
           \quad\text{if}\quad
                  n>\binom{3q-1}{2}.                   \tag{P-F47i}
\]
For $n\le10$ and every even $q$, the nonzero degree-$2n$
invariant can be raised to the $q/2$ power; its Fermat
value stays nonzero. Equation (P-F47b) therefore proves
nonvanishing of $\theta_{3q}^{\wedge n}$ in this range.
These elementary arguments alone do not give general
nonvanishing for $n>10$ and even $q$.
P-F49 below supplies a polynomially bounded nonzero
power for every $n$, and P-F50 proves nonvanishing
for all sufficiently large even $q$, with a polynomial
threshold. The exact low-degree range remains distinct.

\paragraph{Remaining signed-factorization and index gates.}
Equation (P-F47g) asks for a simple regular triple
hypergraph whose perfect-matching factorizations have
nonzero signed sum. An unsigned existence theorem for
such factorizations cannot replace that signed assertion.
The solutions in P-F49--P-F50 give specified actual
nonzero coefficients and a polynomial conductor for
the Fermat degree monoid. They do not
by themselves control every residue-constituent polynomial
or the strict positivity and saturation indices in
GCT VI. Those additional implications require their
own proof. The primary locator for the determinant
constructions and their spanning property is
\url{https://arxiv.org/pdf/1511.02927}, equation~(3.7),
Theorem~3.11, and Remark~3.13, pp.~11--13.


\subsection{The density clause and the size of the cubic families}
\label{sec:cubicdensityscope}

\begin{proposition}[P-F48: the cubic constructions are exponentially sparse in the stated input population]
\label{prop:cubicdensityscope}
For a positive integer $a$, put
$\ell(a)=1+\lfloor\log_2a\rfloor$. For a partition, sum this
quantity over its positive parts; zero parts are omitted.
Write
\[
 L(\lambda,\mu,\pi)=\ell(\lambda)+\ell(\mu)+\ell(\pi).
\]
This is exactly the input-size convention in GCT VI \cite{PGCTVI},
Section~1.1, printed p.~9. The narrow SH(b), Hypothesis~1.6.5(b),
printed p.~25, concerns triples of total length \emph{exactly}
$N$ with $a_{\lambda,\mu}^{\pi}\ne0$. The unified SH(b),
Hypothesis~3.3.2(b), printed p.~51, instead says length
\emph{at most} $N$ and uses the specified input $x$.
For its canonical ordinary-plethysm specification,
Section~3.4, printed p.~53, that length is
$n+L(\lambda,\mu,\pi)$: the group $GL_n$ contributes $n$.
Thus these are different input populations, not just
different cutoffs. Neither clause imposes a fixed inner
shape, fixed outer shape, fixed target shape, or a fixed rank.

Let $\mathcal D_N$ be the set of partition triples of length
exactly $N$ with nonzero ordinary plethysm coefficient, and
put $\mathcal D_{\le N}=\bigcup_{j\le N}\mathcal D_j$.
The counting fractions for narrow triples, using either cutoff,
and an exceptional set
$\mathcal E$ are
\[
 \frac{|\mathcal E\cap\mathcal D_N|}{|\mathcal D_N|},
 \qquad
 \frac{|\mathcal E\cap\mathcal D_{\le N}|}
      {|\mathcal D_{\le N}|},
 \tag{P-F48a}
\]
The first is the narrow clause's natural
conditional reading; the second is its at-most-$N$ variant.
The source states the density in words rather than displaying
a fraction. If instead that wording is read
as density among all partition triples, the estimates below
continue to hold, since that denominator is larger. No
distribution on partitions of a fixed weight or on a chosen
one-parameter family is substituted for the source's population.

For $n\ge4$, set
\[
 g=\gcd(n,3),\quad m=\binom{n+2}{3},\quad
 d=\frac{rn}{g},\quad k=\frac{3r}{g},
 \qquad r\ge1.
\]
The rescaled cubic triples
\[
                  ((d),(3),(k^n))
 \tag{P-F48b}
\]
form a set $\mathcal C$. Their stretching functions are
$f_n(rs)$, in the notation of P-F31--P-F32.
More generally, for arbitrary integers $q,b\ge0$ put
\[
 \begin{split}
  \mu_{n,q}&=(q+3,q^{\,n-1}),\\
  \lambda_{n,r,b}&=(d+b,b^{\,m-1}),\\
  e_{n,q}&=\frac{m(3+nq)}n
          =\frac{(n+1)(n+2)}2+qm,\\
  \pi_{n,r,q,b}&=((k+qd+be_{n,q})^n),
 \end{split}
 \tag{P-F48c}
\]
again omitting all zero parts. These natural inner and outer
determinant twists form a set $\mathcal C^{\det}$ containing
$\mathcal C$. Every one of them has the same stretching
function $f_n(rs)$ as its untwisted counterpart.

For every integer $N\ge12$,
\[
 \begin{split}
 |\mathcal D_N|&\ge 2^{N/2-4},\\
 |\{x\in\mathcal C:L(x)\le N\}|
       &\le N\,2^{N/5},\\
 |\{x\in\mathcal C^{\det}:L(x)\le N\}|
       &\le N(N+1)^3\,2^{N/4}.
 \end{split}
 \tag{P-F48d}
\]
Consequently the conditional density of \emph{any} subset
of $\mathcal C$ is at most
$16N\,2^{-3N/10}$, and that of any subset of
$\mathcal C^{\det}$ is at most
$16N(N+1)^3\,2^{-N/4}$, for both populations in
(P-F48a). In particular these densities are smaller than
$N^{-c}$ for every fixed $c>0$ and all sufficiently large
$N$.

The same conclusion holds separately in the canonical
ordinary-plethysm input population of the unified clause.
Let $\widehat{\mathcal D}_{\le N}$ count admissible
$(h,\lambda,\mu,\pi)$ with nonzero coefficient and
$h+L(\lambda,\mu,\pi)\le N$. Lift the constructed
families using all their admissible ambient ranks
$h\ge n$, including padding above the displayed rank $n$.
For $N\ge13$ their conditional densities in this population
are at most
\[
 32N^2\,2^{-3N/10},\qquad
 32N^2(N+1)^3\,2^{-N/4},
 \tag{P-F48h}
\]
respectively. This statement uses the canonical
ordinary-plethysm specification; it is not a claim
about counting arbitrarily many syntactic descriptions
of a composite subgroup embedding.

The unscaled family $r=1$, including all its twists in
(P-F48c), has zero initial multiplicity and is excluded
from $\mathcal D_N$ altogether. These facts do not prove
SH(b), do not assert that the rescaled functions have zero
strict saturation index, and do not preclude a different
family of positive density.
\end{proposition}

\begin{proof}
\emph{1. A nonzero population of exponential size in every
large length shell.}
For odd $N=2u+1$, choose any integer
$a\in[2^{u-1},2^u-1]$ and take
\[
                         (\lambda,\mu,\pi)=((1),(a),(a)).
\]
Its coefficient is one, since the outer Schur functor is
the identity. Its length is exactly $1+u+u=N$.
There are $2^{u-1}$ such triples.

For even $N=2u+4$, choose
\[
                    \left\lceil\frac{2^{u+1}}3\right\rceil
                              \le a\le 2^u-1
\]
and take $((3),(a),(3a))$. Here $\ell(a)=u$ and
$\ell(3a)=u+2$, so the total length is $2+u+(u+2)=N$.
The coefficient is one: the highest weight $(3a)$ occurs
once in $\operatorname{Sym}^3(\operatorname{Sym}^a V)$.
Indeed its weight space is spanned by the cube of the
highest weight vector of $\operatorname{Sym}^a V$.
This also gives a direct verification in any sufficient
ambient rank, or just in rank one.
For $u\ge4$, the number of choices is at least
$2^u/3-1\ge2^{u-2}$.
Both parity cases therefore give the first inequality in
(P-F48d) for $N\ge12$. In particular the denominator is
nonzero in every such shell.

\emph{2. The determinant twists preserve the actual ray.}
Let $V=\mathbb C^n$, $W=\operatorname{Sym}^3 V$, and
$W_q=W\otimes(\det V)^q$. Then
\[
              W_q=V_{\mu_{n,q}}(GL_n),\qquad \dim W_q=m.
\]
The determinant character of $W$ is
$(\det V)^{3m/n}$. To see this, each diagonal coordinate
occurs with the same total exponent in its monomial
basis, and the sum of those $n$ exponents is $3m$.
The integer $3m/n=(n+1)(n+2)/2$ proves both the formula
and the asserted integrality of $e_{n,q}$.
Consequently
\[
                    \det W_q=(\det V)^{e_{n,q}}.
\]
For every integer $s\ge0$, the usual determinant twist
identity for a highest weight gives
\[
 \begin{split}
 V_{s\lambda_{n,r,b}}(GL(W_q))
   &\simeq \operatorname{Sym}^{sd}(W_q)
                          \otimes(\det W_q)^{sb}\\
   &\simeq \operatorname{Sym}^{sd}(W)
                    \otimes(\det V)^{s(qd+be_{n,q})}.
 \end{split}
 \tag{P-F48e}
\]
The target $V_{s\pi_{n,r,q,b}}(GL_n)$ is the character
$(\det V)^{s(k+qd+be_{n,q})}$. Cancelling the common
determinant twist reduces its multiplicity to that of
$(\det V)^{sk}$ in $\operatorname{Sym}^{sd}(W)$.
The central degree condition is $3sd=nsk$, so that
this is precisely the $SL_n$-invariant multiplicity
$f_n(rs)$. This proves the ray identity for all $s$,
including the boundary cases $q=0$ and $b=0$.

\emph{3. The bit budgets are literal partition budgets.}
For (P-F48b), $d\ge r$ and $k\ge r$, whence
\[
             L((d),(3),(k^n))
                      \ge 2+(n+1)\ell(r).
 \tag{P-F48f}
\]
The target has $n$ positive parts, so $n\le N$ whenever
the total length is at most $N$. For each such $n$,
the number of possible positive $r$ is at most
$2^{N/(n+1)}\le2^{N/5}$.
Summing this bound over at most $N$ values of $n$
proves the second inequality in (P-F48d).

For the larger family define $\ell(0)=0$ only for the
two parameters $q,b$. This convention counts their
omitted zero parts; it does not add zeros to the source's
partition encoding. Put
$u=\ell(r)$, $v=\ell(q)$ and $w=\ell(b)$.
The first part of $\lambda_{n,r,b}$ is at least $r$;
its other $m-1$ positive parts, when $b>0$, equal $b$.
The $n$ parts of $\mu_{n,q}$ are each at least $q$
when $q>0$. The target has $n$ parts each at least $r$.
These are disjoint entries in the three encoded
partitions, and therefore
\[
              L(\lambda_{n,r,b},\mu_{n,q},\pi_{n,r,q,b})
                   \ge (n+1)u+nv+(m-1)w
                   \ge4(u+v+w).
 \tag{P-F48g}
\]
When $q=0$ the omitted contribution on the right is
zero, while the actual inner partition still has its
positive part $3$; when $b=0$ only the first outer
part remains. Thus (P-F48g) is valid also in those
cases without double counting any entry.

There are at most $(N+1)^3$ possible bit-length triples
$(u,v,w)$. For each one, the number of parameter triples
$(r,q,b)$ is at most $2^{u+v+w}\le2^{N/4}$, using the
single zero option when $v=0$ or $w=0$.
Again $n\le N$ because the target has $n$ positive
parts. Summing gives the third inequality in
(P-F48d). Distinct parameter choices that happen to
give the same input only reduce the count.

\emph{4. The density and nonzero-input conclusions.}
For the exact-length density, bound its numerator by
the corresponding at-most-$N$ count in (P-F48d) and
its denominator by the first line there.
For the at-most-$N$ narrow density use the same numerator
bound and $|\mathcal D_{\le N}|\ge|\mathcal D_N|$.
This gives exactly the two displayed exponential
estimates, even if every member of the constructed
family were an SH(b) exception. Exponential decay
times a fixed polynomial is smaller than every
inverse power of $N$.

For the canonical unified ordinary-plethysm inputs,
each underlying triple has at most $N$ admissible
ambient ranks of cost at most $N$.
The stable plethysm coefficient, including its stretched
versions, is unchanged by such padding: specialization
of the symmetric-function character to $h$ variables
preserves the coefficients of all Schur functions of
height at most $h$, including the target.
Multiplying the previous numerator bounds by $N$
therefore covers every padded realization of the family.
The one-row examples in the first part of the proof
are all admissible already at rank one.
Use those of narrow length $N-1$, then add the rank
cost one. They show
\[
        |\widehat{\mathcal D}_{\le N}|
                  \ge 2^{(N-1)/2-4}
                  \ge 2^{N/2-5}
                         \qquad(N\ge13).
\]
Combining this with the enlarged numerator bounds
proves (P-F48h). Any larger population with these
canonical inputs and the same encoding costs has a
larger denominator; no inference about a different
encoding's numerator is made.

Finally P-F32 applied at $s=1$ gives
$f_n(1)=0$ for every $n\ge4$: its exterior parameter
is $k=3/\gcd(n,3)\le3$, for which
$(k-1)(k-2)/2\le1<n$.
Equation (P-F48e) preserves this zero under all the
determinant twists.
If instead $f_n(r)>0$, powers of a nonzero invariant
show $f_n(rs)>0$ for every positive integer $s$.
This last statement concerns realized values only.
Definition~1.2.4 of GCT VI tests every nonzero
constituent polynomial at every positive integer,
including arguments outside that constituent's
residue class. Realized nonvanishing alone thus
does not establish the strict saturation assertion,
and none was used in the counting proof.
\end{proof}

\paragraph{Scope of the remaining density problem.}
The two source versions use different length cutoffs
and, in their canonical ordinary-plethysm specifications,
different rank costs. The proof explicitly treats
each, so their equivalence is not assumed.
Neither changing the conditioning to
rectangular targets nor counting only the parameter
$r$ would test the stated clause. Nor does padding an
ambient rank create new narrow inputs: the source
counts the three partitions and omits their zero parts.
It does create distinct canonical unified inputs,
and the extra factor $N$ in (P-F48h) explicitly covers
them.
The proved conclusion closes this particular
amplification attempt, including its natural determinant
twists. Establishing or refuting SH(b) for the entire
source population remains open here.


\subsection{An actual cubic invariant constructed using a prime in a residue class}
\label{sec:hessian-fermat-prime}

\paragraph{P-F49 (Hessian--Cayley nonvanishing, with exact prime bounds).}
Let $n\ge2$, let $p\equiv1\pmod n$ be prime, and put
$D=p-1$, $h=D/n$. For a cubic $F$ in $n$ variables set
\[
       H_F(x)=6^{-n}\det\left(
                   \frac{\partial^2F}{\partial x_i\partial x_j}
                                      \right)_{1\le i,j\le n}.
                                                        \tag{P-F49a}
\]
There is an explicit polynomial $I_{n,p}$ in the cubic
coefficients, defined over $\mathbb Q$, such that
\[
 \deg I_{n,p}=n(p-1),\qquad
 I_{n,p}(gF)=(\det g)^{3(p-1)}I_{n,p}(F),\qquad
 I_{n,p}(F_n)\ne0,\quad F_n=\sum_{i=1}^n x_i^3.
                                                        \tag{P-F49b}
\]
It is obtained by applying Cayley's degree-$n$ determinant
contraction on forms of degree $D$ to $H_F^h$.
In the actual cubic plethysm notation, $g_n=\gcd(n,3)$,
this proves
\[
              f_n(g_n(p-1))>0,
           \qquad \theta_{3(p-1)}^{\wedge n}\ne0.
                                                        \tag{P-F49c}
\]
Unconditionally one can choose $p\le Cn^{5.18}$ with
an effective absolute constant $C$. For $n\ge4$, the
Dirichlet-GRH least-prime theorem gives the sharper choice
\[
      p\le(\varphi(n)\log n)^2,
      \quad g_n(p-1)<g_n(\varphi(n)\log n)^2,
      \quad\deg I_{n,p}<n^3(\log n)^2.                 \tag{P-F49d}
\]
The supplied Dirichlet statements A-DIR-NV/A-DIR-Z
discharge exactly this GRH premise in the present work.
These are actual point-specific invariant occurrences;
they do not assert a polynomial-size evaluation circuit,
an invariant basis, or a bound on the strict PH2/SH indices.

\paragraph{Cayley's contraction and its normalization.}
For a degree-$D$ form
$G(x)=\sum_{|\alpha|=D}g_\alpha x^\alpha$, let its
symmetric tensor entries be
\[
       G_{i_1\cdots i_D}=g_\alpha\frac{\alpha!}{D!},
       \qquad \alpha!=\prod_{a=1}^n\alpha_a!,
                                                        \tag{P-F49e}
\]
when the index list has multiplicity vector $\alpha$.
Define
\[
 \mathcal C_{D,n}(G)=
 \sum_{\sigma_1,\ldots,\sigma_D\in\mathfrak S_n}
 \left(\prod_{b=1}^D\operatorname{sgn}(\sigma_b)\right)
 \prod_{a=1}^n G_{\sigma_1(a),\ldots,\sigma_D(a)}.
                                                        \tag{P-F49f}
\]
There are $n$ tensor factors in each summand, so this
is a degree-$n$ polynomial in the coefficients of $G$.
Contracting each column of $n$ indices with a determinant
gives
\[
             \mathcal C_{D,n}(gG)
                          =(\det g)^D\mathcal C_{D,n}(G).
                                                        \tag{P-F49g}
\]
This is precisely the invariant in B\"urgisser--Ikenmeyer,
equation~(3.8), printed p.~11, with its transformation
law in Theorem~3.18(1), p.~14. We retain the formula to
make the normalization used below explicit.

\paragraph{The determinant coefficient in prime characteristic.}
For an $n\times n$ matrix $X=(X_{ij})$ of indeterminates,
write $c_L$ for the coefficient of $X^L$ in
$(\det X)^{p-1}$, where every row and column of $L$ sums
to $p-1$. Then
\[
       c_L\equiv\frac{((p-1)!)^n}{\prod_{i,j}L_{ij}!}
                                 \not\equiv0\pmod p.  \tag{P-F49h}
\]
This is Glynn's determinant coefficient identity, stated
with the congruence in Itoh--Shimoyoshi,
Theorem~1.1 and the following Remark, printed p.~1,
\url{https://arxiv.org/pdf/2104.00838}.
Here is the short form of its proof needed for the construction.

Work over the algebraic closure of $\mathbb F_p$ in
\[
        A=\overline{\mathbb F}_p[z_1,\ldots,z_n]
                                      /(z_1^p,\ldots,z_n^p).
\]
Every linear substitution preserves this ideal by
the characteristic-$p$ power identity. The top degree
$n(p-1)$ is a line spanned by $\prod_i z_i^{p-1}$.
Diagonal substitutions act on it by $(\det X)^{p-1}$.
An elementary substitution $z_i\mapsto z_i+a z_j$
acts trivially: every extra term has exponent at least
$p$ in $z_j$ and vanishes. Diagonal and elementary
matrices generate $GL_n$, so the top-line character is
$(\det X)^{p-1}$ for every invertible matrix. Since the
field is infinite this is an identity of polynomials
in the matrix entries. Therefore
\[
 [z_1^{p-1}\cdots z_n^{p-1}]
     \prod_{j=1}^n\left(\sum_i X_{ij}z_i\right)^{p-1}
                          =(\det X)^{p-1}.             \tag{P-F49i}
\]
Expanding the left side proves (P-F49h). Every entry
of $L$ is between zero and $p-1$, so all the factorials
in its denominator are units modulo $p$.

\paragraph{Nonzero evaluation and the exact cubic character.}
Put $G_0=(x_1\cdots x_n)^h$. Its tensor entries vanish
unless every symbol occurs $h$ times, and in that case
equal $(h!)^n/D!$. Consequently (P-F49f) gives the exact
identity over $\mathbb Q$
\[
 \mathcal C_{D,n}(G_0)
       =\left(\frac{(h!)^n}{D!}\right)^n
              [\prod_{i,j}X_{ij}^{h}](\det X)^D.
                                                        \tag{P-F49j}
\]
The coefficient on the right is nonzero modulo $p$
by (P-F49h). Since $h,D<p$, all displayed factorials
are units there; in fact (P-F49j) is congruent to one
modulo $p$. In particular it is nonzero in characteristic
zero. This argument uses the reduction of the specified
monomial evaluation; it does not require reducing the
normalizing factor $6^{-n}$ when $p=3$.

Define $I_{n,p}(F)=\mathcal C_{D,n}(H_F^h)$.
The Hessian determinant has degree $n$ in the cubic
coefficients and degree $n$ in $x$, so $H_F^h$ has
coefficient degree $nh=D$ and variable degree $D$.
The coefficient degree of $I_{n,p}$ is therefore $nD$.
For the convention $(gF)(x)=F(g^t x)$, the chain rule
gives
\[
            H_{gF}(x)=(\det g)^2H_F(g^t x).
                                                        \tag{P-F49k}
\]
After taking the $h$th power and applying (P-F49g),
the determinant exponent is $2hn+D=3D$. This proves
the transformation law in (P-F49b). Finally
$H_{F_n}=x_1\cdots x_n$, so (P-F49j) proves its nonzero
Fermat value. Clearing rational denominators gives an
integral-coefficient invariant if desired, without
changing its degree or its characteristic-zero value.

The resulting invariant is a nonzero element of the
degree-$nD$ invariant space. Dualizing the finite-dimensional
representation preserves the dimension of its invariant
space, and $nD=ns/g_n$ for $s=g_nD$. This proves the
first assertion of (P-F49c); P-F47 proves the second.

\paragraph{The exact analytic premise and the degree bounds.}
Xylouris, \emph{On the least prime in an arithmetic
progression and estimates for the zeros of Dirichlet
$L$-functions}, Theorem~1.1, printed p.~66, proves
$P(a,m)\le C m^{5.18}$ for reduced residue classes,
with an effectively computable absolute constant $C$.
Applying it to $a=1,m=n$ gives the unconditional
polynomial choice in (P-F49b)--(P-F49c).

Bach--Sorenson, \emph{Explicit bounds for primes in
residue classes}, Theorem~5.3, pp.~1731--1732, gives
$P(a,m)<2(m\log m)^2$ under its ERH premise. Its
arithmetic-progression specialization takes
$E=\mathbb Q(\zeta_m)$; the relevant zeta function is
the product of primitive Dirichlet $L$-functions, as
explained in their Section~2, pp.~1719--1720.
Thus the supplied Dirichlet theorem suffices for this
particular specialization; arbitrary number-field
GRH is not being imported.

We use the sharper bound of Lamzouri--Li--Soundararajan,
\emph{Conditional bounds for the least quadratic
non-residue and related problems}, Corollary~1.2,
printed p.~3, which states
$P(a,m)\le(\varphi(m)\log m)^2$ for $m>3$.
Their Section~2, pp.~4--5, uses RH for $\zeta$ and
GRH for primitive Dirichlet $L$-functions. These lie
within the supplied A-DIR-NV/A-DIR-Z scope. Setting
$a=1,m=n\ge4$ gives (P-F49d). For the two smaller
dimensions one may take $p=3$ for $n=2$, and $p=7$
for $n=3$, obtaining stretches $2$ and $18$ respectively.
No optimality in those small cases is asserted.

The primary URLs for the analytic statements are
\begin{itemize}
\item Xylouris, Theorem~1.1, p.~66:
\url{https://www.impan.pl/shop/publication/transaction/download/product/83078}.
\item Bach--Sorenson, Theorem~5.3, pp.~1731--1732:
\url{https://doi.org/10.1090/S0025-5718-96-00763-6}.
\item Lamzouri--Li--Soundararajan, Corollary~1.2,
p.~3, and Section~2, pp.~4--5:
\url{https://arxiv.org/pdf/1309.3595}.
\end{itemize}

\paragraph{What has and has not become constructive.}
The formulas specify a genuine invariant and a proved
nonzero evaluation for every $n$, with a polynomial
degree bound even without the supplied theorem and a
sharper explicit bound using that theorem. The prime
may be found by searching the bounded progression.
However, the displayed Cayley contraction has
$(n!)^{p-1}$ summands before simplification. Its polynomial
degree does not establish a polynomial arithmetic-circuit
size or a positive counting model for multiplicities.
The theorem closes the polynomial-degree Fermat
nonvanishing question posed in P-F47, not the distinct
strict-constituent PH2 or SH question.


\subsection{A polynomial conductor for the Fermat evaluation degrees}
\label{sec:fermat-conductor}

\paragraph{P-F50 (two explicit mixed-Hessian invariants give eventual support).}
For $n\ge2$ let
\[
 \mathscr E_n=\{d\ge0:\ \exists I\in
     \mathbb C[\operatorname{Sym}^3\mathbb C^n]^{SL_n}_d,
                           \ I(F_n)\ne0\},
      \qquad F_n=\sum_i x_i^3.                         \tag{P-F50a}
\]
There are explicitly specified relatively prime positive
integers $u,v$ of polynomial size in $n$ such that
\[
          2n\langle u,v\rangle\ \subseteq\ \mathscr E_n
                         \ \subseteq\ 2n\mathbb Z_{\ge0}.
                                                        \tag{P-F50b}
\]
Here the first inclusion is realized by products of two
specified Hessian--Cayley invariants. In particular
\[
               2nr\in\mathscr E_n
             \quad\text{for all }r\ge (u-1)(v-1).
                                                        \tag{P-F50c}
\]
Unconditionally $(u-1)(v-1)=O(n^{20.72})$ with an
effective absolute implied constant. Under the supplied
Dirichlet theorem there is the explicit bound
\[
 (u-1)(v-1)<\frac14(\varphi(Q)\log Q)^4,
 \qquad Q=2n\,2^{\lceil\log_2(4n)\rceil}<16n^2.
                                                        \tag{P-F50d}
\]
Thus the bound in (P-F50d) is $O(n^8(\log n)^4)$.
In the canonical stretch variable $s$, $g_n=\gcd(n,3)$,
every $s=2g_nr$ with $r\ge(u-1)(v-1)$ has an invariant
nonzero at Fermat and hence $f_n(s)>0$.
No other stretch can be nonzero \emph{at Fermat}.
The whole invariant space can be nonzero at those other
stretches, so this last assertion is not a vanishing
claim for $f_n$.

\paragraph{A mixed-Hessian evaluation identity.}
Use the normalized Hessian $H_F$ and Cayley contraction
$\mathcal C_{D,n}$ of P-F49. Let $p$ be a prime and let
$h,j\ge0$ satisfy
\[
               nh+3j=p-1=:D,\qquad nj<p.
                                                        \tag{P-F50e}
\]
Then
\[
 J_{h,j}(F)=\mathcal C_{D,n}(H_F^h F^j)
                                                        \tag{P-F50f}
\]
is a degree-$nq$ invariant, where $q=nh+j=D-2j$, with
\[
       J_{h,j}(gF)=(\det g)^{3q}J_{h,j}(F),
                       \qquad J_{h,j}(F_n)\ne0.
                                                        \tag{P-F50g}
\]
To prove the nonzero assertion with every normalization
accounted for, first let $G=\sum_{|\alpha|=D}g_\alpha x^\alpha$
have coefficients in $\mathbb Z_{(p)}$. Grouping the
permutation terms of (P-F49f) by their exponent matrices
$L$, and then applying (P-F49h), gives
\[
 \begin{split}
 \mathcal C_{p-1,n}(G)
 &=\sum_L c_L\prod_{i=1}^n
             \left(g_{L_i}\frac{\prod_a L_{ia}!}{(p-1)!}\right)\\
 &\equiv\sum_L\prod_{i=1}^n g_{L_i}
   =[x_1^{p-1}\cdots x_n^{p-1}]G(x)^n\pmod p.
 \end{split}                                           \tag{P-F50h}
\]
Here $L_i$ is row $i$ and every row and column of $L$
sums to $p-1$. All factorial denominators are units.

At Fermat, $H_{F_n}=\prod_i x_i$. Therefore the last
coefficient in (P-F50h), for
$G=(\prod_i x_i)^h(\sum_i x_i^3)^j$, is
\[
 [x_1^D\cdots x_n^D]
       (\prod_i x_i)^{hn}(\sum_i x_i^3)^{jn}
                         =\frac{(nj)!}{(j!)^n}.        \tag{P-F50i}
\]
Indeed each exponent still required after the first
factor is $D-hn=3j$, so each of the $n$ summands must
be chosen $j$ times. The condition $nj<p$ makes this
integer nonzero modulo $p$, proving the evaluation claim.
The polynomial inserted into $\mathcal C$ has variable
degree $nh+3j=D$ and coefficient degree $nh+j=q$.
The Hessian transformation contributes determinant
exponent $2hn$, and the Cayley contraction contributes
$D$, so their sum is $3nh+3j=3q$. Its coefficient
degree is $nq$. This proves (P-F50g).

\paragraph{Two admissible mixed powers from one sufficiently large prime.}
Choose an integer $c$ with
\[
                1\le c<n,
       \qquad\gcd(c,n)=\gcd(6c+1,n)=1.                 \tag{P-F50j}
\]
Such a choice always exists. Modulo a prime divisor
$\ell=2$ or $3$ of $n$, choose any nonzero residue;
$6c+1$ is then automatically nonzero. Modulo each
$\ell\ge5$, exclude the two residues $0,-6^{-1}$;
at least three residues remain. The Chinese remainder
theorem, extended to the prime-power factors of $n$,
gives a residue modulo $n$ satisfying both coprimality
conditions. Its representative is nonzero, so lies
between one and $n-1$. Alternatively it can be found
by testing those $n-1$ integers with gcd computations.

Put
\[
 \ell=2^{\lceil\log_2(4n)\rceil},\quad Q=2n\ell,
 \quad r_0\equiv6c+1\pmod{2n},\quad 1\le r_0<2n,
 \quad a=Q/2+r_0.                                     \tag{P-F50k}
\]
The notation $\ell$ here denotes a power of two, not a
prime. It obeys $4n\le\ell<8n$, hence
$8n^2\le Q<16n^2$. The residue $r_0$ is odd and prime
to $n$. Since $\ell$ is even, $Q/2=n\ell$ is divisible
by $2n$, and consequently $a$ is odd and prime to $n$.
The only additional prime divisor of $Q$ is two.
Thus
\[
          \gcd(a,Q)=1,
                 \qquad4n^2<a<Q.                     \tag{P-F50l}
\]
Choose the least prime $p\equiv a\pmod Q$. In particular
\[
                 p>4n^2,
                 \qquad p\equiv6c+1\pmod{2n}.         \tag{P-F50m}
\]
Set
\[
 \begin{array}{lll}
 j=2c,&h=(p-1-6c)/n,&q=p-1-4c,\\
 j'=2c+n,&h'=h-3,&q'=q-2n.
 \end{array}                                           \tag{P-F50n}
\]
All these numbers are integers. To check (P-F50e) for
both pairs, note that $j'\le3n-2$ and hence
$nj'\le3n^2-2n<p$. Moreover
\[
 p-1-3j'\ge4n^2-9n+6>0\qquad(n\ge2).
\]
The last quadratic is positive because it has positive
leading coefficient and negative discriminant. Thus
$h'>0$ and $h\ge0$; both pairs meet every required
condition. The congruence modulo $2n$ makes $q,q'$ even.
Define
\[
                      u=q/2,\qquad v=q'/2=u-n.
                                                        \tag{P-F50o}
\]
They are positive since $p>4n^2$ and $c<n$.
Furthermore
\[
      u=(p-1)/2-2c\equiv c\pmod n,
          \qquad \gcd(u,v)=\gcd(u,n)=1.                \tag{P-F50p}
\]
The two instances of (P-F50f) therefore give invariants
of degrees $2nu,2nv$, both nonzero at Fermat.

\paragraph{The exact point-specific period obstruction.}
For completeness, let a homogeneous degree-$d$ $SL_n$
invariant $I$ have $I(F_n)\ne0$. The scalar center
$\mu_n\subset SL_n$ first forces $n\mid3d$.
Writing $b=3d/n$, any $g\in GL_n$ can be expressed as
a scalar times an $SL_n$ matrix, and homogeneity gives
\[
                         I(gF)=(\det g)^b I(F).
                                                        \tag{P-F50q}
\]
The diagonal matrix $\operatorname{diag}(\omega,1,\ldots,1)$,
where $\omega$ is a primitive cube root, fixes $F_n$.
Its determinant forces $3\mid b$. A transposition of
two variables also fixes $F_n$ and has determinant
$-1$, forcing $2\mid b$. Thus $6\mid b$, equivalently
$2n\mid d$. This proves the second inclusion in
(P-F50b) directly, without an assumption about the
full Fermat stabilizer. It agrees with the degree
period in B\"urgisser--Ikenmeyer, Proposition~3.24(2),
printed p.~15.

Products of the two constructed invariants stay nonzero
at $F_n$, giving the first inclusion in (P-F50b).
For relatively prime $u,v$, every integer
$t\ge(u-1)(v-1)$ belongs to $\langle u,v\rangle$.
Indeed choose $b\in\{0,\ldots,u-1\}$ with
$bv\equiv t\pmod u$. Then $t-bv$ is a multiple of $u$
and is at least $-(u-1)$, so it is nonnegative.
Setting $a=(t-bv)/u$ writes $t=au+bv$ with both
coefficients nonnegative and proves (P-F50c).

\paragraph{The unconditional and supplied-GRH bounds.}
Xylouris's Theorem~1.1, p.~66, cited in P-F49 gives
$p\le C Q^{5.18}$. Since $u,v<p/2$,
\[
       (u-1)(v-1)<p^2/4
                      \le C^2Q^{10.36}/4=O(n^{20.72}).
                                                        \tag{P-F50r}
\]
For the supplied Dirichlet theorem, apply
Lamzouri--Li--Soundararajan Corollary~1.2, p.~3,
to the reduced class $a\pmod Q$; $Q>3$ is automatic.
It gives $p\le(\varphi(Q)\log Q)^2$, and the same
inequality proves (P-F50d). The exact analytic hypotheses
and primary URLs are those checked in P-F49.

Equations (P-F50b)--(P-F50d) close the question of a
polynomial eventual-support threshold at this specific
Fermat point. They give a family of actual nonzero
cubic plethysm coefficients on all sufficiently large
stretches permitted by that point's period. They do
not claim that every invariant in those spaces is
constructed, and they give no claim of small evaluation
circuits for the Cayley contractions.

Most importantly, eventual positivity of actual values
on the subprogression $2g_n\mathbb Z$ is not a statement
about every constituent polynomial evaluated at every
positive argument after a shift. It does not settle
the strict saturation or coefficient-positivity indices
used in PH2/SH. The distinctions in P-F34 and P-F47
remain in force.

\paragraph{A full-ring support corollary on an infinite family.}
Let $r_0\ge3$ be odd, and put
\[
       k=3r_0,\qquad n=\binom{k-1}{2},\qquad
       C_n=(u-1)(v-1),
       \qquad f_n(s)=\dim\bigl(\operatorname{Sym}^{ns}
                    (\operatorname{Sym}^3\mathbb C^n)\bigr)^{SL_n}.
                                                        \tag{P-F50s}
\]
Here $u,v$ are the explicit integers constructed above for this $n$.
Since $2n=(3r_0-1)(3r_0-2)\equiv2\pmod3$, we have
$\gcd(n,3)=1$, so this is precisely the canonical stretching
grading. The top exterior-power case of (P-F32a) gives
\[
               f_n(r_0)=1,
      \qquad f_n(s)>0\quad\hbox{for every integer }
                           s\ge2C_n+r_0.              \tag{P-F50t}
\]
To verify the first assertion directly, the relevant exterior
power has degree $nr_0=\binom{k}{3}$, hence is the top exterior
power of $\bigwedge^3\mathbb C^k$. Each diagonal weight occurs
$\binom{k-1}{2}=n$ times, so that one-dimensional representation
is $(\det)^n$, with multiplicity one. Let $J$ be a nonzero
invariant in canonical degree $r_0$. For each $t\ge C_n$,
(P-F50c) supplies a nonzero invariant $I_t$ in canonical
degree $2t$. The invariant ring is a subring of the polynomial
ring in the cubic coefficients and is therefore a domain.
Thus $JI_t\ne0$, in degree $r_0+2t$.
Every even $s\ge2C_n+r_0$ is already covered by $I_{s/2}$,
and every odd such $s$ is covered by $JI_{(s-r_0)/2}$.
This proves the second assertion for all actual integer stretches,
including those outside the Fermat point's even progression.

The threshold $2C_n+r_0$ is $O(n^{20.72})$ unconditionally
and $O(n^8\log^4 n)$ under the supplied Dirichlet-GRH premise,
because $r_0=O(\sqrt n)$ and the bounds for $C_n$ were proved
above. This is a full-ring eventual-support statement on the
specified infinite subfamily of dimensions. It is still not
the strict saturation assertion that every nonzero constituent
polynomial is positive at every positive argument after a
uniform shift: actual support tests a constituent only on its
own residue progression. No PH2 or SH conclusion is added.


\subsection{The exact definite-unitary source of the positive theta factor}
\label{sec:definite-theta-factor}

\paragraph{P-F51 (the $E_6$ theta vector from a specified compact source).}
Let $K=E_6$ and $\Theta_K=\sum_{\mu\in K^\vee/K}
\theta_{K+\mu}e_\mu$ be the positive theta vector in P-F46.
There is a positive definite Hermitian space $(V,h)$ of
dimension three over $E=\mathbb Q(\zeta)$, where
$\zeta=e^{2\pi i/3}$, whose underlying trace lattice is
exactly $K$. Put $H=U(V,h)$, so that $H(\mathbb R)=U(3)$.
There are an explicitly specified compact open subgroup
$K_f\subset H(\mathbb A_f)$, a nonnegative compactly
supported automorphic function $f_K$ on
$H(\mathbb Q)\backslash H(\mathbb A)$, and three
Schwartz functions $\Phi_\mu$ such that
\[
 \theta_{K+\mu}(\tau)
   =y^{-3/2}\int_{H(\mathbb Q)\backslash H(\mathbb A)}
       \theta_{\Phi_\mu}(g_\tau,h) f_K(h)\,dh,
 \qquad \tau=x+iy.
 \tag{P-F51a}
\]
The integral is ordinary and absolutely convergent.
The normalization uses the actual measure of a specified
double class, not a genus average or an unstated mass factor.
Its right-hand group is the derived group
$SU(J)=SL_2$ of a split unitary group of signature $(1,1)$.
The Weil action in (P-F51a) is fixed by the canonical
orthogonal Schr\"odinger action; the statement does not
silently choose a central splitting for the full $U(J)$.

The underlying symplectic group contains the exact
seesaw configuration
\[
 \bigl(O(V_{\mathbb Q},B),SL_2\bigr),\qquad
 \bigl(H,U(J)\bigr),\qquad
 H\subset O(V_{\mathbb Q},B),\quad SL_2\subset U(J).
 \tag{P-F51b}
\]
Here the displayed pairs are reductive dual pairs.
The classical theta vector has weight three and the
same finite Weil type as in P-F46; it is modular on
$\Gamma(3)$ component by component. These statements
identify an actual definite-unitary source for the
positive factor, without identifying the additional
operation $\Theta_K\mapsto\delta^{-1}\Theta_K$ or its
multiplicative Borcherds lift with the specific iterated
lifts in the supplied \emph{Remark}.

\paragraph{The Hermitian lattice and the two dual pairs.}
Use the model proved in P-F46:
\[
 A_2=\{(a,b,c)\in\mathbb Z^3:a+b+c=0\},\quad
 v=(2,-1,-1)/3,\qquad
 K=A_2^3+\mathbb Z(v,v,v).
\]
On every $A_2$ factor cyclically permute the three
coordinates, and denote the resulting simultaneous
operator by $\sigma$. It preserves the Euclidean form
$B$, has $\sigma^2+\sigma+1=0$, and satisfies
$\sigma v-v\in A_2$. Thus it preserves the actual glue
lattice $K$, not just $A_2^3$. Let $\zeta$ act by
$\sigma$ on $V=K\otimes\mathbb Q$, and set
\[
 h(u,w)=\frac{1-\zeta^2}{3}B(u,w)
       +\frac{\zeta-\zeta^2}{3}B(u,\sigma w).
 \tag{P-F51c}
\]
As in the direct calculation in (P-F46l), the identities
$B(\sigma u,w)=-B(u,w)-B(u,\sigma w)$ and
$B(\sigma u,\sigma w)=B(u,w)$ give sesquilinearity,
$h(w,u)=\overline{h(u,w)}$, and
\[
       \operatorname{Tr}_{E/\mathbb Q}h=B,
       \qquad h(u,u)=B(u,u)/2.
\]
Consequently $h$ is positive definite, has rank three,
and is even on $K$. Trace duality, applied also to
$au$ for every $a\in\mathbb Z[\zeta]$, shows
$h(K,K)\subset\mathfrak D_E^{-1}$ and that its
inverse-different Hermitian dual
$\{x\in V:h(x,K)\subset\mathfrak D_E^{-1}\}$ is exactly
$K^\vee$. In particular the
quadratic exponents and discriminant cosets are unchanged.

Let $J=E^2$ have the skew-Hermitian form
\[
 j((a,b),(c,d))=a\overline d-b\overline c.
\]
On $\mathcal W=\operatorname{Res}_{E/\mathbb Q}(V\otimes_EJ)$
use the alternating form
$\Omega=\operatorname{Tr}_{E/\mathbb Q}(h\otimes j)$.
Identifying $\mathcal W$ with $V_{\mathbb Q}\oplus V_{\mathbb Q}$
gives exactly
\[
 \Omega((u_1,u_2),(w_1,w_2))
            =B(u_1,w_2)-B(u_2,w_1).
 \tag{P-F51d}
\]
This specifies the trace normalization; no factor of two
is suppressed in the passage to the oscillator representation.
Multiplication of $j$ by $i$ at the real place has one
positive and one negative eigenvalue, so its unitary
group has signature $(1,1)$.

If $A\in SL_2(E)$, then $A^tJA=J$ for the alternating
matrix $J$. The unitary condition is $A^tJ\overline A=J$.
Together these conditions imply $\overline A=A$.
Conversely every $A\in SL_2(\mathbb Q)$ satisfies both.
The same calculation over every $\mathbb Q$-algebra
proves $SU(J)=SL_2$ as algebraic groups, not only on
real points. Its action on the two summands of (P-F51d)
is the usual symplectic action. The inclusions in
(P-F51b) therefore preserve the very same symplectic form.

For completeness, the centralizers can also be checked
directly. After scalar extension to $\mathbb C$, the
unitary group $H$ is $GL_3$, and $V_{\mathbb Q}\otimes\mathbb C$
is the sum of its standard representation and its dual.
The two are nonisomorphic, and each occurs twice in
$\mathcal W\otimes\mathbb C$. Their commuting
endomorphisms are two matrices of size two; preservation
of $\Omega$ makes them inverse transposes. Descending
back to $\mathbb Q$ gives exactly $U(J)$.
The reverse centralizer is $H$ by the same calculation
with dimensions two and three exchanged. Likewise the
standard representation of $O_6$ is absolutely irreducible;
its multiplicity space in $\mathcal W$ has dimension two,
and preservation of its alternating form leaves exactly
$SL_2$. Conversely the centralizer of this $SL_2$ is
$O(V_{\mathbb Q},B)$. This proves both dual-pair assertions
and the stated seesaw configuration.

\paragraph{The splitting and the exact finite level.}
Take the standard additive character of
$\mathbb Q\backslash\mathbb A$, whose real component is
$\psi_\infty(t)=e^{2\pi it}$, and the usual self-dual
measures. The Schr\"odinger model on
$\mathcal S(V_{\mathbb Q}(\mathbb A))$ has the canonical
action
\[
        (\omega(h)\Phi)(u)=\Phi(h^{-1}u),
                   \qquad h\in O(V_{\mathbb Q},B)(\mathbb A).
 \tag{P-F51e}
\]
For the $SL_2$ unipotent $n(b)$ it has action
$\omega(n(b))\Phi(u)=\psi(bB(u,u)/2)\Phi(u)$.
Since the quadratic dimension is six, the symplectic
Weil representation descends from the double cover
to $SL_2$. These actions and the Fourier-transform
operator commute with (P-F51e), because $h$ preserves
$B$ and its self-dual measure. This fixes the action
on $H\times SL_2$ used here. At finite places these are
Kudla, \emph{Notes on the local theta correspondence},
Definition~II.4.2 and Proposition~II.4.3, pp.~37--38,
and the orthogonal splitting on p.~39. The real action
is the corresponding Gaussian Schr\"odinger action,
whose required value is computed below.

The complete finite unitary level is
\[
 K_f=\{h\in H(\mathbb A_f):h\widehat K=\widehat K,
                  \ h\text{ fixes }\widehat K^\vee/\widehat K
                          \text{ pointwise}\}.
 \tag{P-F51f}
\]
This is compact open: it is the kernel of a finite
action of the lattice stabilizer. Away from three,
$K_p=K\otimes\mathbb Z_p$ is self-dual, so the discriminant condition
adds no restriction there. For each class $\mu$ put
\[
 \Phi_\mu(u)=e^{-\pi B(u_\infty,u_\infty)}
                 \mathbf1_{\mu+\widehat K}(u_f).
 \tag{P-F51g}
\]
The Gaussian is $H(\mathbb R)$-invariant and the
finite indicator is $K_f$-invariant. Requiring each
coset to be fixed, rather than merely the lattice,
is what makes this simultaneous assertion immediate.

In the labels $\mu\in\mathbb Z/3\mathbb Z$ with
$q(\mu)=2\mu^2/3$, the classical Weil matrices are
\[
 \rho_K(T)e_\mu=e^{2\pi i(2\mu^2/3)}e_\mu,
 \qquad
 \rho_K(S)e_\mu=\frac{i}{\sqrt3}
                  \sum_\nu e^{-2\pi i\mu\nu/3}e_\nu.
 \tag{P-F51h}
\]
Indeed the first formula follows term by term; the
second is Poisson summation, with the factor
$(-i\tau)^3=i\tau^3$, followed if necessary by
$\nu\mapsto-\nu$, since opposite cosets have equal
theta series. These are the same convention as P-F46.
The finite quadratic module has level three.
Scheithauer, Proposition~4.3, printed p.~14, gives
trivial action of $\Gamma(3)$. His displayed convention
is dual to this one; taking a dual leaves this kernel
assertion unchanged. Thus each component is a genuine
weight-three modular form on $\Gamma(3)$, without
claiming an unspecified scalar level for the full
unitary group.

A full $U(J)$ formulation would additionally select
unitary splitting characters with restrictions
$\chi_V|_{\mathbb A_\mathbb Q^\times}=\varepsilon_{E/\mathbb Q}^3$
and $\chi_J|_{\mathbb A_\mathbb Q^\times}=\varepsilon_{E/\mathbb Q}^2$.
These are the conditions in Harris--Li--Sun,
equation~(1.1.1), p.~3, and Section~1.2, p.~4.
No such choice, nor compatibility with a different
iterated lift's choices, is silently supplied by
(P-F51a); the derived-group and orthogonal splitting
above already suffice for its exact identity.

\paragraph{Selecting the individual lattice class.}
Let $C_K$ be the image of
$H(\mathbb R)K_f$ in $[H]=H(\mathbb Q)\backslash H(\mathbb A)$.
It is a nonempty compact open subset: its preimage is
$H(\mathbb Q)H(\mathbb R)K_f$, which is open, and
$H(\mathbb R)K_f$ is compact. For any fixed Haar
quotient measure its volume $c_K$ is finite and positive.
Set
\[
                   f_K=c_K^{-1}\mathbf1_{C_K}.
 \tag{P-F51i}
\]
It is an automorphic function, is invariant under
$H(\mathbb R)K_f$ on the right, and has compact support.
In particular it is square-integrable. There is no
assertion that this lattice-class selector lies in
one preselected irreducible automorphic representation.

Define the theta kernel in the stated Weil model by
\[
 \theta_{\Phi_\mu}(g,h)=
       \sum_{u\in V_{\mathbb Q}(\mathbb Q)}
                         (\omega(g)\omega(h)\Phi_\mu)(u).
 \tag{P-F51j}
\]
The automorphic theta sum and its integration convention
are the constructions of Harris--Li--Sun,
equations~(1.2.4)--(1.2.5), p.~4, applied to the
explicitly fixed commuting actions above.
Rational $H$-invariance also follows
directly by changing the summation variable.
For $h\in C_K$, write a representative as
$\gamma h_\infty k$ with $\gamma\in H(\mathbb Q)$,
$h_\infty\in H(\mathbb R)$, and $k\in K_f$.
Rational invariance and (P-F51g) give
$\theta_{\Phi_\mu}(g,h)=\theta_{\Phi_\mu}(g,1)$.
Thus its integral against (P-F51i) is exactly
$\theta_{\Phi_\mu}(g,1)$. This proves the normalization
and shows why a genus average has not been substituted
for the requested individual lattice theta series.

Finally take $g_\tau=n(x)m(\sqrt y)$ at infinity and
identity at all finite places. In real quadratic
dimension six the dilation factor is
$(\sqrt y)^{6/2}=y^{3/2}$. The two real operators give
\[
 (\omega(g_\tau)\Phi_\mu)(u)
     =y^{3/2}e^{2\pi i\tau B(u_\infty,u_\infty)/2}
                            \mathbf1_{\mu+\widehat K}(u_f).
\]
The finite condition on a rational vector is exactly
$u\in K+\mu$; this follows coordinate by coordinate in
an integral basis of $K$. Summing proves (P-F51a),
with the exact power of $y$ and the exact exponent.
The Gaussian absolute-value series is convergent,
and on the compact support $C_K$ the kernel is constant;
therefore the integral used in the identity is absolutely
convergent as well.\hfill$\square$

\paragraph{The precise overlap with the supplied theta method.}
This is a source-map statement for the actual factor
$\Theta_{E_6}$ in (P-F46a): it comes from a specified
nonnegative function on a definite $U(3)$ quotient,
with coset level and measure normalization fixed.
The compatible seesaw in (P-F51b) is also explicit.
Its right unitary space is split of signature $(1,1)$;
it is not a pair of totally definite spaces of the
kind named in item~8 of the supplied \emph{Remark}.
This identifies the exact pair involved rather than
asserting that signature differences prohibit all
comparisons. A further seesaw or representation-theoretic
identity would have to match the selected vectors and
splittings to the particular $U(X_j)$ in that note.

Two additional operations remain after this positive
definite factor: multiplication by $\delta^{-1}$ and
the regularized multiplicative lift identified in
P-F46. The former introduces the displayed negative
Fourier powers; the latter outputs the single cubic
discriminant, rather than every invariant-ring graded
piece. The ordinary convergence just proved for
$\Theta_{E_6}$ does not remove the explicit divergence
proved in P-F46 for $g=\Theta_{E_6}/\delta$.
No simultaneous-stretch Hilbert identity, strict-index
bound, or comparison with the supplied iterated lift's
Bessel period is inferred from (P-F51a).

Primary source locators: Stephen Kudla,
\emph{Notes on the local theta correspondence},
Definition~II.4.2 and Proposition~II.4.3, pp.~37--38;
orthogonal splitting, p.~39; seesaw definition,
Section~IV.1, p.~59,
\url{https://www.math.toronto.edu/~skudla/castle.pdf}.
Michael Harris, Jian-Shu Li and Binyong Sun,
\emph{Theta correspondences for close unitary groups},
equation~(1.1.1), p.~3, and equations~(1.2.1)--(1.2.5),
p.~4,
\url{https://webusers.imj-prg.fr/~michael.harris/closeunitaryweb.pdf}.
Nils Scheithauer, \emph{The Weil representation of $SL_2(\mathbb Z)$
and some applications}, Proposition~4.3, p.~14,
\url{https://www.mathematik.tu-darmstadt.de/media/algebra/homepages/scheithauer/publications/Weil_representation.pdf}.


\subsection{The actual central word realizing the quadratic coefficient}
\label{sec:nonstandard-central-word}

\paragraph{P-F52 (an intertwiner inside the displayed nonstandard algebra).}
Use the cubic algebra $B_0$ over $F_0=\mathbb C(q)$, its splitting
field $K=F_0(x)$, and the elements $c,c',P,D,F$ of P-F36 and P-F39.
Here $x^2=D$; the letter $F$ without a subscript continues to denote
the positive Laurent normalization in (P-F39a).
Write $A=Q_1$, $B=Q_2$, and
\[
 f=(q^4+1)(q^4-q^2+1)(q^2+1)/q^5
   =q^5+q+q^{-1}+q^{-5}.
\]
Define the following explicit Laurent polynomial and central word:
\[
 \begin{split}
 S={}&q^{10}+q^{-10}+3(q^8+q^{-8})+(q^6+q^{-6})
       +5(q^4+q^{-4})-2(q^2+q^{-2})+16,\\
 Z={}&(Q_1-Q_2)^2+(S-f^2)1.
 \end{split}                                                    \tag{P-F52a}
\]
In the actual source blocks $U_1,U_2\subset B_0\otimes K$,
left and right multiplication by $Z$ have scalar values $-c,-c'$,
respectively. In particular, this operator is specified by a short
word in the original generators, rather than by assigning arbitrary
eigenvalues to abstract central idempotents.

There is an explicit two-dimensional $F_0$-subspace of $B_0$ on
which multiplication by $Z$ has precisely the matrix $M_0$ of
(P-F39k). Its multiple $FZ$ has matrix $M=FM_0$.
The same construction gives four copies of this matrix on the
eight-dimensional descended two-sided ideal associated with
$U_1\oplus U_2$. These are statements about actual multiplication
in $B_0$, with the original source basis retained.

For every real $q>0$, the same word $Z$ acts as a strictly positive,
self-adjoint operator in the source's actual tensor representation
on $X_q^{\otimes3}$, with the explicit positive tensor form given
below. So does $FZ$. This last assertion concerns the tensor
specialization, whose image can have smaller dimension at exceptional
parameters; the generic $F_0$-linear intertwiners are not asserted to
remain injective at every specialization.

\paragraph{Only the original relation and the original cell element are needed.}
The source gives $A^2=fA$, $B^2=fB$, and the independent
alternating-word basis of lengths at most ten: GCT VII,
Section~7.1.1, pp.~18--19. Put
\[
       C=(A-B)^2=f(A+B)-AB-BA,
       \qquad G=f^2-C,
       \qquad w_j=A(BA)^j.
                                                               \tag{P-F52b}
\]
Direct multiplication gives
$AC=CA=f^2A-ABA$ and $BC=CB=f^2B-BAB$.
Thus $C$ and $G$ are central, and
\[
                    Gw_j=w_{j+1}\qquad(j\ge0).                 \tag{P-F52c}
\]
The actual relation in Figure~2, printed p.~31, has the form
\[
            w_5-Sw_4+a_3w_3+a_2w_2+a_1w_1+a_0w_0=0,          \tag{P-F52d}
\]
where the coefficient of $w_5$ is one and the coefficient of
$w_4$ is exactly $-S$ from (P-F52a). The actual element $u=u_1^1$ in
Figure~7, printed p.~36, has the form
\[
                  u=w_4+cw_3+b_2w_2+b_1w_1+b_0w_0.            \tag{P-F52e}
\]
Its conjugate $u'=u_2^1$ replaces $x$ by $-x$ and therefore has
coefficient $c'$ at $w_3$. This is the same scalar coefficient
used in GCT VIII and P-F39, even though VIII also substitutes
$P$-words for $Q$-words in its corresponding element.

The following full coefficient identity verifies the eigenvector
equation without relying on the simple-block assertion. Set
\[
 \begin{gathered}
 R=-3q^{28}-4q^{24}+2q^{22}-10q^{20}+2q^{18}-4q^{16}-3q^{12},\\
 g_1=\frac{-R-x}{2q^{20}},\qquad
 g_2=\frac{-R+x}{2q^{20}},\\
 g_3=q^4+2+q^{-4},\quad g_4=q^6+2+q^{-6},\quad
 g_5=q^{10}+2+q^{-10},\\
 E_1=g_3+g_4+g_5,\quad E_2=g_3g_4+g_3g_5+g_4g_5,\quad
 E_3=g_3g_4g_5,\\
 L=-R/q^{20},\qquad
 N_0=\frac{(q^4-q^2+1)^2
             (q^{12}+q^{10}+2q^8+2q^4+q^2+1)^2}{q^{16}}.
 \end{gathered}                                                 \tag{P-F52f1}
\]
Expansion gives $R^2-D=4q^{40}N_0$, $g_1+g_2=L$,
$g_1g_2=N_0$, and $L+E_1=S$. All coefficients of the
original relation and the original cell element can be written
as follows; the expressions are direct factorizations of
Figures~2 and 7, including their lower coefficients:
\[
\begin{array}{c|c|c}
 j&[t^j]r(t)&[t^j]p_1(t)\\ \hline
 0&-N_0E_3&g_2E_3\\
 1&N_0E_2+LE_3&-E_3-g_2E_2\\
 2&-N_0E_1-LE_2-E_3&E_2+g_2E_1\\
 3&N_0+LE_1+E_2&-E_1-g_2=c\\
 4&-L-E_1=-S&1\\
 5&1&0
\end{array}                                                     \tag{P-F52f2}
\]
Here $r(G)A=0$ is exactly (P-F52d), and $p_1(G)A=u$
is exactly (P-F52e). For example their constant coefficients
are the factored first rows of the two source figures; the
other rows follow by the same multiplication and collection
of powers of $q$ and $x$. The accompanying exact verifier
transcribes and compares all six relation coefficients and
all five cell coefficients, without interpolation.
Multiplying the polynomials in the table now gives
\[
 \begin{split}
 r(t)&=(t^2-Lt+N_0)(t^3-E_1t^2+E_2t-E_3),\\
 p_1(t)&=(t-g_2)(t^3-E_1t^2+E_2t-E_3),\\
                 (t-g_1)p_1(t)&=r(t).
 \end{split}                                                    \tag{P-F52f3}
\]
Substitution of the central element $G$ proves
$(G-g_1)u=r(G)A=0$ directly. Since $g_1=S+c$,
\[
                 Gu=(S+c)u,\qquad Zu=-cu.                     \tag{P-F52f}
\]
The leading coefficient gives an independent quick check:
in $Gu$ only the reduction of $w_5$ and the term $cw_4$
contribute to $w_4$, with coefficient $S+c$.
Conjugating gives $Zu'=-c'u'$. Both elements are nonzero because
their coefficient at the independent word $w_4$ is one.
Moreover $c-c'=-x/q^{20}\ne0$, so they are distinct; their
belonging to different direct summands also makes them
linearly independent over $K$.

To label the corresponding full blocks, use the split Wedderburn
decomposition in GCT VII, Section~7.1.2, pp.~19--20: $u,u'$
belong to the four-dimensional simple matrix-algebra summands
$U_1,U_2$. The center of a matrix algebra consists of scalars,
as follows by commuting with its diagonal and off-diagonal matrix
units. Thus (P-F52f) gives the same scalar on the full respective
summands, on either side. The eigenvector identity itself has
already been verified independently in (P-F52f1)--(P-F52f3).

This argument does not use the printed two-by-two generator
matrices on p.~20. As printed, both are upper triangular and
would give a common invariant line, inconsistent with the
generic irreducibility claimed there. That display is
quarantined here; no silent correction to it is needed for
(P-F52f), which follows from the complete original relation
and cell-element factorization. The cited split simple-algebra
assertions serve separately to identify the full block labels.

\paragraph{The exact descended intertwiner.}
Let $\tau(x)=-x$ be the nontrivial automorphism of $K/F_0$.
Set
\[
            v_0=u+u',\qquad v_1=x(u-u').                       \tag{P-F52g}
\]
They are fixed by $\tau$ and hence lie in $B_0$. To see this
last assertion explicitly, expand in the $F_0$-word basis:
each fixed coefficient in $K$ belongs to $F_0$.
The coefficient of $w_4$ in $v_0$ is two and in $v_1$ is zero,
while that of $w_3$ in $v_1$ is
$x(c-c')=-D/q^{20}\ne0$. Thus $v_0,v_1$ are independent.

The $F_0$-linear map
\[
       \iota:K\longrightarrow B_0,
           \qquad a\longmapsto au+\tau(a)u'                  \tag{P-F52h}
\]
is injective and sends the basis $(1,x)$ to $(v_0,v_1)$.
Equation (P-F52f) gives
\[
       Z\,\iota(a)=\iota\left(\frac{P+x}{2q^{20}}a\right).
                                                               \tag{P-F52i}
\]
In that basis, actual multiplication by $Z$ is therefore
\[
                 \frac1{2q^{20}}
                       \begin{pmatrix}P&D\\1&P\end{pmatrix}, \tag{P-F52j}
\]
exactly (P-F39k), with no transpose, sign, or scalar discrepancy.
Multiplying $Z$ by the source normalization $F$ gives (P-F39m).
The map $\iota$ intertwines the indicated central action; its
two-dimensional image is not asserted to be a left ideal of $B_0$.

For the actual two-sided ideal, take the four source basis elements
of $U_1$, namely $u_1^1,u_1^{12},u_1^{21},u_1^2$, and their
conjugates in $U_2$. Call them $u_r,u'_r$, respectively.
The eight vectors $u_r+u'_r$, $x(u_r-u'_r)$ are $F_0$-rational
and independent: a relation, projected onto $U_1$ and $U_2$,
has coefficients $a_r+xb_r=0$ and $a_r-xb_r=0$, forcing
$a_r=b_r=0$. Conversely every fixed vector in
$U_1\oplus U_2$ is of this form by writing its $U_1$ coefficients
as $a_r+xb_r$. Their span is therefore the descended two-sided
ideal $W=(U_1\oplus U_2)^\tau$, of dimension eight over $F_0$.
Multiplication by $Z$ acts on it as four copies of (P-F52j).
The $U_i$ here are matrix-algebra summands, of dimension four;
the source's simple left modules $T_i$ have dimension two.
These two roles are not identified.

\paragraph{An explicit positive form for the actual tensor action.}
For $q>0$ set
\[
 [2]=q+q^{-1},\quad [3]=q^2+1+q^{-2},\qquad
 (d_0,d_1,d_2,d_3)
       =\bigl(1,q^{-1}[3],[3][2]^2,q^3[3]^2[2]^2\bigr).
                                                               \tag{P-F52k}
\]
All four entries are positive. Give the source basis
$x_i=f^ix_0$ of $X_q$ the diagonal Gram matrix
$\operatorname{diag}(d_0,d_1,d_2,d_3)$, and use its tensor
product on tensor powers. Here the lowering operator denoted
$f$ in the source's basis notation is distinct from the scalar
Laurent polynomial $f$ in (P-F52a).

The actual rescaled projection matrix $\mathsf P$ in GCT VII,
Figure~1, printed p.~30, is self-adjoint for this tensor form.
The entire check can be organized by $i+j$, since its nonzero
entries preserve that sum. The one-dimensional blocks and the
two-dimensional blocks $i+j=1,5$ have real diagonal entries
and equal opposite off-diagonal entries, with equal Gram weights.
For $i+j=2$, the weights are $(d_2,d_1^2,d_2)$; the only
nontrivial opposite-entry ratio is verified by
\[
 \frac{d_2}{d_1^2}
       =\frac{q^2(q^2+1)^2}{q^4+q^2+1}.
\]
For $i+j=3$, the weights are
$(d_3,d_1d_2,d_1d_2,d_3)$; the displayed outer-to-inner
opposite-entry ratios are $q^{-4}$, and
$d_3/(d_1d_2)=q^4$. Opposite entries between equal-weight
positions are equal. Finally, for $i+j=4$ the weights are
$(d_1d_3,d_2^2,d_1d_3)$, and the required ratio is
\[
 \frac{d_1d_3}{d_2^2}
       =\frac{q^2(q^4+q^2+1)}{(q^2+1)^2}.
\]
Substitution in the entries of Figure~1 gives
$d_id_j\mathsf P_{ij,kl}=d_kd_l\mathsf P_{kl,ij}$
in every block. These identities are Laurent identities,
so also hold at $q=1$, where some off-diagonal entries vanish;
no division by such an entry is being asserted there.
Thus $\mathsf P$ is self-adjoint. The source's complementary
rescaled projection is $\mathsf Q=fI-\mathsf P$ by
equations~(4), (16)--(18), pp.~6--7 and 18, so it too is
self-adjoint. Consequently $Q_1=\mathsf Q\otimes I$ and
$Q_2=I\otimes\mathsf Q$ on $X_q^{\otimes3}$ are
self-adjoint for the displayed positive product form.

Put $t=q^2+q^{-2}=2+v$, so $v\ge0$. A direct expansion gives
\[
 \begin{split}
 S-f^2
 &=3(q^8+q^{-8})-(q^6+q^{-6})
       +3(q^4+q^{-4})-3(q^2+q^{-2})+12\\
 &=3t^4-t^3-9t^2+12
   =3v^4+23v^3+57v^2+48v+16>0.
 \end{split}                                                    \tag{P-F52l}
\]
For every nonzero tensor vector $w$, self-adjointness now gives
\[
 \langle w,Zw\rangle
   =\|(Q_1-Q_2)w\|^2+(S-f^2)\|w\|^2>0.
\]
This proves actual positive-definiteness, without importing a
generic compact-star assertion across an exceptional specialization.
The coefficients of $F$ are positive, so $F(q)>0$ and the same
conclusion holds for $FZ$.\hfill$\square$

\paragraph{What this finite comparison settles.}
The positive quadratic operator of P-F39 is now realized by an
explicit element of the actual nonstandard algebra, with a
source-basis intertwiner and an explicit tensor positive form.
In particular the occurrence of the eigenvalue $-c$ is tied to
the displayed word coefficient of $u_1^1$ through the leading-term
calculation (P-F52d)--(P-F52f); it is not imposed after choosing an
unrelated matrix model. On the descended two-dimensional space,
however, the unscaled determinant still has coefficients $T,N$
from (P-F39c), with their negative Laurent coefficients. Thus the
actual positivity just proved does not imply coefficientwise
positivity. The prescribed scale $F$, the other source coefficients,
their canonical-basis meanings, and any cohomological or theta
intertwiner remain substantive additional data. No uniform positivity
theorem in tensor degree or PH2/SH consequence is claimed here.

Primary source: GCT VII, Section~7.1.1--7.1.2, pp.~18--20,
Figure~1 p.~30, Figure~2 p.~31, Figure~7 p.~36,
\url{https://arxiv.org/pdf/0709.0749}; GCT VIII's matching scalar
coefficient and its normalization are cited and checked in P-F39,
\url{https://arxiv.org/pdf/0709.0751}.

\paragraph{A finite normalization corollary to P-F52.}
Even the combination of positive actual central action and symmetric
Tate--Lefschetz trace and norm does not characterize the source's
prescribed normalization $F$. Indeed put
\[
       H(q)=(q+1+q^{-1})^4
        =19+16(q+q^{-1})+10(q^2+q^{-2})
                    +4(q^3+q^{-3})+(q^4+q^{-4}).               \tag{P-F52m}
\]
For every $q>0$, $HZ$ is an actual strictly positive central
operator on the same tensor space, since $H(q)>0$.
On the descended two-dimensional space its determinant polynomial is
\[
             \det(I+tHZ)=1+H(q)T(q)t+H(q)^2N(q)t^2,             \tag{P-F52n}
\]
where $T,N$ are precisely (P-F39c).
In the symmetric-vector notation $[v]$ of P-F39, direct convolution
gives the complete coefficient certificates
\[
 \begin{split}
 HT=[\alpha^H],\quad
 \alpha^H={}&(390,336,242,160,139,132,128,112,103,88,70,44,23,8,2),\\
 H^2N=[\beta^H],\quad
 \beta^H={}&(169556,164696,152852,139200,127656,118312,109695,\\
 &100480,90558,79776,68121,55832,44025,33832,25905,20000,\\
 &15558,12016,9093,6632,4587,2936,1699,864,377,136,39,8,1).
 \end{split}                                                    \tag{P-F52o}
\]
Every entry is positive, and each list is nonincreasing.
These are exact finite identities: if $h_j$ are the coefficients
in (P-F52m), extended symmetrically and by zero, then
\[
 \alpha_j^H=\sum_i h_i t_{j-i},\qquad
 \beta_j^H=\sum_{i,k}h_ih_kn_{j-i-k},                           \tag{P-F52p}
\]
using the explicitly displayed coefficients $t_j,n_j$ of (P-F39c).
The two coefficient sums at one are $81\cdot44=3564$ and
$81^2\cdot448=2939328$. The telescoping formula (P-F39h)
therefore realizes each as a nonnegative direct sum of the same
Tate-twisted projective-space Lefschetz characters.

Nevertheless $H$ is not divisible by $F$ in
$\mathbb C[q,q^{-1}]$. For a nonzero Laurent polynomial define
its width as its largest exponent minus its smallest exponent.
The extremal coefficients in a product are nonzero products,
so widths add. Here $F$ has width $32$ by (P-F39a), whereas
$H$ has width $8$, excluding $H=FJ$ for any nonzero Laurent
polynomial $J$. In particular these conditions do not even
force the prescribed factor to divide the normalization.

This proves the precise failure of characterization by the stated
finite operator and determinant properties. It does not assert
that $H$ has the source's additional lattice, integral-form, or
canonical-basis compatibility, or that it normalizes all other
coefficients. Those compatibility conditions remain the missing
datum; positivity of this actual action and its two Lefschetz
coefficients alone cannot supply them.\hfill$\square$

\paragraph{The exact spectral-denominator provenance of $F$.}
The additional rule that specifies this source normalization can
itself be identified explicitly. In GCT VII, p.~27,
equations~(32)--(35), the actual two-tensor braid operator
$\widehat R$ has eigenvalues
\[
       q_1=q^{9/2},\qquad q_2=-q^{-3/2},\qquad
       q_3=q^{-11/2},\qquad q_4=-q^{-15/2}.                    \tag{P-F52q}
\]
The positive spectral projector is expressed by Lagrange interpolation:
\[
  P_+=\sum_{i\in\{1,3\}}
       \frac{\prod_{j\ne i}(\widehat R-q_j)}
            {\prod_{j\ne i}(q_i-q_j)}.
\]
The source specifies the product of the two displayed denominators,
not the smallest common denominator after matrix cancellations:
\[
 f_p=\prod_{i\in\{1,3\}}\prod_{j\ne i}(q_i-q_j).
                                                               \tag{P-F52r}
\]
In particular the factors $q_1-q_3$ and $q_3-q_1$ both occur.
Their product is $-q^{-11}(q^{10}-1)^2$. The other four factors
are respectively
\[
 q^{-3/2}(q^6+1),\quad q^{-15/2}(q^{12}+1),\quad
 q^{-11/2}(q^4+1),\quad q^{-15/2}(q^2+1).
\]
Thus their exact product and its ratio to the convenient matrix
normalizer $f=q^{-5}(q^4+1)(q^6+1)$ are
\[
 \begin{split}
 f_p&=-q^{-33}(q^{10}-1)^2(q^6+1)(q^{12}+1)(q^4+1)(q^2+1),\\
 f_p/f&=-q^{-28}(q^{10}-1)^2(q^{12}+1)(q^2+1).
 \end{split}
\]
The source then removes the displayed vanishing factor by defining
$\widehat f_p=-q f_p/((q-1)^2f)$ on printed p.~28.
Consequently the centering used in P-F39 gives exactly
\[
 \begin{split}
 F=q^{11}\widehat f_p
  &=q^{-16}(1+q+\cdots+q^9)^2(1+q^2)(1+q^{12}).
 \end{split}                                                    \tag{P-F52s}
\]
This is a Laurent identity, not an evaluation at a singular
interpolation parameter. The factorization cancels $(q-1)^2$
and extends regularly to $q=1$, where $F(1)=10^2\cdot2\cdot2=400$.
It proves positivity of the coefficients directly. Expanding it
also gives exactly the complete coefficient list in (P-F39a).

Thus the finite provenance of $F$ is fixed once one supplies the
source's particular interpolation-denominator product and its
specified removal and centering factors. The positive central
action alone does not impose that interpolation rule, as the
alternative $H$ proves. This dictionary also does not assert that
the spectral-denominator rule supplies a categorical integral form,
canonical-basis multiplication, or the required geometric
intertwiner for all nonstandard coefficients. Those additional
compatibilities are distinct from the explicit algebraic
provenance just established.\hfill$\square$

\paragraph{An ordinary lambda-operation obstruction for the two coefficient models.}
The separate positive coefficient models cannot be combined into
the first two ordinary exterior-power characters of one graded
vector space, even after replacing $F$ by any common scalar
normalization. More precisely, there is no rational function
$J(q)$ regular and nonzero at $q=1$ and no finite-dimensional
$\mathbb Z$-graded vector space $V$ for which
\[
 \operatorname{ch}_q V=J(q)T(q),\qquad
 \operatorname{ch}_q\bigwedge^2V=J(q)^2N(q).
                                                        \tag{P-F52t}
\]
The same obstruction holds for a virtual graded-vector-space
class with the ordinary special lambda-ring operations.

Indeed evaluation of a graded character at one is its dimension.
Writing $r=44J(1)$ and using $T(1)=44$, $N(1)=448$ from
P-F39, the two equalities in (P-F52t) would give
\[
 r\in\mathbb Z,\quad r\ne0,\qquad
 \frac{r(r-1)}2=448J(1)^2=\frac{28}{121}r^2.
                                                        \tag{P-F52u}
\]
Cancelling $r$ yields $65r=121$, impossible for an integer.
For virtual classes the same calculation applies: their rank
is an integer and the rank of the second lambda operation is
$\binom r2$, as follows by expanding
$(1+t)^{r}$ (including negative integral exponents).
For the prescribed $F$ the numerical discrepancy is particularly
explicit:
\[
 \binom{17600}{2}=154871200\ne71680000=N_1(1).
                                                        \tag{P-F52v}
\]
Thus individual Tate--Lefschetz character realizations of
$T_1$ and $N_1$ do not give this common-object lambda
dictionary.

There is a related obstruction when $q$ is interpreted as
finite-field cardinality and the same Frobenius is required
over all finite extensions. Fix any integer $p>1$. There is
no single finite-dimensional linear operator $A$ over a
characteristic-zero field whose traces satisfy
\[
 \operatorname{Tr}(A^r)=T_1(p^r),\qquad
 \operatorname{Tr}\bigl(\bigwedge^2 A^r\bigr)=N_1(p^r)
                   \quad\hbox{for every }r\ge1.
                                                        \tag{P-F52w}
\]
To prove this, triangularize $A$ over an algebraic closure.
If its diagonal entries are $a_i$, summing
$a_i^r a_j^r$ over $i<j$ gives the elementary identity
\[
 2\operatorname{Tr}\bigl(\bigwedge^2 A^r\bigr)
       =\operatorname{Tr}(A^r)^2-\operatorname{Tr}(A^{2r}).
\]
The proposed traces would therefore make the Laurent
polynomial
\[
       2N_1(q)-T_1(q)^2+T_1(q^2)
                                                        \tag{P-F52x}
\]
vanish at every $q=p^r$. These are infinitely many distinct
nonzero values, so after multiplication by a power of $q$
the resulting polynomial must vanish identically. Evaluating
at one contradicts (P-F52v). No purity or semisimplicity
assumption is used in this contradiction.

These statements exclude the specified ordinary lambda
dictionary and the specified extension-compatible Frobenius
dictionary. A graded trace of an additional endomorphism,
a different relationship between the two coefficient objects,
or a different operation is not identified with an ordinary
graded dimension. In particular the actual central operator
and its determinant remain valid; the obstruction concerns
their proposed common-object cohomological interpretation,
not the algebraic determinant identity.\hfill$\square$


\subsection{Hilbert schemes and the exact positive theta input}

\paragraph{P-F53 (a geometric coefficient model for the identified input).}
Let $S$ be a complex projective K3 surface, and let $S^{[m]}$
be its Hilbert scheme of length-$m$ subschemes, with
$S^{[0]}=\operatorname{Spec}\mathbb C$.
Retain the positive-definite lattice $K=E_6$, its discriminant
cosets $\mu\in K^\vee/K$, and the quadratic form
$Q(x)=(x,x)/2$ of P-F46 and P-F51. Write the actual
Borcherds input from P-F46 as
\[
 g_\mu(\tau)=\frac{\Theta_{K+\mu}(\tau)}{\delta(\tau)}
        =\sum_a c_\mu(a)q^a,
 \qquad \delta(\tau)=q\prod_{r\ge1}(1-q^r)^{24}.
                                                        \tag{P-F53a}
\]
For each rational exponent $a$, define the following finite disjoint union:
\[
 X_{\mu,a}=
 \coprod_{\substack{x\in K+\mu\\
                  m=a+1-Q(x)\in\mathbb Z_{\ge0}}} S^{[m]}.
                                                        \tag{P-F53b}
\]
An empty disjoint union is allowed. Then $X_{\mu,a}$ is a
smooth projective complex scheme (possibly disconnected and
of different dimensions on its components), its odd rational
cohomology vanishes, and
\[
          c_\mu(a)=\dim_{\mathbb Q}H^*(X_{\mu,a},\mathbb Q).
                                                        \tag{P-F53c}
\]
Each $H^j(X_{\mu,a},\mathbb Q)$ is a pure Hodge structure of
weight $j$. The statement concerns purity separately in each
cohomological degree, and does not identify all the cohomology
with one pure weight or assert that it is Tate.

\paragraph{The cited Hilbert-scheme formula and the proof.}
The smoothness and projectivity of $S^{[m]}$ are the
surface Hilbert-scheme theorem recalled in G\"ottsche,
\emph{Hilbert Schemes of Points on Surfaces}, Section~1,
printed p.~484. His Betti-number formula is equation~(2.1),
printed p.~485. In his sign convention the Poincar\'e
polynomial is $p(Y,z)=\sum_j(-1)^jb_j(Y)z^j$.
For a K3 surface the Betti numbers are $(1,0,22,0,1)$,
and the formula specializes to
\[
 \sum_{m\ge0}p(S^{[m]},z)q^m
 =\prod_{r\ge1}
   \frac1{(1-z^{2r-2}q^r)
           (1-z^{2r}q^r)^{22}
           (1-z^{2r+2}q^r)}.
                                                        \tag{P-F53d}
\]
All exponents of $z$ on the right are even. Comparing
coefficients of $q^m z^{2j+1}$ gives
$b_{2j+1}(S^{[m]})=0$.
At $z=1$ the signed polynomial therefore equals the
total cohomology dimension, and
\[
 \sum_{m\ge0}\dim H^*(S^{[m]},\mathbb Q)q^m
                   =\prod_{r\ge1}(1-q^r)^{-24}.
                                                        \tag{P-F53e}
\]
The same Euler-number specialization is explicitly
Corollary~2.4, printed p.~62, of G\"ottsche,
\emph{Invariants of Moduli Spaces and Modular Forms};
the K3 value $e(S)=24$ is stated there on p.~66.

For any fixed $a$, a summand in (P-F53b) requires
$Q(x)\le a+1$. Positive definiteness of $K$ makes that
set of coset vectors finite. For each such vector the
integer $m$ is unique. Hence (P-F53b) is a finite
disjoint union of smooth projective schemes. Finite
disjoint unions preserve smoothness and projectivity;
their rational cohomology is the direct sum of the
component cohomologies. This proves odd-degree
vanishing and the asserted degreewise purity, using
the usual pure Hodge structure on smooth projective
cohomology.

Finally, multiplying (P-F53e) by
$q^{-1}\sum_{x\in K+\mu}q^{Q(x)}$ and taking the
coefficient of $q^a$ gives precisely the direct sum
in (P-F53b), so proves (P-F53c). All coefficient sums
are finite; there is no interchange of conditionally
convergent series. This completes the proof of the
coefficient model.\hfill$\square$

\paragraph{The completed lambda-ring identity.}
This construction also gives an exact, limited lambda-ring
dictionary for the inverse-discriminant operation. Let
$V=H^*(S,\mathbb Q)$ after forgetting its cohomological
grading, so $\dim V=24$. In the completed Grothendieck
ring of finite-dimensional rational vector spaces, put
$\sigma_t(V)=\sum_{j\ge0}[\operatorname{Sym}^jV]t^j$.
Then
\[
 \sum_{m\ge0}[H^*(S^{[m]},\mathbb Q)]q^m
   =\prod_{r\ge1}\sigma_{q^r}(V)
   =\operatorname{Sym}\!\left(V\otimes q\mathbb Q[q]\right)
                                                        \tag{P-F53f}
\]
as classes graded by the exponent of $q$. The last term
denotes the graded class of the symmetric algebra, with
one copy of $V$ in each positive degree. To prove the
identity, the generating function of the dimensions
of symmetric powers of a 24-dimensional vector space
is $(1-t)^{-24}$: a monomial basis counts weak
compositions into 24 parts. Its product over $r$
is (P-F53e). Dimension identifies the Grothendieck
ring of finite-dimensional rational vector spaces
with $\mathbb Z$, so equality of these dimensions
proves the stated class equality in every grade.
Every grade receives only finitely many contributions.
Thus no infinite-dimensional object is being inserted
as one element of the uncompleted ring.

Equation (P-F53f) uses the symmetric operations
$\sigma_t=\lambda_{-t}^{-1}$ in the sense explained in
the Grothendieck plethysm bridge. It is an equality
after forgetting the cohomological grading. This
elementary proof does not assert a canonical
cohomological-ring isomorphism or compatibility with
the nonstandard algebra's multiplication. The much
stronger geometric Heisenberg construction is not
needed for this statement.

\paragraph{The exact principal coefficients and what has been identified.}
The previously checked counts for the $E_6$ cosets
make the nonpositive part completely geometric:
\[
 X_{0,-1}=\operatorname{Spec}\mathbb C,\qquad
 X_{0,0}=S\ \coprod\ 72\operatorname{Spec}\mathbb C,
 \qquad
 X_{\pm,-1/3}=27\operatorname{Spec}\mathbb C.
                                                        \tag{P-F53g}
\]
Indeed the zero coset has one vector of norm zero
and 72 vectors with $Q=1$, while each other coset
has 27 vectors of minimal $Q=2/3$, as proved in
P-F46. Their total dimensions are $1,96,27,27$.
In particular the weight $96/2=48$ of the identified
Borcherds output can be read as half the total
cohomology dimension of $X_{0,0}$.

This supplies smooth projective coefficient spaces
for the \emph{entire particular input} of P-F46,
and expresses its inverse-discriminant factor by
completed symmetric operations on an actual
surface cohomology. Together with P-F51, both
factors now have specified geometric sources.
It does not turn the regularized multiplicative
Borcherds lift into the supplied iterated definite
unitary lift, nor give a cohomological map from
these spaces onto the cubic multiplicity spaces.
The coefficient-growth obstruction of P-F46 still
rules out identifying the whole input with all
$T^{(5)}$ grades by a positive affine reindexing.
The pure Hodge pieces here also need not be Tate:
already $S\subset X_{0,0}$ has its nonzero
$(2,0)$ part. This example therefore adds an
actual Hilbert-scheme and lambda-operation source
to the remaining operation-comparison problem,
without claiming that purity alone solves it.

Primary sources for the cited formulae:
\url{https://arxiv.org/pdf/math/0304302},
Section~1 p.~484 and equation~(2.1) p.~485;
\url{https://rendiconti.dmi.units.it/volumi/41/06.pdf},
Corollary~2.4 p.~62 and the K3 specialization p.~66.


\subsection{An explicit compact-averaging bound for actual cubic support}
\label{sec:cubic-reynolds-support}

\paragraph{P-F54 (all sufficiently large actual degrees, with an explicit polynomial bound).}
For every integer $n\ge33$, there is an explicitly specified
cubic tensor $F_n\in\operatorname{Sym}^3\mathbb C^n$ such that
\[
 \begin{gathered}
 d\ge(10n)^{200},\qquad n\mid3d
 \quad\Longrightarrow\quad
 \text{some homogeneous degree-$d$ $SL_n$ invariant is nonzero at }F_n.
 \end{gathered}                                                   \tag{P-F54a}
\]
In particular, with $g=\gcd(n,3)$ and the actual cubic stretching
function $f_n$ of P-F31,
\[
 f_n(s)>0\qquad\text{for all integers }s\ge
             \left\lceil\frac{g(10n)^{200}}n\right\rceil .        \tag{P-F54b}
\]
The exponent is deliberately coarse. This is an unconditional
bound for actual support, including all allowed residue classes.
It is not a bound for the strict saturation or positivity index:
it does not evaluate a constituent polynomial away from its own
residue progression. The proof gives a compact averaging formula
for a nonzero invariant, not a claim of a polynomial-size
evaluation circuit or a positive counting formula.

\paragraph{1. A sparse exactly balanced tensor.}
Use the Hermitian tensor-product norm on
$W=\operatorname{Sym}^3\mathbb C^n\subset(\mathbb C^n)^{\otimes3}$.
The vectors $e_i^{\otimes3}$ have norm one. For a three-element
set $T$, let $v_T$ be the sum of its six ordered tensor words
divided by $\sqrt6$; the $v_T$ form an orthonormal family.
All indices in the following two lists are in $\mathbb Z/n\mathbb Z$:
\[
 E_i=\{i,i+1,i+3\},\qquad E_{n+i}=\{i,i+4,i+16\}
                      \quad(0\le i<n).                         \tag{P-F54c}
\]
These $2n$ triples are distinct, every vertex occurs in exactly six
of them, and two distinct triples intersect in at most one vertex.
Indeed their possible ordered pair differences are respectively
$\pm1,\pm2,\pm3$ and $\pm4,\pm12,\pm16$. These twelve
residues are distinct when $n\ge33$, and each difference occurs
in exactly one ordered pair within the two base triples. This
also proves that distinct translates or different lists cannot
share two vertices.

Set
\[
 \begin{gathered}
 L=100n^2,\qquad \varepsilon=n^{-20},\qquad
 \eta_j=\exp(2\pi i j/L),\quad0\le j<2n,\\
 F_0=\sum_{i=0}^{n-1}e_i^{\otimes3},\qquad
 H=\sum_{j=0}^{2n-1}\eta_jv_{E_j},\qquad
 F=F_n=F_0+\varepsilon H,\qquad
 N_F=\|F\|^2=n(1+2\varepsilon^2)\le2n.
 \end{gathered}                                                   \tag{P-F54d}
\]
Thus every coefficient is explicitly algebraic. The one-factor
partial trace of $F_0F_0^*$ is the identity matrix. That of
$v_Tv_T^*$ is $\frac13\sum_{i\in T}e_ie_i^*$.
Every cross partial trace between distinct triples is zero,
because a nonzero cross term would require two common tensor
indices. The cross partial traces between $F_0$ and $H$ are
also zero. Consequently the one-factor partial trace of $FF^*$
is exactly $(1+2\varepsilon^2)I_n$.

Let $\rho$ denote the unitary action on $W$. For a traceless
skew-Hermitian matrix $X$, write $A_X=d\rho(X)$ and
$r=\|X\|_{\mathrm F}$, where $\|\cdot\|_{\mathrm F}$ is the
matrix Frobenius norm. The partial-trace identity gives
\[
       \langle F,A_XF\rangle=0,
       \qquad \|A_X\|_{\mathrm{op}}\le3r.                    \tag{P-F54e}
\]
Here the inner product is linear in its second argument.
Directly in the orthonormal cubic monomial basis,
\[
 \|A_XF_0\|^2
   =3\sum_{i\ne j}|X_{ij}|^2+9\sum_i|X_{ii}|^2
   \ge3r^2.
\]
Since $\|H\|=\sqrt{2n}$, it follows that
\[
 \|A_XF\|\ge(\sqrt3-3\varepsilon\sqrt{2n})r\ge r
                        \qquad(n\ge33).                      \tag{P-F54f}
\]
Thus both the exact zero expectation and a quantitative
nondegeneracy bound have been proved, without a balancing
existence theorem.

\paragraph{2. Separation of the finite Fermat symmetries.}
Let $\mathcal K$ be the finite subgroup of $SU(n)$ consisting of
the monomial matrices
$h=\alpha\operatorname{diag}(\zeta_i)P$, where
$\zeta_i^3=1$, $P$ is a permutation matrix, and $\det h=1$.
Such an $h$ satisfies $\rho(h)F_0=\alpha^3F_0$.
Put $k_h=\alpha^{-3}\rho(h)$, a unitary operator fixing $F_0$.
The cube of the chosen $\alpha$ is independent of the common
cube-root ambiguity in this expression.
We claim that
\[
       h\in\mathcal K\setminus Z(SU(n))
       \quad\Longrightarrow\quad
                  \|H-k_hH\|\ge L^{-1}.                     \tag{P-F54g}
\]
If the permutation of variables takes an edge outside the
support in (P-F54c), the difference has a coefficient of
absolute value one. If it preserves the edge set but moves
one edge to another, the corresponding difference of
coefficients is $\eta_i-\zeta\eta_j$ with $i\ne j$ and
$\zeta^3=1$. For $\zeta=1$ its absolute value is at least
$4/L$, by $\sin x\ge2x/\pi$ on $[0,\pi/2]$.
For $\zeta\ne1$, it is at least one: all ratios
$\eta_i/\eta_j$ have arguments of absolute value less than
$4\pi n/L<\pi/6$. These bounds imply (P-F54g) in both cases.

If every edge is fixed individually, every vertex is fixed.
Indeed the indices of the first-list edges incident to vertex
$j$ are $\{j,j-1,j-3\}$. These sets are distinct for different
$j$. A nonzero translation preserving a three-element set
would act on it as a three-cycle, so its nonzero pair
differences would have only two values; the set
$\{0,-1,-3\}$ instead has the six distinct differences
$\pm1,\pm2,\pm3$ when $n\ge33$.
Thus the remaining case is diagonal. Write
$\zeta_i=\exp(2\pi i b_i/3)$, $b_i\in\mathbb F_3$.
If one edge product is nontrivial, the coefficient difference
has absolute value $\sqrt3$. Otherwise, for the cyclic shift
operator $S$ on the vector $(b_i)$,
\[
       (1+S+S^3)b=0,
       \qquad(1+S^4+S^{16})b=0.                              \tag{P-F54h}
\]
In $\mathbb F_3[z]$,
$1+z+z^3=(z-1)(z^2+z+2)$.
The quadratic is irreducible and its roots have fourth
power $-1$ and order eight. At either root,
$1+z^4+z^{16}=1\ne0$; at $z=1$ it vanishes.
The greatest common divisor of the two displayed polynomials
is therefore $z-1$. Bezout's polynomial identity applied to
$S$ gives $(S-1)b=0$, so all $b_i$ agree.
Then $h$ is scalar. This proves (P-F54g) in every case,
including when $3$ divides $n$.

\paragraph{3. A quantitative inverse estimate at Fermat.}
The following elementary estimate will keep the subsequent
global argument from assuming that genericity supplies a
uniform gap. Suppose $U\in SU(n)$, $|\lambda|=1$, and
\[
             \rho_0=\|F_0-\lambda\rho(U)F_0\|
                            \le\frac1{100\sqrt n}.
\]
Then there exist $h\in\mathcal K$, $X\in\mathfrak{su}(n)$,
and $t\in\mathbb R$ such that
\[
 U=h\exp X,\qquad \lambda\rho(h)=e^{it}k_h,
 \qquad \|X\|_{\mathrm F}\le24\rho_0,
 \qquad |t|\le24\rho_0/\sqrt n.                              \tag{P-F54i}
\]
To prove it, write $u_{ij}$ for the entries of $U$.
Since $\langle F_0,\rho(U)F_0\rangle=\sum_{i,j}u_{ij}^3$,
\[
 0\le n-\sum_{i,j}|u_{ij}|^3\le\rho_0^2/2.
\]
In each row let $a_i=\max_j|u_{ij}|$. The row sum of the
squares is one, so its sum of cubes is at most $a_i$.
It follows that $\sum_i(1-a_i)\le\rho_0^2/2$.
Each $a_i>1/\sqrt2$, so the maximizing positions lie in
distinct columns. They define a monomial unitary matrix
$M$, with the same phases $m_i$ at those positions, and
\[
                    \|U-M\|_{\mathrm F}^2
                       =2\sum_i(1-a_i)\le\rho_0^2.
\]
The omitted off-position cubes have total absolute value at
most $\sum_i(1-a_i^2)\le2\sum_i(1-a_i)$.
Also $\sum_i(1-a_i^3)\le3\sum_i(1-a_i)$.
Comparing the dominant-position phases with the entire cubic
sum therefore gives
\[
       \sum_i\bigl(1-\operatorname{Re}(\lambda m_i^3)\bigr)
                   \le3\rho_0^2,
       \qquad \sum_i|1-\lambda m_i^3|^2\le6\rho_0^2.
\]
Choose $\alpha^3=\lambda^{-1}$ and choose cube roots
$\zeta_i$ nearest to $m_i/\alpha$ on the unit circle.
For an angle $|\theta|\le\pi/3$,
$|e^{i\theta}-1|\le|e^{3i\theta}-1|$; hence the monomial
matrix $M'=\alpha\operatorname{diag}(\zeta_i)P$ satisfies
\[
                    \|U-M'\|_{\mathrm F}<4\rho_0.
\]
Let the principal eigenangles of $U^*M'$ be $\theta_j$.
The preceding bound and $|\theta|\le2|e^{i\theta}-1|$
on the relevant small arc imply
$|\sum_j\theta_j|\le8\sqrt n\rho_0$.
Set $\tau=n^{-1}\sum_j\theta_j$ and $h=e^{-i\tau}M'$.
Then $\det h=1$, $h\in\mathcal K$, and
$\|U-h\|_{\mathrm F}\le12\rho_0$.
The principal logarithm of $h^*U$ has Frobenius norm at
most $24\rho_0$. Its trace is a multiple of $2\pi i$,
but its absolute trace is at most
$24\sqrt n\rho_0<2\pi$, so this trace is zero.
This logarithm is the required $X$.
Finally $\lambda\alpha^3e^{-3i\tau}=e^{-3i\tau}$,
so $t=-3\tau$ gives the stated phase bound and proves
(P-F54i).

\paragraph{4. A global overlap gap.}
Put
\[
 z(U)=\frac{\langle F,\rho(U)F\rangle}{N_F},
 \qquad R=\frac{\varepsilon}{16L},
 \qquad \Delta=\frac{R^2}{16n}.
                                                               \tag{P-F54j}
\]
For every $U$ outside the union of the scalar-centre neighborhoods
\[
  \bigcup_{h\in Z(SU(n))}
       \{h\exp X:X\in\mathfrak{su}(n),\ \|X\|_{\mathrm F}<R\},
\]
we claim
\[
                            |z(U)|\le1-\Delta.               \tag{P-F54k}
\]
Choose a phase $\lambda$ minimizing
$r_F=\|F-\lambda\rho(U)F\|$.
If $r_F\ge\varepsilon/(4L)$, then $r_F\ge R/2$ already.
Otherwise
\[
 \rho_0=\|F_0-\lambda\rho(U)F_0\|
       \le r_F+2\varepsilon\sqrt{2n}
       \le5\varepsilon\sqrt n\le\frac1{100\sqrt n}.
\]
Use (P-F54i). Writing $r=\|X\|_{\mathrm F}$, we have
$r\le120\varepsilon\sqrt n$ and $|t|\le120\varepsilon$.
Project $F-\lambda\rho(U)F$ orthogonally to the squarefree
cubic subspace. The vectors $F_0$ and $A_XF_0$ have no
squarefree component, and $k_h$ preserves this subspace.
For a skew-Hermitian operator $A$,
$\|e^A-I-A\|\le\|A\|^2/2$ and
$\|e^A-I\|\le\|A\|$, by their integral remainder formulas.
Thus that projection differs from
$\varepsilon(H-k_hH)$ in norm by at most
\[
 \frac92\sqrt n\,r^2
    +\varepsilon\sqrt{2n}(|t|+3r)
 \le10^5n^{3/2}\varepsilon^2
 \le\frac{\varepsilon}{4L}.                                 \tag{P-F54l}
\]
The last inequality follows from
$4\cdot10^7 n^{7/2}\varepsilon\le1$ for $n\ge33$.
If $h$ were noncentral, (P-F54g) would give
$r_F\ge3\varepsilon/(4L)$, a contradiction.
Hence $h$ is central.

Now $k_h=I$ and the residual is
$\|F-e^{A_X+it}F\|$.
By (P-F54e)--(P-F54f),
\[
       \|(A_X+it)F\|^2
            =\|A_XF\|^2+t^2N_F\ge r^2+nt^2.
\]
For skew-Hermitian $B$, the stronger integral remainder bound
$\|(e^B-I-B)F\|\le\frac12\|B\|\,\|BF\|$ applies.
Here $\|A_X+it\|\le480\varepsilon\sqrt n<1$, so
$r_F\ge\frac12\|(A_X+it)F\|\ge r/2$.
As $U$ lies outside the displayed neighborhoods, $r\ge R$,
and therefore again $r_F\ge R/2$.
Finally
$r_F^2=2N_F(1-|z(U)|)$ and $N_F\le2n$ give
(P-F54k). All smallness inequalities used here follow directly
from $n\ge33$ and $\varepsilon=n^{-20}$; no unquantified
neighborhood or compactness gap has been used.

\paragraph{5. Local modulus and phase estimates.}
For $X\in\mathfrak{su}(n)$ with $r=\|X\|_{\mathrm F}\le R$,
Taylor's formula with the same unitary remainders gives
\[
 |z(e^X)-1|\le\tfrac92 r^2,
 \qquad |\operatorname{Im}z(e^X)|\le\tfrac92 r^3,
 \qquad |\arg z(e^X)|\le10r^3.                               \tag{P-F54m}
\]
The linear term is zero by (P-F54e); the quadratic term
$\langle F,A_X^2F\rangle=-\|A_XF\|^2$ is real.
The second and third Taylor remainder bounds are respectively
$\frac12\|A_X\|^2$ and $\frac16\|A_X\|^3$ after dividing
by $N_F$. Since the real part is at least $1/2$, these also
give the argument bound.

Taking $t=-\arg z(e^X)$ and using the local residual argument
from step~4 gives
\[
                         |z(e^X)|\le1-\frac{r^2}{16n}.
                                                               \tag{P-F54n}
\]
Indeed $\|A_X+it\|\le3r+10r^3<1$; the exact zero
expectation removes the mixed term in its norm, and
(P-F54f) gives residual at least $r/2$.

\paragraph{6. A Haar-volume lower bound without a dimension-free shortcut.}
Normalize Haar measure $\mu$ on $SU(n)$ to total mass one.
For $0<a\le1/(10\sqrt n)$,
\[
 \mu\{e^X:X\in\mathfrak{su}(n),\ \|X\|_{\mathrm F}<a\}
                   \ge(1+8\sqrt n/a)^{-2n^2}.                \tag{P-F54o}
\]
Here is a direct covering proof. A maximal collection of
points of $SU(n)$ separated by Frobenius distance $b$ has
at most $(1+2\sqrt n/b)^{2n^2}$ members: the disjoint
ambient Euclidean balls of radius $b/2$ around them lie
in a ball of radius $\sqrt n+b/2$ in
$\mathbb R^{2n^2}$. The group balls of radius $b$ cover,
and their Haar measures agree by unitary invariance of
Frobenius distance. The identity ball therefore has
measure at least the reciprocal of this number.
Take $b=a/4$, giving the bound
$(1+8\sqrt n/a)^{2n^2}$ exactly.
Every unitary in this identity ball has a principal logarithm
of norm at most $2b=a/2<a$. Its trace is a multiple of $2\pi i$
and has absolute value at most $\sqrt n a<2\pi$, so the
logarithm belongs to $\mathfrak{su}(n)$.
The closed identity ball of radius $b$ is thus contained
in the open logarithm ball of radius $a$, proving (P-F54o)
without an assertion about boundary measure.

\paragraph{7. The complete compact-average comparison.}
Let $d$ be an integer with $d\ge(10n)^{200}$ and $n\mid3d$.
Set
\[
             b_d=d^{-2/5},\qquad a_d=\frac1{10\sqrt d}.
\]
We have $a_d\le b_d\le R$.
For $r\le b_d$, (P-F54m) gives
$d|\arg z(e^X)|\le10d^{-1/5}\le1$.
Thus the real part of $z(e^X)^d$ is nonnegative on this
whole neighborhood. On the smaller neighborhood $r\le a_d$,
\[
 |z(e^X)|\ge1-\frac1{20d},
 \qquad \operatorname{Re}z(e^X)^d\ge\frac14.                \tag{P-F54p}
\]
For the latter use $(1-1/(20d))^d\ge19/20$ and
$\cos(d\arg z)\ge1/2$.

A scalar center element $h=\omega I_n$, $\omega^n=1$,
acts on $W$ by $\omega^3$. Hence
$z(he^X)^d=\omega^{3d}z(e^X)^d=z(e^X)^d$.
All central neighborhoods therefore have the same positive
phase; none are discarded by an absolute-value argument.
Outside their union of radius $b_d$, either (P-F54k) or
(P-F54n) applies and gives
\[
       |z(U)|^d\le
           \exp\left(-\frac{d^{1/5}}{16n}\right).             \tag{P-F54q}
\]
The neighborhoods can be taken disjoint: distinct scalar
centers have Frobenius distance at least $4/\sqrt n$,
whereas $2R<4/\sqrt n$. Combining (P-F54o)--(P-F54q)
and retaining just the smaller identity neighborhood yields
\[
 \operatorname{Re}\int_{SU(n)}z(U)^d\,d\mu(U)
 \ge\frac14(1+80\sqrt{nd})^{-2n^2}
        -\exp\left(-\frac{d^{1/5}}{16n}\right)>0.             \tag{P-F54r}
\]
For clarity, the final strict inequality follows from explicit
elementary bounds. Since $nd\ge1$,
$1+80\sqrt{nd}\le100nd$.
When $d\ge(10n)^{200}$,
$\log(100n)\le\frac1{100}\log d$, so
\[
 \log4+2n^2\log(1+80\sqrt{nd})<5n^2\log d.
\]
Also $d^{1/5}>80n^3\log d$ on this range.
To check the latter at the left endpoint, substitute
$d=(10n)^{200}$ and use $\log(10n)\le10n$;
$10^{40}n^{40}>160000n^4$ suffices.
The function $d^{1/5}/\log d$ is increasing for $d>e^5$,
so the inequality persists for every larger $d$.
This proves the strict positivity in (P-F54r), with the
full dependence on $n$ in the neighborhood volumes included.

\paragraph{8. From the average to an actual polynomial invariant.}
The compact averaging operation is the usual Reynolds
construction; an exact modern primary locator is Aryaman
Maithani, \emph{Splitting the difference: Computations of the
Reynolds operator in classical invariant theory},
\emph{Bulletin of the London Mathematical Society} 58 (2026),
e70175, Theorem~5.6. Theorem~6.1(a) gives the density of
$SU(n)$ in $SL_n(\mathbb C)$:
\url{https://doi.org/10.1112/blms.70175}.
The following short argument records the precise tensor-norm
identity needed for the quantitative estimate above; no degree
bound is being attributed to those general statements.

The Haar average $P_d=\int_{SU(n)}\rho(U)^{\otimes d}d\mu(U)$
is the orthogonal projection onto the invariant subspace:
left invariance makes its image invariant, it fixes every
invariant vector, and inversion invariance makes it
self-adjoint. Therefore
\[
 \int_{SU(n)}z(U)^d\,d\mu(U)
             =\frac{\|P_dF^{\otimes d}\|^2}{N_F^d}.           \tag{P-F54s}
\]
The tensor $F^{\otimes d}$ and its projection lie in
$\operatorname{Sym}^dW$.
By (P-F54r) the projection is nonzero.
The polynomial
$I(v)=\langle P_dF^{\otimes d},v^{\otimes d}\rangle$
is homogeneous of degree $d$, invariant under $SU(n)$,
and satisfies $I(F)=\|P_dF^{\otimes d}\|^2>0$.
It is invariant under $SL_n(\mathbb C)$ as well:
the Lie algebra $\mathfrak{su}(n)$ annihilates it, its
complex span is $\mathfrak{sl}_n(\mathbb C)$, and the
elementary unipotent one-parameter subgroups generate
$SL_n(\mathbb C)$. This proves (P-F54a), and substituting
$d=ns/g$ proves (P-F54b).

The center condition is also necessary: the scalar center
acts on degree-$d$ cubic tensors by $\omega^{3d}$.
Thus this one explicitly specified point has every
sufficiently large degree permitted by the minimal center
condition. No assertion about $n<33$ is needed for the
uniform asymptotic statement, and no generic stabilizer
theorem, GRH, numerical semigroup conductor theorem, or
unquantified stationary-phase limit enters the proof.
The distinction from strict constituent positivity remains
exactly the one stated at the beginning.\hfill$\square$

\paragraph{An exact illustration of the remaining index distinction.}
For any integer $M\ge2$, consider the quasipolynomial
\[
 q_M(s)=M\lfloor s/M\rfloor+1,
 \qquad Q_i(X)=X-i+1\quad(0\le i<M).                       \tag{P-F54t}
\]
It is positive at every actual integer $s\ge0$, and
$Q_i(s)=q_M(s)$ when $s\equiv i\pmod M$.
The bounded periodic function $q_M(s)-s=1-(s\bmod M)$
has minimal period $M$, so these are its constituents.
A shift by a nonnegative integer $c$ permutes the
constituents $Q_i(X+c)$; their smallest constant
coefficient is $c-M+2$. Hence all coefficients are
nonnegative exactly when $c\ge M-2$.
For strict saturation, each constituent is increasing
in $X$, and its smallest value at a positive integer
is $Q_{M-1}(1+c)=c-M+3$. This is strictly positive
exactly when $c\ge M-2$. Thus the literal conventions
of GCT VI, Definitions~1.2.2 and 1.2.4, printed
pp.~13--14, give
\[
                   p(q_M)=s(q_M)=M-2.                    \tag{P-F54u}
\]
See \url{https://arxiv.org/pdf/0704.0229}.
This sequence is not asserted to be a plethysm sequence.
It only demonstrates why even positivity at every actual
nonnegative integer does not by itself bound the two
strict constituent indices.


\subsection{The same theta input and the cubic-fourfold discriminant}

The preceding theta construction also has an exact interpretation for
ordinary cubic fourfolds. The character changes the answer: on the
orthogonal domain the weight-$48$ form must be squared before applying
the invariant-ring isomorphism. This is a fixed-dimensional
compatibility of known constructions, with an explicit discriminant
identity; no historical priority is claimed.

\paragraph{Primary geometric inputs and conventions.}
Let $V=\mathbb C^6$, $\mathscr S_6=\operatorname{Sym}^3V^*$,
and write $\Delta_6$ for its discriminant. Put
\[
 M=U^{\oplus2}\oplus E_8(-1)^{\oplus2}\oplus A_2(-1),
 \qquad K=E_6.
\]
Thus $M$ has signature $(2,20)$ and embeds primitively in
$L=II_{2,26}$ with complement $K(-1)$, exactly as in P-F46.
Let $D^+$ be one connected component of its type-IV domain and
$\Gamma^+=\widetilde O^+(M)$ the group preserving that component
and acting trivially on $A_M=M^\vee/M$.

Looijenga, \emph{The period map for cubic fourfolds}, Section~1
printed p.~3 and Section~2 p.~7, identifies the group for the
cubic-fourfold period map with this discriminant kernel, allowing
both components before restriction to $D^+$. Theorem~3.1,
printed p.~11, gives an isomorphism of algebras
\[
 \mathcal P_6:\bigoplus_{k\ge0}
 H^0(\dot D^+,\mathcal A(k))^{\Gamma^+}
       \ \xrightarrow{\ \sim\ }
 \mathbb C[\mathscr S_6]^{SL(V)},
 \qquad \deg\mathcal P_6(h)=2\operatorname{wt}(h).
 \tag{P-F55a}
\]
Here $\dot D^+$ and its section spaces are precisely the
arrangement-complement spaces in that theorem. In particular this
notation retains its automorphic section convention. The theorem
identifies the line bundles over the locus of cubics with at most
ADE singularities. Laza, \emph{The moduli space of cubic fourfolds
via the period map}, Theorem~1.1, printed p.~1, identifies this
image with the complement of $\mathcal H_\infty$, and the smooth
image with the complement of
$\mathcal H_\infty\cup\mathcal H_\Delta$.
The source URLs are
\url{https://arxiv.org/pdf/0705.0951} and
\url{https://arxiv.org/pdf/0705.0949}.

Let $g=\Theta_K/\delta$ denote the vector-valued form, with the
discriminant-coset identification and Weil convention of (P-F46a).
Write $\Psi=\Psi_M(g)$. By the exact theta contraction of P-F46,
this is, up to a nonzero scalar, the $E_6$ quasi-pullback of
$\Phi_{12}$. Laza, Notation~6.4 printed p.~24 and
Proposition~6.6 p.~25, uses the same primitive embedding and
quasi-pullback: its weight is $48$ and its orders on
$\mathcal H_\Delta,\mathcal H_\infty$ are $1,27$ respectively.

\begin{corollary}[P-F55: an exact second invariant image and its cyclic compatibility]
\label{cor:fourfold-theta-discriminant}
The multiplier of $\Psi$ on $\Gamma^+$ is $\det$, and
\[
          \mathcal P_6(\Psi^2)=c_6\Delta_6,
                    \qquad c_6\in\mathbb C^*.
 \tag{P-F55b}
\]
The square has automorphic weight $96$, ordinary coefficient
degree $192$, and canonical stretching degree $96$ in the
actual cubic plethysm algebra
\[
       T^{(6)}=\bigoplus_{s\ge0}
        \left(\operatorname{Sym}^{2s}
                    (\operatorname{Sym}^3\mathbb C^6)\right)^{SL_6}.
 \tag{P-F55c}
\]
Its $GL_6$ change-of-variables character is $\det^{96}$.
The same identification respects all products in
$\mathbb C[\Psi^2]\cong\mathbb C[\Delta_6]$.

For a cubic $f$ in five variables and an additional variable $w$,
there is a nonzero constant $C$, depending only on the two
discriminant normalizations, such that the following polynomial
identity holds in $a$ and all coefficients of $f$:
\[
           \Delta_6(a w^3-f)=C a^{32}\Delta_5(f)^2.
 \tag{P-F55d}
\]
With the discriminant normalization of Bus\'e--Jouanolou
Definition~4.6, the constant is $C=3^{33}$.
At $a=1$ this agrees, up to scalar normalization, with restriction
of the named orthogonal form to the cubic-threefold ball:
\[
        \Psi^2|_{\mathbb B}\doteq\phi^2,
        \qquad \mathcal P_5(\phi^2)\doteq\Delta_5^2.
 \tag{P-F55e}
\]
Here $\phi$ and $\mathcal P_5$ are those of P-F41, and
$\doteq$ denotes equality up to a nonzero scalar. This last
statement concerns the named forms and their powers; it does not
assert that the affine substitution $f\mapsto w^3-f$ is a graded
homomorphism between the full invariant rings.
\end{corollary}

\begin{proof}
We first check the character on the exact group of (P-F55a).
Gritsenko, \emph{24 faces of the Borcherds modular form
$\Phi_{12}$}, printed p.~2, states
\[
             \Phi_{12}\in M_{12}(O^+(II_{2,26}),\det).
\]
The source is \url{https://arxiv.org/pdf/1203.6503}.
The gluing subgroup for
$M\oplus K(-1)\subset L$ is the graph of the discriminant
isometry used in P-F46. For $\gamma\in\Gamma^+$,
$\gamma\oplus\mathrm{id}_K$ acts trivially on this graph, hence
extends to an isometry $\widetilde\gamma$ of $L$.
It preserves the chosen domain component, since it preserves
the positive-plane component containing $D^+$, and
$\det\widetilde\gamma=\det\gamma$.
In the quasi-pullback one divides $\Phi_{12}$ by the product
of one linear form for each pair of opposite roots in $K(-1)$,
and then restricts to $D^+$. There are $72/2=36$ such factors.
The extension fixes each of these linear forms, not merely
their collection, because its restriction to $K(-1)$ is the
identity. Thus the quotient and its nonzero restriction have
character $\det\gamma$. This proves the multiplier assertion
and the trivial character of $\Psi^2$.

For clarity, passing to one domain component changes no invariant
section ring in Looijenga's theorem. His $D$ has two components
interchanged by an element of the discriminant kernel
(Section~2, p.~7). A section invariant under the stabilizer
of $D^+$ extends uniquely to the other component by such an
element; independence of the chosen element follows from
invariance under that stabilizer. This identifies the ring in
the cited theorem with the ring in (P-F55a).

The form $\Psi^2$ is holomorphic on the full orthogonal domain,
so is an allowed section in the arrangement-complement ring.
The isomorphism (P-F55a) sends it to a nonzero homogeneous
invariant $I$ of coefficient degree $2\cdot96=192$.
We next identify this polynomial without assuming a divisor
multiplicity convention on an orbifold quotient.

The discriminant in $\mathscr S_6$, whose dimension is $56$,
has a dense locus of cubics with a unique ordinary node.
Here is a direct verification of the density. The incidence
\[
  \mathscr I=\{(F,p)\in\mathscr S_6\times\mathbb P^5:
                                  dF(p)=0\}
\]
is a rank-$50$ vector bundle over $\mathbb P^5$. At
$p=[1:0:\cdots:0]$ its six independent conditions remove
the coefficients of $x_0^3$ and $x_0^2x_i$, $1\le i\le5$.
It is therefore irreducible of dimension $55$, and projection
to $\mathscr S_6$ is proper. Requiring two distinct singular
points imposes twelve independent conditions: after moving
them to the first two coordinate points, the two lists of
six degree-three monomials are disjoint. The corresponding
ordered-pair incidence has dimension $10+(56-12)=54$.
The condition that the quadratic part at the marked point be
degenerate is a proper closed condition on $\mathscr I$.
Properness follows, for example, from
\[
        F_0=x_0\sum_{i=1}^5x_i^2+\sum_{i=1}^5x_i^3.
 \tag{P-F55f}
\]
If $x_0=0$, its gradient equations force all $x_i=0$.
Otherwise set $x_0=1$. Each $x_i$ is then $0$ or $-2/3$,
and $\sum_i x_i^2=0$ forces all to be zero. Hence its only
singular point is $[1:0:\cdots:0]$, with nondegenerate affine
quadratic part $\sum_i x_i^2$.
The degenerate-quadratic locus and the closure of the
two-point locus have images of dimension at most $54$.
Their complement in the discriminant is the asserted dense
ordinary-node locus. In particular the incidence image has
dimension $55$, also following from the zero-dimensional
fiber just exhibited.

These nodal cubics have ADE singularities and are stable by
Looijenga's Theorem~3.2, p.~11. Their periods lie on
$\mathcal H_\Delta$ and outside $\mathcal H_\infty$ by
the smooth and ADE image descriptions above.
Since $\Psi$ vanishes on $\mathcal H_\Delta$, the
line-bundle identification in Theorem~3.1 shows that $I$
vanishes on this dense nodal locus.

Bus\'e--Jouanolou, \emph{A computational approach to the
discriminant of homogeneous polynomials}, Definition~4.6 and
Proposition~4.7(i), printed p.~40, and Theorem~4.26, p.~53,
show that over $\mathbb C$ the cubic discriminant is a prime
polynomial cutting out this hypersurface, with coefficient
degree $6(3-1)^5=192$; see
\url{https://arxiv.org/pdf/1210.4697}.
Thus $\Delta_6$ divides $I$ in the polynomial coefficient
ring. Their degrees agree, so $I=c_6\Delta_6$, with
$c_6\ne0$. This proves (P-F55b).
The same source, Proposition~4.13, pp.~42--43, gives
\[
          \Delta_6(F\circ h)=\det(h)^{3\cdot2^5}\Delta_6(F),
                       \qquad h\in GL_6,
 \tag{P-F55g}
\]
which is the claimed character.
The center of $SL_6$ permits only even invariant coefficient
degrees, and the canonical regrading is $d_0=6/\gcd(6,3)=2$.
The actual Cartan identification in P-F19 therefore gives
(P-F55c) with degree $192/2=96$. The algebra isomorphism
respects multiplication. Since a nonconstant polynomial is
transcendental over $\mathbb C$ inside a polynomial ring,
the subalgebras generated by the two named elements are
polynomial algebras and are identified as asserted.

The cyclic discriminant formula is already known, and can
be applied directly. Bus\'e--Jouanolou, Example~4.10,
printed p.~41, gives in their normalization
\[
 \operatorname{Disc}(a w^d+h)
   =d^{(d-1)^{n-1}+(-1)^n}
       a^{(d-1)^{n-1}}\operatorname{Disc}(h)^{d-1},
 \tag{P-F55h}
\]
where $h$ is a degree-$d$ form in $n-1$ variables.
Substitute $n=6$, $d=3$ and $h=-f$.
Since $\deg\Delta_5=80$ is even, Proposition~4.7(i)
gives $\Delta_5(-f)=\Delta_5(f)$. Thus (P-F55h)
is exactly (P-F55d) with $C=3^{33}$ in that normalization;
arbitrary nonzero normalizations change only $C$.

Finally P-F41 and P-F46 identify the restriction of this very
quasi-pullback to the cyclic-cubic-threefold ball with $\phi$,
up to scalar, and identify its invariant image with $\Delta_5$.
Squaring these identities gives (P-F55e). On the algebraic
side (P-F55d) at $a=1$ gives the same squared discriminant.
The gradings are consistent: the restricted square has weight
$96$, five-variable coefficient degree $(5/3)96=160$, and
five-variable stretching degree $160/5=32$.
These differ from the six-variable coefficient and stretching
degrees $192,96$, which is why no equal-degree full-ring
substitution assertion has been made.
\end{proof}

\paragraph{Scope for the theta comparison.}
P-F51 realizes the definite-$U(3)$ theta factor of the same
input $g$ used here, and P-F53 models its coefficient dimensions.
P-F55 identifies another actual invariant output and the exact
cyclic restriction of that output. The operation still includes
division by $\delta$, regularized multiplicative lifting, and
the prescribed quasi-pullback or restriction. It supplies no
identification of these operations with the supplied Remark's
iterated totally definite unitary lifts, and no basis or positive
dimension formula for all graded pieces of either invariant
ring. In particular it gives no growing-dimension PH2/SH bound.


\subsection{Disposition of the main request}
The literal homogeneous PH1 and the old strong zero-index assertions
are refuted by the cited counterexamples. A shifted two-dimensional
model for the PH1 counterexample is proved in P-F1; the general
positive-rational-form transfer is P-F2. P-F11--P-F12 prove the
quantitative implication from revised PH3 to revised PH2, and under
PH3 bound the size of a shifted model in the expanded group parameters.
P-F31 refutes the weaker positive-form clause, so this implication
cannot settle PH2 from the printed PH3 premise. Uniform construction
and direct PH2 or saturation bounds remain separate targets. P-F13 makes reconstruction and positive-form verification
finite under a stated denominator bound. P-F14--P-F15 refute the
printed rank-only PH3, PH2 and saturation-index bounds for general
subgroup restriction. P-F16--P-F18 instead prove finite bounds for
fixed semisimple embeddings and fixed inner dimension in ordinary
plethysm. P-F19 proves that the optimal canonical envelope has an
exponential subsequence, refuting the stronger canonical clause of
ordinary plethysm PH3. P-F30--P-F31 prove actual exponential-order
Hilbert poles and thereby refute the numerical positive-form clause
as well. PH2 and the narrower saturation bounds remain unresolved;
P-F32 gives a direct exterior-duality reduction and a polynomial
lower bound on their indices, without settling the upper bounds.
P-F33 blocks a general period-and-reciprocity shortcut; P-F34 gives
a cubic-specific Fourier test with one explicit average-polynomial
estimate still missing; P-F35 computes the four-variable indices.
P-F36 proves that a direct nonstandard projector or character
realization must account for a nontrivial quadratic descent datum.
P-F39 computes the exact positive normalized trace and norm and
constructs an actual positive operator, without identifying its
coefficient models with the nonstandard modules. P-F52 realizes
that operator by an explicit central word in the actual algebra,
with a descended source-basis intertwiner and a positive tensor
form. The remaining comparison concerns the prescribed
normalization, products and geometric or theta coefficient models.
P-F37 and P-F41
supply an actual automorphic algebra and a named theta output:
the weight-$48$ form equals the degree-$80$ cubic discriminant.
P-F46 identifies its exact positive theta input and gluing,
while proving that direct Fourier-coefficient reindexing
cannot give the actual graded multiplicities.
P-F51 constructs the exact $E_6$ factor from a specified
nonnegative function on a definite $U(3)$ quotient, through
an ordinary convergent theta integral with all normalizations.
The comparison still has to identify division by the elliptic
discriminant and the regularized multiplicative operation.
P-F53 supplies actual Hilbert-scheme cohomology spaces for
every coefficient of that full input and a completed symmetric-
operation identity for its inverse-discriminant factor. The
remaining task is comparison of operations and coefficient maps.
P-F55 gives a second actual invariant image from the same input:
the square of the orthogonal form is the degree-$192$
six-variable cubic discriminant, with an exact cyclic restriction
to the square of the earlier five-variable discriminant.
P-F42 bounds the actual average-polynomial norm, with a constant
that still exceeds the index-refutation threshold. Its refined
bound already puts the highest fifth of the coefficient indices
below that threshold, reducing the sufficient gate to the lower
coefficients. P-F40 and
P-F44 expose local and global signed cancellations, and give
an explicit remainder whose coefficient estimates remain open.
P-F43 independently resolves the quarantined binary-degree-twelve
divisibility claim by a full parameter-system counterexample.
P-F47--P-F50 identify exterior-volume powers with actual Fermat
evaluations, construct polynomial-degree nonzero invariants and
prove a polynomial conductor for that point's permitted degrees,
and full eventual support on a specified infinite dimension family.
The supplied Dirichlet theorem sharpens both quantitative bounds.
P-F54 goes further for actual support: one explicit balanced
cubic has a nonzero invariant in every coefficient degree
$d\ge(10n)^{200}$ permitted by the scalar center, for every
$n\ge33$. The compact-average proof controls all dimensions,
phases and global gaps explicitly.
These pointwise occurrences do not control the strict constituent
indices. The rescaled cubic families and their determinant twists
also have exponentially small density in the SH(b) input populations.
PLE-F7--PLE-F8 identify the actual canonical multiplicity algebra on
the obstructing ray as a normal surface with pure even Tate ordinary
and intersection cohomology and a positive multiplicative basis.
Its graded Hilbert function still violates homogeneous PH1.
The known geometric route to
canonical positivity can be expressed by P-F6, but its nonstandard
realization, efficiency bounds and comparison with the relevant
coordinate-ring quotients have not been constructed here.
The supplied RH/Dirichlet-GRH statements do not state any of those
constructions. Reproving an existing finite-field purity theorem does
not enlarge its domain of applicability to an unidentified object.
This is a precise account of a missing deduction, not a logical
independence theorem between RH and a GCT conjecture.


% Agent fragment. All finding IDs begin PL-. Bibliography locators are inline.
\section{The plethystic bridge and the actual content of the archive}
\label{sec:plethysm-bridge}

\subsection{Scope, conventions, and known dictionary}
The archive evidence here consists of the supplied transcriptions, including their
deletions and uncertainty marks, not images of every manuscript page. A claim made
by the supplied study is not thereby a checked claim of this paper. In particular,
the study's order-two check of cote 133, scan 19 fails equivariantly, as proved
below. No assertion of historical priority is made for the elementary consequences
labelled PL below.

We use ``pre-$\lambda$-ring'' for an additive-to-multiplicative map
$\lambda_t:R\longrightarrow 1+tR[[t]]$ with first coefficient $x$, and
``$\lambda$-ring'' for one also satisfying $\lambda_t(1)=1+t$ and the
universal product and iteration identities. This is the convention of SGA 6, Expos\'e V,
Definitions 2.1 and 2.4; some authors call the latter a special $\lambda$-ring.
The vector-bundle Grothendieck ring of a quasi-compact scheme has the latter
structure, with $\lambda^i[E]=[\bigwedge^iE]$ (SGA 6, Expos\'e VI,
Theorem 3.2). The free $\lambda$-ring on one generator is the symmetric-function
ring $\Lambda=\mathbb Z[e_1,e_2,\ldots]$; see the attached Borger,
\emph{$\Lambda$-rings and the field with one element}, \S1.1, the paragraph
beginning ``The functor $W^*$ has a right adjoint''. The universal map to a
$\lambda$-ring $R$ sending the generator to $x$ evaluates $e_i$ as $\lambda^i(x)$.

The characteristic-zero dictionary is exact:
\[
 R=\bigoplus_{n\geq0}K_0(\operatorname{Rep}_{\mathbb C}(\mathfrak S_n))
 \xrightarrow{\operatorname{ch}}\Lambda,
 \qquad [\lambda]\longmapsto s_\lambda,
 \quad [\mathbf1_{\mathfrak S_n}]\longmapsto h_n,
 \quad [\operatorname{sgn}_n]\longmapsto e_n.
\]
Here the multiplication on $R$ is induction from
$\mathfrak S_a\times\mathfrak S_b$ to $\mathfrak S_{a+b}$. It is not the
inner tensor product. For the latter, within one degree,
\[
 \operatorname{ch}([\mu]\otimes[\nu])=s_\mu*s_\nu
   =\sum_{\lambda\vdash n}g(\lambda,\mu,\nu)s_\lambda.
\]
For composition of polynomial Schur functors,
\[
 \operatorname{ch}(S_\alpha\circ S_\beta)=s_\alpha[s_\beta]
   =\sum_{\lambda\vdash |\alpha||\beta|}p^\lambda_{\alpha,\beta}s_\lambda.
\]
The first multiplication adds degrees, the inner tensor product preserves the
common degree, and plethysm multiplies degrees. Exact standard locators are
Macdonald, \emph{Symmetric Functions and Hall Polynomials}, second edition,
I~(7.3), I~(7.9), and Appendix A~(7.3)--(7.4); the last formula already proves
ordinary plethysm coefficients nonnegative. The positive categorical structure
is made explicit in Baez--Moeller--Trimble, \emph{Schur Functors and Categorified
Plethysm}, Theorem 6.11. None of these statements supplies a polynomial-time
algorithm or an Ehrhart model.

For a finite group $G$, let $B(G)$ be its Burnside ring. There are two different
operations to distinguish:
\[
 \epsilon^i(X)=[\{A\subset X:|A|=i\}],\qquad
 \sigma^i(X)=[\operatorname{Sym}^iX].
\]
Both extend to pre-$\lambda$ structures (with the appropriate generating series)
by disjoint-union decompositions. Define the exterior-compatible structure
$\lambda^{\mathrm{op}}$ instead by
\[
 \lambda^{\mathrm{op}}_t(x)=\sigma_{-t}(x)^{-1}.
\]
This generally differs from $\epsilon$. Linearization
$L:B(G)\to R_{\mathbb C}(G)$, $[X]\mapsto[\mathbb C[X]]$, respects
$\sigma$ and $\lambda^{\mathrm{op}}$, but not $\epsilon$ in general.
These conventions agree with Bouc--R\"okaeus,
\emph{A note on the $\lambda$-structure on the Burnside ring},
\S1 and Theorem 1.1, equation (1.2). The following proofs isolate the relevant
small counterexamples rather than reprove that general theorem.

\subsection{PL-01: a false identity in the supplied reading of scan 19}
\paragraph{Statement PL-01.}
The identity $2\epsilon^2(X)+X=X^2$ is false in $B(G)$.
For $G=C_2$ and its free two-point orbit $E$, put $x=[E]$. Then
\[
 x^2=2x,\qquad \epsilon^2(x)=1,\qquad
 \sigma^2(x)=x+1,\qquad \lambda^{\mathrm{op},2}(x)=x-1.
\]
Thus neither the subset operation nor the exterior-compatible operation makes
all three order-two equalities transcribed on scan 19 true at once. The
cardinality equality remains true.

\begin{proof}
The diagonal and off-diagonal pairs in $E\times E$ are two free orbits, proving
$x^2=2x$. Its unique two-element subset is fixed, proving $\epsilon^2(x)=1$.
The two repeated multisets form a free orbit and the multiset containing both
elements is fixed; hence $\sigma^2(x)=x+1$. The coefficient of $t^2$ in
$\sigma_t(x)\lambda^{\mathrm{op}}_{-t}(x)=1$ gives
$\lambda^{\mathrm{op},2}(x)=x^2-\sigma^2(x)=x-1$.
The fixed-point homomorphism $\Phi_G:B(G)\to\mathbb Z$ takes $x$ to zero
and $1$ to one. Therefore it takes $2\epsilon^2(x)+x$ to two and $x^2$ to
zero, so these classes differ. It also distinguishes
$2\lambda^{\mathrm{op},2}(x)+x=3x-2$ from $2x$.
Applying cardinality instead gives $2\binom{2}{2}+2=2^2$.
\end{proof}

The supplied \texttt{SA17\_STUDY\_cote133\_lambda.md}, section
``Scan 19: $\lambda$- and $\sigma$-operations'', claims that with subsets and
multisets the two descriptions agree. PL-01 refutes precisely that assertion.
The correction preserves the true identity
$\operatorname{Sym}^2X\simeq\epsilon^2X\sqcup X$, and the true opposite-series
identity, but uses different symbols for their two exterior-like operations.

\subsection{PL-02: what survives linearization, and why specialness matters}
\paragraph{Statement PL-02.}
Linearization preserves symmetric powers and their opposite-series exterior
operations, but the subset operations need not be special. Moreover, the
opposite-series structure itself need not be special: $G=C_2\times C_2$ gives
an explicit counterexample. Thus the free special $\lambda$-ring universal
property cannot simply be applied to these Burnside structures.

\begin{proof}
The degree-$i$ monomial basis of $\operatorname{Sym}^i\mathbb C[X]$ is indexed
by multisets of size $i$ in $X$, and the $G$-actions agree. Thus
$L(\sigma^i(X))=[\operatorname{Sym}^i\mathbb C[X]]$.
Disjoint unions give the addition law and extend this identity to virtual
classes. Applying $L$ to inverse series proves compatibility with
$\lambda^{\mathrm{op}}$. By contrast, in the example PL-01,
$L(\epsilon^2(x))=1$ whereas $\bigwedge^2\mathbb C[E]=\operatorname{sgn}$.

In every special $\lambda$-ring the second Adams operation
$\psi^2(z)=z^2-2\lambda^2(z)$ is multiplicative. At order two this follows
directly from the universal identity
\[
 \lambda^2(uv)=u^2\lambda^2(v)+v^2\lambda^2(u)
                       -2\lambda^2(u)\lambda^2(v).
\]
For the subset pre-structure on $B(C_2)$ one has
$\psi^2(x)=2x-2$, so
$\psi^2(x)^2=4$, while additivity gives
$\psi^2(x^2)=\psi^2(2x)=4x-4$. Their $C_2$-fixed-point marks are
$4$ and $-4$, so they differ.

For the opposite-series structure take $G=C_2\times C_2$, choose distinct
subgroups $H,K$ of order two, and let $a=[G/H]$, $b=[G/K]$.
Both sets have two elements and their actions factor through $C_2$. Hence
$a^2=2a$, $\sigma^2(a)=a+1$, and
$\psi^2(a)=2\sigma^2(a)-a^2=2$; likewise $\psi^2(b)=2$.
The product $(G/H)\times(G/K)$ is the regular four-element $G$-set $G$,
since its stabilizer is $H\cap K=1$. Write $r=[G]$. A $G$-fixed multiset on
the transitive set $G$ has constant multiplicity at all four points, so no
multiset of size two is fixed. Consequently
\[
 \Phi_G(\psi^2(r))=
  2\Phi_G(\sigma^2(r))-\Phi_G(r^2)=0,
 \qquad
 \Phi_G(\psi^2(a)\psi^2(b))=4.
\]
Multiplicativity fails, proving the claim.
\end{proof}

\subsection{PL-03: the characteristic-zero meaning of scan 28}
\paragraph{Statement PL-03.}
Let $V$ be a complex vector space, regarded universally as a representation of
$GL(V)$. Give $V^{\otimes n}$ its commuting place-permutation action of
$\mathfrak S_n$. Put $\rho_i=[\mathbf1_{\mathfrak S_i}]$, and multiply the
$\rho_i$ by induction. For $\alpha=(\alpha_1,\ldots,\alpha_n)$ with
$\sum i\alpha_i=n$, put $\mu(\alpha)=(1^{\alpha_1}2^{\alpha_2}\cdots)$.
Then the formula on cote 133, scan 28 has the following correct form in the
representation Grothendieck group:
\[
 [V^{\otimes n}]=
 \sum_{\alpha}m_{\mu(\alpha)}[V]\,
          \rho_1^{\alpha_1}\cdots\rho_n^{\alpha_n}.
 \tag{PL-03a}
\]
Here $m_\mu[V]$ means the monomial symmetric function expressed as an
integral polynomial in the exterior powers of $V$; it need not be effective.
In degree two the identity is
\[
 [V^{\otimes2}]=\psi^2[V]\,\rho_2+[\bigwedge^2V]\,\rho_1^2,
 \qquad \psi^2[V]=[\operatorname{Sym}^2V]-[\bigwedge^2V].
 \tag{PL-03b}
\]
The coefficient of $\rho_1^n$ is $e_n[V]$, and that of $\rho_n$ is
$p_n[V]=\psi^n[V]$. These coefficient assertions use the induction ring in
characteristic zero, where the $\rho_i$ freely generate a polynomial ring.

\begin{proof}
Restrict $GL(V)$ to a diagonal torus with eigencharacters $x_1,\ldots,x_r$.
For each exponent vector $(b_1,\ldots,b_r)$ with sum $n$, the tensors having
exactly $b_i$ copies of the $i$th basis vector span a torus-weight space of
weight $x_1^{b_1}\cdots x_r^{b_r}$. As an $\mathfrak S_n$-module this weight
space is the permutation module on ordered words of that content, namely
\[
 \operatorname{Ind}_{\mathfrak S_{b_1}\times\cdots\times\mathfrak S_{b_r}}
              ^{\mathfrak S_n}\mathbf1.
\]
Group exponent vectors by their unordered positive parts $\mu$.
The sum of their torus weights is $m_\mu(x)$, while their common permutation
module is the induction product $\prod_i\rho_i^{\alpha_i}$.
This proves (PL-03a) on the torus. Characters of polynomial $GL(V)$ modules
are symmetric polynomials, with the Schur characters linearly independent,
so torus restriction is injective on their Grothendieck group. Thus it proves
the asserted identity there. By universality it is the corresponding
polynomial identity in exterior powers in characteristic zero.
Finally $m_{(1^n)}=e_n$, $m_{(n)}=p_n$, and
$m_{(2)}=e_1^2-2e_2$; these give the coefficient statements and (PL-03b).
\end{proof}

For $\dim V\geq2$, $\psi^2[V]$ is not the class of a genuine polynomial
$GL(V)$ representation: its irreducible decomposition has coefficient $-1$
on $\bigwedge^2V$. Thus the terms of this particular monomial/induction
expansion need not be direct summands individually. In contrast the full
Schur--Weyl expansion
$V^{\otimes n}=\bigoplus_\lambda S_\lambda V\otimes[\lambda]$
is a direct-sum decomposition. Replacing one basis by another is exactly
where cancellation enters; the archive formula is not a positive combinatorial
formula for arbitrary plethysm multiplicities. This also gives the concrete
counterexample $\psi^2(s_1)=p_2=s_2-s_{11}$ to the inference that Adams
operations or Frobenius lifts must preserve Schur positivity.

\subsection{PL-04: the modular qualification cannot be omitted}
\paragraph{Statement PL-04.}
For the exact Grothendieck groups of finite-dimensional representations over
$k=\mathbb F_2$, the induction ring
$\bigoplus_nG_0(k[\mathfrak S_n])$ is not a polynomial ring freely generated
by the trivial classes $\rho_n$. Indeed $\rho_1^2=2\rho_2$.

\begin{proof}
The induction product $\rho_1^2$ is the regular $k[C_2]$-module.
Writing its generator as $g$, its invariant line is $k(1+g)$. In
characteristic two both this line and the one-dimensional quotient are
trivial modules. The short exact sequence therefore gives
$[k[C_2]]=2[k]$ in $G_0$, which is the asserted relation. A polynomial ring
on formal generators has no such relation, so the natural map from that
polynomial ring is not injective.
\end{proof}

The split Grothendieck group is a different object: the regular module here
is indecomposable, and the displayed short exact sequence does not split.
This distinction must be fixed before answering any question about arbitrary
local coefficient rings. The supplied study discusses scan 29, which was not
itself attached; PL-04 checks the quoted mathematical assertion, not that
page's wording. It does not classify all local rings for which a chosen
representation Grothendieck ring has polynomial generators.

\subsection{PL-05: positive sets cannot supply every positive representation}
\paragraph{Statement PL-05.}
An effective $G$-set linearizes to a character nonnegative on every element
of $G$. In particular, for $n\geq2$ the sign representation of $\mathfrak S_n$
is not the linearization of an effective finite $\mathfrak S_n$-set, even
though its class is positive in the representation semiring. Consequently a
construction that uses only effective sets, disjoint union, Cartesian product,
and linearization cannot produce all Schur-positive representations.

\begin{proof}
The permutation matrix of $g$ on $\mathbb C[X]$ has trace $|X^g|\geq0$.
The sign character has value $-1$ on a transposition, which rules out such
an $X$. Disjoint union and Cartesian product preserve the class of effective
sets, so adding these constructions does not remove the obstruction.
\end{proof}

This is a limitation on that proposed positive-set realization, not a
nonexistence theorem for positive formulas for Kronecker coefficients.
Restriction along the diagonal $G\to G\times G$ turns the external tensor
product into the internal one; cote 2, scan 17 writes exactly this operation.
For $G=\mathfrak S_n$ it is the operation whose irreducible multiplicities
are Kronecker coefficients. Nothing on the attached page identifies their
irreducible multiplicity spaces with explicitly countable sets.

\subsection{PL-06: the exact extra datum needed for the Chow obstruction}
Write $\widehat{\operatorname{Ch}}_n(W)\subset\operatorname{Sym}^nW$ for
the affine cone of products of $n$ linear vectors. The known
Hermite--Hadamard--Howe map is
\[
 h_{d,n}:\operatorname{Sym}^d(\operatorname{Sym}^nW^*)
       \longrightarrow\operatorname{Sym}^n(\operatorname{Sym}^dW^*).
\]
Its kernel is the degree-$d$ Chow ideal (Landsberg, attached survey,
Proposition 7.6). Its target is the degree-$d$ piece of the normalization's
coordinate ring (ibid., Propositions 7.11--7.12). If $\dim W=n$, the sum
of these targets is also the polynomial-representation part of the orbit
coordinate ring (ibid., \S7.5, using Theorem 3.7). Hence it is necessary
to distinguish orbit, orbit closure, and normalization.

\paragraph{Statement PL-06.}
Let $a_\lambda$ and $b_\lambda$ be the multiplicities of $S_\lambda W^*$
in the source and target of $h_{d,n}$. There is a linear map
$H_\lambda:\mathbb C^{a_\lambda}\to\mathbb C^{b_\lambda}$ such that
\[
 m_{\widehat{\operatorname{Ch}}_n,d}(\lambda)=\operatorname{rank}H_\lambda,
 \qquad
 m_{I_d}(\lambda)=a_\lambda-\operatorname{rank}H_\lambda.
 \tag{PL-06}
\]
In particular $0\leq m_{\widehat{\operatorname{Ch}}_n,d}(\lambda)
\leq\min(a_\lambda,b_\lambda)$. Knowledge of the two complete lists of
plethysm multiplicities does not determine these ranks for an arbitrary
equivariant map with that source and target. Even for the natural map,
the equality of source and target characters when $d=n$ does not force
injectivity.

\begin{proof}
Complete reducibility decomposes the two modules into
$\bigoplus_\lambda S_\lambda W^*\otimes\mathbb C^{a_\lambda}$ and
$\bigoplus_\lambda S_\lambda W^*\otimes\mathbb C^{b_\lambda}$.
Schur's lemma identifies an equivariant map with the sum of
$1_{S_\lambda W^*}\otimes H_\lambda$. Its image and kernel on each
summand have the multiplicities in (PL-06). Proposition 7.6 identifies its
image with the coordinate-ring quotient, proving the formula.
When $a_\lambda,b_\lambda>0$, the same source and target admit a zero map
and a rank-one multiplicity-space map, so their characters alone do not
force a rank. For the natural $h_{5,5}$ on $W=\mathbb C^5$, Landsberg's Proposition 7.9(1),
attributed there to Ikenmeyer--Mkrtchyan, supplies a nonzero kernel; for
example the module labelled $(14,7,2,2)$ occurs in it. Yet source and
target are the same module. This proves the final assertion by the cited
counterexample.
\end{proof}

Foulkes' multiplicity comparison and injectivity of a particular
Hermite--Hadamard--Howe map are therefore different assertions. A proof of
ordinary Schur positivity, or even an oracle giving all plethysm coefficients,
still needs control of $H_\lambda$ to produce these orbit-closure multiplicities.
For general GCT orbit closures the analogous missing datum is restriction
of functions from the ambient representation to the closure, including
boundary extension. This is a precise remaining gate, not a proof that a
future purity argument cannot control those maps.

\subsection{PL-07: the tensor-power bridge does not give saturation}
\paragraph{Statement PL-07.}
Kronecker coefficients fail the naive saturation implication
\[
 g(N\lambda,N\mu,N\nu)>0\ \Longrightarrow\ g(\lambda,\mu,\nu)>0.
\]
Specifically,
\[
 g((1,1),(1,1),(1,1))=0,\qquad
 g((2,2),(2,2),(2,2))=1.
\]

\begin{proof}
For $\mathfrak S_2$ the first invariant space is the invariant subspace of
the sign representation cubed, again sign, and is zero.
For $\mathfrak S_4$, let $M$ be its permutation representation on the
three partitions of four letters into two unordered pairs, and let
$U=M/\mathbf1$. On the conjugacy classes
$1,(12),(12)(34),(123),(1234)$, of sizes $1,6,3,8,6$, the numbers of
fixed pairings are $3,1,3,0,1$. Thus $U$ has character
$(2,0,2,-1,0)$. Its character has norm
$(4+3\cdot4+8)/24=1$, so $U$ is irreducible. The hook-length dimensions
of the five Specht modules for $\mathfrak S_4$ are $1,3,2,3,1$
(Macdonald, Chapter I, (7.6)(ii), and Section 5, Example 2), so it is
the module labelled $(2,2)$. The dimension of its tensor-cube invariant
space is the average of the cubed character:
\[
 \frac{1}{24}(1\cdot2^3+6\cdot0^3+3\cdot2^3+8\cdot(-1)^3+6\cdot0^3)
 =1.
\]
This proves both values.
\end{proof}

Ordinary representation positivity and the entire characteristic-zero
$\lambda$-ring dictionary already hold in this counterexample. Consequently
they cannot imply naive saturation. This leaves untouched carefully
qualified saturation-index hypotheses in GCT, whose exact quantifiers must
be considered separately.

\subsection{PL-08: two quadratic qualifications on scans 6 and 8}
For comparison, the usual Grothendieck--Witt ring of a scheme over a field
of characteristic different from two is known to be a special
$\lambda$-ring (Zibrowius, \emph{The $\gamma$-filtration on the Witt ring of
a scheme}, Theorem 2). The analogous statement for symmetric representations
of an affine algebraic group is Zibrowius, \emph{Symmetric representation
rings are $\lambda$-rings}, main theorem, p.~1056. These results concern
the Grothendieck--Witt definitions in those works; one must not silently
substitute the split group completion on an arbitrary locally ringed topos.
\paragraph{Statement PL-08.}
Exterior powers of underlying modules do not define rank-compatible
$\lambda$ operations on the group of strictly nondegenerate quadratic forms
in the generality proposed on cote 133, scan 6, when $2$ is not invertible.
Also, even when $2$ is invertible, the usual exterior-power operations on
Grothendieck--Witt groups need not descend to Witt groups, answering the
question on scan 8 negatively in that generality.

\begin{proof}
Over $\mathbb F_2$ the polar bilinear form of a quadratic form is alternating.
A nonsingular alternating form has even dimension: if $v\ne0$,
nonsingularity supplies $w$ with nonzero pairing, their span is a
nondegenerate plane, and induction on its orthogonal complement proves
the assertion. The hyperbolic quadratic plane $q(x,y)=xy$ has nonsingular
polar form. Its second exterior power as a vector space has dimension one,
so it cannot carry a strictly nondegenerate quadratic form. More generally
all ranks in the Grothendieck group generated by such forms are even,
whereas rank compatibility would require its second exterior operation
to have rank $\binom22=1$. This is a contradiction.

For the descent question work over $\mathbb R$, with symmetric bilinear
forms and the usual exterior-power form. The hyperbolic plane
$h=\langle1\rangle+\langle-1\rangle$ is zero in the Witt group.
Its exterior square is $\langle-1\rangle$. This is nonzero in the Witt
group: signature is additive, invariant under isometry, and zero on every
hyperbolic form, but its value on this line is $-1$. Thus a descended
operation would send the zero Witt class both to zero (from the zero
form) and to a nonzero class (from $h$). It is not well-defined.
\end{proof}

\subsection{PLE-F5: an invariant-ring realization of the obstructing ray}
\paragraph{Statement PLE-F5.}
Let $\zeta$ be a primitive sixth root of unity, and let $\Gamma$ act on
$\mathbb C[x,y]$ through the substitutions
\[
 a(x,y)=(\zeta x,\zeta^{-1}y),\qquad b(x,y)=(iy,ix).
\]
Then $\Gamma$ has order $12$, and
\[
 \mathbb C[x,y]^\Gamma\cong
 \mathbb C[U,V,Z]/(Z^2-UV^2-4U^4),
 \quad U=x^2y^2,\quad V=x^6-y^6,\quad Z=xy(x^6+y^6).
 \tag{PLE-F5a}
\]
Every invariant has even ordinary degree. Dividing degrees by two gives
$\deg U=2$, $\deg V=3$, $\deg Z=4$, and its Hilbert series and Hilbert
function are
\[
 \operatorname{Hilb}(t)=\frac{1+t^4}{(1-t^2)(1-t^3)},\qquad
 \dim\bigl(\mathbb C[x,y]^\Gamma\bigr)_s
      =\left\lfloor\frac{s}{3}\right\rfloor+\mathbf1_{2\mid s}.
 \tag{PLE-F5b}
\]
Consequently this function agrees with the plethysm/Kronecker ray
$\phi(s)$ in Kahle--Micha\l ek, \emph{Obstructions to combinatorial formulas
for plethysm}, Theorem 5 and the paragraph after Remark 13.
The dimension comparison here is strengthened in PLE-F8, which
constructs an explicit graded-algebra isomorphism with the
highest-weight multiplicity algebra of the cubic-plethysm ray.
No compatible Kronecker-algebra identification is asserted.
The invariant ring is normal, but its degree support
$\{0,2,3,4,\ldots\}$ is not saturated in the group $\mathbb Z$ it generates.

\begin{proof}
One checks $a^6=1$, $b^2=a^3$, and $bab^{-1}=a^{-1}$.
Every group element therefore has the form $a^j$ or $a^jb$, $0\leq j<6$.
The former substitutions are diagonal, the latter off-diagonal, and all
twelve are distinct, proving the order.
A monomial $x^py^q$ is invariant under $a$ exactly when $p-q$ is divisible
by six. If $p\geq q$, it is $(xy)^q(x^6)^{(p-q)/6}$; if $q\geq p$ use
$y^6$ instead. Hence
\[
 \mathbb C[x,y]^{\langle a\rangle}
   =\mathbb C[A,B,C]/(AB-C^6),
 \quad A=x^6,\ B=y^6,\ C=xy.
\]
To justify that there is no further relation, reduce a polynomial modulo
$AB-C^6$ to a linear combination of monomials $A^rC^q$ ($r,q\geq0$)
and $B^rC^q$ ($r\geq1,q\geq0$); their images are distinct monomials in
$x,y$, so they are linearly independent.
The remaining involution sends $(A,B,C)$ to $(-B,-A,-C)$.
Putting $V=A-B$, $D=A+B$ gives
\[
 \mathbb C[A,B,C]/(AB-C^6)
   =\mathbb C[C,V,D]/(D^2-V^2-4C^6),
\]
where the involution fixes $V$ and negates $C,D$.
Because the relation is monic in $D$, every class uniquely has the form
$f(C,V)+Dg(C,V)$. It is invariant exactly when $f$ is even in $C$ and
$g$ is odd in $C$. The invariant ring therefore is the free
$\mathbb C[U=C^2,V]$-module with basis $1,Z=CD$; its multiplication is
specified by
$Z^2=C^2D^2=UV^2+4U^4$. This proves (PLE-F5a).

The two basis elements have half-degrees zero and four, proving the
Hilbert series in (PLE-F5b). Independently,
\[
 \sum_{s\geq0}\left\lfloor\frac{s}{3}\right\rfloor t^s
       =\frac{t^3}{(1-t)(1-t^3)},\qquad
 \sum_{s\geq0}\mathbf1_{2\mid s}t^s=\frac1{1-t^2}.
\]
Adding and cancelling $1-t$ gives the same Hilbert series. Finally the
residue classes $s=0,1,\ldots,5\pmod6$ give the periodic correction
$(3,-1,1,0,2,-2)/3$ in $\phi(s)=(s+r(s))/3$, which is the cited ray.

Normality has a direct proof. Put $A=\mathbb C[x,y]$ and $B=A^\Gamma$.
If $f\in\operatorname{Frac}(B)$ is integral over $B$, the same monic
equation shows that it is integral over $A$. Since $A$ is integrally
closed in $\operatorname{Frac}(A)$, we have $f\in A$. As a fraction of
elements of $B$, it is fixed by $\Gamma$, so $f\in B$. Thus $B$ is
normal. Its displayed Hilbert function is zero in degree one and positive
in every degree at least two. Its support generates $\mathbb Z$ because
it contains two and three, but is not saturated because $2\cdot1$ belongs
to the support and $1$ does not.
\end{proof}

The invariant-ring realization shows directly why a positive Hilbert
function and an ordinary homogeneous Ehrhart function are distinct
requirements. The ring is normal and has a positive finite free-module
basis over $\mathbb C[U,V]$, but the associated ray still violates
homogeneous Ehrhart reciprocity and its degree support is not saturated.
These are properties of this weighted grading; normality of a toric
semigroup algebra with its full semigroup grading is a different assertion.
No RH or GRH assumption was used in this calculation.
Positive combinatorial formulas already cover all plethysm coefficients
with two-row output and the corresponding equal-rectangle Kronecker family:
Pak--Panova--Swanson, \emph{A combinatorial interpretation for certain
plethysm and Kronecker coefficients}, arXiv:2511.02312v3 (July 2026),
Theorems 1.1--1.2 and Lemma 2.1. Thus PLE-F5 is an explicit invariant-ring
realization of one ray, not a claim to a first positive formula for that ray.

\subsection{PLE-F6: exactly which rescalings repair homogeneous counting}
\paragraph{Statement PLE-F6.}
For an integer $q\ge1$, define
$\phi_q(s)=\phi(qs)$ for integers $s\ge0$, with $\phi$ as in PLE-F5.
There is a rational polytope $Q$ such that
$\phi_q(s)=\#(sQ\cap\mathbb Z^D)$ for every $s\ge0$ if and only if
$\gcd(q,6)>1$. More precisely, when $q$ is even the rational interval
$Q=[0,q/3]$ works. When $q$ is odd and divisible by three, the interval
$Q=[1/2,q/3+1/2]$ works. Hence two is the smallest positive rescaling
that repairs homogeneous Ehrhart counting on this ray.

\begin{proof}
If $q$ is even, the formula in PLE-F5 gives
\[
 \phi_q(s)=\left\lfloor\frac{qs}{3}\right\rfloor+1
 =\#\bigl([0,qs/3]\cap\mathbb Z\bigr),
\]
including $s=0$. If $q$ is odd and divisible by three, put $k=q/3$.
Then
\[
 \phi_q(s)=ks+\mathbf1_{2\mid s}
 =\left\lfloor(k+1/2)s\right\rfloor
  -\left\lceil s/2\right\rceil+1,
\]
which is the claimed interval count, again including $s=0$.

It remains to exclude odd $q$ not divisible by three. Suppose a rational
polytope $Q$ gave the count. Its dimension would be one. Indeed,
dimension zero gives at most one point. In dimension $d\ge2$, choose
$d+1$ affinely independent rational vertices and clear their denominators
by $M$, obtaining integral vertices $v_0,\ldots,v_d$ in $MQ$. The distinct
integral points
$tv_0+\sum_{j=1}^d k_j(v_j-v_0)$ with $k_j\ge0$ and
$\sum_jk_j\le t$ all belong to $MtQ$. There are
$\binom{t+d}{d}$ of them, contradicting
$\phi_q(Mt)=qMt/3+O(1)$.
The quasi-polynomial extension of the formula for $\phi_q$ to negative
integers therefore obeys one-dimensional Ehrhart reciprocity, as stated
in Kahle--Micha\l ek, Remark 6. At $s=1$ its interior count would be
\[
 \#\bigl(\operatorname{relint}(Q)\cap\mathbb Z^D\bigr)
 =-\phi_q(-1)
 =-\left\lfloor-\frac q3\right\rfloor
 =\left\lceil\frac q3\right\rceil.
\]
But its closed count is
$\phi_q(1)=\lfloor q/3\rfloor$ since $q$ is odd. These two integers
differ by one since $3\nmid q$, contradicting inclusion of the relative
interior in $Q$. No rational homogeneous model exists in this case.
The two constructive cases and this exclusion exhaust all $q\ge1$.
\end{proof}

Kahle--Micha\l ek, Question 15 and its following paragraph, already
consider rescaling a stretching function into an Ehrhart function and
note a positive result in degree one. PLE-F6 records the exact allowed
rescalings for this particular ray. It does not restore the original
$q=1$ instance of PH1, and it does not provide a uniform polynomial bound
for rescaling arbitrary plethysm data.

\subsection{PLE-F7: pure Tate cohomology on the same obstructing cone}
\paragraph{Statement PLE-F7.}
Let $q$ be a prime power with $q\equiv1\pmod{12}$ and let
\[
 X=\operatorname{Spec}
 \mathbb F_q[U,V,Z]/(Z^2-UV^2-4U^4),
 \qquad \deg(U,V,Z)=(2,3,4).
\]
This is a normal affine surface with the Hilbert function $\phi$ of PLE-F5.
For every prime $\ell\nmid6q$, its geometric compactly supported
cohomology is
\[
 H_c^j(X_{\overline{\mathbb F}_q},\mathbb Q_\ell)
 \cong
 \begin{cases}\mathbb Q_\ell(-2)&j=4,\\0&j\ne4.\end{cases}
 \tag{PLE-F7}
\]
The identification respects Frobenius, so this cohomology is even,
Tate, and pure of the expected weight. Nevertheless the Hilbert function
is not a homogeneous rational Ehrhart function and the degree support
is not saturated. With the perverse normalization
$IC_X=j_{!*}\mathbb Q_{\ell,X_{\rm sm}}[2]$, one also has
$IC_X\simeq\mathbb Q_{\ell,X}[2]$. Consequently
$H_c^j(X_{\overline{\mathbb F}_q},IC_X)$ is zero except for
$j=2$, where it is $\mathbb Q_\ell(-2)$.
In fact the displayed hypersurface has
$q^{2r}$ points over $\mathbb F_{q^r}$ for every $r\ge1$ for every odd
$q$, even without $q\equiv1\pmod{12}$.

\begin{proof}
For $q\equiv1\pmod{12}$, choose $\zeta,i\in\mathbb F_q$ of orders
six and four. The substitutions in PLE-F5 define the same group $\Gamma$
of order twelve. Every step in the invariant-ring calculation in that
proof is valid over $\mathbb F_q$: the cyclic weights are distinct and
two is invertible. Thus $X=\mathbb A^2/\Gamma$ with the displayed
coordinate ring, free over $\mathbb F_q[U,V]$ with basis $1,Z$.
The same integral-closure argument proves normality, and the free basis
proves the same Hilbert function. In particular $X$ has dimension two.

We give the cohomological comparison explicitly, so no identification
between graded multiplicities and cohomological degrees is implicit.
The quotient map $\pi:\mathbb A^2\to X$ is finite: each element $a$ of
the polynomial ring satisfies the monic equation
$\prod_{g\in\Gamma}(T-ga)$ with invariant coefficients, and the ring
is generated by its two integral coordinate functions.
Each geometric fiber is a single $\Gamma$-orbit as a set. To see the
last assertion, two distinct finite orbits in $\overline{\mathbb F}_q^2$
can be separated by a polynomial that is zero on the first and one on
the second, by the Chinese remainder theorem for their distinct maximal
ideals. Averaging that polynomial over $\Gamma$ produces an invariant
with the same values, contradicting equality of their quotient points.
Surjectivity follows from the integral extension. The same explicit
invariant presentation is valid after extending the field.

Put $\Lambda_n=\mathbb Z/\ell^n\mathbb Z$. For a finite morphism,
the stalk of $\pi_*\Lambda_n$ is the space of functions on the finite
geometric fiber and $R^j\pi_*\Lambda_n=0$ for $j>0$
(Stacks Project, Proposition 59.55.2, tag \texttt{03QP}). Since the fibers
are single orbits, the unit map identifies
\[
 \Lambda_n\simeq(\pi_*\Lambda_n)^\Gamma
 =\operatorname{im}\left(e:\pi_*\Lambda_n\to\pi_*\Lambda_n\right),
 \qquad e=\frac1{12}\sum_{g\in\Gamma}g.
\]
Here $12$ is invertible in $\Lambda_n$, and $e^2=e$, so this is a direct
summand of sheaves. Finite maps are proper, and the composition theorem
for cohomology with proper support (Milne, \emph{Lectures on \'Etale
Cohomology}, version 2.21, Proposition 18.4(d)) therefore gives
\[
 H_c^j(X_{\overline{\mathbb F}_q},\Lambda_n)
 =eH_c^j(X_{\overline{\mathbb F}_q},\pi_*\Lambda_n)
 =H_c^j(\mathbb A^2_{\overline{\mathbb F}_q},\Lambda_n)^\Gamma.
\]
Affine-space acyclicity and Poincar\'e duality give only
$H_c^4(\mathbb A^2,\Lambda_n)=\Lambda_n(-2)$
(Milne, Example 16.3 and Theorem 24.1).
The trace-normalized class of the origin generates the top group after
the Tate twist. Every element of $\Gamma$ fixes the origin, hence fixes
this generator. Taking the inverse limit in $n$ and tensoring with
$\mathbb Q_\ell$ proves (PLE-F7). All maps are defined over
$\mathbb F_q$, so the assertion is Frobenius-equivariant; geometric
Frobenius acts by $q^2$ on $\mathbb Q_\ell(-2)$.

We verify the intersection-complex assertion separately. The partial
derivatives of $Z^2-UV^2-4U^4$ are
$2Z,-2UV,-V^2-16U^3$. Since the characteristic is not two, their
common zero on $X$ is only the origin $o$. Thus
$j:X\setminus\{o\}\hookrightarrow X$ is the smooth locus.
The fiber of $\pi$ over $o$ consists only of the origin in $\mathbb A^2$:
$U=x^2y^2=0$ and $V=x^6-y^6=0$ force $x=y=0$.
Apply the preceding finite-pushforward and averaging argument to the
localization triangles for the pairs $(X,\{o\})$ and
$(\mathbb A^2,\{0\})$. It gives
\[
 R\Gamma_{\{o\}}(X_{\overline{\mathbb F}_q},\mathbb Q_\ell)
 \simeq
 R\Gamma_{\{0\}}(\mathbb A^2_{\overline{\mathbb F}_q},
                 \mathbb Q_\ell)^\Gamma
 \simeq \mathbb Q_\ell(-2)[-4].
\]
The second isomorphism is purity for the smooth pair
$(\{0\},\mathbb A^2)$ (Milne, Theorem 16.1); the generator is again
the class of the origin and is fixed by $\Gamma$.
Let $i:\{o\}\hookrightarrow X$ and $K=\mathbb Q_{\ell,X}[2]$.
Since a geometric point has no higher cohomology, the local calculation
and the constant stalk give
\[
 i^*K=\mathbb Q_\ell[2],\qquad
 i^!K=\mathbb Q_\ell(-2)[-2].
\]
On the smooth open locus $K$ is a local system shifted by its dimension;
at $o$ its stalk and costalk satisfy the middle-perversity conditions.
Thus $K$ is perverse (BBD, \emph{Faisceaux pervers},
(4.0.1)--(4.0.2)). Moreover $H^0(i^*K)=H^0(i^!K)=0$.
The intermediate-extension criterion (BBD, Corollaries 1.4.24--1.4.25)
therefore identifies $K$ with $j_{!*}\mathbb Q_{\ell,X\setminus\{o\}}[2]$:
it has no nonzero perverse subobject or quotient supported at $o$.
The identification restricts to the identity on the smooth locus;
uniqueness of intermediate extension makes it Frobenius-equivariant.
The stated intersection-cohomology groups now follow by this explicit
shift, independently of the point count below.

There is also an elementary, independent point count. Let $Q$ be any odd
prime power and let $\chi$ be the quadratic character of $\mathbb F_Q$,
extended by $\chi(0)=0$. For $c\ne0$,
\[
 \sum_{v\in\mathbb F_Q}\chi(v^2+c)=-1.
\]
Indeed the number of pairs $(v,w)$ with $w^2-v^2=c$ is $Q-1$,
because $(w-v)(w+v)=c$ is parametrized by an arbitrary nonzero first
factor and two is invertible; counting first by $v$ gives $Q$ plus the
displayed sum. Therefore
\[
 \#X(\mathbb F_Q)
 =\sum_{u,v\in\mathbb F_Q}\bigl(1+\chi(uv^2+4u^4)\bigr)
 =Q^2-\sum_{u\ne0}\chi(u)=Q^2.
\]
The $u=0$ character contribution is zero; for $u\ne0$ use
$uv^2+4u^4=u(v^2+4u^3)$ and the preceding identity. Taking $Q=q^r$
proves the stated point counts. The Hilbert-function and support
failures are those already proved in PLE-F5 and PLE-F6.
\end{proof}

Thus the pure Tate compactly supported cohomology of this affine
quotient, including its intersection cohomology, even with normality
and a positive finite free-module Hilbert description, cannot supply
the missing homogeneous counting assertion.
The ordinary compactly supported cohomology in PLE-F7 has only one
nonzero degree, and the intersection complex is the constant sheaf
shifted by two. The grading whose pieces have dimensions $\phi(s)$ is
different data. PLE-F8 identifies the graded algebra with actual
plethysm multiplicities and supplies a positive multiplicative basis.
It does not identify those graded pieces with cohomological groups
carrying all the additional structure in the categorical criterion
P-F6. This example does not refute a proposed construction with
those additional properties.

\subsection{Archive locator inventory and the scope of the bridge}
\begin{description}
\item[Cote 133, scan 6.] Quadratic Grothendieck--Witt groups, and explicit
questions about the universal product and iteration formulas. This is not
a proof that the quadratic group is a special $\lambda$-ring. The page's
warning about noninvertible $2$ matters; PL-08 gives a counterexample to
the unrestricted exterior-power assertion.
\item[Cote 133, scan 8.] The hyperbolic ideal and the question whether the
pre-$\lambda$ operations descend to its quotient. A question on the sheet
cannot be used as an affirmative theorem; PL-08 answers it negatively
for the standard operations over $\mathbb R$.
\item[Cote 133, scan 9.] Stiefel--Whitney operations and universal tensor
formulas; the assertion that the total class factors through the Witt
quotient is visibly deleted in the transcription.
\item[Cote 133, scan 19.] Subsets, multisets, their proposed comparison, and
an open product-formula question. PL-01--PL-02 give the needed equivariant
correction.
\item[Cote 133, scan 28.] Tensor powers with symmetric-group action, the
induction product, and monomial symmetric functions in the exterior powers.
PL-03 gives its characteristic-zero meaning; PL-04 explains why integral
or modular variants cannot silently use the same freeness.
\item[Cote 2, scan 17.] External tensor products, pullback, projection
formula, and restriction along the diagonal. This supplies the categorical
operation for the inner tensor product, not its numerical multiplicities.
\item[Cote 43, scan 37.] The ring of augmentation-one formal series,
convolution, $\lambda$ operations, and Frobenius--Verschiebung formulas in
characteristic $p$. The last map to $p$-typical Witt vectors carries a
question mark. Formal Frobenius structure is not coefficient positivity;
$p_2=s_2-s_{11}$ is an explicit distinction.
\item[Cote 119, scan 217 (site foliation 216).] A retrospective attribution
of $K$-theory and $\lambda$ operations to the Riemann--Roch work and their
links with symmetric-group operations; no plethysm positivity theorem is
stated on this sheet.
\item[Cote 12, scan 41.] A proposed determination of $\lambda$ operations
through finite operator groups, not a worked-out coefficient formula.
\item[Cote 16, scan 38.] Lefschetz operators, primitive parts, and
hyperplane restriction and pushforward identities. Its capital $\Lambda$ is a
Lefschetz operator, not a plethystic exterior-power operation.
\item[Cote 4, scan 7.] Crystalline/\'etale vanishing cohomology, torsion,
and unit Frobenius roots, with a correction marked false in the margin.
This excerpt does not identify any plethysm multiplicity with that
cohomology.
\item[Cote 79, scan 51.] Scalar torsor actions on triangulated finite
polyhedral structures. Its $\lambda$ is a scalar/action parameter.
\item[Cote 89, scan 18.] Reconstruction of a vector space from punctured
lines and wedge/torsor data. Its $\Lambda$ is a coefficient ring. Neither
of these last two sheets supplies $\lambda$-ring positivity.
\end{description}

Thus the attached pages contain substantial algebra of operations, including
the correct source of tensor, Schur, and plethystic constructions. They do
not state a positive lattice-point rule, a saturation theorem, or an
asymptotic determinant--permanent obstruction. The meaningful next target
is a realization of the relevant multiplicity spaces together with their
restriction maps and quantitative complexity bounds; the formal algebra
alone has the counterexamples above.


\subsection{The actual cubic-plethysm multiplicity algebra}

\begin{proposition}[PLE-F8: the obstructing ray is the binary-dihedral algebra]
\label{prop:actualplethysmalgebra}
Let $V=\mathbb C^2$, let $W=\operatorname{Sym}^3V$, and write
$S=\operatorname{Sym}(W)=\mathbb C[a,b,c,d]$ using the basis
$a=x^3$, $b=x^2y$, $c=xy^2$, $d=y^3$ of $W$.
Here $a,b,c,d$ are algebraically independent generators of
$\operatorname{Sym}(W)$; the monomials in $x,y$ label the basis
vectors of $W$.
The torus $\operatorname{diag}(t,t^{-1})$ assigns these variables
weights $3,1,-1,-3$, and the raising operator is the derivation
\[
 E=a\frac{\partial}{\partial b}
      +2b\frac{\partial}{\partial c}
      +3c\frac{\partial}{\partial d}.
\]
Define the graded algebra
\[
 T=\bigoplus_{s\ge0}T_s,\qquad
 T_s=\ker(E)\cap(S_{4s})_{2s},
 \tag{PLE-F8a}
\]
where the first subscript on $S$ is ordinary polynomial degree and
the second is torus weight. Then $T_s$ is the highest-weight
multiplicity space for
$S^{(7s,5s)}V$ in $\operatorname{Sym}^{4s}(\operatorname{Sym}^3V)$,
and there is an explicit graded-algebra isomorphism
\[
 T\simeq
 \mathbb C[U,V_0,Z]/(Z^2-UV_0^2-4U^4),
 \qquad \deg(U,V_0,Z)=(2,3,4).
 \tag{PLE-F8b}
\]
Thus the ring in PLE-F5 is the actual characteristic-zero plethysm
multiplicity algebra on this ray, rather than merely another algebra
with the same Hilbert function.

The algebra has a homogeneous basis
$\{U^iV_0^jZ^\epsilon:i,j\ge0,\ \epsilon\in\{0,1\}\}$
with nonnegative integral multiplication constants. It admits the
flat integral model
\[
 \mathcal T=\mathbb Z[U,V_0,Z]/(Z^2-UV_0^2-4U^4).
 \tag{PLE-F8c}
\]
For every finite field $\mathbb F_q$ with $q\equiv1\pmod{12}$
and prime $\ell\nmid6q$, PLE-F7 computes the ordinary and
intersection cohomology of its fiber: the only nonzero ordinary
compactly supported group is $H_c^4=\mathbb Q_\ell(-2)$, and
$IC=\mathbb Q_\ell[2]$. Nevertheless its Hilbert function is the
Kahle--Micha\l ek ray, which is not a homogeneous rational Ehrhart
function.
\end{proposition}

\begin{proof}
The action $E(x)=0$, $E(y)=x$ on $V$ gives exactly the displayed
derivation on $S$. Multiplication adds ordinary degrees and torus
weights, and the kernel of a derivation is a subalgebra, so (PLE-F8a)
is an algebra. Each irreducible polynomial $GL_2$ representation of
highest weight $(p,q)$ is $\det^q\otimes\operatorname{Sym}^{p-q}V$.
In its monomial weight basis the raising operator kills exactly its
highest vector, of torus weight $p-q$. All constituents of $S_{4s}$
have central weight $12s$, since a central scalar acts on $W$ with
weight three. Consequently a highest vector of torus weight $2s$
has $p+q=12s$, $p-q=2s$, hence $(p,q)=(7s,5s)$.
Complete reducibility in characteristic zero now identifies the
dimension and the space in (PLE-F8a) with the stated multiplicity
space, after choosing the usual highest line in each irreducible.
This identification respects the canonical Cartan multiplication.
Choose compatible highest vectors
$v_s=x^{2s}\otimes(x\wedge y)^{\otimes5s}$ in
$S^{(7s,5s)}V\simeq\operatorname{Sym}^{2s}V\otimes(\det V)^{5s}$.
Evaluation $F\mapsto F(v_s)$ identifies
$\operatorname{Hom}_{GL_2}(S^{(7s,5s)}V,S_{4s})$ with $T_s$.
Normalize the Cartan injection
$S^{(7(s+t),5(s+t))}V\to
 S^{(7s,5s)}V\otimes S^{(7t,5t)}V$
to send $v_{s+t}$ to $v_s\otimes v_t$.
Composing it with $F\otimes G$ and multiplication in $S$ sends
$v_{s+t}$ to $F(v_s)G(v_t)$. Hence evaluation is an isomorphism
of graded algebras for the canonical multiplication of GCT VI,
Section 4.1.2, equation (4.5), printed p.~69, constructed in
equations (4.2)--(4.3), printed p.~67.
For the $\mathfrak{sl}_2$ representation facts, see
Borcherds--Haiman--Reshetikhin--Serganova--Johnson-Freyd,
\emph{Lectures in Lie Groups and Quantum Groups}, Part I,
Example 5.2.0.5 and Proposition 5.2.0.7, printed pp.~68--69,
and Theorem 7.6.1.3, printed p.~167; the passage to $GL_2$ here
adds the scalar central character just computed.
Kahle--Micha\l ek, Theorem 5, therefore gives for $s\ge1$
\[
 \dim T_s=\phi(s)=\left\lfloor\frac{s}{3}\right\rfloor
                      +\mathbf1_{\{2\mid s\}}.
 \tag{PLE-F8d}
\]
Their notation $(7s,5s,0)$ is the same partition with a trailing
zero. Specializing the Schur-character identity to two variables
preserves this coefficient: Schur polynomials with more than two
rows vanish, and those with at most two rows are linearly independent.
At $s=0$, both sides are one directly.

We construct the isomorphism. Put
\[
 h=ac-b^2,\qquad
 j=a^2d-3abc+2b^3,
\]
and
\[
 \delta=a^2d^2-6abcd+4ac^3+4b^3d-3b^2c^2.
\]
Direct application of $E$ gives $E(a)=E(h)=E(j)=0$, and polynomial
expansion gives the identity
\[
 j^2+4h^3=a^2\delta.
 \tag{PLE-F8e}
\]
Applying $E$ to this identity proves $a^2E(\delta)=0$; since
$S$ is a domain, $E(\delta)=0$.
The pairs (ordinary degree, torus weight) of $a,h,j,\delta$ are
respectively $(1,3),(2,2),(3,3),(4,0)$.
Define actual elements of $S$ by
\[
 u=h^2\delta,\qquad
 v=aj\delta^2,\qquad
 w=a^2h\delta^3,\qquad z=w-2u^2.
\]
The pairs for $u,v,z$ are $(8,4),(12,6),(16,8)$, so they belong
to $T_2,T_3,T_4$, respectively. From (PLE-F8e) one obtains
\[
 v^2+4uw=a^4\delta^5,\qquad
 w^2=u(v^2+4uw).
\]
It follows that
\[
 z^2=(w-2u^2)^2=uv^2+4u^4.
 \tag{PLE-F8f}
\]
Thus $U\mapsto u$, $V_0\mapsto v$, $Z\mapsto z$ defines a graded
homomorphism from the right side of (PLE-F8b) into $T$.

We prove its injectivity, including the algebraic independence needed
for that step. Suppose $F(u,v)=0$ for a nonzero polynomial
$F(X,Y)\in\mathbb C[X,Y]$. Remove the largest common power of $X$
from $F$. Since $u\ne0$ and $S$ is a domain, the remaining polynomial
still gives a relation, and now $F(0,Y)\ne0$.
For instance, $u=4$ at $a=c=1$, $b=d=0$, verifying $u\ne0$.
Specialize $a=1$, $b=c=0$. Then $h=0$, $j=d$, $\delta=d^2$,
so $u=0$ and $v=d^5$. The relation would give
$F(0,d^5)=0$ in $\mathbb C[d]$, impossible for a nonzero
$F(0,Y)$. Hence $u,v$ are algebraically independent.

In the field $K=\mathbb C(U,V_0)$, the element
$UV_0^2+4U^4=U(V_0^2+4U^3)$ has valuation one at the prime $U$.
It is therefore not a square, since every square has even valuation.
Consequently $Z^2-UV_0^2-4U^4$ is irreducible over $K$.
Algebraic independence embeds $K$ in $\operatorname{Frac}(S)$ via
$U,V_0\mapsto u,v$, and the element $z$ satisfies this irreducible
quadratic. The resulting map from the degree-two field extension of
$K$ is injective. Restricting to its polynomial subalgebra proves
injectivity of the homomorphism in (PLE-F8b). Equivalently, no
relation $A(u,v)+zB(u,v)=0$ is possible: if $B\ne0$, it would put
$z$ in $K$ and make the displayed element a square, and if $B=0$,
algebraic independence forces $A=0$.

The source of this homomorphism is free over $\mathbb C[U,V_0]$
with basis $1,Z$, because its relation is monic quadratic in $Z$.
Its Hilbert series is
\[
 \frac{1+t^4}{(1-t^2)(1-t^3)}.
\]
PLE-F5 computed its coefficient in degree $s$ as $\phi(s)$.
By (PLE-F8d), the injective homomorphism is therefore an isomorphism
on each finite-dimensional graded piece, hence on the graded algebra.
This proves (PLE-F8b) without presupposing a presentation theorem
for the full algebra of binary-cubic covariants.

The same monic relation shows that the monomials stated in the
proposition form a basis over $\mathbb Z$ for (PLE-F8c).
They span freely over every coefficient ring. In multiplying two
such basis monomials, the only needed reduction is
$Z^2=UV_0^2+4U^4$, whose coefficients are nonnegative integers.
Thus the multiplication constants have the claimed positivity.
Freeness over $\mathbb Z$ also proves flatness of $\mathcal T$,
and all its fibers have the same graded Hilbert function.
Its complex fiber is $T$ by the proved isomorphism. The finite-field
statement is exactly PLE-F7 for the identical equation and grading;
the non-Ehrhart assertion is Kahle--Micha\l ek, Theorem 5 and
Remark 6, as used in P-F1 and PLE-F6. This completes the proof.
\end{proof}

\paragraph{What passes to finite characteristic, and what does not.}
The finite-field scheme here is an explicit flat model of the
characteristic-zero multiplicity algebra. It is not asserted to be
the kernel of the same Lie-algebra raising operator after reduction
modulo every prime. In characteristic $p$, a derivation kills
$p$th powers, so that kernel can acquire additional elements and
cannot be substituted for highest-weight multiplicities without
further arguments. The characteristic-zero isomorphism and the
finite-field geometric model are the claims proved above.

This result removes the possible objection that the binary-dihedral
example only reproduced the numerical Hilbert function of an
unrelated algebra. On this actual plethysm multiplicity algebra,
normality, a positive multiplicative basis, and a pure Tate model
still do not force the literal homogeneous PH1 conclusion.
The coordinate-ring degree and the cohomological degree remain
different gradings. The richer compatibility and polyhedral
requirements of a canonical-basis or crystal construction have not
been proved or refuted by this basis alone. Nor does this example
refute the narrow quantitative PH2 or PH3 bounds: it already has
the displayed positive rational form with denominator weights two
and three. No historical priority is asserted for these explicit
covariant formulas or their algebraic consequences.


% Bridge sections by /root/bridges.  All source locators are explicit in text.
% Bibliography keys and evidence URLs are listed in bridges_notes.md.

\section{Schubert $K$-theory, alternating positivity, and saturation}
\label{sec:bridge-k}

\paragraph{B3-KNOWN: the exact bridge.}
For the complex Grassmannian $\operatorname{Gr}(r,N)$, let
$b_\lambda=[\mathcal O_{\Omega_\lambda}]$ be its Schubert structure-sheaf
basis.  Write
\[
 b_\lambda b_\mu=\sum_\nu c^{K,\nu}_{\lambda\mu}b_\nu.
\]
Grothendieck polynomials are the representatives introduced by Lascoux and
Sch\"utzenberger; they are not polynomials written in the supplied
Grothendieck manuscript pages.  Buch's \emph{A Littlewood--Richardson rule
for the $K$-theory of Grassmannians}, Theorem~5.4 and Corollary~5.5,
express $(-1)^{|\nu|-|\lambda|-|\mu|}c^{K,\nu}_{\lambda\mu}$ as a
set-valued-tableau count.  His Introduction, immediately after the
definition of the product, identifies $c^{K,\nu}_{\lambda\mu}$ with the
ordinary Littlewood--Richardson coefficient when
$|\nu|=|\lambda|+|\mu|$, and says it is zero when
$|\nu|<|\lambda|+|\mu|$ \cite{buch2002}.
For a complex flag variety, Brion's \emph{Positivity in the Grothendieck
group of complex flag varieties}, Theorem~1, proves
\[
 (-1)^{\operatorname{codim}X_w-\operatorname{codim}Y}
 c_Y^w\geq0,
 \qquad [\mathcal O_Y]=\sum_w c_Y^w[\mathcal O_{X_w}],
\]
when $Y$ has rational singularities.  The product version follows from
Theorem~1 and Lemma~2, as explained directly after Theorem~1
\cite{brion2002}.  These hypotheses and the choice of basis matter.

The classical dilation equivalence says, for every integer $N\geq1$,
\[
 c_{N\lambda,N\mu}^{N\nu}>0
 \quad\Longleftrightarrow\quad c_{\lambda\mu}^{\nu}>0.
\]
An exact primary locator is Knutson--Tao, \emph{The honeycomb model of
$GL_n(\mathbb C)$ tensor products I}, \S6, Theorem~3 and its following
corollary converting honeycombs to tensor-product multiplicities,
with the forward semigroup property explained in its Introduction
\cite{kt1999}.  This is the additional integrality result in the
Littlewood--Richardson model, not a consequence of nonnegative
coefficients alone.

\subsection*{B3.1. The sign change and its precise computational scope}
\paragraph{Statement.}
Put $e_\lambda=(-1)^{|\lambda|}b_\lambda$.  The multiplication constants
in this basis are the nonnegative integers
$a_{\lambda\mu}^{\nu}=(-1)^{|\nu|-|\lambda|-|\mu|}
c^{K,\nu}_{\lambda\mu}$.  On the degree-equality locus
$|\nu|=|\lambda|+|\mu|$, an oracle for exact $K$-theoretic constants is
also an oracle for exact Littlewood--Richardson constants, with the same
three partitions as input.  In particular, changing the basis does not
produce a plethysm formula.
\paragraph{Proof.}
Substitute $b_\rho=(-1)^{|\rho|}e_\rho$ into the product:
\[
 e_\lambda e_\mu
 =\sum_\nu(-1)^{|\lambda|+|\mu|+|\nu|}
      c^{K,\nu}_{\lambda\mu}e_\nu
 =\sum_\nu a_{\lambda\mu}^{\nu}e_\nu.
\]
The two sign exponents differ by the even integer
$2(|\lambda|+|\mu|)$.  Nonnegativity is the cited Buch theorem.
On degree equality the cited identification gives
$a_{\lambda\mu}^{\nu}=c^{K,\nu}_{\lambda\mu}=c_{\lambda\mu}^{\nu}$,
so one oracle call proves the computational assertion.  Every operation
in this argument is a product or a sign change; it supplies no identity
for a Schur-functor composition.  This last sentence describes the scope
of the construction, not a formal independence theorem.\hfill$\square$

\subsection*{B3.2. Geometric alternating positivity does not preserve dilation support}
\paragraph{Statement.}
The sign-corrected $K$-theoretic coefficients do not satisfy the classical
dilation equivalence.  In fact
\[
 a_{(1),(1)}^{(2,1)}=1,
 \qquad a_{(2),(2)}^{(4,2)}=0.
\]
Both triples fit in the Schubert rectangle of $\operatorname{Gr}(2,6)$.
This refutes even the implication from nonvanishing at scale one to
nonvanishing at scale two, for these coefficients.  It does not refute
the separate, opposite implication on its own.
\paragraph{Proof.}
Use Lenart's Pieri formula in Buch, Corollary~4.7, together with
conjugation of partitions (Buch, Corollary~6.7; equivalently the
Grassmannian-duality identity in the Introduction).  In horizontal-strip
form it says
\[
 c^{K,\nu}_{\lambda,(p)}=
 (-1)^{|\nu/\lambda|-p}
 \binom{r(\nu/\lambda)-1}{|\nu/\lambda|-p}
\]
when $\nu/\lambda$ is a horizontal strip, where $r$ is its number of
nonempty rows, and the coefficient is zero otherwise.  For
$\lambda=(1)$, $p=1$, $\nu=(2,1)$, the two added boxes occupy two
different columns and two rows.  The coefficient is
$(-1)^1\binom11=-1$, also recorded in Buch, Example~5.6.
For $\lambda=(2)$, $p=2$, $\nu=(4,2)$, the four added boxes occupy
columns $3,4$ in the first row and columns $1,2$ in the second.  Hence
the strip is horizontal, $r=2$, and the coefficient is
$(-1)^2\binom12=0$.  Applying the defining sign correction proves both
values.\hfill$\square$

\paragraph{Obstruction gate for bridge 3.}
To transfer a $K$-theoretic computation into determinant-versus-permanent
GCT, one must identify a particular coordinate-ring multiplicity (or a
map detecting its failure to extend across the boundary) with the
specified $K$-theoretic coefficient, with all signs and degrees fixed.
Neither the name ``Grothendieck'' nor Brion's theorem provides that
identification.  B3.2 shows that even an available geometric positivity
theorem cannot simply be used as the Littlewood--Richardson saturation
step of the flip.  The supplied archive page cote~119, scan~217
(site page~216), discusses the $K$-invariant and $\lambda_i$ operations;
it contains neither the later Lascoux--Sch\"utzenberger polynomials nor
Brion's theorem.

\section{The metric inequality: what its positivity actually controls}
\label{sec:bridge-metric}

\paragraph{B4-KNOWN: dictionary and exact locators.}
For a real matrix $A=(a_{ij})$, define
\[
 B(A)=\max_{\varepsilon_i,\delta_j\in\{-1,1\}}
          \sum_{ij}a_{ij}\varepsilon_i\delta_j,
 \qquad
 S(A)=\max_{\|u_i\|=\|v_j\|=1}
          \sum_{ij}a_{ij}\langle u_i,v_j\rangle.
\]
The real Grothendieck inequality is $B(A)\leq S(A)\leq K_GB(A)$,
where one universal constant works for all dimensions.  The historical
primary locator is Grothendieck's 1953 \emph{R\'esum\'e de la th\'eorie
m\'etrique des produits tensoriels topologiques}, \S4, no.~2,
Theorems~1--2; its tensor-norm formulation is transported to the displayed
finite-matrix formulation in Alon--Naor, Introduction, programs~(1)--(2)
\cite{groth1953,alon-naor2006}.  Alon--Naor, Theorem~4.3, supplies
polynomial-time randomized rounding with expected objective at least
$2\log(1+\sqrt2)/\pi$ times the SDP objective.  Their Lemma~2.1 relates this
optimization to the cut norm.

Raghavendra--Steurer, \emph{Towards Computing the Grothendieck Constant},
Theorem~1.1, proves that a factor strictly better than $K_G$ for this
\emph{bipartite quadratic optimization problem} is NP-hard under the
Unique Games Conjecture.  This is not a threshold for every problem
involving tensors \cite{rs2009}.  In a two-party XOR game with a specified
target parity $f(i,j)$ and question distribution $\pi(i,j)$, set
$a_{ij}=\pi(i,j)(-1)^{f(i,j)}$.  Tsirelson's vector characterization
identifies the classical bias with $B(A)$ and quantum bias with $S(A)$.
For an accessible exact research-paper locator, see Cleve--H\o yer--Toner--Watrous,
\emph{Consequences and Limits of Nonlocal Strategies}, revised version
of 11 January 2010, Theorem~4 (Tsirelson), Proposition~5, and Theorem~7
\cite{chtw2010}.  Bias is twice winning probability minus one; this
qualification avoids conflating a bias ratio with a probability ratio.

\subsection*{B4.1. A zero-versus-positive norm realization already carries the decision problem}
\paragraph{Statement.}
For a rational real matrix $A$,
\[
 B(A)=0\quad\Longleftrightarrow\quad S(A)=0
 \quad\Longleftrightarrow\quad \|A\|_{\mathrm{cut}}=0
 \quad\Longleftrightarrow\quad A=0.
\]
Consequently, suppose a nonnegative integer function $m(z)$ has an
explicit polynomial-time construction of a polynomial-size rational
matrix $A_z$, with polynomial-bit entries, such that
$m(z)>0$ if and only if $B(A_z)>0$ (or the analogous condition with
$S$ or the cut norm).  Then deciding $m(z)>0$ is in deterministic
polynomial time, by inspecting entries alone.
\paragraph{Proof.}
The maximum of a bilinear form over two cubes is attained at vertices:
fix one cube variable at a time and move it to an endpoint that does not
decrease the objective.  Therefore $B(A)$ is also the maximum over
$x_i,y_j\in[-1,1]$.  For any entry $a_{ij}$, choose $x$ and $y$ supported
only at $i$ and $j$, with their signs chosen to make the objective
$|a_{ij}|$.  Thus $B(A)\geq\max_{ij}|a_{ij}|$.
Since scalar signs are admissible unit vectors, $S(A)\geq B(A)$.
A singleton rectangle gives
$\|A\|_{\mathrm{cut}}\geq\max_{ij}|a_{ij}|$.
All three quantities vanish for the zero matrix.  These observations
prove the equivalences.  The algorithm constructs $A_z$ and checks
whether any rational numerator is nonzero, taking polynomial time in
the asserted output size.\hfill$\square$

\subsection*{B4.2. A constant-factor estimate cannot generally decide a one-unit multiplicity gap}
\paragraph{Statement.}
Fix $\alpha>1$.  Knowing only a number $v$ satisfying
$c\leq v\leq\alpha c$ for an unknown positive integer $c$ does not
determine $c$, even when $c$ is promised to be one of two consecutive
integers.  Hence such an estimate by itself cannot certify every strict
comparison between two multiplicities.
\paragraph{Proof.}
Choose an integer $c\geq1/(\alpha-1)$.  The common output $v=c+1$
is permitted when the true value is $c$, since
$c+1\leq\alpha c$, and also when it is $c+1$.  An algorithm supplied
only with this guarantee and this output receives identical information
in the two cases.  Against a comparison value $c$, strict inequality
holds in one case and fails in the other.\hfill$\square$

\paragraph{Obstruction gate for bridge 4.}
The established passage is from a bilinear objective to a semidefinite
relaxation and a rounding guarantee.  GCT needs an explicit map from a
specified multiplicity or orbit-closure question to an optimization
instance with a preserved, adequately separated threshold.  B4.1 locates
the whole positivity difficulty in any proposed zero-preserving
matrix construction, and B4.2 records the precision needed for a
multiplicity obstruction.  None of the supplied cote/scan pages is the
1953 \emph{R\'esum\'e}.  No claim that RH or a subexponential SAT
algorithm resolves the Unique Games Conjecture is made here.

\section{Projective tensor norms and algebraic degenerations}
\label{sec:bridge-norm}

\paragraph{B5-KNOWN: exact objects.}
For finite-dimensional real or complex Hilbert spaces $V_1,\ldots,V_d$,
the projective tensor norm is
\[
 \|T\|_\pi=
 \inf_{T=\sum_j v_{1j}\otimes\cdots\otimes v_{dj}}
        \sum_j\prod_{i=1}^d\|v_{ij}\|.
\]
Normalizing each nonzero factor gives exactly the tensor nuclear norm
of Friedland--Lim, \emph{Nuclear Norm of Higher-Order Tensors},
equation~(1) and \S2.  Their Theorem~8.1 states NP-hardness for real
$d$-tensors for $d\geq3$; Theorem~8.10 supplies the complex statement
for $d\geq4$ in the arXiv v3 numbering.  These are norm-computation
statements, not tensor-rank lower bounds.  Base field matters for the
norm as well as for rank \cite{friedland-lim2018}.
Tensor rank counts decomposable summands without a norm bound.
Border rank asks membership in the Zariski closure of the tensors of
bounded rank.  For the matrix-multiplication tensor, a decomposition
$T_{\mathrm{MM}}=\sum_{j=1}^r u_j\otimes v_j\otimes w_j$ means precisely
$AB=\sum_{j=1}^r u_j(A)v_j(B)w_j$: substitution in the tensor identity
gives an algorithm using $r$ products of scalar linear forms.
The norm additionally depends on Hilbert metrics, which general linear
changes of bases do not preserve.

\subsection*{B5.1. Scalar changes rule out an unnormalized norm obstruction for a full linear orbit}
\paragraph{Statement.}
For $a\neq0$ and $T\neq0$,
\[
 R(aT)=R(T),\qquad \underline R(aT)=\underline R(T),
 \qquad \|aT\|_\pi=|a|\|T\|_\pi.
\]
Moreover, $aT$ belongs to the same
$\prod_i\operatorname{GL}(V_i)$-orbit as $T$.  In particular a
nontrivial absolute norm threshold cannot be constant on these orbits
or distinguish their scalar-invariant closure membership.
\paragraph{Proof.}
Multiplication of the first factor of every decomposable summand by
$a$ sends a rank-$r$ decomposition of $T$ to one of $aT$; multiplication
by $a^{-1}$ proves the reverse rank inequality.  The linear
automorphism $T\mapsto aT$ preserves each rank-at-most-$r$ locus, hence
its Zariski closure, proving the border-rank identity.  Applying the
same change to every decomposition in the definition of $\|\cdot\|_\pi$
gives $\|aT\|_\pi\leq|a|\|T\|_\pi$; the inverse change proves
equality.  Finally $(a\operatorname{id}_{V_1},\operatorname{id}_{V_2},
\ldots,\operatorname{id}_{V_d})$ is the required group element.
As $|a|$ ranges over positive reals, the norm crosses any fixed positive
threshold while the orbit is unchanged.\hfill$\square$

\subsection*{B5.2. The exact bounded-operator version that does carry an obstruction}
\paragraph{Statement.}
For linear maps $A_i:V_i\to W_i$ between finite-dimensional Hilbert
spaces,
\[
 \|(A_1\otimes\cdots\otimes A_d)T\|_\pi
 \leq\left(\prod_i\|A_i\|_{\mathrm{op}}\right)\|T\|_\pi.
\]
Therefore $\|S\|_\pi>\|T\|_\pi$ proves that $S$ is not in the
Euclidean closure of the tensors obtained from $T$ by maps satisfying
$\|A_i\|_{\mathrm{op}}\leq1$.  This statement does not impose a
restriction on a full $\operatorname{GL}$ orbit.
\paragraph{Proof.}
Apply the tensor product of the maps to an arbitrary decomposition of
$T$.  Each summand has projective cost at most
$\prod_i\|A_i\|_{\mathrm{op}}$ times its original cost, by the
defining operator-norm inequalities.  Sum and take the infimum.
When all maps are contractions, the resulting tensor has nuclear norm
at most $\|T\|_\pi$.  The norm is continuous in finite dimensions, so
the inequality persists under Euclidean limits.  B5.1 supplies the
unbounded scalar elements of the full general linear group, explaining
why that group is not covered.\hfill$\square$

\subsection*{B5.3. Even a natural normalization gives no uniform rank gap}
\paragraph{Statement.}
For every fixed order $d\geq2$ there are tensors $T_\epsilon$ of rank
and border rank two such that
\[
 \frac{\|T_\epsilon\|_\pi}{\|T_\epsilon\|_F}\longrightarrow1
 \quad(\epsilon\downarrow0),
\]
the same ratio as for a nonzero decomposable tensor.  All
$T_\epsilon$, $\epsilon>0$, lie in one general linear orbit.  Thus a
uniform positive additive gap in this ratio cannot separate rank one
from rank two.
\paragraph{Proof.}
Take
$T_\epsilon=(e_1\otimes e_1+\epsilon e_2\otimes e_2)\otimes
e_1\otimes\cdots\otimes e_1$, with every displayed $e_i$ unit and
orthogonal when indices differ.  It has a two-term decomposition.
Flattening between its first factor and all remaining factors has
matrix rank two for $\epsilon>0$.  Flattenings of a rank-one tensor
have matrix rank at most one, and their $2$-by-$2$ minors continue to
vanish on its Zariski closure.  This proves both asserted ranks.
The displayed decomposition gives
$\|T_\epsilon\|_\pi\leq1+\epsilon$.
The multilinear functional
\[
 L(v_1,\ldots,v_d)=
   \langle v_1,v_2\rangle_{\mathrm{bil}}
       \prod_{i=3}^d(v_i)_1
\]
has absolute value at most $\prod_i\|v_i\|$; over $\mathbb C$ the
first pairing is the bilinear coordinate pairing, for which the same
Cauchy--Schwarz bound holds.  Hence its linear extension satisfies
$|L(T)|\leq\|T\|_\pi$ by applying it to every decomposition.
Since $L(T_\epsilon)=1+\epsilon$, equality follows.  Orthogonality
gives $\|T_\epsilon\|_F=\sqrt{1+\epsilon^2}$, proving the limit.
The invertible diagonal map with entries $1,\epsilon$ on the first
factor sends $T_1$ to $T_\epsilon$.\hfill$\square$

\paragraph{Obstruction gate for bridge 5.}
A nuclear-norm bound can be used for a norm-constrained tensor model,
as B5.2 proves.  An algebraic orbit-closure obstruction requires a
quantity compatible with the allowed general linear action and with
degeneration, or an explicit theorem converting the bounded model into
the algebraic model with controlled loss.  Neither NP-hardness of
computing $\|\cdot\|_\pi$ nor a faster Boolean counting algorithm
supplies such a conversion.  The tensor operations in cote~2,
scan~17 are categorical pullback, pushforward, external tensor product,
and restriction along the diagonal; no Banach norm is specified there.

\section{Integral group schemes, Hilbert and Quot schemes, and fixed witnesses}
\label{sec:bridge-integral}

\paragraph{B7-KNOWN: what the foundations provide.}
SGA~3, Expos\'e~I (Demazure), \S4.2 describes affine group schemes by
their coordinate algebras, and \S4.5 constructs
\[
 \operatorname{GL}_{m,\mathbb Z}
 =\operatorname{Spec}\mathbb Z[t_{ij},\det(t)^{-1}].
\]
Thus integral tensor representations and their orbit maps can be
formed before choosing a characteristic \cite{sga3-I}.
Grothendieck, FGA, Bourbaki Expos\'e~221,
\emph{Les sch\'emas de Hilbert}, Theorems~3.1--3.2, represents the
functor of flat coherent quotients of a coherent sheaf on a projective
scheme over a noetherian base, with fixed Hilbert polynomial, by a
projective scheme.  The Hilbert scheme is the special case of a
quotient of the structure sheaf, immediately following Proposition~3.8.
Remark~3.9 explicitly reduces the universal construction to schemes
over $\mathbb Z$ and warns of components confined to a positive
characteristic \cite{fga221}.  These are the relevant exact historical
locators; they are not among the supplied cote/scan transcriptions.

\subsection*{B7.1. A characteristic-zero separation survives outside an explicit finite set}
\paragraph{Statement.}
Let $W$ be a finite free $\mathbb Z$-module with an algebraic action of
$G=\operatorname{GL}_{m,\mathbb Z}$, and let $x,y\in W$.
If $y$ does not belong to $\overline{G_{\overline{\mathbb Q}}x}$,
then there is $f\in\mathbb Z[W]$ with $f(y)\neq0$ whose restriction
to the orbit map $g\mapsto gx$ is identically zero over $\mathbb Z$.
For every prime $p\nmid f(y)$,
\[
 \overline y\notin
 \overline{G_{\overline{\mathbb F}_p}\overline x}.
\]
Consequently noncontainment of the corresponding orbit closures
survives at all these primes.  This assertion concerns one fixed pair
and one fixed witness; it asserts no uniform bound in a growing GCT
family.
\paragraph{Proof.}
Put $R=\mathbb Z[W]$ and
$A=\mathbb Z[t_{ij},\det(t)^{-1}]$.  The integral orbit map induces a
homomorphism $\phi_x:R\to A$.  Let
$I_{\mathbb Q}=\ker(\phi_x\otimes\mathbb Q)$.
Flatness of field extensions implies
$I_{\mathbb Q}\otimes\overline{\mathbb Q}
=\ker(\phi_x\otimes\overline{\mathbb Q})$.
This kernel is exactly the ideal of the orbit closure: a polynomial
vanishes on the orbit if and only if its pullback vanishes as a regular
function on the reduced group $G_{\overline{\mathbb Q}}$.
There is consequently $h$ in the extended ideal with $h(y)\neq0$.
Write $h$ as a finite $\overline{\mathbb Q}$-linear combination of
elements of $I_{\mathbb Q}$.  Evaluation is linear, so at least one of
those elements, say $f_0\in I_{\mathbb Q}$, has $f_0(y)\neq0$.
Clear its rational denominators to get $f\in R$.
Now $\phi_x(f)$ is zero in $A\otimes\mathbb Q$.
The ring $A$ is a localization of a polynomial ring over $\mathbb Z$,
hence torsion-free over $\mathbb Z$; therefore $\phi_x(f)=0$ in $A$.
Reduce this identity modulo $p$ and then extend to
$\overline{\mathbb F}_p$.  The reduced polynomial vanishes on the
reduced orbit, hence on its Zariski closure.  But it takes the nonzero
value $f(y)\bmod p$ at $\overline y$ if $p\nmid f(y)$, proving the
assertion.  Finally, for a group-stable closed set $X$,
$\overline{Gy}\subseteq X$ implies $y\in X$, and $y\in X$ implies
$\overline{Gy}\subseteq X$.  Apply this with the orbit closure of
$x$ in each field.\hfill$\square$

\subsection*{B7.2. Fixed-degree separation has an explicit rank criterion and stable reductions}
\paragraph{Statement.}
In B7.1, fix $d\geq0$.  There are an explicit integer matrix $M_d$
and an integer row $v_d$ such that, over any algebraically closed field
$k$, a polynomial of degree at most $d$ separates $y_k$ from
$\overline{G_kx_k}$ if and only if
\[
 \operatorname{rank}_k\begin{pmatrix}M_d\\v_d\end{pmatrix}
    >\operatorname{rank}_k M_d.                    \tag{B7.2}
\]
One can choose a nonzero integer $D_d$, as a product of at most two
nonzero rank-size minors, so that both ranks, and therefore this
separation condition, agree with characteristic zero whenever
$p\nmid D_d$.
\paragraph{Proof.}
Choose the monomials of degree at most $d$ as a basis of
$R_{\leq d}$.  Each of their pullbacks under $\phi_x$ is a regular
function on $G$, hence an integral polynomial in $t_{ij}$ divided by
a power of $\det(t)$.  There are finitely many basis monomials, so a
single power $\det(t)^K$ clears all these denominators.  Expand
$\det(t)^K\phi_x$ on this basis in ordinary monomials of $t_{ij}$,
using the finite union of all their supports as rows.  Its coefficient
matrix is $M_d$.  Let $v_d$ be evaluation of the original basis
monomials at $y$.  Over any field, the determinant polynomial is
nonzero and the polynomial ring is an integral domain.  Thus clearing
its power does not change whether a pullback is zero after localization.
The separating polynomials are exactly vectors $c$ such that
$M_dc=0$ and $v_dc\neq0$.

For a matrix over a field, a row vanishes on its kernel if and only if
it lies in its row space.  To see the nontrivial implication without an
extra hypothesis, a linear functional that vanishes on $\ker M_d$
factors through $\operatorname{im}M_d$; extend that functional to the
whole target by extending a basis.  This gives the desired row-space
expression.  Appending $v_d$ raises rank precisely when it fails to lie
in that row space.  This proves (B7.2).

For each of the two integer matrices choose a nonzero minor of size
its rational rank; if that rank is zero, use the empty minor $1$.
Let $D_d$ be their product.  A prime not dividing $D_d$ preserves these
two nonzero minors, so each reduced rank is at least its rational rank.
Every larger minor is the zero integer, so each reduced rank is at most
its rational rank.  Both ranks are equal as asserted.  Rank does not
change upon field extension, so the argument also applies to algebraic
closures.\hfill$\square$

\subsection*{B7.3. Bad-prime noncontainment need not lift}
\paragraph{Statement and proof.}
Let $\mathbb G_m$ act on $\mathbb A^1_{\mathbb Z}$ by scaling,
and take $x=2$, $y=1$.  Over $\overline{\mathbb Q}$ the orbit of $2$
is the nonzero locus, whose closure is all of $\mathbb A^1$; therefore
$y$ belongs to it.  In characteristic two, $x$ becomes zero, its orbit
closure is $\{0\}$, and $y=1$ lies outside.  This is a counterexample
to lifting arbitrary finite-characteristic noncontainment without a
good-reduction certificate.\hfill$\square$

\paragraph{Obstruction gate for bridge 7.}
B7.2 reduces a bounded-degree geometric separation to two concrete
linear algebra computations, and supplies a precise good-prime
certificate.  The unresolved GCT part is to obtain a sufficiently
small separating degree and a succinct, efficiently usable description
of the relevant matrices or representation summands for the growing
determinant/permanent family.  The monomial space has
$\binom{\operatorname{rank}W+d}{d}$ columns, so explicit expansion
alone does not supply a polynomial-time flip.  In positive
characteristic, one must also specify whether one is comparing formal
characters, Weyl filtrations, composition factors, or actual direct
summands; characteristic-zero semisimplicity cannot simply be reused.
Hilbert and Quot representability records families and boundaries, but
does not assert the needed degree bounds or multiplicity inequalities.


\section{Which analytic hypothesis is actually used?}
\label{sec:complexity-grh}

In this section \(A_{\rm Dir}\) denotes precisely the supplied theorem
\texttt{LFunction\_eq\_zero\_re} in \texttt{ALPOGE\_Main.lean}, whose type is
\texttt{DirichletLZerosOnLine} in \texttt{ALPOGE\_Statements.lean}.
It quantifies over Dirichlet characters of every positive modulus and zeros in
\(0<\Re(s)<1\).  The stronger nonvanishing statement
\texttt{LFunction\_ne\_zero\_of\_half\_lt\_re} has type
\texttt{DirichletLNonvanishingRightOfHalf}.  Neither type quantifies over number
fields, their Dedekind zeta functions, or Hecke characters of such fields.
The following consequences use the supplied theorems as cited premises; they
do not assert that the missing imported Lean modules were rebuilt in this run.
The cited premise is Alp\"oge's supplied note and formal statements
\cite{CAlpoge26}.

\subsection{The part that passes to number fields}

\begin{proposition}[C-GRH-01: the cyclotomic-field transfer]
Assume \(A_{\rm Dir}\).  Let \(K\) be a subfield of a cyclotomic field
\(\mathbb Q(\zeta_N)\).  All zeros of \(\zeta_K(s)\) in
\(0<\Re(s)<1\) lie on \(\Re(s)=1/2\).
\end{proposition}
\begin{proof}
Put \(G=\operatorname{Gal}(K/\mathbb Q)\), a finite abelian group, and let
\(X_K\) be its character group, identified with the primitive Dirichlet
characters inducing the characters of this quotient of
\((\mathbb Z/N\mathbb Z)^\times\).  We verify the factorization
\[
 \zeta_K(s)=\prod_{\chi\in X_K}L(s,\chi).
\]
At a rational prime \(p\), let \(I\subseteq G\) be inertia, let
\(e=|I|\), and let the Frobenius element of \(G/I\) have order \(f\).
There are \(g=|G|/(ef)\) primes above \(p\), each of residue degree \(f\).
The Dedekind local factor is therefore \((1-p^{-fs})^{-g}\).
The Dirichlet factors that have a local factor at \(p\) are exactly those
characters trivial on \(I\).  On the cyclic subgroup generated by Frobenius,
their values run through all \(f\)-th roots of unity, each \(g\) times.
Thus their product is also \((1-p^{-fs})^{-g}\).  Equality of the absolutely
convergent Euler products proves the displayed identity for \(\Re(s)>1\),
and meromorphic continuation proves it everywhere.  In the open critical
strip each factor is holomorphic; the possible principal-character pole is
at \(s=1\), outside this strip.  A zero of their finite product must be a
zero of one factor, and \(A_{\rm Dir}\) puts that zero on the critical line.
\end{proof}

This is a genuine enlargement of the directly named family, but not an
extension to arbitrary number fields.  The standard Kronecker--Weber theorem
identifies these fields with the finite abelian extensions of \(\mathbb Q\);
the proposition was stated for cyclotomic subfields so that its proof needs
only the explicit character dictionary.  This dictionary is also explained
in Bach \cite[\S6, p.~369]{CBach90}.

\subsection{Koiran's theorem and its remaining gate}

The exact result needed is Koiran's \emph{Hilbert's Nullstellensatz is in the
polynomial hierarchy}, DIMACS long version, Theorem~2: under the paper's
analytic hypothesis, \(\mathrm{HN}\in\mathrm{AM}\).  Here HN asks whether
integer polynomial equations have a common complex zero.  The analytic
ingredient is \S5, Theorem~9 and Corollary~1, for an \emph{arbitrary}
irreducible \(f\in\mathbb Z[X]\).  If \(d=\deg f\), \(\Delta\) is its
discriminant, and \(W_f(p)\) is its number of roots modulo \(p\), Theorem~9
bounds the discrepancy
\[
 \left|\sum_{\substack{p\le x\\p\nmid\Delta}}
                 (1-W_f(p))\right|
 =O\!\left(x^{1/2}\log(|\Delta|x^d)\right).
\]
Koiran explicitly says that an effective prime ideal theorem suffices.  The
field is \(\mathbb Q(\alpha)\) for a primitive algebraic solution coordinate;
it need not be abelian.  The 8 August 1997 erratum corrects the quantitative
bound in Theorem~8, without changing Theorems~1 or~2
\cite[Theorem~2; \S5, Theorem~9, Corollary~1; erratum]{CKoiran96}.

Consequently \(A_{\rm Dir}\) discharges this input for cyclotomic subfields,
by C-GRH-01, but the supplied statement does not instantiate the theorem for
general HN.  This is a \emph{scope gap}, not a proof of the negation of
\(A_{\rm Dir}\Rightarrow\mathrm{HN}\in\mathrm{AM}\).  The exact unresolved
gate is a uniform effective prime-ideal estimate of the above strength for
the arbitrary primitive-element fields arising in the reduction.  One may
use a weaker suitable analytic statement instead of full Dedekind GRH:
Rojas \cite[Theorem~1 and Theorem~2]{CRojas03} gives the chain
\[
 \text{Dedekind GRH}\Longrightarrow\mathrm{DZH}
 \Longrightarrow\mathrm{GIPIT}\Longrightarrow\mathrm{HN}\in\mathrm{AM}.
\]
Neither DZH nor GIPIT for these arbitrary fields is supplied by a theorem
about Dirichlet characters alone.

\begin{proposition}[C-GRH-03: an HN instance outside the abelian-field restriction]
The feasible Hilbert--Nullstellensatz instance \(X^3-2=0\) has no
solution in any finite abelian extension of \(\mathbb Q\).
Consequently, feasible HN instances cannot all be handled by restricting
their solution fields to those covered by C-GRH-01.
\end{proposition}
\begin{proof}
The polynomial has a real root, so the instance is feasible.
Eisenstein's criterion at \(2\) proves that it is irreducible over
\(\mathbb Q\).  For any of its roots \(\alpha\), therefore,
\([\mathbb Q(\alpha):\mathbb Q]=3\).
If this degree-three extension were Galois, the irreducible polynomial
would split in it.  Taking the ratio of two distinct roots would put
a primitive cube root of unity \(\zeta_3\) in \(\mathbb Q(\alpha)\).
But \([\mathbb Q(\zeta_3):\mathbb Q]=2\), contradicting the tower law.
Thus no root field is Galois.  Every intermediate field of a finite
abelian Galois extension of \(\mathbb Q\) is Galois, since every
subgroup of an abelian group is normal.  No root \(\alpha\) can belong
to such an extension, proving the assertion.
\end{proof}

This example excludes the proposed restriction of solution fields;
it does not prove that Dirichlet GRH cannot have other consequences
for HN.

\subsection{Noether normalization versus identity testing}

GCT~V uses h.s.o.p. for a homogeneous system of parameters, s.s.o.p.
for a small system of parameters, and e.s.o.p. for an explicit system
of parameters; explicitness includes polynomial-time construction
\cite[\S1.2]{CGCTV}.
There are two separate inputs in GCT~V.  Its Theorem~2.10 invokes Koiran.
This is used for the conditional PH bound for constructing an h.s.o.p.
(Theorem~3.4), and for the conditional EXPH bound for constructing or
verifying an s.s.o.p. of the succinct determinant variety
(Theorem~4.10).  These inherit the number-field gate just described.
Theorem~4.1, concerning an h.s.o.p. for that succinct variety, has the further
explicit hypothesis that its defining equations can be computed in
exponential time.  Even full GRH would not erase that additional hypothesis
\cite[Theorems~2.10, 3.4, 4.1, 4.10]{CGCTV}.

The derandomization results use different hypotheses.  Theorem~4.9(a) states
that a polynomial-time construction of an e.s.o.p. for
\(\Delta[\det,m]\) is equivalent to \emph{strengthened} black-box
derandomization of symbolic determinant identity testing.  Strengthened
means a hitting set that also hits nonzero limits of the prescribed
small-circuit family.  Theorems~1.7 and 4.7 obtain a quasi-polynomial
construction from subexponential \emph{border} determinantal lower bounds
for the permanent.  Theorem~5.11 gives the corresponding statements for
general explicit varieties, using strengthened low-degree PIT or an
explicit integral multilinear family with the stated arithmetic border
lower bound.  No GRH hypothesis substitutes for these circuit or hitting-set
hypotheses.  Conversely, these equivalences and reductions are already
proved without GRH.  The remaining gate is the specified hitting set or
border lower bound, not an unproved occurrence of the word ``GRH''
\cite[Theorems~1.7, 4.7, 4.9, 5.11]{CGCTV}.

\subsection{Miller and Bach: the direct discharge}

\begin{proposition}[C-GRH-02: the witness bounds under the supplied theorem]
Assume \(A_{\rm Dir}\) and the supplied Riemann-hypothesis theorem.
The conditional hypotheses in Miller's deterministic primality test and in
Bach's rational-field bounds are discharged.  In particular an odd
composite integer \(m\) has a Miller compositeness witness at most
\(2(\ln m)^2\).  The integers below \(3(\ln m)^2\) that are coprime to
\(m\) generate \((\mathbb Z/m\mathbb Z)^\times\).
\end{proposition}
\begin{proof}
Bach's \S6 explicitly restricts its analysis to Dirichlet characters and
the Riemann zeta function.  Theorem~2, p.~372, bounds the first positive
integer outside any proper subgroup, allowing a nonunit, by
\(2(\ln m)^2\).  Theorem~3 bounds the first \emph{unit} outside any
proper subgroup by \(3(\ln m)^2\).  The consequences for witnesses and
generators are stated on pp.~372--373.  Their analytic assumptions are
exactly covered by the supplied theorems, so these cited results apply.
For the generator assertion specifically, the subgroup generated by the
specified units cannot be proper: Theorem~3 would supply a unit outside
it already among its generators, a contradiction.
\end{proof}

Miller's original polynomial-time conclusion is Theorem~2, with the
analytic hypothesis defined in his Appendix \cite{CMiller76}; Lenstra's
simplified criterion needs only a certain family of quadratic Dirichlet
characters \cite[Theorem~1]{CLenstra79}.  Ordinary deterministic
polynomial-time primality testing was already unconditional by AKS
\cite[Theorems~4.1 and 5.1]{CAKS04}.  What the supplied theorem adds here
is the validity of these specific small-witness guarantees.  Bach's
Theorem~4 concerns Hecke characters over general number fields and retains
its broader analytic gate; it must not be bundled with Theorems~2 and~3.

\section{The supplied counting theorem: quantitative consequences and limits}
\label{sec:complexity-eth}

Write \(A_{\#\mathrm{SAT}}\) for OpenAI's supplied Theorem~1.1 in
\path{LIT_OpenAI_ETH_is_False/main_full_version.tex}.  Its precise
bound is, for each fixed \(k\ge2\),
\[
 (n+m)^{O_k(1)}2^{O_k(n\log\log n/\sqrt{\log n})},
\]
with polynomial space.  For every fixed \(D\), when \(m\le Dn\), the
exponent improves to \(O_{k,D}(n/\sqrt{\log n})\).  The source is
\cite[Theorem~1.1; Appendix~D]{COAI26}.  These are bit-operation bounds,
not unit-cost bounds for symbolic arithmetic.

\subsection{ETH, SETH, counting, and the direction of reductions}

\begin{proposition}[C-SAT-01: the quantifiers that refute SETH]
\(A_{\#\mathrm{SAT}}\) refutes the usual deterministic ETH, \(\#\mathrm{ETH}\),
SETH, and their bounded-error randomized versions.  It also refutes the
usual Gap-ETH and NSETH.  A uniform bound in growing \(k\) is not needed
for this statement.
\end{proposition}
\begin{proof}
For fixed \(k\), the exponent divided by \(n\) tends to zero because
\(\log\log n/\sqrt{\log n}\to0\).  Testing whether the exact count is
positive decides satisfiability, and testing whether it is zero decides
unsatisfiability.  Taking \(k=3\) contradicts ETH and \(\#\mathrm{ETH}\).
Fix, for example, \(\varepsilon=1/2\).  For every fixed \(k\) its counting
algorithm, with its own constants and threshold, eventually runs in
\(2^{(1-\varepsilon)n}(n+m)^{O_k(1)}\) time; the finitely many smaller
lengths can be absorbed into its constant.  This is precisely the
negation of the SETH quantifier
\(\forall\varepsilon>0\ \exists k\) for which no such algorithm exists.
A deterministic algorithm is also a bounded-error randomized algorithm,
and also a nondeterministic one.  Its unsatisfiability version therefore
contradicts NSETH.  On the Gap-ETH promise, satisfiability distinguishes
the two promised cases, giving the last assertion.
\end{proof}

The conclusions of C-SAT-01 are already stated in the supplied paper,
\S1.1.1--1.1.2; the short proof above records the quantifier issue explicitly.
An implication \(\mathrm{ETH}\Rightarrow H\) remains a true conditional
theorem when ETH is false.  It no longer establishes \(H\), but it also
does not establish \(\neg H\).  The same applies to ETH-based
\(f(k)n^{o(k)}\) exclusions and to SETH-based fine-grained lower bounds.
For example, refuting SETH does not refute the Orthogonal Vectors
Conjecture merely because SETH implies that conjecture.  Nor does a
fixed-width guarantee settle the stronger hypothesis about unrestricted
CNF-SAT with width growing with \(n\).  These distinctions are already
made in \cite[\S1.1.3]{COAI26}.

\begin{proposition}[C-COUNT-01: a precise counting transfer]
Let \(V_x\) be a given Boolean verifier with \(w\) free input bits and
\(s\) gates of fan-in at most two, with all other inputs fixed by \(x\).
Under \(A_{\#\mathrm{SAT}}\), its number of accepting inputs can be computed
in
\[
 (|x|+w+s)^{O(1)}
 2^{O((w+s+1)/\sqrt{\log(w+s+2)})}
\]
bit operations and polynomial space.  The constants are absolute for
the fixed gate basis.  The construction is parsimonious.
\end{proposition}
\begin{proof}
Introduce one variable for each free input and one for each gate value.
For a two-input gate computing \(z=g(a,b)\), there are four forbidden
triples \((a,b,z)\), one for each choice of \((a,b)\).  For each forbidden
triple, insert the clause that is false precisely at that triple.
These clauses have width at most three.  A unary or constant gate has
at most four such clauses too, after substituting fixed bits.  Repeated
inputs simply produce shorter clauses or tautologies.  Pin the output
to one with a unit clause.  The formula has \(N=w+s\) variables and
at most \(4s+1\le5(N+1)\) clauses.  Once its free inputs are assigned,
the circuit's topological order determines all gate variables uniquely.
The clauses hold exactly for that assignment to gate values, and the
output condition holds exactly for accepting inputs.  This proves
parsimony.  The bounded exceptional cases \(N<20\) are enumerated.
Otherwise the fixed-density part of the supplied theorem, with \(k=3\)
and a fixed density bound, proves the stated time and space estimates.
\end{proof}

If a representation-theoretic coefficient already has a positive counting
formula with a verifier of \(B\) Boolean inputs and gates, C-COUNT-01
computes it in \(\operatorname{poly}(B)2^{O(B/\sqrt{\log B})}\).
The existence of the formula is an explicit hypothesis: this observation
does not construct a \(\#\mathrm P\) interpretation for plethysm or
Kronecker coefficients.  If \(B=O(L^a)\) in an input parameter \(L\),
the certified exponent is \(O(L^a/\sqrt{\log L})\), not automatically
\(o(L)\), and not \(O(\log L)\).  Thus polynomial reductions preserve
polynomial-time membership but need not preserve subexponential time
in the original parameter.  A target-to-SAT reduction is required to
obtain an algorithm for that target; the SAT-to-target reductions used
to prove hardness point in the other direction.

\begin{corollary}[C-COUNT-03: the parsimonious completeness form]
Counting satisfying assignments of width-at-most-three CNF formulas is
\(\#\mathrm P\)-complete under polynomial-time parsimonious reductions.
This statement does not require \(A_{\#\mathrm{SAT}}\).
\end{corollary}
\begin{proof}
Use the witness definition of \(\#\mathrm P\): a function counts the
strings \(y\in\{0,1\}^{p(|x|)}\) accepted by a fixed polynomial-time
verifier on \((x,y)\).  For the equivalent accepting-computation
definition, pad the unused nondeterministic choices by zeros and have
the verifier reject nonzero padding; this preserves the count.
Let \(T\) bound the verifier's running time and input length, and freeze a halting
configuration thereafter.  A standard one-tape simulation may enlarge
\(T\) polynomially.  Encode each tape cell by a fixed number of bits
for its symbol, head marker and finite control state.  In each time
step a cell's new data depend only on it and its two neighbors.
The finitely many local update rules have constant-size fan-in-two
Boolean circuits.  A tableau with \(T\) layers of \(O(T)\) cells
therefore gives an \(O(T^2)\)-gate circuit, constructed in polynomial
time, with \(x\) fixed and the witness bits free.  An OR tree over
the final cells tests the accepting state.  Starting configurations
are valid, so behavior of the local rules on invalid configurations
is irrelevant.  The clauses in C-COUNT-01 now give a width-three
formula with exactly one satisfying extension per accepting witness.
Conversely a satisfying assignment is itself a polynomial-length,
polynomially checkable witness.  This proves both hardness and membership.
\end{proof}

\begin{corollary}[C-SIZE-01: a concrete width and size accounting]
Under \(A_{\#\mathrm{SAT}}\), a CNF with \(n\) declared variables and
\(L\) literal occurrences can be counted in
\(\operatorname{poly}(n+L)2^{O((n+L)/\sqrt{\log(n+L+2)})}\), using
polynomial space, regardless of its width.  In particular this is
\(2^{o(n)}\operatorname{poly}(n+L)\) when
\(L=O(n(\log n)^a)\) for a fixed \(0\le a<1/2\).
\end{corollary}
\begin{proof}
Use binary OR trees for clauses, a binary AND tree for their conjunction,
and at most one negation gate per negative occurrence.  The resulting
verifier has \(n\) free inputs and \(O(L+1)\) gates; empty clauses and
empty conjunctions are handled directly.  C-COUNT-01 gives the first
bound.  Under the additional assumption,
\(\log(n+L+2)=\Theta(\log(n+2))\) and the exponent is
\(O(n/(\log n)^{1/2-a})=o(n)\), proving the second assertion.
\end{proof}

This also shows concretely why the cited theorem alone does not give the
standard uniform variable-width strong-SAT conclusion: general instances
need not satisfy this incidence bound.

\subsection{A signed, weighted permanent consequence}

The supplied Boolean contraction theorem is stronger than the decision
consequence alone.  Its exact interface is Theorem~D.1 in
\texttt{quantitative-reduction.tex}: the coefficient ring is
\((\mathbb Z/2^{P_0}\mathbb Z)[V]/(V^\ell)\), with \(P_0,\ell=O(s)\);
there are \(O(s)\) Boolean wires; maximum degree is at most
\(\lceil s^{9/10}\rceil\); and local entries have the stated uniform
\(\operatorname{poly}(s+\mathcal L)2^{O(s^{9/10})}\) time and polynomial
space descriptions.  The conclusion is
\(\operatorname{poly}(s+\mathcal L)2^{O(s/\sqrt{\log s})}\) time and
polynomial space.  We denote this cited result by \(A_{\rm contract}\).

\begin{proposition}[C-COUNT-02: signed sparse permanent]
Assume \(A_{\rm contract}\).  Let \(A\) be an \(n\times n\) integer
matrix with \(E\) nonzero entries, each of absolute value at most \(2^b\),
where \(b\ge0\).  Put \(S=\max\{2,E+n(b+1)\}\).  Its exact permanent
can be computed deterministically in
\[
 (S+\mathcal L)^{O(1)}2^{O(S/\sqrt{\log S})}
\]
time and polynomial space, where \(\mathcal L\) is its input-description
length.  In particular, bounded integer weights and \(E=O(n)\) give
\(2^{O(n/\sqrt{\log n})}\) time up to a polynomial factor.
\end{proposition}
\begin{proof}
Return zero if a row or column has no nonzero entry; handle \(n=0\)
by the empty permanent, equal to one.  Otherwise introduce a Boolean
selection variable for each supported edge of the row-column bipartite
graph.  At each row and each column impose that exactly one incident
edge is selected.  Realize such a constraint of degree \(d\) by a chain
of \(d\) ternary tensors
\[
 T_j(u,x,v)=\begin{cases}1&u+x=v\text{ as integers},\\0&\text{otherwise},\end{cases}
 \qquad u,x,v\in\{0,1\},
\]
with initial state pinned to zero and final state pinned to one.
The state after each step is the number of selections so far.  There
is a unique state assignment for a selection with exactly one edge,
and no state assignment otherwise.  On a supported edge insert a
binary diagonal tensor with values \(1\) on \((0,0)\), \(a_{ij}\)
on \((1,1)\), and zero on the two unequal pairs.
The contraction is now precisely \(\operatorname{perm}(A)\): every
perfect matching contributes the product of its weights once.

All vertices have degree at most three, and the graph has \(O(E+n)\)
edges and vertices.  Subdivide every edge by a binary equality tensor.
The old and the new vertices are its two bipartition classes; equal
bits extend uniquely across a subdivision, so the contraction is
unchanged.  The size remains \(O(E+n)\).

Let \(d_i\) be the number of supported entries in row \(i\).  The
number of perfect matchings is at most \(\prod_i d_i\le2^E\), since
\(d_i\le2^{d_i}\) and \(\sum_i d_i=E\).  Each matching weight has
absolute value at most \(2^{nb}\), whence
\(|\operatorname{perm}(A)|\le2^{E+nb}\).
Use modulus \(2^P\) with \(P=E+nb+2=O(S)\), and take \(\ell=1\).
The local tensors can be evaluated in polynomial time and space,
including their integer residues.  They satisfy every hypothesis of
\(A_{\rm contract}\) with size parameter \(S\); the bounded small
exceptions to the degree inequality can be enumerated directly.
After contraction, choose the representative in
\([-2^{P-1},2^{P-1})\).  The displayed absolute bound is at most
\(2^{P-2}\), so this representative is the exact signed permanent.
The stated costs follow from the contraction theorem.
\end{proof}

For dense bounded-entry matrices this construction has \(S=\Theta(n^2)\).
It certifies an exponent \(O(n^2/\sqrt{\log n})\), which is larger
than the \(O(n)\) exponent given by the \(2^{n-1}\)-term formula in
Landsberg \cite[\S8.2, equation~(22)]{CLandsberg13}.
It therefore supplies no improved dense-permanent bound.
The signed weighted statement is an extension of the unweighted
perfect-matching case already given in \cite[Corollary~1.3]{COAI26}.
It concerns evaluation on integer inputs, not polynomial-size
arithmetic circuits for the symbolic permanent, and not membership
or nonmembership in a determinant orbit closure.

\subsection{A bounded finite-field PIT transfer}

\begin{proposition}[C-PIT-01]
Assume \(A_{\#\mathrm{SAT}}\).  Let a division-free arithmetic circuit
over a specified prime field \(\mathbb F_p\) have \(r\) inputs,
\(s\) arithmetic gates, and output polynomial of total degree \(D<p\).
Put \(b=\lceil\log_2 p\rceil\) and
\(B=\max\{2,(s+r+1)b^2\}\).  Its polynomial identity can be tested
deterministically in \(\operatorname{poly}(B)2^{O(B/\sqrt{\log B})}\)
bit operations and polynomial space.
\end{proposition}
\begin{proof}
Use \(rb\) Boolean input bits to encode a field point uniquely, rejecting
any coordinate encoding at least \(p\).  Schoolbook addition,
multiplication and modular reduction on \(b\)-bit integers have
fan-in-two Boolean circuits of \(O(b^2)\) gates.  Replacing every
arithmetic gate by these circuits, and testing whether the output is
nonzero, gives a Boolean verifier of \(O(B)\) inputs and gates.
C-COUNT-01 counts its accepting inputs within the claimed bounds.
The zero polynomial has no accepting input.  A nonzero polynomial
of degree \(D<p\) cannot vanish on all of \(\mathbb F_p^r\).
For completeness, the latter assertion follows by induction on \(r\):
write the polynomial in its last variable; if it vanishes at every
point then, at every fixed choice of the other variables, this
univariate polynomial of degree less than \(p\) has \(p\) roots,
so all its coefficients vanish.  Apply the induction hypothesis to
each coefficient, whose individual degrees are still less than \(p\).
Thus the exact count is positive if and only if the polynomial is
nonzero, completing the test.
\end{proof}

C-PIT-01 makes the model boundary explicit.  It supplies neither a
polynomial-time hitting set, nor characteristic-zero identity testing
with arbitrary coefficient growth, nor a hitting set for the
approximative families in GCT~V.  The exponent depends on Boolean
evaluation size.  Furthermore \(2^{O(B/\sqrt{\log B})}\) is not the
\(\bigcap_{\varepsilon>0}\operatorname{NTIME}(2^{B^\varepsilon})\)
bound called NSUBEXP in arithmetic hardness-versus-randomness results;
one must not substitute one convention for the other.

\subsection{What the contraction method says about positivity}

The source uses collision packing (Lemma~D.3), integral local
coordinates (Lemma~D.4), and tapered Hermite recovery (Lemma~D.5).
At depth \(d\), a network of \(M=O(s8^{-d})\) wires uses
\(u=\Theta((\log s)/2^d)\) retained valuation levels.  Weighted modes
are discarded because their coefficients are divisible by the
working power of two, and the remaining exponent counts sum as
\(O(s/\sqrt{\log s})\).  These are divisibility and interpolation
statements; the source proves no ordered-sign or Schur-positive
interpretation for them \cite[Appendix~D.4]{COAI26}.

\begin{proposition}[C-SIGN-01: divisibility is not a positivity certificate]
The existence of an \(A_{\rm contract}\)-admissible network does not
imply nonnegativity of its integer contraction.  Nor does nonnegativity
of its canonical residue modulo \(2^P\) imply nonnegativity of the
integer it represents.
\end{proposition}
\begin{proof}
Take a network with one edge, a unary tensor constantly equal to one
at one endpoint, and the unary tensor \((-1,0)\) at the other.  Its
contraction is \(-1\); its entries, degrees, size and entry costs
satisfy the theorem for every fixed \(P\), after increasing the
size parameter if necessary.  Its canonical residue is \(2^P-1\),
a positive integer.  More structurally, a nonzero finite ring cannot
carry an additive order with \(1>0\): repeated addition would give
\(0=2^P\cdot1>0\) in \(\mathbb Z/2^P\mathbb Z\).
Thus a separate positivity interpretation is indispensable.
\end{proof}

The theorem's evaluation capacity can therefore pass to a proposed
representation-theoretic counting or tensor model once that model and
its resource bounds are proved.  It does not build that model,
establish saturation, construct a canonical basis, identify an
intersection-cohomology stalk, or compare the permanent and
determinant multiplicities.  Those are the outstanding mathematical
interfaces.

\subsection{Circuit lower bounds and already considered algebraic bridges}

The supplied paper already proves the quantitative consequence
\[
 \exists L\in\mathrm E^{\mathrm{NP}}\ \exists c>0:
 C_L(n)>c n\sqrt{\log n}\quad\text{for infinitely many }n,
\]
where gates have fan-in at most two and unrestricted fan-out
\cite[Corollary~1.4; Appendix~C.1; \texttt{quantitative-circuits.tex}]{COAI26}.
The language is in \(\mathrm E^{\mathrm{NP}}\), the bound is
superlinear rather than superpolynomial, and the model is Boolean.
This supplies none of the three missing changes needed for a direct
GCT separation: an appropriate algebraic family, a sufficiently large
arithmetic lower bound, and a border/orbit-closure certificate.

A later paper already addresses this bridge.  Chukhin--Gajulapalli--Golovnev--Mihajlin,
ECCC TR26-313, Theorem~1, proves the independent disjunction
\[
 \mathrm{NE}\not\subseteq\mathrm{SIZE}[2^n/(2n)]
 \quad\text{or}\quad
 \mathrm{VP}^{0}_{\mathbb Z}\ne\mathrm{VNP}^{0}_{\mathbb Z}.
\]
Here \(0\) means constant-free circuits of polynomial formal degree,
a different model from arbitrary complex constants.  Their Theorem~6
turns sufficiently small integer arithmetic circuits for the permanent
into perfectly complete input-oblivious Merlin--Arthur protocols for
sparse 3-UNSAT, with running time \(s(Dn)n^D\), for a fixed \(D\)
and the size bound \(s\) satisfying the theorem's uniformity conditions.
Their discussion on p.~9 explicitly observes that the supplied
\(2^{O(n/\sqrt{\log n})}\) deterministic solver is an empty-message
protocol, and compares it with the stronger
\(2^{O(n/\log n)}\) protocol gate \cite[Theorems~1 and~6; pp.~8--9]{CCGGM26}.
Thus the exact remaining contrapositive gate is a lower bound against
the specified protocols at a \emph{faster} time scale.  Refuting ETH
does not establish this gate.

The same source gives a concrete positive zero-test: Theorem~5 uses
the sum of squares of coefficients of a syntactically set-multilinear
integer circuit.  With \(d\) variable blocks of size at most \(m\),
\(s\) gates, and every gate polynomial having coefficient
\(\ell_1\)-norm at most \(2^L\), its running time is
\(m^{d-\lceil d/6\rceil}(dsL\log(2m))^{O(1)}\).
This promised positive-scalar test supplies no plethysm formula
\cite[Theorem~5, \S8.2]{CCGGM26}.

\paragraph{Status of the major separations.}
None of the deductions in this section proves \(\mathrm P=\mathrm{NP}\),
\(\mathrm P\ne\mathrm{NP}\), \(\mathrm{VP}=\mathrm{VNP}\), or
\(\mathrm{VP}\ne\mathrm{VNP}\) in the characteristic-zero model of
the question.  This sentence reports the scope of the proved deductions;
it is not a claimed logical independence theorem.  In particular,
\(\exp(O(n/\sqrt{\log n}))\) is a superpolynomial \emph{bound},
so it does not certify polynomial time, and an algorithmic upper
bound does not supply a superpolynomial lower bound for the same
problem.  The remaining gates are recorded individually in
\texttt{QUESTIONS.tsv}; none is declared refuted merely because the
supplied theorem has a different type.


\part{Unchecked targets and remaining gates}
\label{part:open}
\section{What has not been proved}
There is no unchecked finding offered as a theorem in this part. The
following are the unresolved targets of the task. Their state is a
current gate, not a proved negative answer. A precise counterexample to
one proposed transfer is recorded separately from a proof that no other
transfer can exist.

\begin{description}
\item[O-PH2-SH: direct bounds for irreducible-inner plethysm.]
Both clauses of ordinary plethysm PH3 are now refuted: P-F19 gives
an exponential canonical parameter degree, and P-F30--P-F31 prove
actual simple poles of exponential order in the same cubic family.
The complete support classification closes the singular-boundary
and possible-cancellation gates of P-F23--P-F29. The first PH3
clause therefore cannot serve as the premise for a proof of PH2,
although the conditional implication P-F11--P-F12 remains valid.

The unresolved question is a direct proof or refutation of the
narrower polynomial positivity-index and saturation-index bounds,
with the source's exact parameter restrictions and constituent
polynomial definitions. A large quasipolynomial period does not by
itself force either index to be large. P-F32 supplies a precise
exterior-power complement identity and an infinite family with
first positive stretch $\Theta(\sqrt n)$; it gives only a polynomial
index lower bound. It does not settle SH(a), its separate density
clause SH(b), or PH2. P-F14--P-F15 already refute the broader
rank-only subgroup-restriction versions, a different formulation.

P-F34 makes a direct sufficient gate precise. Let $B_n$ be the
average constituent polynomial, determined by the principal part
of $H_n$ at one, and let $U_n$ be its coefficient $\ell^1$ norm.
For a specified primitive root, the proved Fourier amplitude satisfies
$\log_2|A_n|=n^4/24+O(n^3)$. Positivity or saturation
after a shift forces $B_n$ at that shift (or one more) to dominate
this amplitude. The additional bound
$\limsup (\log_2 U_n)/n^4<1/24$ would therefore refute polynomial
PH2 and SH(a). This bound remains unproved; the Fourier inequality
is necessary, not an equivalence to the index conjectures.
P-F42 proves an unconditional actual-family upper bound
$\limsup(\log_2 U_n)/n^4\le\log_2(5)/12$; its constant
is larger than $1/24$. Its coefficient-sensitive refinement
(P-F42z)--(P-F42ah) proves that the highest fifth of the
coefficient indices already has logarithmic leading constant
at most $\log_2(5)/60<1/24$. The sufficient norm gate is
equivalent to the same strict bound for the complementary
lower coefficients only. P-F44 identifies the exact global
primitive-Klein contribution as a signed cyclotomic trace and
gives an explicit complementary balanced-cone expression.
Its odd-$D$ contribution cancels, and even-$D$ signs can mix.
The remaining task is therefore a signed coefficient estimate,
not a bound on a single local amplitude.
The refined bound also removes the high-degree tail uniformly
on every polynomially bounded interval. Thus any polynomially
bounded strict index must make the signed lower polynomial in
(P-F42ai)--(P-F42ak) reach almost the full Fourier amplitude.
P-F47 identifies exterior-volume powers with nonvanishing at the
Fermat cubic. P-F49 constructs such invariants in polynomial
degree, with a sharper explicit bound from the supplied Dirichlet
theorem. P-F50 gives a polynomial eventual-support threshold on
every stretch allowed by that point's period. Combining this
with the odd top-exterior-power invariant gives polynomial
full-ring eventual support on the infinite family in
(P-F50s)--(P-F50t). These results
control realized point evaluations; the source's strict indices
also test constituent polynomials away from their residue classes.
P-F54 now proves unconditional polynomial eventual actual support
for every $n\ge33$, by a specified exactly balanced cubic and
quantified compact averaging. Every coefficient degree
$d\ge(10n)^{200}$ with $n\mid3d$ occurs at this point. Thus
the missing strict-index estimate cannot be attributed to an
unproved eventual-support bound on this cubic family; its
remaining content is the off-residue constituent tests.
P-F48 proves that the cubic family and its natural determinant
twists are exponentially sparse in each of the two stated SH(b)
input populations, so this amplification does not settle SH(b).
P-F38 excludes deriving this bound from small matrix entries, bounded
coordinates and cubic-scale dimension alone, even after normalizing by
the row rank. The actual cubic weight system and Weyl alternation remain
necessary extra structure. P-F33 shows why normal rational Gorenstein geometry with fixed
negative $a$-invariant and fixed dimension cannot alone supply
such a bound: its genuine arbitrarily large periods coexist with
zero indices and an unbounded average constant term. P-F35 gives
the exact calibration $p(f_4)=s(f_4)=23$ and minimal period $120$.

\item[O-COUNT: uniform effective shifted counting.]
P-F2 constructs a single shifted polytope from explicit positive Hilbert
data. P-F12 controls their existence and size under a bounded
positive-form premise in expanded parameters. P-F31 shows that the
proposed numerical PH3 bound cannot supply that premise in general. P-F13 reconstructs
and verifies candidate data from a finite prefix under an explicit
denominator promise. A procedure constructing them with the original
partition-input complexity bounds remains to be given. A small object
whose existence is known separately for every instance is not automatically
a polynomial-time constructible object. The homogeneous version is already
refuted by P-K1; no revision should silently restore that false assertion.

\item[O-NSQ: nonstandard geometric realization.]
For the particular bases and products in GCT VII--VIII, one must construct
the category, identify the basis, control parity or shifts, establish
product and specialization compatibility, and supply size bounds.
P-F6 is a sufficient criterion after these data exist. The determinant
and theta-lift passages of the supplied note do not themselves identify
the relevant Hom spaces. P-F5 refutes the simplest attempted substitution
of ordinary interlacing for plethysm restriction.
P-F36 additionally proves that the actual cubic example has a
quadratic coefficient extension that survives every finite power
regrading. Its separate conjugate blocks cannot descend through
an ordinary rational projector, and its displayed raw coefficient
cannot be a finite graded Laurent character. A comparison must
retain that extension and its descent datum, or explicitly recover
the blocks from a conjugate sum. Arbitrary coefficient systems
and the source's relaxed minimal-polynomial conjectures remain
outside this obstruction.
P-F39 supplies an actual rational symmetrizable positive operator
for the two conjugate scalars, the exact source normalization of
its trace and norm, and positive Tate--Lefschetz models for those
two normalized coefficients. Before normalization its Laurent
coefficients have negative entries. P-F52 now constructs the
missing actual-algebra intertwiner for this operator: the short
central word $(Q_1-Q_2)^2+S-f^2$ gives precisely its matrix,
and an explicit positive tensor form proves positivity at every
real $q>0$. The complete relation and cell-element factorization
prove the scalar equation independently of the source's
inconsistent printed two-by-two generator matrices. Comparison
of the coefficient models with the geometric or theta objects,
their products, and the prescribed normalization remains open.
The explicit alternative normalization in (P-F52m)--(P-F52p)
shows that positive actual action and positive Lefschetz trace
and norm do not even force the source factor to divide the scale.
The source's actual spectral-interpolation rule does give an
exact positive factorization of that scale in (P-F52q)--(P-F52s).
Moreover (P-F52t)--(P-F52x) refute two literal common-object
dictionaries: the trace and norm cannot be the first two ordinary
lambda characters under any common scalar normalization, and
the prescribed normalized pair cannot be simultaneous traces of
one Frobenius and its exterior square over every finite extension.
These are precise excluded comparisons, not a general exclusion
of geometric or theta constructions.

\item[O-THETA: comparison with an identified invariant.]
P-F37 supplies the actual meromorphic automorphic ring of
five-variable cubics. P-F41 identifies its weight-$48$
Borcherds form with the ordinary degree-$80$ cubic discriminant,
including all products in the generated polynomial subalgebra.
The specified construction uses a regularized multiplicative
orthogonal theta lift, an $E_6$ quasi-pullback and restriction.
The next comparison must identify this named output with the
supplied iterated unitary lift or determinant, or produce an
identified larger graded collection with quantitative bounds.
The supplied note already mentions orthogonal groups; the group
name alone neither establishes nor excludes this comparison.
P-F46 identifies the exact positive input as the vector-valued
series $\Theta_{E_6}/\delta$, including the discriminant-group
gluing and principal coefficients $1,96,27,27$. Its ordinary
Petersson norm diverges. Its positive Fourier coefficient counts
grow faster than the actual $T^{(5)}$ Hilbert coefficients under
every positive affine reindexing. Thus the remaining comparison
must identify the regularized multiplicative operation or another
explicit coefficient-space construction; direct reindexing of
these input counts is refuted. Hofmann's unitary Borcherds
construction identifies the same form as a multiplicative
unitary lift of this input; the remaining comparison concerns
that exact regularized lift and the supplied iterated definite
unitary operations, not whether the output can be called unitary.
P-F51 further realizes the exact positive $E_6$ factor as an
ordinary convergent theta integral from a specified nonnegative
function on a definite $U(3)$ quotient. Its lattice cosets,
finite level, splitting on $U(3)\times SL_2$, individual-class
measure and $y^{-3/2}$ normalization are fixed. The compatible
right unitary group is split $U(1,1)$. The remaining comparison
must match the supplied definite chains to these specified
vectors and account for division by $\delta$ and the regularized
multiplicative operation; the positive factor itself is identified.
P-F53 realizes all coefficients of that exact input by finite
disjoint unions of K3 Hilbert schemes, with even cohomology
and purity in each cohomological degree. Its completed
lambda-ring identity supplies a geometric source for the
inverse-discriminant coefficient series. This still does not
identify the required regularized operation or a map to all
actual multiplicity spaces.
P-F55 identifies the square of the same orthogonal theta output
with the six-variable cubic discriminant, of coefficient degree
$192$, and matches its cyclic restriction with the square of
the five-variable discriminant. This adds an exact fixed-dimensional
compatibility; it does not supply the missing comparison with
the supplied iterated operations or an asymptotic family.
P-F45 proves that the full cyclic-cubic-cover eigenspace
classifying domain is a ball exactly for $n=4,5$. Beyond
this range an asymptotic construction must specify what
replaces that full-domain model; selected substructures,
special period images and other theta constructions remain open.

\item[O-OBSTRUCTION: an actual family of separating equations.]
In the no-occurrence range, prove a strict inequality between the actual
coordinate-ring multiplicities, or construct an explicit polynomial
vanishing on the determinant closure and not on the padded permanent.
The ambient plethysm and Kronecker multiplicities do not determine the
ranks of the quotient maps. P-F8 and PL-06 isolate this rank/comparison
requirement. No such asymptotic obstruction family is constructed here.

\item[O-DIR-HN: arbitrary solution fields in Nullstellensatz.]
C-GRH-01 proves only the cyclotomic-field transfer. Koiran's algorithm uses
a uniform prime-ideal/root-modulo-prime estimate for arbitrary primitive
solution fields. Rojas's DZH or GIPIT conditions are alternative sufficient
analytic inputs. The present supplied Dirichlet theorem does not state
these inputs. No formal nonimplication between the analytic conjectures
is proved.

\item[O-PIT: the derandomization required by GCT V.]
The exact algorithm or hardness hypothesis of each theorem must still be
provided. The supplied counting algorithm yields C-PIT-01 in a bounded
finite-field circuit model. It does not supply the characteristic-zero
black-box, border or strengthened symbolic determinant PIT assertions
used elsewhere. GRH does not replace those separate hypotheses.

\item[O-INTEGRAL: uniform bounds for specialization.]
B7.1 proves good-prime survival for a fixed integral separator; B7.2 gives
an exact rank test in any fixed degree and integer minors excluding bad
primes. An efficient GCT transfer needs a uniform separator-degree bound,
succinct rank data and a usable prime bound. Neither Hilbert-scheme
representability nor the pointwise specialization theorem supplies these
quantitative statements.

\item[O-COMPLEXITY: Boolean and algebraic separation.]
The ETH theorem refutes the precise exponential hypotheses treated in
C-SAT-01. Its supplied circuit corollary is superlinear for
$\mathsf E^{\mathsf{NP}}$, not superpolynomial for $\mathsf{NP}$.
The current constant-free integer Boolean--arithmetic disjunction cited
in the complexity section does not choose a branch and does not identify
the usual complex-field $\mathsf{VP}$ with its constant-free integer
variant. The stated upper bounds and conditional bridges do not settle
$\mathsf P$ versus $\mathsf{NP}$ or the target GCT obstruction.
\end{description}

\section{Approaches tried and the next discriminating steps}
The qualitative-sign route was tested by actual determinant projectors
and a virtual-character counterexample (P-F3--P-F4). Its next step is an
identity with the intended multiplicity projector together with a size
bound. The standard-branching route was tested on $\operatorname{Sym}^2$
(P-F5); a replacement must use the actual Schur-functor homomorphism.

The purity route was tested at three increasingly specific levels:
semisimple categorical structure (P-F6), pure finite schemes encoding
\#3SAT (P-F7), and the actual canonical plethysm algebra carrying the
obstructing Hilbert function, a positive multiplicative basis and pure
intersection cohomology (PLE-F7--PLE-F8). The next step is an
equivariant graded comparison theorem, not another proof that an unrelated
cohomology group is pure. P-F10 separately tests the missing quantitative
control of a grading.

The Hilbert-series route made positive progress: P-F1 constructs the
optimal shifted model on the obstructing ray; PLE-F5 realizes its normal
invariant ring; PLE-F6 classifies all rescalings that repair homogeneous
Ehrhart realization; P-F2 treats every explicit positive rational form;
P-F9 and P-F11 supply coefficient-positivity bounds; P-F12 connects the
last bound to the numerical positive-form clause of revised PH3.
P-F13 gives finite reconstruction under a denominator promise;
P-F16--P-F18 give finite fixed-embedding and fixed-dimension bounds;
P-F19 refutes polynomial growth of the canonical parameter bound.
P-F30--P-F31 subsequently close the numerical positive-form question
negatively, using a complete support theorem and a fixed-component
calculation. The next step is a direct PH2 or saturation argument,
together with uniform construction of a suitable revised model.
The ordinary pole is proved by these additional arguments, not
inferred from the h.s.o.p. bound alone.
P-F20 confirms that distinction with universal denominators and
fixed-target covariants. Its multivariate counterexample defeats a
transfer based only on quasi-polynomiality from an affine slice to
simultaneous rays. P-F21--P-F22 then defeat generic normal,
Cohen--Macaulay, rational and Gorenstein attempts to repair the
parameter-degree inference. P-F23--P-F24 made the geometric gate precise: the known chain
component is smooth, and the other fixed components had to be
controlled. P-F26 identifies the entire smooth fixed locus with
the chain orbit; P-F27 excludes every graph-support boundary
competitor, including rank-deficient supports and lifts of larger
two-power order. P-F29 handles the complementary cube-and-squarefree
case. P-F30 excludes all remaining mixed supports, and P-F31 applies
the noncancellation criterion. No competing nonzero component remains. The binary determinant subfamily is separately treated
in P-F25: its realized holes are bounded by thirteen, but the
source's strict saturation definition also tests off-residue
constituent values and is not settled by that table.

The archive route distinguishes the three products on representations,
separates the two Burnside operations, checks the transcribed identities
on small groups, and treats the modular and quadratic-form caveats
(PL-01--PL-08). Its obstruction gate is the actual rank of the
Foulkes--Howe multiplicity-space map. A Schur-positive character of its
source and target does not compute that rank.

The norm and approximation routes were tested against scaling, exact
zero detection, adjacent integer multiplicities, and an explicit rank-two
family approaching rank one in normalized norm (B4.1--B5.3). They do give
valid contraction-model obstructions under the stated operator bounds.
A full algebraic-orbit obstruction requires an invariant and a certified
separating gap that survive the allowed degenerations.

The arithmetic and ETH routes track the exact analytic family, all input
and reduction sizes, the direction of reductions, arithmetic precision,
and the class in which a lower bound is obtained. Their next steps are
specified in O-DIR-HN, O-PIT and O-COMPLEXITY. These are not resolved by
relabeling the supplied subexponential bound as polynomial time.

The exterior-volume route is now an exact Fermat-evaluation
dictionary (P-F47). A prime-characteristic determinant identity,
applied through normalized Hessians and Cayley's contraction,
gives the actual invariants of P-F49--P-F50. The remaining
strict-index step cannot be replaced by this point-specific
conductor. P-F48 separately closes the proposed density
amplification by ordinary determinant twists, with the exact
rank and partition-size conventions of each source clause.


\clearpage\begin{thebibliography}{99}

\bibitem{PGCTVI} K. D. Mulmuley, \emph{Geometric Complexity Theory VI: the flip via saturated and positive integer programming}, arXiv:0704.0229v4, supplied revised version. Problems 1.1.1--1.1.4; Definitions 1.2.2, 1.2.4, 1.2.7; Theorems 1.6.1--1.6.3; Hypotheses 1.6.4--1.6.7 and 3.3.2--3.3.4; Lemma 3.1.6; Theorems 3.4.1, 4.1.1, 4.4.1; Sections 4.1.1--4.1.2, equations (4.2)--(4.8). \url{https://arxiv.org/abs/0704.0229v4}.
\bibitem{PGCTVII} K. D. Mulmuley, \emph{Geometric Complexity Theory VII: Nonstandard quantum groups for the plethysm problem}, arXiv:0709.0749. Theorem 1.2; Conjectures 1.3--1.4. \url{https://arxiv.org/abs/0709.0749}.
\bibitem{PGCTVIII} K. D. Mulmuley, \emph{Geometric Complexity Theory VIII: On canonical bases for the nonstandard quantum groups}, arXiv:0709.0751. Conjectures 2.6--2.9 and 3.4--3.7. \url{https://arxiv.org/abs/0709.0751}.
\bibitem{PKM} T. Kahle and M. Micha\l ek, \emph{Obstructions to combinatorial formulas for plethysm}, Electron. J. Combin. 25(1) (2018), P1.41. Theorem 5, Remark 6, Proposition 11, the paragraph after Remark 13, Question 15 and its following paragraph. \url{https://www.combinatorics.org/ojs/index.php/eljc/article/download/v25i1p41/pdf}.
\bibitem{PBOR} E. Briand, R. Orellana and M. Rosas, \emph{Reduced Kronecker coefficients and counter-examples to Mulmuley's strong saturation conjecture SH}, arXiv:0810.3163. Theorem 2.4; Section 2.5; Mulmuley's appendix. \url{https://arxiv.org/abs/0810.3163}.
\bibitem{PIMW} C. Ikenmeyer, K. D. Mulmuley and M. Walter, \emph{On vanishing of Kronecker coefficients}, arXiv:1507.02955. Theorem 1.3. \url{https://arxiv.org/abs/1507.02955}.
\bibitem{PFI} N. Fischer and C. Ikenmeyer, \emph{The computational complexity of plethysm coefficients}, arXiv:2002.00788. Theorems 3.1--3.2. \url{https://arxiv.org/abs/2002.00788}.
\bibitem{PBIP} P. B\"urgisser, C. Ikenmeyer and G. Panova, \emph{No occurrence obstructions in geometric complexity theory}, arXiv:1604.06431. Theorem 1.4. \url{https://arxiv.org/abs/1604.06431}.
\bibitem{PWeilII} P. Deligne, \emph{La conjecture de Weil. II}, Publ. Math. IH\'ES 52 (1980), 137--252. Theorem 3.3.1, Corollaries 3.3.4--3.3.6 and 3.3.9, paragraph 3.3.11. \url{https://www.numdam.org/item/PMIHES_1980__52__137_0.pdf}.
\bibitem{PBBD} A. A. Beilinson, J. Bernstein and P. Deligne, \emph{Faisceaux pervers}, Ast\'erisque 100 (1982). Corollaries 1.4.24--1.4.25, p.~55; equations (4.0.1)--(4.0.2), p.~102; Theorems 5.3.8, 5.4.5, 6.2.5. \url{https://www.sas.rochester.edu/mth/sites/doug-ravenel/otherpapers/BBD-pervers-2x.pdf}.
\bibitem{PHaines} T. J. Haines, \emph{A proof of the Kazhdan--Lusztig purity theorem via the decomposition theorem of BBD}, author note. Theorem 2.0.1, p.~2; Theorem 5.0.4, p.~7. \url{https://math.umd.edu/~tjh/KL_purity1.pdf}.
\bibitem{PMilne} J. S. Milne, \emph{Lectures on \'Etale Cohomology}, version 2.21. Theorem 16.1; Example 16.3; Proposition 18.4(d); Theorem 24.1. \url{https://www.jmilne.org/math/CourseNotes/LEC.pdf}.
\bibitem{PStacks} The Stacks Project, Proposition 59.55.2, tag 03QP, finite-morphism direct images and geometric stalks. \url{https://stacks.math.columbia.edu/tag/03QP}.
\bibitem{PStacksCM} The Stacks Project, Lemma 10.104.2, tag 02JN, the regular-parameter criterion in Cohen--Macaulay local rings. \url{https://stacks.math.columbia.edu/tag/02JN}.
\bibitem{PWeyl} R. Borcherds, M. Haiman, N. Reshetikhin and V. Serganova, notes by T. Johnson-Freyd, \emph{Lectures in Lie Groups and Quantum Groups}, Part I (2016). Example 5.2.0.5 and Proposition 5.2.0.7, pp.~68--69; Theorem 6.1.1.30, pp.~102--103; Theorem 7.6.1.3, p.~167. \url{https://math.berkeley.edu/~serganov/261/part1.pdf}.
\bibitem{PDiscriminant} C. McCrory, T. Shifrin and R. Varley, \emph{Siegel modular forms generated by invariants of cubic hypersurfaces}, arXiv:alg-geom/9403005v3. Section 3.1 and Section 3.1.2(i), printed p.~15: discriminant invariance, degree and singular-hypersurface locus. \url{https://arxiv.org/pdf/alg-geom/9403005}.
\bibitem{PBuse} L. Bus\'e and J.-P. Jouanolou, \emph{A computational approach to the discriminant of homogeneous polynomials}, arXiv:1210.4697v1. Definition 4.6 and Proposition 4.7(i), p.~40; Example 4.12, p.~42; Proposition 4.13, pp.~42--43. \url{https://arxiv.org/pdf/1210.4697}.
\bibitem{PUniversalDenom} H. Derksen, \emph{Universal denominators of Hilbert series}, J. Algebra 285 (2005), 586--607. Definition 1.5; Theorem 1.10; Example 4.7, pp.~597--598. \url{https://sites.lsa.umich.edu/hderksen/wp-content/uploads/sites/614/2018/05/A.I.a.19.pdf}.
\bibitem{PMacdonald} I. G. Macdonald, \emph{Symmetric Functions and Hall Polynomials}, second edition, Oxford, 1995. I (7.3), (7.6)(ii), (7.9); I \S5 Example 2; Appendix A (7.3)--(7.4). \url{https://math.berkeley.edu/~corteel/MATH249/macdonald.pdf}.
\bibitem{PBorger} J. Borger, \emph{$\Lambda$-rings and the field with one element}, arXiv:0906.3146, supplied file \texttt{lrfoe13.tex}. Section 1.1, paragraph on the right adjoint of the big Witt functor. \url{https://arxiv.org/abs/0906.3146}.
\bibitem{PSGA6} P. Berthelot, A. Grothendieck and L. Illusie, eds., \emph{Th\'eorie des intersections et th\'eor\`eme de Riemann--Roch}, SGA 6, LNM 225, 1971. Expos\'e V, Definitions 2.1, 2.4; Expos\'e VI, Theorem 3.2. Original-page retrieval was not completed in this run; the check register distinguishes this background citation from directly inspected primary pages.
\bibitem{PBouc} S. Bouc and K. R\"okaeus, \emph{A note on the $\lambda$-structure on the Burnside ring}. Section 1 and Theorem 1.1, equation (1.2). \url{https://www.lamfa.u-picardie.fr/bouc/Burnside2.pdf}.
\bibitem{PBaez} J. Baez, J. Moeller and T. Trimble, \emph{Schur Functors and Categorified Plethysm}, arXiv:2106.00190. Theorem 6.11. \url{https://arxiv.org/abs/2106.00190}.
\bibitem{PZib1} M. Zibrowius, \emph{Symmetric representation rings are $\lambda$-rings}, New York J. Math. 21 (2015), 1055--1092. Main theorem, p.~1056; Theorem 3.5. \url{https://nyjm.albany.edu/j/2015/21-49v.pdf}.
\bibitem{PZib2} M. Zibrowius, \emph{The gamma-filtration on the Witt ring of a scheme}, arXiv:1411.3198v2. Theorem 2; Introduction, p.~2. \url{https://arxiv.org/abs/1411.3198}.
\bibitem{PPPS} I. Pak, G. Panova and J. P. Swanson, \emph{A combinatorial interpretation for certain plethysm and Kronecker coefficients}, arXiv:2511.02312v3, July 2026 version. Theorems 1.1--1.2 and Lemma 2.1. \url{https://arxiv.org/abs/2511.02312}.
\bibitem{PArchive} A. Grothendieck, supplied manuscript transcriptions, cotes/scans 133/6, 133/8, 133/9, 133/19, 133/28, 2/17, 43/37, 119/217, 12/41, 16/38, 4/7, 79/51 and 89/18; filenames and hashes in the source manifest. The associated supplied study \texttt{SA17\_STUDY\_cote133\_lambda.md} is treated separately from the transcriptions and is corrected where PL-01 applies.


% Items to append inside the final paper's thebibliography environment.
\bibitem{brion2002} M. Brion, \emph{Positivity in the Grothendieck group of complex flag varieties}, J. Algebra 258 (2002), 137--159. \url{https://arxiv.org/abs/math/0105254}. Theorem 1 and Lemma 2.
\bibitem{buch2002} A. S. Buch, \emph{A Littlewood--Richardson rule for the $K$-theory of Grassmannians}, Acta Math. 189 (2002), 37--78. \url{https://arxiv.org/abs/math/0004137}. Locators in arXiv v2: Corollary 4.7, Theorem 5.4, Corollary 5.5, Example 5.6, Corollary 6.7.
\bibitem{kt1999} A. Knutson and T. Tao, \emph{The honeycomb model of $GL_n(\mathbb C)$ tensor products I: proof of the saturation conjecture}, J. Amer. Math. Soc. 12 (1999), 1055--1090. \url{https://arxiv.org/abs/math/9807160}. Section 6, Theorem 3 and following Corollary.
\bibitem{groth1953} A. Grothendieck, \emph{R\'esum\'e de la th\'eorie m\'etrique des produits tensoriels topologiques}, Bol. Soc. Mat. S\~ao Paulo 8 (1953), 1--79. Section 4, no. 2, Theorems 1--2. \url{https://webusers.imj-prg.fr/~leila.schneps/grothendieckcircle/AG/AG-22.pdf}.
\bibitem{alon-naor2006} N. Alon and A. Naor, \emph{Approximating the cut-norm via Grothendieck's inequality}, SIAM J. Comput. 35 (2006), 787--803. Lemma 2.1, Proposition 2.2, Theorem 4.3. \url{https://web.math.princeton.edu/~nalon/PDFS/cutnorm3.pdf}.
\bibitem{rs2009} P. Raghavendra and D. Steurer, \emph{Towards computing the Grothendieck constant}, Proc. SODA (2009), 525--534. Theorem 1.1. \url{https://www.dsteurer.org/paper/grothendieck.pdf}.
\bibitem{chtw2010} R. Cleve, P. H\o yer, B. Toner and J. Watrous, \emph{Consequences and limits of nonlocal strategies}, arXiv:quant-ph/0404076v2, 11 January 2010. Theorem 4 (Tsirelson), Proposition 5, Theorem 7. \url{https://arxiv.org/abs/quant-ph/0404076}.
\bibitem{friedland-lim2018} S. Friedland and L.-H. Lim, \emph{Nuclear norm of higher-order tensors}, Math. Comp. 87 (2018), 1255--1281. \url{https://arxiv.org/abs/1410.6072}. Equation (1), Section 2, Theorems 8.1 and 8.10 in arXiv v3.
\bibitem{sga3-I} M. Demazure, \emph{Structures alg\'ebriques. Cohomologie des groupes}, SGA 3, Expos\'e I, Sections 4.2 and 4.5. \url{https://webusers.imj-prg.fr/~patrick.polo/SGA3/origExp1.pdf}.
\bibitem{fga221} A. Grothendieck, \emph{Techniques de construction et th\'eor\`emes d'existence en g\'eom\'etrie alg\'ebrique IV: les sch\'emas de Hilbert}, S\'eminaire Bourbaki, Expos\'e 221 (May 1961), 249--276. Theorems 3.1--3.2, Proposition 3.8 and subsequent Hilbert-scheme definition, Remark 3.9. \url{https://www.numdam.org/article/SB_1960-1961__6__249_0.pdf}.


% Append these entries inside the final paper's thebibliography environment.
\bibitem{CAlpoge26} L. Alp\"oge, \emph{Remark}, supplied four-page note, 20 November 2026. \url{https://alpo.ge/remark.pdf}. Supplied \texttt{ALPOGE\_Main.lean}: \texttt{riemannHypothesis}, \texttt{riemannZeta\_eq\_zero\_re}, \texttt{LFunction\_ne\_zero\_of\_half\_lt\_re}, \texttt{LFunction\_eq\_zero\_re}, \texttt{LFunction\_eq\_zero\_re\_of\_isPrimitive}; types in \texttt{ALPOGE\_Statements.lean}. Used as cited premises as requested, without claiming a Lean rebuild.
\bibitem{CKoiran96} P. Koiran, \emph{Hilbert's Nullstellensatz is in the polynomial hierarchy}, J. Complexity 12 (1996), 273--286, \url{https://doi.org/10.1006/jcom.1996.0019}. Exact numbering here is that of the DIMACS long version, Technical Report 96-27: Theorem 2; Section 5, Theorem 9 and Corollary 1. \url{https://perso.ens-lyon.fr/pascal.koiran/Publis/DIMACS.96-27.ps}. Erratum, 8 August 1997: \url{https://perso.ens-lyon.fr/pascal.koiran/Publis/hn.erratum.ps}.
\bibitem{CRojas03} J. M. Rojas, \emph{Dedekind zeta functions and the complexity of Hilbert's Nullstellensatz}, arXiv:math/0301111v2 (2003). Theorems 1 and 2. \url{https://arxiv.org/abs/math/0301111}.
\bibitem{CGCTV} K. D. Mulmuley, \emph{Geometric Complexity Theory V: Efficient algorithms for Noether Normalization}, arXiv:1209.5993v5. Theorems 1.7, 1.9, 2.10, 3.4, 4.1, 4.7, 4.9, 4.10, 5.11, 5.14. \url{https://arxiv.org/abs/1209.5993}. The supplied LaTeX and the retrieved PDF use the numbering cited here.
\bibitem{CMiller76} G. L. Miller, \emph{Riemann's hypothesis and tests for primality}, J. Comput. System Sci. 13 (1976), 300--317. Theorem 2 and Appendix. \url{https://www.cs.cmu.edu/~glmiller/Publications/Papers/Mi76.pdf}.
\bibitem{CLenstra79} H. W. Lenstra, Jr., \emph{Miller's primality test}, Inform. Process. Lett. 8 (1979), 86--88. Theorem 1. \url{https://pub.math.leidenuniv.nl/~lenstrahw/PUBLICATIONS/1979e/art.pdf}.
\bibitem{CBach90} E. Bach, \emph{Explicit bounds for primality testing and related problems}, Math. Comp. 55 (1990), 355--380. Section 6, Theorems 2--3, pp.~369--373; Section 7, Theorem 4. \url{https://doi.org/10.1090/S0025-5718-1990-1023756-8}.
\bibitem{CAKS04} M. Agrawal, N. Kayal and N. Saxena, \emph{PRIMES is in P}, Ann. of Math. 160 (2004), 781--793. Theorems 4.1 and 5.1 (also Theorem 5.3 for the improved bound). \url{https://annals.math.princeton.edu/wp-content/uploads/annals-v160-n2-p12.pdf}.
\bibitem{COAI26} OpenAI, \emph{The Exponential Time Hypothesis is False}, OpenAI Math Release preprint, 17 October 2026, revised 25 October 2026. Theorem 1.1, Corollaries 1.2--1.4, Sections 1.1.1--1.1.3, Appendix C.1, Appendix D, and the supplied files \texttt{quantitative-reduction.tex}, \texttt{quantitative-contraction.tex}, \texttt{quantitative-circuits.tex}. \url{https://github.com/openai/math/blob/main/preprints/The-Exponential-Time-Hypothesis-is-False-October-25-2026/paper.pdf}.
\bibitem{CCGGM26} N. Chukhin, K. Gajulapalli, A. Golovnev and I. Mihajlin, \emph{A Little Derandomization Goes a Long Way: Unconditional Circuit Lower Bounds using Low-Degree PIT}, arXiv:2611.07972v1, 5 November 2026; ECCC TR26-313, published 9 November 2026. Theorems 1, 5 and 6; Sections 5 and 8.2. \url{https://arxiv.org/abs/2611.07972}; \url{https://eccc.weizmann.ac.il/report/2026/313/}.
\bibitem{CLandsberg13} J. M. Landsberg, \emph{Geometric complexity theory: an introduction for geometers}, arXiv:1305.7387v3, 11 December 2013. Section 8.2, equation (22), p.~28. \url{https://arxiv.org/abs/1305.7387}.
\bibitem{CBarvinokGF} A. Barvinok, \emph{The Complexity of Generating Functions for Integer Points in Polyhedra and Beyond}. Theorem 2.5 (Brion's identity for rational polyhedra), printed p.~9. Author-hosted text: \url{https://sites.lsa.umich.edu/barvinok/wp-content/uploads/sites/1434/2025/05/lecture.pdf}.


\end{thebibliography}\end{document}